\chapt{सुन्दरकाण्डः}

\sect{प्रथमः सर्गः}

\twolineshloka
{ततो रावणनीतायाः सीतायाः शत्रुकर्शनः}
{इयेष पदमन्वेष्टुं चारणाचरिते पथि} %||5-1-1||

\twolineshloka
{अथ वैदूर्यवर्णेषु शाद्वलेषु महाबलः}
{धीरः सलिलकल्पेषु विचचार यथासुखम्} %||5-1-2||

\twolineshloka
{द्विजान्वित्रासयन्धीमानुरसा पादपान्हरन्}
{मृगांश्च सुबहून्निघ्नन्प्रवृद्ध इव केसरी} %||5-1-3||

\twolineshloka
{नीललोहितमाञ्जिष्ठपद्मवर्णैः सितासितैः}
{स्वभावविहितैश्चित्रैर्धातुभिः समलंकृतम्} %||5-1-4||

\twolineshloka
{कामरूपिभिराविष्टमभीक्ष्णं सपरिच्छदैः}
{यक्षकिंनरगन्धर्वैर्देवकल्पैश्च पन्नगैः} %||5-1-5||

\twolineshloka
{स तस्य गिरिवर्यस्य तले नागवरायुते}
{तिष्ठन्कपिवरस्तत्र ह्रदे नाग इवाबभौ} %||5-1-6||

\twolineshloka
{स सूर्याय महेन्द्राय पवनाय स्वयम्भुवे}
{भूतेभ्यश्चाञ्जलिं कृत्वा चकार गमने मतिम्} %||5-1-7||

\twolineshloka
{अञ्जलिं प्राङ्मुखः कुर्वन्पवनायात्मयोनये}
{ततो हि ववृधे गन्तुं दक्षिणो दक्षिणां दिशम्} %||5-1-8||

\twolineshloka
{प्लवंगप्रवरैर्दृष्टः प्लवने कृतनिश्चयः}
{ववृधे रामवृद्ध्यर्थं समुद्र इव पर्वसु} %||5-1-9||

\twolineshloka
{निष्प्रमाण शरीरः सँल्लिलङ्घयिषुरर्णवम्}
{बाहुभ्यां पीडयामास चरणाभ्यां च पर्वतम्} %||5-1-10||

\twolineshloka
{स चचालाचलाश्चारु मुहूर्तं कपिपीडितः}
{तरूणां पुष्पिताग्राणां सर्वं पुष्पमशातयत्} %||5-1-11||

\twolineshloka
{तेन पादपमुक्तेन पुष्पौघेण सुगन्धिना}
{सर्वतः संवृतः शैलो बभौ पुष्पमयो यथा} %||5-1-12||

\twolineshloka
{तेन चोत्तमवीर्येण पीड्यमानः स पर्वतः}
{सलिलं संप्रसुस्राव मदं मत्त इव द्विपः} %||5-1-13||

\threelineshloka
{पीड्यमानस्तु बलिना महेन्द्रस्तेन पर्वतः}
{रीतिर्निर्वर्तयामास काञ्चनाञ्जनराजतीः}
{मुमोच च शिलाः शैलो विशालाः समनःशिलाः} %||5-1-14||

\twolineshloka
{गिरिणा पीड्यमानेन पीड्यमानानि सर्वशः}
{गुहाविष्टानि भूतानि विनेदुर्विकृतैः स्वरैः} %||5-1-15||

\twolineshloka
{स महासत्त्वसंनादः शैलपीडानिमित्तजः}
{पृथिवीं पूरयामास दिशश्चोपवनानि च} %||5-1-16||

\twolineshloka
{शिरोभिः पृथुभिः सर्पा व्यक्तस्वस्तिकलक्षणैः}
{वमन्तः पावकं घोरं ददंशुर्दशनैः शिलाः} %||5-1-17||

\twolineshloka
{तास्तदा सविषैर्दष्टाः कुपितैस्तैर्महाशिलाः}
{जज्वलुः पावकोद्दीप्ता विभिदुश्च सहस्रधा} %||5-1-18||

\twolineshloka
{यानि चौषधजालानि तस्मिञ्जातानि पर्वते}
{विषघ्नान्यपि नागानां न शेकुः शमितुं विषम्} %||5-1-19||

\twolineshloka
{भिद्यतेऽयं गिरिर्भूतैरिति मत्वा तपस्विनः}
{त्रस्ता विद्याधरास्तस्मादुत्पेतुः स्त्रीगणैः सह} %||5-1-20||

\twolineshloka
{पानभूमिगतं हित्वा हैममासनभाजनम्}
{पात्राणि च महार्हाणि करकांश्च हिरण्मयान्} %||5-1-21||

\twolineshloka
{लेह्यानुच्चावचान्भक्ष्यान्मांसानि विविधानि च}
{आर्षभाणि च चर्माणि खड्गांश्च कनकत्सरून्} %||5-1-22||

\twolineshloka
{कृतकण्ठगुणाः क्षीबा रक्तमाल्यानुलेपनाः}
{रक्ताक्षाः पुष्कराक्षाश्च गगनं प्रतिपेदिरे} %||5-1-23||

\twolineshloka
{हारनूपुरकेयूर पारिहार्य धराः स्त्रियः}
{विस्मिताः सस्मितास्तस्थुराकाशे रमणैः सह} %||5-1-24||

\twolineshloka
{दर्शयन्तो महाविद्यां विद्याधरमहर्षयः}
{सहितास्तस्थुराकाशे वीक्षां चक्रुश्च पर्वतम्} %||5-1-25||

\twolineshloka
{शुश्रुवुश्च तदा शब्दमृषीणां भावितात्मनाम्}
{चारणानां च सिद्धानां स्थितानां विमलेऽम्बरे} %||5-1-26||

\twolineshloka
{एष पर्वतसंकाशो हनूमान्मारुतात्मजः}
{तितीर्षति महावेगं समुद्रं मकरालयम्} %||5-1-27||

\twolineshloka
{रामार्थं वानरार्थं च चिकीर्षन्कर्म दुष्करम्}
{समुद्रस्य परं पारं दुष्प्रापं प्राप्तुमिच्छति} %||5-1-28||

\twolineshloka
{दुधुवे च स रोमाणि चकम्पे चाचलोपमः}
{ननाद च महानादं सुमहानिव तोयदः} %||5-1-29||

\twolineshloka
{आनुपूर्व्याच्च वृत्तं च लाङ्गूलं रोमभिश्चितम्}
{उत्पतिष्यन्विचिक्षेप पक्षिराज इवोरगम्} %||5-1-30||

\twolineshloka
{तस्य लाङ्गूलमाविद्धमतिवेगस्य पृष्ठतः}
{ददृशे गरुडेनेव ह्रियमाणो महोरगः} %||5-1-31||

\twolineshloka
{बाहू संस्तम्भयामास महापरिघसंनिभौ}
{ससाद च कपिः कट्यां चरणौ संचुकोप च} %||5-1-32||

\twolineshloka
{संहृत्य च भुजौ श्रीमांस्तथैव च शिरोधराम्}
{तेजः सत्त्वं तथा वीर्यमाविवेश स वीर्यवान्} %||5-1-33||

\twolineshloka
{मार्गमालोकयन्दूरादूर्ध्वप्रणिहितेक्षणः}
{रुरोध हृदये प्राणानाकाशमवलोकयन्} %||5-1-34||

\threelineshloka
{पद्भ्यां दृढमवस्थानं कृत्वा स कपिकुञ्जरः}
{निकुञ्च्य कर्णौ हनुमानुत्पतिष्यन्महाबलः}
{वानरान्वानरश्रेष्ठ इदं वचनमब्रवीत्} %||5-1-35||

\twolineshloka
{यथा राघवनिर्मुक्तः शरः श्वसनविक्रमः}
{गच्छेत्तद्वद्गमिष्यामि लङ्कां रावणपालिताम्} %||5-1-36||

\twolineshloka
{न हि द्रक्ष्यामि यदि तां लङ्कायां जनकात्मजाम्}
{अनेनैव हि वेगेन गमिष्यामि सुरालयम्} %||5-1-37||

\twolineshloka
{यदि वा त्रिदिवे सीतां न द्रक्ष्यामि कृतश्रमः}
{बद्ध्वा राक्षसराजानमानयिष्यामि रावणम्} %||5-1-38||

\twolineshloka
{सर्वथा कृतकार्योऽहमेष्यामि सह सीतया}
{आनयिष्यामि वा लङ्कां समुत्पाट्य सरावणाम्} %||5-1-39||

\twolineshloka
{एवमुक्त्वा तु हनुमान्वानरान्वानरोत्तमः}
{उत्पपाताथ वेगेन वेगवानविचारयन्} %||5-1-40||

\twolineshloka
{समुत्पतति तस्मिंस्तु वेगात्ते नगरोहिणः}
{संहृत्य विटपान्सर्वान्समुत्पेतुः समन्ततः} %||5-1-41||

\twolineshloka
{स मत्तकोयष्टिभकान्पादपान्पुष्पशालिनः}
{उद्वहन्नूरुवेगेन जगाम विमलेऽम्बरे} %||5-1-42||

\twolineshloka
{ऊरुवेगोद्धता वृक्षा मुहूर्तं कपिमन्वयुः}
{प्रस्थितं दीर्घमध्वानं स्वबन्धुमिव बान्धवाः} %||5-1-43||

\twolineshloka
{तमूरुवेगोन्मथिताः सालाश्चान्ये नगोत्तमाः}
{अनुजग्मुर्हनूमन्तं सैन्या इव महीपतिम्} %||5-1-44||

\twolineshloka
{सुपुष्पिताग्रैर्बहुभिः पादपैरन्वितः कपिः}
{हनुमान्पर्वताकारो बभूवाद्भुतदर्शनः} %||5-1-45||

\twolineshloka
{सारवन्तोऽथ ये वृक्षा न्यमज्जँल्लवणाम्भसि}
{भयादिव महेन्द्रस्य पर्वता वरुणालये} %||5-1-46||

\twolineshloka
{स नानाकुसुमैः कीर्णः कपिः साङ्कुरकोरकैः}
{शुशुभे मेघसंकाशः खद्योतैरिव पर्वतः} %||5-1-47||

\twolineshloka
{विमुक्तास्तस्य वेगेन मुक्त्वा पुष्पाणि ते द्रुमाः}
{अवशीर्यन्त सलिले निवृत्ताः सुहृदो यथा} %||5-1-48||

\twolineshloka
{लघुत्वेनोपपन्नं तद्विचित्रं सागरेऽपतत्}
{द्रुमाणां विविधं पुष्पं कपिवायुसमीरितम्} %||5-1-49||

\twolineshloka
{पुष्पौघेणानुबद्धेन नानावर्णेन वानरः}
{बभौ मेघ इवोद्यन्वै विद्युद्गणविभूषितः} %||5-1-50||

\twolineshloka
{तस्य वेगसमुद्भूतैः पुष्पैस्तोयमदृश्यत}
{ताराभिरभिरामाभिरुदिताभिरिवाम्बरम्} %||5-1-51||

\twolineshloka
{तस्याम्बरगतौ बाहू ददृशाते प्रसारितौ}
{पर्वताग्राद्विनिष्क्रान्तौ पञ्चास्याविव पन्नगौ} %||5-1-52||

\twolineshloka
{पिबन्निव बभौ चापि सोर्मिजालं महार्णवम्}
{पिपासुरिव चाकाशं ददृशे स महाकपिः} %||5-1-53||

\twolineshloka
{तस्य विद्युत्प्रभाकारे वायुमार्गानुसारिणः}
{नयने विप्रकाशेते पर्वतस्थाविवानलौ} %||5-1-54||

\twolineshloka
{पिङ्गे पिङ्गाक्षमुख्यस्य बृहती परिमण्डले}
{चक्षुषी संप्रकशेते चन्द्रसूर्याविव स्थितौ} %||5-1-55||

\twolineshloka
{मुखं नासिकया तस्य ताम्रया ताम्रमाबभौ}
{संध्यया समभिस्पृष्टं यथा सूर्यस्य मण्डलम्} %||5-1-56||

\twolineshloka
{लाङ्गलं च समाविद्धं प्लवमानस्य शोभते}
{अम्बरे वायुपुत्रस्य शक्रध्वज इवोच्छ्रितः} %||5-1-57||

\twolineshloka
{लाङ्गूलचक्रेण महाञ्शुक्लदंष्ट्रोऽनिलात्मजः}
{व्यरोचत महाप्राज्ञः परिवेषीव भास्करः} %||5-1-58||

\twolineshloka
{स्फिग्देशेनाभिताम्रेण रराज स महाकपिः}
{महता दारितेनेव गिरिर्गैरिकधातुना} %||5-1-59||

\twolineshloka
{तस्य वानरसिंहस्य प्लवमानस्य सागरम्}
{कक्षान्तरगतो वायुर्जीमूत इव गर्जति} %||5-1-60||

\twolineshloka
{खे यथा निपतत्युल्का उत्तरान्ताद्विनिःसृता}
{दृश्यते सानुबन्धा च तथा स कपिकुञ्जरः} %||5-1-61||

\twolineshloka
{पतत्पतंगसंकाशो व्यायतः शुशुभे कपिः}
{प्रवृद्ध इव मातंगः कक्ष्यया बध्यमानया} %||5-1-62||

\twolineshloka
{उपरिष्टाच्छरीरेण छायया चावगाढया}
{सागरे मारुताविष्टा नौरिवासीत्तदा कपिः} %||5-1-63||

\twolineshloka
{यं यं देशं समुद्रस्य जगाम स महाकपिः}
{स स तस्याङ्गवेगेन सोन्माद इव लक्ष्यते} %||5-1-64||

\twolineshloka
{सागरस्योर्मिजालानामुरसा शैलवर्ष्मणाम्}
{अभिघ्नंस्तु महावेगः पुप्लुवे स महाकपिः} %||5-1-65||

\twolineshloka
{कपिवातश्च बलवान्मेघवातश्च निःसृतः}
{सागरं भीमनिर्घोषं कम्पयामासतुर्भृशम्} %||5-1-66||

\twolineshloka
{विकर्षन्नूर्मिजालानि बृहन्ति लवणाम्भसि}
{अत्यक्रामन्महावेगस्तरङ्गान्गणयन्निव} %||5-1-67||

\twolineshloka
{प्लवमानं समीक्ष्याथ भुजङ्गाः सागरालयाः}
{व्योम्नि तं कपिशार्दूलं सुपर्णमिति मेनिरे} %||5-1-68||

\twolineshloka
{दशयोजनविस्तीर्णा त्रिंशद्योजनमायता}
{छाया वानरसिंहस्य जले चारुतराभवत्} %||5-1-69||

\twolineshloka
{श्वेताभ्रघनराजीव वायुपुत्रानुगामिनी}
{तस्य सा शुशुभे छाया वितता लवणाम्भसि} %||5-1-70||

\twolineshloka
{प्लवमानं तु तं दृष्ट्वा प्लवगं त्वरितं तदा}
{ववृषुः पुष्पवर्षाणि देवगन्धर्वदानवाः} %||5-1-71||

\twolineshloka
{तताप न हि तं सूर्यः प्लवन्तं वानरेश्वरम्}
{सिषेवे च तदा वायू रामकार्यार्थसिद्धये} %||5-1-72||

\twolineshloka
{ऋषयस्तुष्टुवुश्चैनं प्लवमानं विहायसा}
{जगुश्च देवगन्धर्वाः प्रशंसन्तो महौजसं} %||5-1-73||

\twolineshloka
{नागाश्च तुष्टुवुर्यक्षा रक्षांसि विबुधाः खगाः}
{प्रेक्ष्याकाशे कपिवरं सहसा विगतक्लमम्} %||5-1-74||

\twolineshloka
{तस्मिन्प्लवगशार्दूले प्लवमाने हनूमति}
{इक्ष्वाकुकुलमानार्थी चिन्तयामास सागरः} %||5-1-75||

\twolineshloka
{साहाय्यं वानरेन्द्रस्य यदि नाहं हनूमतः}
{करिष्यामि भविष्यामि सर्ववाच्यो विवक्षताम्} %||5-1-76||

\twolineshloka
{अहमिक्ष्वाकुनाथेन सगरेण विवर्धितः}
{इक्ष्वाकुसचिवश्चायं नावसीदितुमर्हति} %||5-1-77||

\twolineshloka
{तथा मया विधातव्यं विश्रमेत यथा कपिः}
{शेषं च मयि विश्रान्तः सुखेनातिपतिष्यति} %||5-1-78||

\twolineshloka
{इति कृत्वा मतिं साध्वीं समुद्रश्छन्नमम्भसि}
{हिरण्यनाभं मैनाकमुवाच गिरिसत्तमम्} %||5-1-79||

\twolineshloka
{त्वमिहासुरसंघानां पातालतलवासिनाम्}
{देवराज्ञा गिरिश्रेष्ठ परिघः संनिवेशितः} %||5-1-80||

\twolineshloka
{त्वमेषां ज्ञातवीर्याणां पुनरेवोत्पतिष्यताम्}
{पातालस्याप्रमेयस्य द्वारमावृत्य तिष्ठसि} %||5-1-81||

\twolineshloka
{तिर्यगूर्ध्वमधश्चैव शक्तिस्ते शैलवर्धितुम्}
{तस्मात्संचोदयामि त्वामुत्तिष्ठ नगसत्तम} %||5-1-82||

\twolineshloka
{स एष कपिशार्दूलस्त्वामुपर्येति वीर्यवान्}
{हनूमान्रामकार्यार्थं भीमकर्मा खमाप्लुतः} %||5-1-83||

\twolineshloka
{तस्य साह्यं मया कार्यमिक्ष्वाकुकुलवर्तिनः}
{मम इक्ष्वाकवः पूज्याः परं पूज्यतमास्तव} %||5-1-84||

\twolineshloka
{कुरु साचिव्यमस्माकं न नः कार्यमतिक्रमेत्}
{कर्तव्यमकृतं कार्यं सतां मन्युमुदीरयेत्} %||5-1-85||

\twolineshloka
{सलिलादूर्ध्वमुत्तिष्ठ तिष्ठत्वेष कपिस्त्वयि}
{अस्माकमतिथिश्चैव पूज्यश्च प्लवतां वरः} %||5-1-86||

\twolineshloka
{चामीकरमहानाभ देवगन्धर्वसेवित}
{हनूमांस्त्वयि विश्रान्तस्ततः शेषं गमिष्यति} %||5-1-87||

\twolineshloka
{काकुत्स्थस्यानृशंस्यं च मैथिल्याश्च विवासनम्}
{श्रमं च प्लवगेन्द्रस्य समीक्ष्योत्थातुमर्हसि} %||5-1-88||

\twolineshloka
{हिरण्यनाभो मैनाको निशम्य लवणाम्भसः}
{उत्पपात जलात्तूर्णं महाद्रुमलतायुतः} %||5-1-89||

\twolineshloka
{स सागरजलं भित्त्वा बभूवात्युत्थितस्तदा}
{यथा जलधरं भित्त्वा दीप्तरश्मिर्दिवाकरः} %||5-1-90||

\twolineshloka
{शातकुम्भमयैः शृङ्गैः सकिंनरमहोरगैः}
{आदित्योदयसंकाशैरालिखद्भिरिवाम्बरम्} %||5-1-91||

\twolineshloka
{तस्य जाम्बूनदैः शृङ्गैः पर्वतस्य समुत्थितैः}
{आकाशं शस्त्रसंकाशमभवत्काञ्चनप्रभम्} %||5-1-92||

\twolineshloka
{जातरूपमयैः शृङ्गैर्भ्राजमानैः स्वयं प्रभैः}
{आदित्यशतसंकाशः सोऽभवद्गिरिसत्तमः} %||5-1-93||

\twolineshloka
{तमुत्थितमसंगेन हनूमानग्रतः स्थितम्}
{मध्ये लवणतोयस्य विघ्नोऽयमिति निश्चितः} %||5-1-94||

\twolineshloka
{स तमुच्छ्रितमत्यर्थं महावेगो महाकपिः}
{उरसा पातयामास जीमूतमिव मारुतः} %||5-1-95||

\twolineshloka
{स तदा पातितस्तेन कपिना पर्वतोत्तमः}
{बुद्ध्वा तस्य कपेर्वेगं जहर्ष च ननन्द च} %||5-1-96||

\threelineshloka
{तमाकाशगतं वीरमाकाशे समवस्थितम्}
{प्रीतो हृष्टमना वाक्यमब्रवीत्पर्वतः कपिम्}
{मानुषं धरयन्रूपमात्मनः शिखरे स्थितः} %||5-1-97||

\twolineshloka
{दुष्करं कृतवान्कर्म त्वमिदं वानरोत्तम}
{निपत्य मम शृङ्गेषु विश्रमस्व यथासुखम्} %||5-1-98||

\twolineshloka
{राघावस्य कुले जातैरुदधिः परिवर्धितः}
{स त्वां रामहिते युक्तं प्रत्यर्चयति सागरः} %||5-1-99||

\twolineshloka
{कृते च प्रतिकर्तव्यमेष धर्मः सनातनः}
{सोऽयं तत्प्रतिकारार्थी त्वत्तः संमानमर्हति} %||5-1-100||

\threelineshloka
{त्वन्निमित्तमनेनाहं बहुमानात्प्रचोदितः}
{योजनानां शतं चापि कपिरेष समाप्लुतः}
{तव सानुषु विश्रान्तः शेषं प्रक्रमतामिति} %||5-1-101||

\threelineshloka
{तिष्ठ त्वं हरिशार्दूल मयि विश्रम्य गम्यताम्}
{तदिदं गन्धवत्स्वादु कन्दमूलफलं बहु}
{तदास्वाद्य हरिश्रेष्ठ विश्रान्तोऽनुगमिष्यसि} %||5-1-102||

\twolineshloka
{अस्माकमपि संबन्धः कपिमुख्यस्त्वयास्ति वै}
{प्रख्यातस्त्रिषु लोकेषु महागुणपरिग्रहः} %||5-1-103||

\twolineshloka
{वेगवन्तः प्लवन्तो ये प्लवगा मारुतात्मज}
{तेषां मुख्यतमं मन्ये त्वामहं कपिकुञ्जर} %||5-1-104||

\twolineshloka
{अतिथिः किल पूजार्हः प्राकृतोऽपि विजानता}
{धर्मं जिज्ञासमानेन किं पुनर्यादृशो भवान्} %||5-1-105||

\twolineshloka
{त्वं हि देववरिष्ठस्य मारुतस्य महात्मनः}
{पुत्रस्तस्यैव वेगेन सदृशः कपिकुञ्जर} %||5-1-106||

\twolineshloka
{पूजिते त्वयि धर्मज्ञ पूजां प्राप्नोति मारुतः}
{तस्मात्त्वं पूजनीयो मे शृणु चाप्यत्र कारणम्} %||5-1-107||

\twolineshloka
{पूर्वं कृतयुगे तात पर्वताः पक्षिणोऽभवन्}
{तेऽपि जग्मुर्दिशः सर्वा गरुडानिलवेगिनः} %||5-1-108||

\twolineshloka
{ततस्तेषु प्रयातेषु देवसंघाः सहर्षिभिः}
{भूतानि च भयं जग्मुस्तेषां पतनशङ्कया} %||5-1-109||

\twolineshloka
{ततः क्रुद्धः सहस्राक्षः पर्वतानां शतक्रतुः}
{पक्षांश्चिच्छेद वज्रेण तत्र तत्र सहस्रशः} %||5-1-110||

\twolineshloka
{स मामुपगतः क्रुद्धो वज्रमुद्यम्य देवराट्}
{ततोऽहं सहसा क्षिप्तः श्वसनेन महात्मना} %||5-1-111||

\twolineshloka
{अस्मिँल्लवणतोये च प्रक्षिप्तः प्लवगोत्तम}
{गुप्तपक्षः समग्रश्च तव पित्राभिरक्षितः} %||5-1-112||

\twolineshloka
{ततोऽहं मानयामि त्वां मान्यो हि मम मारुतः}
{त्वया मे ह्येष संबन्धः कपिमुख्य महागुणः} %||5-1-113||

\twolineshloka
{अस्मिन्नेवंगते कार्ये सागरस्य ममैव च}
{प्रीतिं प्रीतमना कर्तुं त्वमर्हसि महाकपे} %||5-1-114||

\twolineshloka
{श्रमं मोक्षय पूजां च गृहाण कपिसत्तम}
{प्रीतिं च बहुमन्यस्व प्रीतोऽस्मि तव दर्शनात्} %||5-1-115||

\twolineshloka
{एवमुक्तः कपिश्रेष्ठस्तं नगोत्तममब्रवीत्}
{प्रीतोऽस्मि कृतमातिथ्यं मन्युरेषोऽपनीयताम्} %||5-1-116||

\twolineshloka
{त्वरते कार्यकालो मे अहश्चाप्यतिवर्तते}
{प्रतिज्ञा च मया दत्ता न स्थातव्यमिहान्तरा} %||5-1-117||

\twolineshloka
{इत्युक्त्वा पाणिना शैलमालभ्य हरिपुंगवः}
{जगामाकाशमाविश्य वीर्यवान्प्रहसन्निव} %||5-1-118||

\twolineshloka
{स पर्वतसमुद्राभ्यां बहुमानादवेक्षितः}
{पूजितश्चोपपन्नाभिराशीर्भिरनिलात्मजः} %||5-1-119||

\twolineshloka
{अथोर्ध्वं दूरमुत्पत्य हित्वा शैलमहार्णवौ}
{पितुः पन्थानमास्थाय जगाम विमलेऽम्बरे} %||5-1-120||

\twolineshloka
{भूयश्चोर्ध्वगतिं प्राप्य गिरिं तमवलोकयन्}
{वायुसूनुर्निरालम्बे जगाम विमलेऽम्बरे} %||5-1-121||

\twolineshloka
{तद्द्वितीयं हनुमतो दृष्ट्वा कर्म सुदुष्करम्}
{प्रशशंसुः सुराः सर्वे सिद्धाश्च परमर्षयः} %||5-1-122||

\twolineshloka
{देवताश्चाभवन्हृष्टास्तत्रस्थास्तस्य कर्मणा}
{काञ्चनस्य सुनाभस्य सहस्राक्षश्च वासवः} %||5-1-123||

\twolineshloka
{उवाच वचनं धीमान्परितोषात्सगद्गदम्}
{सुनाभं पर्वतश्रेष्ठं स्वयमेव शचीपतिः} %||5-1-124||

\twolineshloka
{हिरण्यनाभशैलेन्द्रपरितुष्टोऽस्मि ते भृशम्}
{अभयं ते प्रयच्छामि तिष्ठ सौम्य यथासुखम्} %||5-1-125||

\twolineshloka
{साह्यं कृतं ते सुमहद्विक्रान्तस्य हनूमतः}
{क्रमतो योजनशतं निर्भयस्य भये सति} %||5-1-126||

\twolineshloka
{रामस्यैष हि दौत्येन याति दाशरथेर्हरिः}
{सत्क्रियां कुर्वता शक्या तोषितोऽस्मि दृढं त्वया} %||5-1-127||

\twolineshloka
{ततः प्रहर्षमलभद्विपुलं पर्वतोत्तमः}
{देवतानां पतिं दृष्ट्वा परितुष्टं शतक्रतुम्} %||5-1-128||

\twolineshloka
{स वै दत्तवरः शैलो बभूवावस्थितस्तदा}
{हनूमांश्च मुहूर्तेन व्यतिचक्राम सागरम्} %||5-1-129||

\twolineshloka
{ततो देवाः सगन्धर्वाः सिद्धाश्च परमर्षयः}
{अब्रुवन्सूर्यसंकाशां सुरसां नागमातरम्} %||5-1-130||

\twolineshloka
{अयं वातात्मजः श्रीमान्प्लवते सागरोपरि}
{हनूमान्नाम तस्य त्वं मुहूर्तं विघ्नमाचर} %||5-1-131||

\twolineshloka
{राक्षसं रूपमास्थाय सुघोरं पर्वतोपमम्}
{दंष्ट्राकरालं पिङ्गाक्षं वक्त्रं कृत्वा नभःस्पृशम्} %||5-1-132||

\twolineshloka
{बलमिच्छामहे ज्ञातुं भूयश्चास्य पराक्रमम्}
{त्वां विजेष्यत्युपायेन विषदं वा गमिष्यति} %||5-1-133||

\twolineshloka
{एवमुक्ता तु सा देवी दैवतैरभिसत्कृता}
{समुद्रमध्ये सुरसा बिभ्रती राक्षसं वपुः} %||5-1-134||

\twolineshloka
{विकृतं च विरूपं च सर्वस्य च भयावहम्}
{प्लवमानं हनूमन्तमावृत्येदमुवाच ह} %||5-1-135||

\twolineshloka
{मम भक्षः प्रदिष्टस्त्वमीश्वरैर्वानरर्षभ}
{अहं त्वां भक्षयिष्यामि प्रविशेदं ममाननम्} %||5-1-136||

\twolineshloka
{एवमुक्तः सुरसया प्राञ्जलिर्वानरर्षभः}
{प्रहृष्टवदनः श्रीमानिदं वचनमब्रवीत्} %||5-1-137||

\twolineshloka
{रामो दाशरथिर्नाम प्रविष्टो दण्डकावनम्}
{लक्ष्मणेन सह भ्रात्रा वैदेह्या चापि भार्यया} %||5-1-138||

\twolineshloka
{अस्य कार्यविषक्तस्य बद्धवैरस्य राक्षसैः}
{तस्य सीता हृता भार्या रावणेन यशस्विनी} %||5-1-139||

\twolineshloka
{तस्याः सकाशं दूतोऽहं गमिष्ये रामशासनात्}
{कर्तुमर्हसि रामस्य साह्यं विषयवासिनि} %||5-1-140||

\twolineshloka
{अथ वा मैथिलीं दृष्ट्वा रामं चाक्लिष्टकारिणम्}
{आगमिष्यामि ते वक्त्रं सत्यं प्रतिशृणोमि ते} %||5-1-141||

\twolineshloka
{एवमुक्ता हनुमता सुरसा कामरूपिणी}
{अब्रवीन्नातिवर्तेन्मां कश्चिदेष वरो मम} %||5-1-142||

\twolineshloka
{एवमुक्तः सुरसया क्रुद्धो वानरपुंगवः}
{अब्रवीत्कुरु वै वक्त्रं येन मां विषहिष्यसे} %||5-1-143||

\twolineshloka
{इत्युक्त्वा सुरसां क्रुद्धो दशयोजनमायतः}
{दशयोजनविस्तारो बभूव हनुमांस्तदा} %||5-1-144||

\twolineshloka
{तं दृष्ट्वा मेघसंकाशं दशयोजनमायतम्}
{चकार सुरसाप्यास्यं विंशद्योजनमायतम्} %||5-1-145||

\twolineshloka
{हनुमांस्तु ततः क्रुद्धस्त्रिंशद्योजनमायतः}
{चकार सुरसा वक्त्रं चत्वारिंशत्तथोच्छ्रितम्} %||5-1-146||

\twolineshloka
{बभूव हनुमान्वीरः पञ्चाशद्योजनोच्छ्रितः}
{चकार सुरसा वक्त्रं षष्टियोजनमायतम्} %||5-1-147||

\twolineshloka
{तथैव हनुमान्वीरः सप्ततिं योजनोच्छ्रितः}
{चकार सुरसा वक्त्रमशीतिं योजनायतम्} %||5-1-148||

\twolineshloka
{हनूमानचल प्रख्यो नवतिं योजनोच्छ्रितः}
{चकार सुरसा वक्त्रं शतयोजनमायतम्} %||5-1-149||

\twolineshloka
{तद्दृष्ट्वा व्यादितं त्वास्यं वायुपुत्रः स बुद्धिमान्}
{दीर्घजिह्वं सुरसया सुघोरं नरकोपमम्} %||5-1-150||

\twolineshloka
{स संक्षिप्यात्मनः कायं जीमूत इव मारुतिः}
{तस्मिन्मुहूर्ते हनुमान्बभूवाङ्गुष्ठमात्रकः} %||5-1-151||

\twolineshloka
{सोऽभिपत्याशु तद्वक्त्रं निष्पत्य च महाजवः}
{अन्तरिक्षे स्थितः श्रीमानिदं वचनमब्रवीत्} %||5-1-152||

\twolineshloka
{प्रविष्टोऽस्मि हि ते वक्त्रं दाक्षायणि नमोऽस्तु ते}
{गमिष्ये यत्र वैदेही सत्यं चास्तु वचस्तव} %||5-1-153||

\twolineshloka
{तं दृष्ट्वा वदनान्मुक्तं चन्द्रं राहुमुखादिव}
{अब्रवीत्सुरसा देवी स्वेन रूपेण वानरम्} %||5-1-154||

\twolineshloka
{अर्थसिद्ध्यै हरिश्रेष्ठ गच्छ सौम्य यथासुखम्}
{समानय च वैदेहीं राघवेण महात्मना} %||5-1-155||

\twolineshloka
{तत्तृतीयं हनुमतो दृष्ट्वा कर्म सुदुष्करम्}
{साधु साध्विति भूतानि प्रशशंसुस्तदा हरिम्} %||5-1-156||

\twolineshloka
{स सागरमनाधृष्यमभ्येत्य वरुणालयम्}
{जगामाकाशमाविश्य वेगेन गरुणोपमः} %||5-1-157||

\twolineshloka
{सेविते वारिधारिभिः पतगैश्च निषेविते}
{चरिते कैशिकाचार्यैरैरावतनिषेविते} %||5-1-158||

\twolineshloka
{सिंहकुञ्जरशार्दूलपतगोरगवाहनैः}
{विमानैः संपतद्भिश्च विमलैः समलंकृते} %||5-1-159||

\twolineshloka
{वज्राशनिसमाघातैः पावकैरुपशोभिते}
{कृतपुण्यैर्महाभागैः स्वर्गजिद्भिरलंकृते} %||5-1-160||

\twolineshloka
{बहता हव्यमत्यन्तं सेविते चित्रभानुना}
{ग्रहनक्षत्रचन्द्रार्कतारागणविभूषिते} %||5-1-161||

\twolineshloka
{महर्षिगणगन्धर्वनागयक्षसमाकुले}
{विविक्ते विमले विश्वे विश्वावसुनिषेविते} %||5-1-162||

\twolineshloka
{देवराजगजाक्रान्ते चन्द्रसूर्यपथे शिवे}
{विताने जीवलोकस्य विततो ब्रह्मनिर्मिते} %||5-1-163||

\twolineshloka
{बहुशः सेविते वीरैर्विद्याधरगणैर्वरैः}
{कपिना कृष्यमाणानि महाभ्राणि चकाशिरे} %||5-1-164||

\twolineshloka
{प्रविशन्नभ्रजालानि निष्पतंश्च पुनः पुनः}
{प्रावृषीन्दुरिवाभाति निष्पतन्प्रविशंस्तदा} %||5-1-165||

\twolineshloka
{प्लवमानं तु तं दृष्ट्वा सिंहिका नाम राक्षसी}
{मनसा चिन्तयामास प्रवृद्धा कामरूपिणी} %||5-1-166||

\twolineshloka
{अद्य दीर्घस्य कालस्य भविष्याम्यहमाशिता}
{इदं हि मे महत्सत्त्वं चिरस्य वशमागतम्} %||5-1-167||

\twolineshloka
{इति संचिन्त्य मनसा छायामस्य समक्षिपत्}
{छायायां संगृहीतायां चिन्तयामास वानरः} %||5-1-168||

\twolineshloka
{समाक्षिप्तोऽस्मि सहसा पङ्गूकृतपराक्रमः}
{प्रतिलोमेन वातेन महानौरिव सागरे} %||5-1-169||

\twolineshloka
{तिर्यगूर्ध्वमधश्चैव वीक्षमाणस्ततः कपिः}
{ददर्श स महासत्त्वमुत्थितं लवणाम्भसि} %||5-1-170||

\twolineshloka
{कपिराज्ञा यदाख्यातं सत्त्वमद्भुतदर्शनम्}
{छायाग्राहि महावीर्यं तदिदं नात्र संशयः} %||5-1-171||

\twolineshloka
{स तां बुद्ध्वार्थतत्त्वेन सिंहिकां मतिमान्कपिः}
{व्यवर्धत महाकायः प्रावृषीव बलाहकः} %||5-1-172||

\twolineshloka
{तस्य सा कायमुद्वीक्ष्य वर्धमानं महाकपेः}
{वक्त्रं प्रसारयामास पातालाम्बरसंनिभम्} %||5-1-173||

\twolineshloka
{स ददर्श ततस्तस्या विकृतं सुमहन्मुखम्}
{कायमात्रं च मेधावी मर्माणि च महाकपिः} %||5-1-174||

\twolineshloka
{स तस्या विवृते वक्त्रे वज्रसंहननः कपिः}
{संक्षिप्य मुहुरात्मानं निष्पपात महाबलः} %||5-1-175||

\twolineshloka
{आस्ये तस्या निमज्जन्तं ददृशुः सिद्धचारणाः}
{ग्रस्यमानं यथा चन्द्रं पूर्णं पर्वणि राहुणा} %||5-1-176||

\twolineshloka
{ततस्तस्य नखैस्तीक्ष्णैर्मर्माण्युत्कृत्य वानरः}
{उत्पपाताथ वेगेन मनःसंपातविक्रमः} %||5-1-177||

\twolineshloka
{तां हतां वानरेणाशु पतितां वीक्ष्य सिंहिकाम्}
{भूतान्याकाशचारीणि तमूचुः प्लवगर्षभम्} %||5-1-178||

\twolineshloka
{भीममद्य कृतं कर्म महत्सत्त्वं त्वया हतम्}
{साधयार्थमभिप्रेतमरिष्टं प्लवतां वर} %||5-1-179||

\twolineshloka
{यस्य त्वेतानि चत्वारि वानरेन्द्र यथा तव}
{धृतिर्दृष्टिर्मतिर्दाक्ष्यं स कर्मसु न सीदति} %||5-1-180||

\twolineshloka
{स तैः संभावितः पूज्यः प्रतिपन्नप्रयोजनः}
{जगामाकाशमाविश्य पन्नगाशनवत्कपिः} %||5-1-181||

\twolineshloka
{प्राप्तभूयिष्ठ पारस्तु सर्वतः प्रतिलोकयन्}
{योजनानां शतस्यान्ते वनराजिं ददर्श सः} %||5-1-182||

\twolineshloka
{ददर्श च पतन्नेव विविधद्रुमभूषितम्}
{द्वीपं शाखामृगश्रेष्ठो मलयोपवनानि च} %||5-1-183||

\twolineshloka
{सागरं सागरानूपान्सागरानूपजान्द्रुमान्}
{सागरस्य च पत्नीनां मुखान्यपि विलोकयन्} %||5-1-184||

\twolineshloka
{स महामेघसंकाशं समीक्ष्यात्मानमात्मना}
{निरुन्धन्तमिवाकाशं चकार मतिमान्मतिम्} %||5-1-185||

\twolineshloka
{कायवृद्धिं प्रवेगं च मम दृष्ट्वैव राक्षसाः}
{मयि कौतूहलं कुर्युरिति मेने महाकपिः} %||5-1-186||

\twolineshloka
{ततः शरीरं संक्षिप्य तन्महीधरसंनिभम्}
{पुनः प्रकृतिमापेदे वीतमोह इवात्मवान्} %||5-1-187||

\fourlineindentedshloka
{स चारुनानाविधरूपधारी}
{ परं समासाद्य समुद्रतीरम्}
{परैरशक्यप्रतिपन्नरूपः}
{ समीक्षितात्मा समवेक्षितार्थः} %||5-1-188||

\fourlineindentedshloka
{ततः स लम्बस्य गिरेः समृद्धे}
{ विचित्रकूटे निपपात कूटे}
{सकेतकोद्दालकनालिकेरे}
{ महाद्रिकूटप्रतिमो महात्मा} %||5-1-189||

\fourlineindentedshloka
{स सागरं दानवपन्नगायुतं}
{ बलेन विक्रम्य महोर्मिमालिनम्}
{निपत्य तीरे च महोदधेस्तदा}
{ ददर्श लङ्काममरावतीमिव} %||5-2-190||


{॥इत्यार्षे श्रीमद्रामायणे वाल्मीकीये आदिकाव्ये सुन्दरकाण्डे प्रथमः सर्गः॥}

\sect{द्वितीयः सर्गः}

\twolineshloka
{स सागरमनाधृष्यमतिक्रम्य महाबलः}
{त्रिकूटशिखरे लङ्कां स्थितां स्वस्थो ददर्श ह} %||5-2-1||

\twolineshloka
{ततः पादपमुक्तेन पुष्पवर्षेण वीर्यवान्}
{अभिवृष्टः स्थितस्तत्र बभौ पुष्पमयो यथा} %||5-2-2||

\twolineshloka
{योजनानां शतं श्रीमांस्तीर्त्वाप्युत्तमविक्रमः}
{अनिश्वसन्कपिस्तत्र न ग्लानिमधिगच्छति} %||5-2-3||

\twolineshloka
{शतान्यहं योजनानां क्रमेयं सुबहून्यपि}
{किं पुनः सागरस्यान्तं संख्यातं शतयोजनम्} %||5-2-4||

\twolineshloka
{स तु वीर्यवतां श्रेष्ठः प्लवतामपि चोत्तमः}
{जगाम वेगवाँल्लङ्कां लङ्घयित्वा महोदधिम्} %||5-2-5||

\twolineshloka
{शाद्वलानि च नीलानि गन्धवन्ति वनानि च}
{गण्डवन्ति च मध्येन जगाम नगवन्ति च} %||5-2-6||

\twolineshloka
{शैलांश्च तरुसंछन्नान्वनराजीश्च पुष्पिताः}
{अभिचक्राम तेजस्वी हनुमान्प्लवगर्षभः} %||5-2-7||

\twolineshloka
{स तस्मिन्नचले तिष्ठन्वनान्युपवनानि च}
{स नगाग्रे च तां लङ्कां ददर्श पवनात्मजः} %||5-2-8||

\twolineshloka
{सरलान्कर्णिकारांश्च खर्जूरांश्च सुपुष्पितान्}
{प्रियालान्मुचुलिन्दांश्च कुटजान्केतकानपि} %||5-2-9||

\twolineshloka
{प्रियङ्गून्गन्धपूर्णांश्च नीपान्सप्तच्छदांस्तथा}
{असनान्कोविदारांश्च करवीरांश्च पुष्पितान्} %||5-2-10||

\twolineshloka
{पुष्पभारनिबद्धांश्च तथा मुकुलितानपि}
{पादपान्विहगाकीर्णान्पवनाधूतमस्तकान्} %||5-2-11||

\twolineshloka
{हंसकारण्डवाकीर्णा वापीः पद्मोत्पलायुताः}
{आक्रीडान्विविधान्रम्यान्विविधांश्च जलाशयान्} %||5-2-12||

\twolineshloka
{संततान्विविधैर्वृक्षैः सर्वर्तुफलपुष्पितैः}
{उद्यानानि च रम्याणि ददर्श कपिकुञ्जरः} %||5-2-13||

\twolineshloka
{समासाद्य च लक्ष्मीवाँल्लङ्कां रावणपालिताम्}
{परिखाभिः सपद्माभिः सोत्पलाभिरलंकृताम्} %||5-2-14||

\twolineshloka
{सीतापहरणार्थेन रावणेन सुरक्षिताम्}
{समन्ताद्विचरद्भिश्च राक्षसैरुग्रधन्विभिः} %||5-2-15||

\twolineshloka
{काञ्चनेनावृतां रम्यां प्राकारेण महापुरीम्}
{अट्टालकशताकीर्णां पताकाध्वजमालिनीम्} %||5-2-16||

\twolineshloka
{तोरणैः काञ्चनैर्दिव्यैर्लतापङ्क्तिविचित्रितैः}
{ददर्श हनुमाँल्लङ्कां दिवि देवपुरीमिव} %||5-2-17||

\twolineshloka
{गिरिमूर्ध्नि स्थितां लङ्कां पाण्डुरैर्भवनैः शुभैः}
{ददर्श स कपिः श्रीमान्पुरमाकाशगं यथा} %||5-2-18||

\twolineshloka
{पालितां राक्षसेन्द्रेण निर्मितां विश्वकर्मणा}
{प्लवमानामिवाकाशे ददर्श हनुमान्पुरीम्} %||5-2-19||

\twolineshloka
{संपूर्णां राक्षसैर्घोरैर्नागैर्भोगवतीमिव}
{अचिन्त्यां सुकृतां स्पष्टां कुबेराध्युषितां पुरा} %||5-2-20||

\twolineshloka
{दंष्ट्रिभिर्बहुभिः शूरैः शूलपट्टिशपाणिभिः}
{रक्षितां राक्षसैर्घोरैर्गुहामाशीविषैरपि} %||5-2-21||

\twolineshloka
{वप्रप्राकारजघनां विपुलाम्बुनवाम्बराम्}
{शतघ्नीशूलकेशान्तामट्टालकवतंसकाम्} %||5-2-22||

\threelineshloka
{द्वारमुत्तरमासाद्य चिन्तयामास वानरः}
{कैलासशिखरप्रख्यमालिखन्तमिवाम्बरम्}
{ध्रियमाणमिवाकाशमुच्छ्रितैर्भवनोत्तमैः} %||5-2-23||

\twolineshloka
{तस्याश्च महतीं गुप्तिं सागरं च निरीक्ष्य सः}
{रावणं च रिपुं घोरं चिन्तयामास वानरः} %||5-2-24||

\twolineshloka
{आगत्यापीह हरयो भविष्यन्ति निरर्थकाः}
{न हि युद्धेन वै लङ्का शक्या जेतुं सुरैरपि} %||5-2-25||

\twolineshloka
{इमां तु विषमां दुर्गां लङ्कां रावणपालिताम्}
{प्राप्यापि स महाबाहुः किं करिष्यति राघवः} %||5-2-26||

\twolineshloka
{अवकाशो न सान्त्वस्य राक्षसेष्वभिगम्यते}
{न दानस्य न भेदस्य नैव युद्धस्य दृश्यते} %||5-2-27||

\twolineshloka
{चतुर्णामेव हि गतिर्वानराणां महात्मनाम्}
{वालिपुत्रस्य नीलस्य मम राज्ञश्च धीमतः} %||5-2-28||

\twolineshloka
{यावज्जानामि वैदेहीं यदि जीवति वा न वा}
{तत्रैव चिन्तयिष्यामि दृष्ट्वा तां जनकात्मजाम्} %||5-2-29||

\twolineshloka
{ततः स चिन्तयामास मुहूर्तं कपिकुञ्जरः}
{गिरिशृङ्गे स्थितस्तस्मिन्रामस्याभ्युदये रतः} %||5-2-30||

\twolineshloka
{अनेन रूपेण मया न शक्या रक्षसां पुरी}
{प्रवेष्टुं राक्षसैर्गुप्ता क्रूरैर्बलसमन्वितैः} %||5-2-31||

\twolineshloka
{उग्रौजसो महावीर्यो बलवन्तश्च राक्षसाः}
{वञ्चनीया मया सर्वे जानकीं परिमार्गिता} %||5-2-32||

\twolineshloka
{लक्ष्यालक्ष्येण रूपेण रात्रौ लङ्का पुरी मया}
{प्रवेष्टुं प्राप्तकालं मे कृत्यं साधयितुं महत्} %||5-2-33||

\twolineshloka
{तां पुरीं तादृशीं दृष्ट्वा दुराधर्षां सुरासुरैः}
{हनूमांश्चिन्तयामास विनिःश्वस्य मुहुर्मुहुः} %||5-2-34||

\twolineshloka
{केनोपायेन पश्येयं मैथिलीं जनकात्मजाम्}
{अदृष्टो राक्षसेन्द्रेण रावणेन दुरात्मना} %||5-2-35||

\twolineshloka
{न विनश्येत्कथं कार्यं रामस्य विदितात्मनः}
{एकामेकश्च पश्येयं रहिते जनकात्मजाम्} %||5-2-36||

\twolineshloka
{भूताश्चार्थो विपद्यन्ते देशकालविरोधिताः}
{विक्लवं दूतमासाद्य तमः सूर्योदये यथा} %||5-2-37||

\twolineshloka
{अर्थानर्थान्तरे बुद्धिर्निश्चितापि न शोभते}
{घातयन्ति हि कार्याणि दूताः पण्डितमानिनः} %||5-2-38||

\twolineshloka
{न विनश्येत्कथं कार्यं वैक्लव्यं न कथं भवेत्}
{लङ्घनं च समुद्रस्य कथं नु न वृथा भवेत्} %||5-2-39||

\twolineshloka
{मयि दृष्टे तु रक्षोभी रामस्य विदितात्मनः}
{भवेद्व्यर्थमिदं कार्यं रावणानर्थमिच्छतः} %||5-2-40||

\twolineshloka
{न हि शक्यं क्वचित्स्थातुमविज्ञातेन राक्षसैः}
{अपि राक्षसरूपेण किमुतान्येन केनचित्} %||5-2-41||

\twolineshloka
{वायुरप्यत्र नाज्ञातश्चरेदिति मतिर्मम}
{न ह्यस्त्यविदितं किंचिद्राक्षसानां बलीयसाम्} %||5-2-42||

\twolineshloka
{इहाहं यदि तिष्ठामि स्वेन रूपेण संवृतः}
{विनाशमुपयास्यामि भर्तुरर्थश्च हीयते} %||5-2-43||

\twolineshloka
{तदहं स्वेन रूपेण रजन्यां ह्रस्वतां गतः}
{लङ्कामभिपतिष्यामि राघवस्यार्थसिद्धये} %||5-2-44||

\twolineshloka
{रावणस्य पुरीं रात्रौ प्रविश्य सुदुरासदाम्}
{विचिन्वन्भवनं सर्वं द्रक्ष्यामि जनकात्मजाम्} %||5-2-45||

\threelineshloka
{इति संचिन्त्य हनुमान्सूर्यस्यास्तमयं कपिः}
{आचकाङ्क्षे तदा वीरा वैदेह्या दर्शनोत्सुकः}
{पृषदंशकमात्रः सन्बभूवाद्भुतदर्शनः} %||5-2-46||

\twolineshloka
{प्रदोषकाले हनुमांस्तूर्णमुत्पत्य वीर्यवान्}
{प्रविवेश पुरीं रम्यां सुविभक्तमहापथम्} %||5-2-47||

\twolineshloka
{प्रासादमालाविततां स्तम्भैः काञ्चनराजतैः}
{शातकुम्भमयैर्जालैर्गन्धर्वनगरोपमाम्} %||5-2-48||

\twolineshloka
{सप्तभौमाष्टभौमैश्च स ददर्श महापुरीम्}
{तलैः स्फाटिकसंपूर्णैः कार्तस्वरविभूषितैः} %||5-2-49||

\twolineshloka
{वैदूर्यमणिचित्रैश्च मुक्ताजालविभूषितैः}
{तलैः शुशुभिरे तानि भवनान्यत्र रक्षसाम्} %||5-2-50||

\twolineshloka
{काञ्चनानि विचित्राणि तोरणानि च रक्षसाम्}
{लङ्कामुद्द्योतयामासुः सर्वतः समलंकृताम्} %||5-2-51||

\twolineshloka
{अचिन्त्यामद्भुताकारां दृष्ट्वा लङ्कां महाकपिः}
{आसीद्विषण्णो हृष्टश्च वैदेह्या दर्शनोत्सुकः} %||5-2-52||

\fourlineindentedshloka
{स पाण्डुरोद्विद्धविमानमालिनीं}
{ महार्हजाम्बूनदजालतोरणाम्}
{यशस्विनां रावणबाहुपालितां}
{ क्षपाचरैर्भीमबलैः समावृताम्} %||5-2-53||

\fourlineindentedshloka
{चन्द्रोऽपि साचिव्यमिवास्य कुर्वं}
{स्तारागणैर्मध्यगतो विराजन्}
{ज्योत्स्नावितानेन वितत्य लोक}
{मुत्तिष्ठते नैकसहस्ररश्मिः} %||5-2-54||

\fourlineindentedshloka
{शङ्खप्रभं क्षीरमृणालवर्ण}
{मुद्गच्छमानं व्यवभासमानम्}
{ददर्श चन्द्रं स कपिप्रवीरः}
{ पोप्लूयमानं सरसीव हंसं} %||5-3-55||


{॥इत्यार्षे श्रीमद्रामायणे वाल्मीकीये आदिकाव्ये सुन्दरकाण्डे द्वितीयः सर्गः॥}

\sect{तृतीयः सर्गः}

\twolineshloka
{स लम्बशिखरे लम्बे लम्बतोयदसंनिभे}
{सत्त्वमास्थाय मेधावी हनुमान्मारुतात्मजः} %||5-3-1||

\twolineshloka
{निशि लङ्कां महासत्त्वो विवेश कपिकुञ्जरः}
{रम्यकाननतोयाढ्यां पुरीं रावणपालिताम्} %||5-3-2||

\twolineshloka
{शारदाम्बुधरप्रख्यैर्भवनैरुपशोभिताम्}
{सागरोपमनिर्घोषां सागरानिलसेविताम्} %||5-3-3||

\twolineshloka
{सुपुष्टबलसंगुप्तां यथैव विटपावतीम्}
{चारुतोरणनिर्यूहां पाण्डुरद्वारतोरणाम्} %||5-3-4||

\twolineshloka
{भुजगाचरितां गुप्तां शुभां भोगवतीमिव}
{तां सविद्युद्घनाकीर्णां ज्योतिर्मार्गनिषेविताम्} %||5-3-5||

\twolineshloka
{चण्डमारुतनिर्ह्रादां यथेन्द्रस्यामरावतीम्}
{शातकुम्भेन महता प्राकारेणाभिसंवृताम्} %||5-3-6||

\twolineshloka
{किङ्किणीजालघोषाभिः पताकाभिरलंकृताम्}
{आसाद्य सहसा हृष्टः प्राकारमभिपेदिवान्} %||5-3-7||

\twolineshloka
{विस्मयाविष्टहृदयः पुरीमालोक्य सर्वतः}
{जाम्बूनदमयैर्द्वारैर्वैदूर्यकृतवेदिकैः} %||5-3-8||

\twolineshloka
{मणिस्फटिक मुक्ताभिर्मणिकुट्टिमभूषितैः}
{तप्तहाटकनिर्यूहै राजतामलपाण्डुरैः} %||5-3-9||

\twolineshloka
{वैदूर्यतलसोपानैः स्फाटिकान्तरपांसुभिः}
{चारुसंजवनोपेतैः खमिवोत्पतितैः शुभैः} %||5-3-10||

\twolineshloka
{क्रौञ्चबर्हिणसंघुष्टे राजहंसनिषेवितैः}
{तूर्याभरणनिर्घोषैः सर्वतः प्रतिनादिताम्} %||5-3-11||

\twolineshloka
{वस्वोकसाराप्रतिमां समीक्ष्य नगरीं ततः}
{खमिवोत्पतितां लङ्कां जहर्ष हनुमान्कपिः} %||5-3-12||

\twolineshloka
{तां समीक्ष्य पुरीं लङ्कां राक्षसाधिपतेः शुभाम्}
{अनुत्तमामृद्धियुतां चिन्तयामास वीर्यवान्} %||5-3-13||

\twolineshloka
{नेयमन्येन नगरी शक्या धर्षयितुं बलात्}
{रक्षिता रावणबलैरुद्यतायुधधारिभिः} %||5-3-14||

\twolineshloka
{कुमुदाङ्गदयोर्वापि सुषेणस्य महाकपेः}
{प्रसिद्धेयं भवेद्भूमिर्मैन्दद्विविदयोरपि} %||5-3-15||

\twolineshloka
{विवस्वतस्तनूजस्य हरेश्च कुशपर्वणः}
{ऋक्षस्य केतुमालस्य मम चैव गतिर्भवेत्} %||5-3-16||

\twolineshloka
{समीक्ष्य तु महाबाहो राघवस्य पराक्रमम्}
{लक्ष्मणस्य च विक्रान्तमभवत्प्रीतिमान्कपिः} %||5-3-17||

\twolineshloka
{तां रत्नवसनोपेतां कोष्ठागारावतंसकाम्}
{यन्त्रागारस्तनीमृद्धां प्रमदामिव भूषिताम्} %||5-3-18||

\twolineshloka
{तां नष्टतिमिरां दीपैर्भास्वरैश्च महागृहैः}
{नगरीं राक्षसेन्द्रस्य ददर्श स महाकपिः} %||5-3-19||

\twolineshloka
{प्रविष्टः सत्त्वसंपन्नो निशायां मारुतात्मजः}
{स महापथमास्थाय मुक्तापुष्पविराजितम्} %||5-3-20||

\threelineshloka
{हसितोद्घुष्टनिनदैस्तूर्यघोष पुरः सरैः}
{वज्राङ्कुशनिकाशैश्च वज्रजालविभूषितैः}
{गृहमेधैः पुरी रम्या बभासे द्यौरिवाम्बुदैः} %||5-3-21||

\threelineshloka
{प्रजज्वाल तदा लङ्का रक्षोगणगृहैः शुभैः}
{सिताभ्रसदृशैश्चित्रैः पद्मस्वस्तिकसंस्थितैः}
{वर्धमानगृहैश्चापि सर्वतः सुविभाषितैः} %||5-3-22||

\twolineshloka
{तां चित्रमाल्याभरणां कपिराजहितंकरः}
{राघवार्थं चरञ्श्रीमान्ददर्श च ननन्द च} %||5-3-23||

\twolineshloka
{शुश्राव मधुरं गीतं त्रिस्थानस्वरभूषितम्}
{स्त्रीणां मदसमृद्धानां दिवि चाप्सरसामिव} %||5-3-24||

\threelineshloka
{शुश्राव काञ्चीनिनादं नूपुराणां च निःस्वनम्}
{सोपाननिनदांश्चैव भवनेषु महात्मनम्}
{आस्फोटितनिनादांश्च क्ष्वेडितांश्च ततस्ततः} %||5-3-25||

\twolineshloka
{स्वाध्याय निरतांश्चैव यातुधानान्ददर्श सः}
{रावणस्तवसंयुक्तान्गर्जतो राक्षसानपि} %||5-3-26||

\twolineshloka
{राजमार्गं समावृत्य स्थितं रक्षोबलं महत्}
{ददर्श मध्यमे गुल्मे राक्षसस्य चरान्बहून्} %||5-3-27||

\twolineshloka
{दीक्षिताञ्जटिलान्मुण्डान्गोऽजिनाम्बरवाससः}
{दर्भमुष्टिप्रहरणानग्निकुण्डायुधांस्तथा} %||5-3-28||

\twolineshloka
{कूटमुद्गरपाणींश्च दण्डायुधधरानपि}
{एकाक्षानेककर्णांश्च चलल्लम्बपयोधरान्} %||5-3-29||

\threelineshloka
{करालान्भुग्नवक्त्रांश्च विकटान्वामनांस्तथा}
{धन्विनः खड्गिनश्चैव शतघ्नी मुसलायुधान्}
{परिघोत्तमहस्तांश्च विचित्रकवचोज्ज्वलान्} %||5-3-30||

\twolineshloka
{नातिष्ठूलान्नातिकृशान्नातिदीर्घातिह्रस्वकान्}
{विरूपान्बहुरूपांश्च सुरूपांश्च सुवर्चसः} %||5-3-31||

\twolineshloka
{शक्तिवृक्षायुधांश्चैव पट्टिशाशनिधारिणः}
{क्षेपणीपाशहस्तांश्च ददर्श स महाकपिः} %||5-3-32||

\twolineshloka
{स्रग्विणस्त्वनुलिप्तांश्च वराभरणभूषितान्}
{तीक्ष्णशूलधरांश्चैव वज्रिणश्च महाबलान्} %||5-3-33||

\twolineshloka
{शतसाहस्रमव्यग्रमारक्षं मध्यमं कपिः}
{प्राकारावृतमत्यन्तं ददर्श स महाकपिः} %||5-3-34||

\twolineshloka
{त्रिविष्टपनिभं दिव्यं दिव्यनादविनादितम्}
{वाजिहेषितसंघुष्टं नादितं भूषणैस्तथा} %||5-3-35||

\twolineshloka
{रथैर्यानैर्विमानैश्च तथा गजहयैः शुभैः}
{वारणैश्च चतुर्दन्तैः श्वेताभ्रनिचयोपमैः} %||5-3-36||

\twolineshloka
{भूषितं रुचिरद्वारं मत्तैश्च मृगपक्षिभिः}
{राक्षसाधिपतेर्गुप्तमाविवेश गृहं कपिः} %||5-4-37||


{॥इत्यार्षे श्रीमद्रामायणे वाल्मीकीये आदिकाव्ये सुन्दरकाण्डे तृतीयः सर्गः॥}

\sect{चतुर्थः सर्गः}

\fourlineindentedshloka
{ततः स मध्यं गतमंशुमन्तं}
{ ज्योत्स्नावितानं महदुद्वमन्तम्}
{ददर्श धीमान्दिवि भानुमन्तं}
{ गोष्ठे वृषं मत्तमिव भ्रमन्तम्} %||5-4-1||

\fourlineindentedshloka
{लोकस्य पापानि विनाशयन्तं}
{ महोदधिं चापि समेधयन्तम्}
{भूतानि सर्वाणि विराजयन्तं}
{ ददर्श शीतांशुमथाभियान्तम्} %||5-4-2||

\fourlineindentedshloka
{या भाति लक्ष्मीर्भुवि मन्दरस्था}
{ तथा प्रदोषेषु च सागरस्था}
{तथैव तोयेषु च पुष्करस्था}
{ रराज सा चारुनिशाकरस्था} %||5-4-3||

\fourlineindentedshloka
{हंसो यथा राजतपञ्जुरस्थः}
{ सिंहो यथा मन्दरकन्दरस्थः}
{वीरो यथा गर्वितकुञ्जरस्थ}
{श्चन्द्रोऽपि बभ्राज तथाम्बरस्थः} %||5-4-4||

\fourlineindentedshloka
{स्थितः ककुद्मानिव तीक्ष्णशृङ्गो}
{ महाचलः श्वेत इवोच्चशृङ्गः}
{हस्तीव जाम्बूनदबद्धशृङ्गो}
{ विभाति चन्द्रः परिपूर्णशृङ्गः} %||5-4-5||

\fourlineindentedshloka
{प्रकाशचन्द्रोदयनष्टदोषः}
{ प्रवृद्धरक्षः पिशिताशदोषः}
{रामाभिरामेरितचित्तदोषः}
{ स्वर्गप्रकाशो भगवान्प्रदोषः} %||5-4-6||

\fourlineindentedshloka
{तन्त्री स्वनाः कर्णसुखाः प्रवृत्ताः}
{ स्वपन्ति नार्यः पतिभिः सुवृत्ताः}
{नक्तंचराश्चापि तथा प्रवृत्ता}
{ विहर्तुमत्यद्भुतरौद्रवृत्ताः} %||5-4-7||

\fourlineindentedshloka
{मत्तप्रमत्तानि समाकुलानि}
{ रथाश्वभद्रासनसंकुलानि}
{वीरश्रिया चापि समाकुलानि}
{ ददर्श धीमान्स कपिः कुलानि} %||5-4-8||

\fourlineindentedshloka
{परस्परं चाधिकमाक्षिपन्ति}
{ भुजांश्च पीनानधिविक्षिपन्ति}
{मत्तप्रलापानधिविक्षिपन्ति}
{ मत्तानि चान्योन्यमधिक्षिपन्ति} %||5-4-9||

\fourlineindentedshloka
{रक्षांसि वक्षांसि च विक्षिपन्ति}
{ गात्राणि कान्तासु च विक्षिपन्ति}
{ददर्श कान्ताश्च समालपन्ति}
{ तथापरास्तत्र पुनः स्वपन्ति} %||5-4-10||

\fourlineindentedshloka
{महागजैश्चापि तथा नदद्भिः}
{ सूपूजितैश्चापि तथा सुसद्भिः}
{रराज वीरैश्च विनिःश्वसद्भि}
{र्ह्रदो भुजङ्गैरिव निःश्वसद्भिः} %||5-4-11||

\fourlineindentedshloka
{बुद्धिप्रधानान्रुचिराभिधाना}
{न्संश्रद्दधानाञ्जगतः प्रधानान्}
{नानाविधानान्रुचिराभिधाना}
{न्ददर्श तस्यां पुरि यातुधानान्} %||5-4-12||

\fourlineindentedshloka
{ननन्द दृष्ट्वा स च तान्सुरूपा}
{न्नानागुणानात्मगुणानुरूपान्}
{विद्योतमानान्स च तान्सुरूपा}
{न्ददर्श कांश्चिच्च पुनर्विरूपान्} %||5-4-13||

\fourlineindentedshloka
{ततो वरार्हाः सुविशुद्धभावा}
{स्तेषां स्त्रियस्तत्र महानुभावाः}
{प्रियेषु पानेषु च सक्तभावा}
{ ददर्श तारा इव सुप्रभावाः} %||5-4-14||

\fourlineindentedshloka
{श्रिया ज्वलन्तीस्त्रपयोपगूढा}
{ निशीथकाले रमणोपगूढाः}
{ददर्श काश्चित्प्रमदोपगूढा}
{ यथा विहंगाः कुसुमोपगूडाः} %||5-4-15||

\fourlineindentedshloka
{अन्याः पुनर्हर्म्यतलोपविष्टा}
{स्तत्र प्रियाङ्केषु सुखोपविष्टाः}
{भर्तुः प्रिया धर्मपरा निविष्टा}
{ ददर्श धीमान्मनदाभिविष्टाः} %||5-4-16||

\fourlineindentedshloka
{अप्रावृताः काञ्चनराजिवर्णाः}
{ काश्चित्परार्ध्यास्तपनीयवर्णाः}
{पुनश्च काश्चिच्छशलक्ष्मवर्णाः}
{ कान्तप्रहीणा रुचिराङ्गवर्णाः} %||5-4-17||

\fourlineindentedshloka
{ततः प्रियान्प्राप्य मनोऽभिरामा}
{न्सुप्रीतियुक्ताः प्रसमीक्ष्य रामाः}
{गृहेषु हृष्टाः परमाभिरामा}
{ हरिप्रवीरः स ददर्श रामाः} %||5-4-18||

\fourlineindentedshloka
{चन्द्रप्रकाशाश्च हि वक्त्रमाला}
{ वक्राक्षिपक्ष्माश्च सुनेत्रमालाः}
{विभूषणानां च ददर्श मालाः}
{ शतह्रदानामिव चारुमालाः} %||5-4-19||

\fourlineindentedshloka
{न त्वेव सीतां परमाभिजातां}
{ पथि स्थिते राजकुले प्रजाताम्}
{लतां प्रफुल्लामिव साधुजातां}
{ ददर्श तन्वीं मनसाभिजाताम्} %||5-4-20||

\fourlineindentedshloka
{सनातने वर्त्मनि संनिविष्टां}
{ रामेक्षणीं तां मदनाभिविष्टाम्}
{भर्तुर्मनः श्रीमदनुप्रविष्टां}
{ स्त्रीभ्यो वराभ्यश्च सदा विशिष्टाम्} %||5-4-21||

\fourlineindentedshloka
{उष्णार्दितां सानुसृतास्रकण्ठीं}
{ पुरा वरार्होत्तमनिष्ककण्ठीम्}
{सुजातपक्ष्मामभिरक्तकण्ठीं}
{ वने प्रवृत्तामिव नीलकण्ठीम्} %||5-4-22||

\fourlineindentedshloka
{अव्यक्तलेखामिव चन्द्रलेखां}
{ पांसुप्रदिग्धामिव हेमलेखाम्}
{क्षतप्ररूढामिव बाणलेखां}
{ वायुप्रभिन्नामिव मेघलेखाम्} %||5-4-23||

\fourlineindentedshloka
{सीतामपश्यन्मनुजेश्वरस्य}
{ रामस्य पत्नीं वदतां वरस्य}
{बभूव दुःखाभिहतश्चिरस्य}
{ प्लवंगमो मन्द इवाचिरस्य} %||5-5-24||


{॥इत्यार्षे श्रीमद्रामायणे वाल्मीकीये आदिकाव्ये सुन्दरकाण्डे चतुर्थः सर्गः॥}

\sect{पञ्चमः सर्गः}

\twolineshloka
{स निकामं विनामेषु विचरन्कामरूपधृक्}
{विचचार कपिर्लङ्कां लाघवेन समन्वितः} %||5-5-1||

\twolineshloka
{आससादाथ लक्ष्मीवान्राक्षसेन्द्रनिवेशनम्}
{प्राकारेणार्कवर्णेन भास्वरेणाभिसंवृतम्} %||5-5-2||

\twolineshloka
{रक्षितं राक्षसैर्भीमैः सिंहैरिव महद्वनम्}
{समीक्षमाणो भवनं चकाशे कपिकुञ्जरः} %||5-5-3||

\twolineshloka
{रूप्यकोपहितैश्चित्रैस्तोरणैर्हेमभूषितैः}
{विचित्राभिश्च कक्ष्याभिर्द्वारैश्च रुचिरैर्वृतम्} %||5-5-4||

\twolineshloka
{गजास्थितैर्महामात्रैः शूरैश्च विगतश्रमैः}
{उपस्थितमसंहार्यैर्हयैः स्यन्दनयायिभिः} %||5-5-5||

\twolineshloka
{सिंहव्याघ्रतनुत्राणैर्दान्तकाञ्चनराजतैः}
{घोषवद्भिर्विचित्रैश्च सदा विचरितं रथैः} %||5-5-6||

\twolineshloka
{बहुरत्नसमाकीर्णं परार्ध्यासनभाजनम्}
{महारथसमावासं महारथमहासनम्} %||5-5-7||

\twolineshloka
{दृश्यैश्च परमोदारैस्तैस्तैश्च मृगपक्षिभिः}
{विविधैर्बहुसाहस्रैः परिपूर्णं समन्ततः} %||5-5-8||

\twolineshloka
{विनीतैरन्तपालैश्च रक्षोभिश्च सुरक्षितम्}
{मुख्याभिश्च वरस्त्रीभिः परिपूर्णं समन्ततः} %||5-5-9||

\twolineshloka
{मुदितप्रमदा रत्नं राक्षसेन्द्रनिवेशनम्}
{वराभरणनिर्ह्रादैः समुद्रस्वननिःस्वनम्} %||5-5-10||

\twolineshloka
{तद्राजगुणसंपन्नं मुख्यैश्च वरचन्दनैः}
{भेरीमृदङ्गाभिरुतं शङ्खघोषविनादितम्} %||5-5-11||

\twolineshloka
{नित्यार्चितं पर्वहुतं पूजितं राक्षसैः सदा}
{समुद्रमिव गम्भीरं समुद्रमिव निःस्वनम्} %||5-5-12||

\twolineshloka
{महात्मानो महद्वेश्म महारत्नपरिच्छदम्}
{महाजनसमाकीर्णं ददर्श स महाकपिः} %||5-5-13||

\twolineshloka
{विराजमानं वपुषा गजाश्वरथसंकुलम्}
{लङ्काभरणमित्येव सोऽमन्यत महाकपिः} %||5-5-14||

\twolineshloka
{गृहाद्गृहं राक्षसानामुद्यानानि च वानरः}
{वीक्षमाणो ह्यसंत्रस्तः प्रासादांश्च चचार सः} %||5-5-15||

\twolineshloka
{अवप्लुत्य महावेगः प्रहस्तस्य निवेशनम्}
{ततोऽन्यत्पुप्लुवे वेश्म महापार्श्वस्य वीर्यवान्} %||5-5-16||

\twolineshloka
{अथ मेघप्रतीकाशं कुम्भकर्णनिवेशनम्}
{विभीषणस्य च तथा पुप्लुवे स महाकपिः} %||5-5-17||

\threelineshloka
{महोदरस्य च तथा विरूपाक्षस्य चैव हि}
{विद्युज्जिह्वस्य भवनं विद्युन्मालेस्तथैव च}
{वज्रदंष्ट्रस्य च तथा पुप्लुवे स महाकपिः} %||5-5-18||

\twolineshloka
{शुकस्य च महावेगः सारणस्य च धीमतः}
{तथा चेन्द्रजितो वेश्म जगाम हरियूथपः} %||5-5-19||

\twolineshloka
{जम्बुमालेः सुमालेश्च जगाम हरियूथपः}
{रश्मिकेतोश्च भवनं सूर्यशत्रोस्तथैव च} %||5-5-20||

\twolineshloka
{धूम्राक्षस्य च संपातेर्भवनं मारुतात्मजः}
{विद्युद्रूपस्य भीमस्य घनस्य विघनस्य च} %||5-5-21||

\twolineshloka
{शुकनाभस्य वक्रस्य शठस्य विकटस्य च}
{ह्रस्वकर्णस्य दंष्ट्रस्य रोमशस्य च रक्षसः} %||5-5-22||

\twolineshloka
{युद्धोन्मत्तस्य मत्तस्य ध्वजग्रीवस्य नादिनः}
{विद्युज्जिह्वेन्द्रजिह्वानां तथा हस्तिमुखस्य च} %||5-5-23||

\twolineshloka
{करालस्य पिशाचस्य शोणिताक्षस्य चैव हि}
{क्रममाणः क्रमेणैव हनूमान्मारुतात्मजः} %||5-5-24||

\twolineshloka
{तेषु तेषु महार्हेषु भवनेषु महायशाः}
{तेषामृद्धिमतामृद्धिं ददर्श स महाकपिः} %||5-5-25||

\twolineshloka
{सर्वेषां समतिक्रम्य भवनानि समन्ततः}
{आससादाथ लक्ष्मीवान्राक्षसेन्द्रनिवेशनम्} %||5-5-26||

\threelineshloka
{रावणस्योपशायिन्यो ददर्श हरिसत्तमः}
{विचरन्हरिशार्दूलो राक्षसीर्विकृतेक्षणाः}
{शूलमुद्गरहस्ताश्च शक्तो तोमरधारिणीः} %||5-5-27||

{ददर्श विविधान्गुल्मांस्तस्य रक्षःपतेर्गृहे॥\devanumber{28}॥} %||5-5-28|| (Check)

\addtocounter{shlokacount}{1}

\twolineshloka
{रक्ताञ्श्वेतान्सितांश्चैव हरींश्चैव महाजवान्}
{कुलीनान्रूपसंपन्नान्गजान्परगजारुजान्} %||5-5-29||

\twolineshloka
{निष्ठितान्गजशिखायामैरावतसमान्युधि}
{निहन्तॄन्परसैन्यानां गृहे तस्मिन्ददर्श सः} %||5-5-30||

\twolineshloka
{क्षरतश्च यथा मेघान्स्रवतश्च यथा गिरीन्}
{मेघस्तनितनिर्घोषान्दुर्धर्षान्समरे परैः} %||5-5-31||

\twolineshloka
{सहस्रं वाहिनीस्तत्र जाम्बूनदपरिष्कृताः}
{हेमजालैरविच्छिन्नास्तरुणादित्यसंनिभाः} %||5-5-32||

\twolineshloka
{ददर्श राक्षसेन्द्रस्य रावणस्य निवेशने}
{शिबिका विविधाकाराः स कपिर्मारुतात्मजः} %||5-5-33||

\twolineshloka
{लतागृहाणि चित्राणि चित्रशालागृहाणि च}
{क्रीडागृहाणि चान्यानि दारुपर्वतकानपि} %||5-5-34||

\twolineshloka
{कामस्य गृहकं रम्यं दिवागृहकमेव च}
{ददर्श राक्षसेन्द्रस्य रावणस्य निवेशने} %||5-5-35||

\twolineshloka
{स मन्दरतलप्रख्यं मयूरस्थानसंकुलम्}
{ध्वजयष्टिभिराकीर्णं ददर्श भवनोत्तमम्} %||5-5-36||

\twolineshloka
{अनन्तरत्ननिचयं निधिजालं समन्ततः}
{धीरनिष्ठितकर्मान्तं गृहं भूतपतेरिव} %||5-5-37||

\twolineshloka
{अर्चिर्भिश्चापि रत्नानां तेजसा रावणस्य च}
{विरराजाथ तद्वेश्म रश्मिमानिव रश्मिभिः} %||5-5-38||

\twolineshloka
{जाम्बूनदमयान्येव शयनान्यासनानि च}
{भाजनानि च शुभ्राणि ददर्श हरियूथपः} %||5-5-39||

\twolineshloka
{मध्वासवकृतक्लेदं मणिभाजनसंकुलम्}
{मनोरममसंबाधं कुबेरभवनं यथा} %||5-5-40||

\twolineshloka
{नूपुराणां च घोषेण काञ्चीनां निनदेन च}
{मृदङ्गतलघोषैश्च घोषवद्भिर्विनादितम्} %||5-5-41||

\twolineshloka
{प्रासादसंघातयुतं स्त्रीरत्नशतसंकुलम्}
{सुव्यूढकक्ष्यं हनुमान्प्रविवेश महागृहम्} %||5-6-42||


{॥इत्यार्षे श्रीमद्रामायणे वाल्मीकीये आदिकाव्ये सुन्दरकाण्डे पञ्चमः सर्गः॥}

\sect{षष्ठः सर्गः}

\fourlineindentedshloka
{स वेश्मजालं बलवान्ददर्श}
{ व्यासक्तवैदूर्यसुवर्णजालम्}
{यथा महत्प्रावृषि मेघजालं}
{ विद्युत्पिनद्धं सविहंगजालम्} %||5-6-1||

\fourlineindentedshloka
{निवेशनानां विविधाश्च शालाः}
{ प्रधानशङ्खायुधचापशालाः}
{मनोहराश्चापि पुनर्विशाला}
{ ददर्श वेश्माद्रिषु चन्द्रशालाः} %||5-6-2||

\fourlineindentedshloka
{गृहाणि नानावसुराजितानि}
{ देवासुरैश्चापि सुपूजितानि}
{सर्वैश्च दोषैः परिवर्जितानि}
{ कपिर्ददर्श स्वबलार्जितानि} %||5-6-3||

\fourlineindentedshloka
{तानि प्रयत्नाभिसमाहितानि}
{ मयेन साक्षादिव निर्मितानि}
{महीतले सर्वगुणोत्तराणि}
{ ददर्श लङ्काधिपतेर्गृहाणि} %||5-6-4||

\fourlineindentedshloka
{ततो ददर्शोच्छ्रितमेघरूपं}
{ मनोहरं काञ्चनचारुरूपम्}
{रक्षोऽधिपस्यात्मबलानुरूपं}
{ गृहोत्तमं ह्यप्रतिरूपरूपम्} %||5-6-5||

\fourlineindentedshloka
{महीतले स्वर्गमिव प्रकीर्णं}
{ श्रिया ज्वलन्तं बहुरत्नकीर्णम्}
{नानातरूणां कुसुमावकीर्णं}
{ गिरेरिवाग्रं रजसावकीर्णम्} %||5-6-6||

\fourlineindentedshloka
{नारीप्रवेकैरिव दीप्यमानं}
{ तडिद्भिरम्भोदवदर्च्यमानम्}
{हंसप्रवेकैरिव वाह्यमानं}
{ श्रिया युतं खे सुकृतां विमानम्} %||5-6-7||

\fourlineindentedshloka
{यथा नगाग्रं बहुधातुचित्रं}
{ यथा नभश्च ग्रहचन्द्रचित्रम्}
{ददर्श युक्तीकृतमेघचित्रं}
{ विमानरत्नं बहुरत्नचित्रम्} %||5-6-8||

\fourlineindentedshloka
{मही कृता पर्वतराजिपूर्णा}
{ शैलाः कृता वृक्षवितानपूर्णाः}
{वृक्षाः कृताः पुष्पवितानपूर्णाः}
{ पुष्पं कृतं केसरपत्रपूर्णम्} %||5-6-9||

\fourlineindentedshloka
{कृतानि वेश्मानि च पाण्डुराणि}
{ तथा सुपुष्पा अपि पुष्करिण्यः}
{पुनश्च पद्मानि सकेसराणि}
{ धन्यानि चित्राणि तथा वनानि} %||5-6-10||

\fourlineindentedshloka
{पुष्पाह्वयं नाम विराजमानं}
{ रत्नप्रभाभिश्च विवर्धमानम्}
{वेश्मोत्तमानामपि चोच्चमानं}
{ महाकपिस्तत्र महाविमानम्} %||5-6-11||

\fourlineindentedshloka
{कृताश्च वैदूर्यमया विहंगा}
{ रूप्यप्रवालैश्च तथा विहंगाः}
{चित्राश्च नानावसुभिर्भुजंगा}
{ जात्यानुरूपास्तुरगाः शुभाङ्गाः} %||5-6-12||

\fourlineindentedshloka
{प्रवालजाम्बूनदपुष्पपक्षाः}
{ सलीलमावर्जितजिह्मपक्षाः}
{कामस्य साक्षादिव भान्ति पक्षाः}
{ कृता विहंगाः सुमुखाः सुपक्षाः} %||5-6-13||

\fourlineindentedshloka
{नियुज्यमानाश्च गजाः सुहस्ताः}
{ सकेसराश्चोत्पलपत्रहस्ताः}
{बभूव देवी च कृता सुहस्ता}
{ लक्ष्मीस्तथा पद्मिनि पद्महस्ता} %||5-6-14||

\fourlineindentedshloka
{इतीव तद्गृहमभिगम्य शोभनं}
{ सविस्मयो नगमिव चारुशोभनम्}
{पुनश्च तत्परमसुगन्धि सुन्दरं}
{ हिमात्यये नगमिव चारुकन्दरम्} %||5-6-15||

\fourlineindentedshloka
{ततः स तां कपिरभिपत्य पूजितां}
{ चरन्पुरीं दशमुखबाहुपालिताम्}
{अदृश्य तां जनकसुतां सुपूजितां}
{ सुदुःखितां पतिगुणवेगनिर्जिताम्} %||5-6-16||

\fourlineindentedshloka
{ततस्तदा बहुविधभावितात्मनः}
{ कृतात्मनो जनकसुतां सुवर्त्मनः}
{अपश्यतोऽभवदतिदुःखितं मनः}
{ सुचक्षुषः प्रविचरतो महात्मनः} %||5-7-17||


{॥इत्यार्षे श्रीमद्रामायणे वाल्मीकीये आदिकाव्ये सुन्दरकाण्डे षष्ठः सर्गः॥}

\sect{सप्तमः सर्गः}

\twolineshloka
{तस्यालयवरिष्ठस्य मध्ये विपुलमायतम्}
{ददर्श भवनश्रेष्ठं हनूमान्मारुतात्मजः} %||5-7-1||

\twolineshloka
{अर्धयोजनविस्तीर्णमायतं योजनं हि तत्}
{भवनं राक्षसेन्द्रस्य बहुप्रासादसंकुलम्} %||5-7-2||

\twolineshloka
{मार्गमाणस्तु वैदेहीं सीतामायतलोचनाम्}
{सर्वतः परिचक्राम हनूमानरिसूदनः} %||5-7-3||

\twolineshloka
{चतुर्विषाणैर्द्विरदैस्त्रिविषाणैस्तथैव च}
{परिक्षिप्तमसंबाधं रक्ष्यमाणमुदायुधैः} %||5-7-4||

\twolineshloka
{राक्षसीभिश्च पत्नीभी रावणस्य निवेशनम्}
{आहृताभिश्च विक्रम्य राजकन्याभिरावृतम्} %||5-7-5||

\twolineshloka
{तन्नक्रमकराकीर्णं तिमिंगिलझषाकुलम्}
{वायुवेगसमाधूतं पन्नगैरिव सागरम्} %||5-7-6||

\twolineshloka
{या हि वैश्वरणे लक्ष्मीर्या चेन्द्रे हरिवाहने}
{सा रावणगृहे सर्वा नित्यमेवानपायिनी} %||5-7-7||

\twolineshloka
{या च राज्ञः कुबेरस्य यमस्य वरुणस्य च}
{तादृशी तद्विशिष्टा वा ऋद्धी रक्षो गृहेष्विह} %||5-7-8||

\twolineshloka
{तस्य हर्म्यस्य मध्यस्थं वेश्म चान्यत्सुनिर्मितम्}
{बहुनिर्यूह संकीर्णं ददर्श पवनात्मजः} %||5-7-9||

\twolineshloka
{ब्रह्मणोऽर्थे कृतं दिव्यं दिवि यद्विश्वकर्मणा}
{विमानं पुष्पकं नाम सर्वरत्नविभूषितम्} %||5-7-10||

\twolineshloka
{परेण तपसा लेभे यत्कुबेरः पितामहात्}
{कुबेरमोजसा जित्वा लेभे तद्राक्षसेश्वरः} %||5-7-11||

\twolineshloka
{ईहा मृगसमायुक्तैः कार्यस्वरहिरण्मयैः}
{सुकृतैराचितं स्तम्भैः प्रदीप्तमिव च श्रिया} %||5-7-12||

\twolineshloka
{मेरुमन्दरसंकाशैरुल्लिखद्भिरिवाम्बरम्}
{कूटागारैः शुभाकारैः सर्वतः समलंकृतम्} %||5-7-13||

\twolineshloka
{ज्वलनार्कप्रतीकाशं सुकृतं विश्वकर्मणा}
{हेमसोपानसंयुक्तं चारुप्रवरवेदिकम्} %||5-7-14||

\threelineshloka
{जालवातायनैर्युक्तं काञ्चनैः स्थाटिकैरपि}
{इन्द्रनीलमहानीलमणिप्रवरवेदिकम्}
{विमानं पुष्पकं दिव्यमारुरोह महाकपिः} %||5-7-15||

\twolineshloka
{तत्रस्थः स तदा गन्धं पानभक्ष्यान्नसंभवम्}
{दिव्यं संमूर्छितं जिघ्रन्रूपवन्तमिवानिलम्} %||5-7-16||

\twolineshloka
{स गन्धस्तं महासत्त्वं बन्धुर्बन्धुमिवोत्तमम्}
{इत एहीत्युवाचेव तत्र यत्र स रावणः} %||5-7-17||

\twolineshloka
{ततस्तां प्रस्थितः शालां ददर्श महतीं शुभाम्}
{रावणस्य मनःकान्तां कान्तामिव वरस्त्रियम्} %||5-7-18||

\twolineshloka
{मणिसोपानविकृतां हेमजालविराजिताम्}
{स्फाटिकैरावृततलां दन्तान्तरितरूपिकाम्} %||5-7-19||

\twolineshloka
{मुक्ताभिश्च प्रवालैश्च रूप्यचामीकरैरपि}
{विभूषितां मणिस्तम्भैः सुबहुस्तम्भभूषिताम्} %||5-7-20||

\twolineshloka
{समैरृजुभिरत्युच्चैः समन्तात्सुविभूषितैः}
{स्तम्भैः पक्षैरिवात्युच्चैर्दिवं संप्रस्थितामिव} %||5-7-21||

\twolineshloka
{महत्या कुथयास्त्रीणं पृथिवीलक्षणाङ्कया}
{पृथिवीमिव विस्तीर्णां सराष्ट्रगृहमालिनीम्} %||5-7-22||

\twolineshloka
{नादितां मत्तविहगैर्दिव्यगन्धाधिवासिताम्}
{परार्ध्यास्तरणोपेतां रक्षोऽधिपनिषेविताम्} %||5-7-23||

\twolineshloka
{धूम्रामगरुधूपेन विमलां हंसपाण्डुराम्}
{चित्रां पुष्पोपहारेण कल्माषीमिव सुप्रभाम्} %||5-7-24||

\twolineshloka
{मनःसंह्लादजननीं वर्णस्यापि प्रसादिनीम्}
{तां शोकनाशिनीं दिव्यां श्रियः संजननीमिव} %||5-7-25||

\twolineshloka
{इन्द्रियाणीन्द्रियार्थैस्तु पञ्च पञ्चभिरुत्तमैः}
{तर्पयामास मातेव तदा रावणपालिता} %||5-7-26||

\twolineshloka
{स्वर्गोऽयं देवलोकोऽयमिन्द्रस्येयं पुरी भवेत्}
{सिद्धिर्वेयं परा हि स्यादित्यमन्यत मारुतिः} %||5-7-27||

\twolineshloka
{प्रध्यायत इवापश्यत्प्रदीपांस्तत्र काञ्चनान्}
{धूर्तानिव महाधूर्तैर्देवनेन पराजितान्} %||5-7-28||

\twolineshloka
{दीपानां च प्रकाशेन तेजसा रावणस्य च}
{अर्चिर्भिर्भूषणानां च प्रदीप्तेत्यभ्यमन्यत} %||5-7-29||

\twolineshloka
{ततोऽपश्यत्कुथासीनं नानावर्णाम्बरस्रजम्}
{सहस्रं वरनारीणां नानावेषविभूषितम्} %||5-7-30||

\twolineshloka
{परिवृत्तेऽर्धरात्रे तु पाननिद्रावशं गतम्}
{क्रीडित्वोपरतं रात्रौ सुष्वाप बलवत्तदा} %||5-7-31||

\twolineshloka
{तत्प्रसुप्तं विरुरुचे निःशब्दान्तरभूषणम्}
{निःशब्दहंसभ्रमरं यथा पद्मवनं महत्} %||5-7-32||

\twolineshloka
{तासां संवृतदन्तानि मीलिताक्षाणि मारुतिः}
{अपश्यत्पद्मगन्धीनि वदनानि सुयोषिताम्} %||5-7-33||

\twolineshloka
{प्रबुद्धानीव पद्मानि तासां भूत्वा क्षपाक्षये}
{पुनःसंवृतपत्राणि रात्राविव बभुस्तदा} %||5-7-34||

\twolineshloka
{इमानि मुखपद्मानि नियतं मत्तषट्पदाः}
{अम्बुजानीव फुल्लानि प्रार्थयन्ति पुनः पुनः} %||5-7-35||

\twolineshloka
{इति वामन्यत श्रीमानुपपत्त्या महाकपिः}
{मेने हि गुणतस्तानि समानि सलिलोद्भवैः} %||5-7-36||

\twolineshloka
{सा तस्य शुशुभे शाला ताभिः स्त्रीभिर्विराजिता}
{शारदीव प्रसन्ना द्यौस्ताराभिरभिशोभिता} %||5-7-37||

\twolineshloka
{स च ताभिः परिवृतः शुशुभे राक्षसाधिपः}
{यथा ह्युडुपतिः श्रीमांस्ताराभिरभिसंवृतः} %||5-7-38||

\twolineshloka
{याश्च्यवन्तेऽम्बरात्ताराः पुण्यशेषसमावृताः}
{इमास्ताः संगताः कृत्स्ना इति मेने हरिस्तदा} %||5-7-39||

\twolineshloka
{ताराणामिव सुव्यक्तं महतीनां शुभार्चिषाम्}
{प्रभावर्णप्रसादाश्च विरेजुस्तत्र योषिताम्} %||5-7-40||

\twolineshloka
{व्यावृत्तगुरुपीनस्रक्प्रकीर्णवरभूषणाः}
{पानव्यायामकालेषु निद्रापहृतचेतसः} %||5-7-41||

\twolineshloka
{व्यावृत्ततिलकाः काश्चित्काश्चिदुद्भ्रान्तनूपुराः}
{पार्श्वे गलितहाराश्च काश्चित्परमयोषितः} %||5-7-42||

\twolineshloka
{मुखा हारवृताश्चान्याः काश्चित्प्रस्रस्तवाससः}
{व्याविद्धरशना दामाः किशोर्य इव वाहिताः} %||5-7-43||

\twolineshloka
{सुकुण्डलधराश्चान्या विच्छिन्नमृदितस्रजः}
{गजेन्द्रमृदिताः फुल्ला लता इव महावने} %||5-7-44||

\twolineshloka
{चन्द्रांशुकिरणाभाश्च हाराः कासांचिदुत्कटाः}
{हंसा इव बभुः सुप्ताः स्तनमध्येषु योषिताम्} %||5-7-45||

\twolineshloka
{अपरासां च वैदूर्याः कादम्बा इव पक्षिणः}
{हेमसूत्राणि चान्यासां चक्रवाका इवाभवन्} %||5-7-46||

\twolineshloka
{हंसकारण्डवाकीर्णाश्चक्रवाकोपशोभिताः}
{आपगा इव ता रेजुर्जघनैः पुलिनैरिव} %||5-7-47||

\twolineshloka
{किङ्किणीजालसंकाशास्ता हेमविपुलाम्बुजाः}
{भावग्राहा यशस्तीराः सुप्ता नद्य इवाबभुः} %||5-7-48||

\twolineshloka
{मृदुष्वङ्गेषु कासांचित्कुचाग्रेषु च संस्थिताः}
{बभूवुर्भूषणानीव शुभा भूषणराजयः} %||5-7-49||

\twolineshloka
{अंशुकान्ताश्च कासांचिन्मुखमारुतकम्पिताः}
{उपर्युपरि वक्त्राणां व्याधूयन्ते पुनः पुनः} %||5-7-50||

\twolineshloka
{ताः पाताका इवोद्धूताः पत्नीनां रुचिरप्रभाः}
{नानावर्णसुवर्णानां वक्त्रमूलेषु रेजिरे} %||5-7-51||

\twolineshloka
{ववल्गुश्चात्र कासांचित्कुण्डलानि शुभार्चिषाम्}
{मुखमारुतसंसर्गान्मन्दं मन्दं सुयोषिताम्} %||5-7-52||

\twolineshloka
{शर्करासवगन्धः स प्रकृत्या सुरभिः सुखः}
{तासां वदननिःश्वासः सिषेवे रावणं तदा} %||5-7-53||

\twolineshloka
{रावणाननशङ्काश्च काश्चिद्रावणयोषितः}
{मुखानि स्म सपत्नीनामुपाजिघ्रन्पुनः पुनः} %||5-7-54||

\twolineshloka
{अत्यर्थं सक्तमनसो रावणे ता वरस्त्रियः}
{अस्वतन्त्राः सपत्नीनां प्रियमेवाचरंस्तदा} %||5-7-55||

\twolineshloka
{बाहूनुपनिधायान्याः पारिहार्य विभूषिताः}
{अंशुकानि च रम्याणि प्रमदास्तत्र शिश्यिरे} %||5-7-56||

\twolineshloka
{अन्या वक्षसि चान्यस्यास्तस्याः काचित्पुनर्भुजम्}
{अपरा त्वङ्कमन्यस्यास्तस्याश्चाप्यपरा भुजौ} %||5-7-57||

\twolineshloka
{ऊरुपार्श्वकटीपृष्ठमन्योन्यस्य समाश्रिताः}
{परस्परनिविष्टाङ्ग्यो मदस्नेहवशानुगाः} %||5-7-58||

\twolineshloka
{अन्योन्यस्याङ्गसंस्पर्शात्प्रीयमाणाः सुमध्यमाः}
{एकीकृतभुजाः सर्वाः सुषुपुस्तत्र योषितः} %||5-7-59||

\twolineshloka
{अन्योन्यभुजसूत्रेण स्त्रीमालाग्रथिता हि सा}
{मालेव ग्रथिता सूत्रे शुशुभे मत्तषट्पदा} %||5-7-60||

\twolineshloka
{लतानां माधवे मासि फुल्लानां वायुसेवनात्}
{अन्योन्यमालाग्रथितं संसक्तकुसुमोच्चयम्} %||5-7-61||

\twolineshloka
{व्यतिवेष्टितसुस्कन्थमन्योन्यभ्रमराकुलम्}
{आसीद्वनमिवोद्धूतं स्त्रीवनं रावणस्य तत्} %||5-7-62||

\twolineshloka
{उचितेष्वपि सुव्यक्तं न तासां योषितां तदा}
{विवेकः शक्य आधातुं भूषणाङ्गाम्बरस्रजाम्} %||5-7-63||

\twolineshloka
{रावणे सुखसंविष्टे ताः स्त्रियो विविधप्रभाः}
{ज्वलन्तः काञ्चना दीपाः प्रेक्षन्तानिमिषा इव} %||5-7-64||

\twolineshloka
{राजर्षिपितृदैत्यानां गन्धर्वाणां च योषितः}
{रक्षसां चाभवन्कन्यास्तस्य कामवशं गताः} %||5-7-65||

\fourlineindentedshloka
{न तत्र काचित्प्रमदा प्रसह्य}
{ वीर्योपपन्नेन गुणेन लब्धा}
{न चान्यकामापि न चान्यपूर्वा}
{ विना वरार्हां जनकात्मजां तु} %||5-7-66||

\fourlineindentedshloka
{न चाकुलीना न च हीनरूपा}
{ नादक्षिणा नानुपचार युक्ता}
{भार्याभवत्तस्य न हीनसत्त्वा}
{ न चापि कान्तस्य न कामनीया} %||5-7-67||

\fourlineindentedshloka
{बभूव बुद्धिस्तु हरीश्वरस्य}
{ यदीदृशी राघवधर्मपत्नी}
{इमा यथा राक्षसराजभार्याः}
{ सुजातमस्येति हि साधुबुद्धेः} %||5-7-68||

\fourlineindentedshloka
{पुनश्च सोऽचिन्तयदार्तरूपो}
{ ध्रुवं विशिष्टा गुणतो हि सीता}
{अथायमस्यां कृतवान्महात्मा}
{ लङ्केश्वरः कष्टमनार्यकर्म} %||5-8-69||


{॥इत्यार्षे श्रीमद्रामायणे वाल्मीकीये आदिकाव्ये सुन्दरकाण्डे सप्तमः सर्गः॥}

\sect{अष्टमः सर्गः}

\twolineshloka
{तत्र दिव्योपमं मुख्यं स्फाटिकं रत्नभूषितम्}
{अवेक्षमाणो हनुमान्ददर्श शयनासनम्} %||5-8-1||

\twolineshloka
{तस्य चैकतमे देशे सोऽग्र्यमाल्यविभूषितम्}
{ददर्श पाण्डुरं छत्रं ताराधिपतिसंनिभम्} %||5-8-2||

\twolineshloka
{बालव्यजनहस्ताभिर्वीज्यमानं समन्ततः}
{गन्धैश्च विविधैर्जुष्टं वरधूपेन धूपितम्} %||5-8-3||

\twolineshloka
{परमास्तरणास्तीर्णमाविकाजिनसंवृतम्}
{दामभिर्वरमाल्यानां समन्तादुपशोभितम्} %||5-8-4||

\twolineshloka
{तस्मिञ्जीमूतसंकाशं प्रदीप्तोत्तमकुण्डलम्}
{लोहिताक्षं महाबाहुं महारजतवाससं} %||5-8-5||

\twolineshloka
{लोहितेनानुलिप्ताङ्गं चन्दनेन सुगन्धिना}
{संध्यारक्तमिवाकाशे तोयदं सतडिद्गुणम्} %||5-8-6||

\twolineshloka
{वृतमाभरणैर्दिव्यैः सुरूपं कामरूपिणम्}
{सवृक्षवनगुल्माढ्यं प्रसुप्तमिव मन्दरम्} %||5-8-7||

\twolineshloka
{क्रीडित्वोपरतं रात्रौ वराभरणभूषितम्}
{प्रियं राक्षसकन्यानां राक्षसानां सुखावहम्} %||5-8-8||

\twolineshloka
{पीत्वाप्युपरतं चापि ददर्श स महाकपिः}
{भास्करे शयने वीरं प्रसुप्तं राक्षसाधिपम्} %||5-8-9||

\twolineshloka
{निःश्वसन्तं यथा नागं रावणं वानरोत्तमः}
{आसाद्य परमोद्विग्नः सोऽपासर्पत्सुभीतवत्} %||5-8-10||

\twolineshloka
{अथारोहणमासाद्य वेदिकान्तरमाश्रितः}
{सुप्तं राक्षसशार्दूलं प्रेक्षते स्म महाकपिः} %||5-8-11||

\twolineshloka
{शुशुभे राक्षसेन्द्रस्य स्वपतः शयनोत्तमम्}
{गन्धहस्तिनि संविष्टे यथाप्रस्रवणं महत्} %||5-8-12||

\twolineshloka
{काञ्चनाङ्गदनद्धौ च ददर्श स महात्मनः}
{विक्षिप्तौ राक्षसेन्द्रस्य भुजाविन्द्रध्वजोपमौ} %||5-8-13||

\twolineshloka
{ऐरावतविषाणाग्रैरापीडितकृतव्रणौ}
{वज्रोल्लिखितपीनांसौ विष्णुचक्रपरिक्षितौ} %||5-8-14||

\twolineshloka
{पीनौ समसुजातांसौ संगतौ बलसंयुतौ}
{सुलक्षण नखाङ्गुष्ठौ स्वङ्गुलीतललक्षितौ} %||5-8-15||

\twolineshloka
{संहतौ परिघाकारौ वृत्तौ करिकरोपमौ}
{विक्षिप्तौ शयने शुभ्रे पञ्चशीर्षाविवोरगौ} %||5-8-16||

\twolineshloka
{शशक्षतजकल्पेन सुशीतेन सुगन्धिना}
{चन्दनेन परार्ध्येन स्वनुलिप्तौ स्वलंकृतौ} %||5-8-17||

\twolineshloka
{उत्तमस्त्रीविमृदितौ गन्धोत्तमनिषेवितौ}
{यक्षपन्नगगन्धर्वदेवदानवराविणौ} %||5-8-18||

\twolineshloka
{ददर्श स कपिस्तस्य बाहू शयनसंस्थितौ}
{मन्दरस्यान्तरे सुप्तौ महार्ही रुषिताविव} %||5-8-19||

\twolineshloka
{ताभ्यां स परिपूर्णाभ्यां भुजाभ्यां राक्षसाधिपः}
{शुशुभेऽचलसंकाशः शृङ्गाभ्यामिव मन्दरः} %||5-8-20||

\twolineshloka
{चूतपुंनागसुरभिर्बकुलोत्तमसंयुतः}
{मृष्टान्नरससंयुक्तः पानगन्धपुरःसरः} %||5-8-21||

\twolineshloka
{तस्य राक्षससिंहस्य निश्चक्राम मुखान्महान्}
{शयानस्य विनिःश्वासः पूरयन्निव तद्गृहम्} %||5-8-22||

\twolineshloka
{मुक्तामणिविचित्रेण काञ्चनेन विराजता}
{मुकुटेनापवृत्तेन कुण्डलोज्ज्वलिताननम्} %||5-8-23||

\twolineshloka
{रक्तचन्दनदिग्धेन तथा हारेण शोभिता}
{पीनायतविशालेन वक्षसाभिविराजितम्} %||5-8-24||

\twolineshloka
{पाण्डुरेणापविद्धेन क्षौमेण क्षतजेक्षणम्}
{महार्हेण सुसंवीतं पीतेनोत्तमवाससा} %||5-8-25||

\twolineshloka
{माषराशिप्रतीकाशं निःश्वसन्तं भुजङ्गवत्}
{गाङ्गे महति तोयान्ते प्रसुतमिव कुञ्जरम्} %||5-8-26||

\twolineshloka
{चतुर्भिः काञ्चनैर्दीपैर्दीप्यमानैश्चतुर्दिशम्}
{प्रकाशीकृतसर्वाङ्गं मेघं विद्युद्गणैरिव} %||5-8-27||

\twolineshloka
{पादमूलगताश्चापि ददर्श सुमहात्मनः}
{पत्नीः स प्रियभार्यस्य तस्य रक्षःपतेर्गृहे} %||5-8-28||

\twolineshloka
{शशिप्रकाशवदना वरकुण्डलभूषिताः}
{अम्लानमाल्याभरणा ददर्श हरियूथपः} %||5-8-29||

\twolineshloka
{नृत्तवादित्रकुशला राक्षसेन्द्रभुजाङ्कगाः}
{वराभरणधारिण्यो निषन्ना ददृशे कपिः} %||5-8-30||

\twolineshloka
{वज्रवैदूर्यगर्भाणि श्रवणान्तेषु योषिताम्}
{ददर्श तापनीयानि कुण्डलान्यङ्गदानि च} %||5-8-31||

\twolineshloka
{तासां चन्द्रोपमैर्वक्त्रैः शुभैर्ललितकुण्डलैः}
{विरराज विमानं तन्नभस्तारागणैरिव} %||5-8-32||

\twolineshloka
{मदव्यायामखिन्नास्ता राक्षसेन्द्रस्य योषितः}
{तेषु तेष्ववकाशेषु प्रसुप्तास्तनुमध्यमाः} %||5-8-33||

\twolineshloka
{काचिद्वीणां परिष्वज्य प्रसुप्ता संप्रकाशते}
{महानदीप्रकीर्णेव नलिनी पोतमाश्रिता} %||5-8-34||

\twolineshloka
{अन्या कक्षगतेनैव मड्डुकेनासितेक्षणा}
{प्रसुप्ता भामिनी भाति बालपुत्रेव वत्सला} %||5-8-35||

\twolineshloka
{पटहं चारुसर्वाङ्गी पीड्य शेते शुभस्तनी}
{चिरस्य रमणं लब्ध्वा परिष्वज्येव कामिनी} %||5-8-36||

\twolineshloka
{काचिदंशं परिष्वज्य सुप्ता कमललोचना}
{निद्रावशमनुप्राप्ता सहकान्तेव भामिनी} %||5-8-37||

\twolineshloka
{अन्या कनकसंकाशैर्मृदुपीनैर्मनोरमैः}
{मृदङ्गं परिपीड्याङ्गैः प्रसुप्ता मत्तलोचना} %||5-8-38||

\twolineshloka
{भुजपार्श्वान्तरस्थेन कक्षगेण कृशोदरी}
{पणवेन सहानिन्द्या सुप्ता मदकृतश्रमा} %||5-8-39||

\twolineshloka
{डिण्डिमं परिगृह्यान्या तथैवासक्तडिण्डिमा}
{प्रसुप्ता तरुणं वत्समुपगूह्येव भामिनी} %||5-8-40||

\twolineshloka
{काचिदाडम्बरं नारी भुजसंभोगपीडितम्}
{कृत्वा कमलपत्राक्षी प्रसुप्ता मदमोहिता} %||5-8-41||

\twolineshloka
{कलशीमपविद्ध्यान्या प्रसुप्ता भाति भामिनी}
{वसन्ते पुष्पशबला मालेव परिमार्जिता} %||5-8-42||

\twolineshloka
{पाणिभ्यां च कुचौ काचित्सुवर्णकलशोपमौ}
{उपगूह्याबला सुप्ता निद्राबलपराजिता} %||5-8-43||

\twolineshloka
{अन्या कमलपत्राक्षी पूर्णेन्दुसदृशानना}
{अन्यामालिङ्ग्य सुश्रोणी प्रसुप्ता मदविह्वला} %||5-8-44||

\twolineshloka
{आतोद्यानि विचित्राणि परिष्वज्य वरस्त्रियः}
{निपीड्य च कुचैः सुप्ताः कामिन्यः कामुकानिव} %||5-8-45||

\twolineshloka
{तासामेकान्तविन्यस्ते शयानां शयने शुभे}
{ददर्श रूपसंपन्नामपरां स कपिः स्त्रियम्} %||5-8-46||

\twolineshloka
{मुक्तामणिसमायुक्तैर्भूषणैः सुविभूषिताम्}
{विभूषयन्तीमिव च स्वश्रिया भवनोत्तमम्} %||5-8-47||

\twolineshloka
{गौरीं कनकवर्णाभामिष्टामन्तःपुरेश्वरीम्}
{कपिर्मन्दोदरीं तत्र शयानां चारुरूपिणीम्} %||5-8-48||

\threelineshloka
{स तां दृष्ट्वा महाबाहुर्भूषितां मारुतात्मजः}
{तर्कयामास सीतेति रूपयौवनसंपदा}
{हर्षेण महता युक्तो ननन्द हरियूथपः} %||5-8-49||

\fourlineindentedshloka
{आस्ह्पोटयामास चुचुम्ब पुच्छं}
{ ननन्द चिक्रीड जगौ जगाम}
{स्तम्भानरोहन्निपपात भूमौ}
{ निदर्शयन्स्वां प्रकृतिं कपीनाम्} %||5-9-50||


{॥इत्यार्षे श्रीमद्रामायणे वाल्मीकीये आदिकाव्ये सुन्दरकाण्डे अष्टमः सर्गः॥}

\sect{नवमः सर्गः}

\twolineshloka
{अवधूय च तां बुद्धिं बभूवावस्थितस्तदा}
{जगाम चापरां चिन्तां सीतां प्रति महाकपिः} %||5-9-1||

\twolineshloka
{न रामेण वियुक्ता सा स्वप्तुमर्हति भामिनी}
{न भोक्तुं नाप्यलंकर्तुं न पानमुपसेवितुम्} %||5-9-2||

\threelineshloka
{नान्यं नरमुपस्थातुं सुराणामपि चेश्वरम्}
{न हि रामसमः कश्चिद्विद्यते त्रिदशेष्वपि}
{अन्येयमिति निश्चित्य पानभूमौ चचार सः} %||5-9-3||

\twolineshloka
{क्रीडितेनापराः क्लान्ता गीतेन च तथा पराः}
{नृत्तेन चापराः क्लान्ताः पानविप्रहतास्तथा} %||5-9-4||

\twolineshloka
{मुरजेषु मृदङ्गेषु पीठिकासु च संस्थिताः}
{तथास्तरणमुख्य्येषु संविष्टाश्चापराः स्त्रियः} %||5-9-5||

\twolineshloka
{अङ्गनानां सहस्रेण भूषितेन विभूषणैः}
{रूपसंलापशीलेन युक्तगीतार्थभाषिणा} %||5-9-6||

\twolineshloka
{देशकालाभियुक्तेन युक्तवाक्याभिधायिना}
{रताभिरतसंसुप्तं ददर्श हरियूथपः} %||5-9-7||

\twolineshloka
{तासां मध्ये महाबाहुः शुशुभे राक्षसेश्वरः}
{गोष्ठे महति मुख्यानां गवां मध्ये यथा वृषः} %||5-9-8||

\twolineshloka
{स राक्षसेन्द्रः शुशुभे ताभिः परिवृतः स्वयम्}
{करेणुभिर्यथारण्यं परिकीर्णो महाद्विपः} %||5-9-9||

\twolineshloka
{सर्वकामैरुपेतां च पानभूमिं महात्मनः}
{ददर्श कपिशार्दूलस्तस्य रक्षःपतेर्गृहे} %||5-9-10||

\twolineshloka
{मृगाणां महिषाणां च वराहाणां च भागशः}
{तत्र न्यस्तानि मांसानि पानभूमौ ददर्श सः} %||5-9-11||

\twolineshloka
{रौक्मेषु च विशलेषु भाजनेष्वर्धभक्षितान्}
{ददर्श कपिशार्दूल मयूरान्कुक्कुटांस्तथा} %||5-9-12||

\twolineshloka
{वराहवार्ध्राणसकान्दधिसौवर्चलायुतान्}
{शल्यान्मृगमयूरांश्च हनूमानन्ववैक्षत} %||5-9-13||

\threelineshloka
{कृकरान्विविधान्सिद्धांश्चकोरानर्धभक्षितान्}
{महिषानेकशल्यांश्च छागांश्च कृतनिष्ठितान्}
{लेख्यमुच्चावचं पेयं भोज्यानि विविधानि च} %||5-9-14||

\twolineshloka
{तथाम्ललवणोत्तंसैर्विविधै रागषाडवैः}
{हार नूपुरकेयूरैरपविद्धैर्महाधनैः} %||5-9-15||

\twolineshloka
{पानभाजनविक्षिप्तैः फलैश्च विविधैरपि}
{कृतपुष्पोपहारा भूरधिकं पुष्यति श्रियम्} %||5-9-16||

\twolineshloka
{तत्र तत्र च विन्यस्तैः सुश्लिष्टैः शयनासनैः}
{पानभूमिर्विना वह्निं प्रदीप्तेवोपलक्ष्यते} %||5-9-17||

\twolineshloka
{बहुप्रकारैर्विविधैर्वरसंस्कारसंस्कृतैः}
{मांसैः कुशलसंयुक्तैः पानभूमिगतैः पृथक्} %||5-9-18||

\threelineshloka
{दिव्याः प्रसन्ना विविधाः सुराः कृतसुरा अपि}
{शर्करासवमाध्वीकाः पुष्पासवफलासवाः}
{वासचूर्णैश्च विविधैर्मृष्टास्तैस्तैः पृथक्पृथक्} %||5-9-19||

\threelineshloka
{संतता शुशुभे भूमिर्माल्यैश्च बहुसंस्थितैः}
{हिरण्मयैश्च करकैर्भाजनैः स्फाटिकैरपि}
{जाम्बूनदमयैश्चान्यैः करकैरभिसंवृता} %||5-9-20||

\twolineshloka
{राजतेषु च कुम्भेषु जाम्बूनदमयेषु च}
{पानश्रेष्ठं तदा भूरि कपिस्तत्र ददर्श ह} %||5-9-21||

\twolineshloka
{सोऽपश्यच्छातकुम्भानि शीधोर्मणिमयानि च}
{राजतानि च पूर्णानि भाजनानि महाकपिः} %||5-9-22||

\twolineshloka
{क्वचिदर्धावशेषाणि क्वचित्पीतानि सर्वशः}
{क्वचिन्नैव प्रपीतानि पानानि स ददर्श ह} %||5-9-23||

\twolineshloka
{क्वचिद्भक्ष्यांश्च विविधान्क्वचित्पानानि भागशः}
{क्वचिदन्नावशेषाणि पश्यन्वै विचचार ह} %||5-9-24||

\twolineshloka
{क्वचित्प्रभिन्नैः करकैः क्वचिदालोडितैर्घटैः}
{क्वचित्संपृक्तमाल्यानि जलानि च फलानि च} %||5-9-25||

\twolineshloka
{शयनान्यत्र नारीणां शून्यानि बहुधा पुनः}
{परस्परं समाश्लिष्य काश्चित्सुप्ता वराङ्गनाः} %||5-9-26||

\twolineshloka
{काचिच्च वस्त्रमन्यस्या अपहृत्योपगुह्य च}
{उपगम्याबला सुप्ता निद्राबलपराजिता} %||5-9-27||

\twolineshloka
{तासामुच्छ्वासवातेन वस्त्रं माल्यं च गात्रजम्}
{नात्यर्थं स्पन्दते चित्रं प्राप्य मन्दमिवानिलम्} %||5-9-28||

\twolineshloka
{चन्दनस्य च शीतस्य शीधोर्मधुरसस्य च}
{विविधस्य च माल्यस्य पुष्पस्य विविधस्य च} %||5-9-29||

\threelineshloka
{बहुधा मारुतस्तत्र गन्धं विविधमुद्वहन्}
{स्नानानां चन्दनानां च धूपानां चैव मूर्छितः}
{प्रववौ सुरभिर्गन्धो विमाने पुष्पके तदा} %||5-9-30||

\twolineshloka
{श्यामावदातास्तत्रान्याः काश्चित्कृष्णा वराङ्गनाः}
{काश्चित्काञ्चनवर्णाङ्ग्यः प्रमदा राक्षसालये} %||5-9-31||

\twolineshloka
{तासां निद्रावशत्वाच्च मदनेन विमूर्छितम्}
{पद्मिनीनां प्रसुप्तानां रूपमासीद्यथैव हि} %||5-9-32||

\twolineshloka
{एवं सर्वमशेषेण रावणान्तःपुरं कपिः}
{ददर्श सुमहातेजा न ददर्श च जानकीम्} %||5-9-33||

\twolineshloka
{निरीक्षमाणश्च ततस्ताः स्त्रियः स महाकपिः}
{जगाम महतीं चिन्तां धर्मसाध्वसशङ्कितः} %||5-9-34||

\twolineshloka
{परदारावरोधस्य प्रसुप्तस्य निरीक्षणम्}
{इदं खलु ममात्यर्थं धर्मलोपं करिष्यति} %||5-9-35||

\twolineshloka
{न हि मे परदाराणां दृष्टिर्विषयवर्तिनी}
{अयं चात्र मया दृष्टः परदारपरिग्रहः} %||5-9-36||

\twolineshloka
{तस्य प्रादुरभूच्चिन्तापुनरन्या मनस्विनः}
{निश्चितैकान्तचित्तस्य कार्यनिश्चयदर्शिनी} %||5-9-37||

\twolineshloka
{कामं दृष्ट्वा मया सर्वा विश्वस्ता रावणस्त्रियः}
{न तु मे मनसः किंचिद्वैकृत्यमुपपद्यते} %||5-9-38||

\twolineshloka
{मनो हि हेतुः सर्वेषामिन्द्रियाणां प्रवर्तते}
{शुभाशुभास्ववस्थासु तच्च मे सुव्यवस्थितम्} %||5-9-39||

\twolineshloka
{नान्यत्र हि मया शक्या वैदेही परिमार्गितुम्}
{स्त्रियो हि स्त्रीषु दृश्यन्ते सदा संपरिमार्गणे} %||5-9-40||

\twolineshloka
{यस्य सत्त्वस्य या योनिस्तस्यां तत्परिमार्ग्यते}
{न शक्यं प्रमदा नष्टा मृगीषु परिमार्गितुम्} %||5-9-41||

\twolineshloka
{तदिदं मार्गितं तावच्छुद्धेन मनसा मया}
{रावणान्तःपुरं सरं दृश्यते न च जानकी} %||5-9-42||

\twolineshloka
{देवगन्धर्वकन्याश्च नागकन्याश्च वीर्यवान्}
{अवेक्षमाणो हनुमान्नैवापश्यत जानकीम्} %||5-9-43||

\twolineshloka
{तामपश्यन्कपिस्तत्र पश्यंश्चान्या वरस्त्रियः}
{अपक्रम्य तदा वीरः प्रध्यातुमुपचक्रमे} %||5-10-44||


{॥इत्यार्षे श्रीमद्रामायणे वाल्मीकीये आदिकाव्ये सुन्दरकाण्डे नवमः सर्गः॥}

\sect{दशमः सर्गः}

\fourlineindentedshloka
{स तस्य मध्ये भवनस्य वानरो}
{ लतागृहांश्चित्रगृहान्निशागृहान्}
{जगाम सीतां प्रति दर्शनोत्सुको}
{ न चैव तां पश्यति चारुदर्शनाम्} %||5-10-1||

\fourlineindentedshloka
{स चिन्तयामास ततो महाकपिः}
{ प्रियामपश्यन्रघुनन्दनस्य ताम्}
{ध्रुवं नु सीता म्रियते यथा न मे}
{ विचिन्वतो दर्शनमेति मैथिली} %||5-10-2||

\fourlineindentedshloka
{सा राक्षसानां प्रवरेण बाला}
{ स्वशीलसंरक्षण तत्परा सती}
{अनेन नूनं प्रतिदुष्टकर्मणा}
{ हता भवेदार्यपथे परे स्थिता} %||5-10-3||

\fourlineindentedshloka
{विरूपरूपा विकृता विवर्चसो}
{ महानना दीर्घविरूपदर्शनाः}
{समीक्ष्य सा राक्षसराजयोषितो}
{ भयाद्विनष्टा जनकेश्वरात्मजा} %||5-10-4||

\fourlineindentedshloka
{सीतामदृष्ट्वा ह्यनवाप्य पौरुषं}
{ विहृत्य कालं सह वानरैश्चिरम्}
{न मेऽस्ति सुग्रीवसमीपगा गतिः}
{ सुतीक्ष्णदण्डो बलवांश्च वानरः} %||5-10-5||

\twolineshloka
{दृष्टमन्तःपुरं सर्वं दृष्ट्वा रावणयोषितः}
{न सीता दृश्यते साध्वी वृथा जातो मम श्रमः} %||5-10-6||

\twolineshloka
{किं नु मां वानराः सर्वे गतं वक्ष्यन्ति संगताः}
{गत्वा तत्र त्वया वीर किं कृतं तद्वदस्व नः} %||5-10-7||

\twolineshloka
{अदृष्ट्वा किं प्रवक्ष्यामि तामहं जनकात्मजाम्}
{ध्रुवं प्रायमुपेष्यन्ति कालस्य व्यतिवर्तने} %||5-10-8||

\twolineshloka
{किं वा वक्ष्यति वृद्धश्च जाम्बवानङ्गदश्च सः}
{गतं पारं समुद्रस्य वानराश्च समागताः} %||5-10-9||

\twolineshloka
{अनिर्वेदः श्रियो मूलमनिर्वेदः परं सुखम्}
{भूयस्तावद्विचेष्यामि न यत्र विचयः कृतः} %||5-10-10||

\twolineshloka
{अनिर्वेदो हि सततं सर्वार्थेषु प्रवर्तकः}
{करोति सफलं जन्तोः कर्म यच्च करोति सः} %||5-10-11||

\twolineshloka
{तस्मादनिर्वेद कृतं यत्नं चेष्टेऽहमुत्तमम्}
{अदृष्टांश्च विचेष्यामि देशान्रावणपालितान्} %||5-10-12||

\twolineshloka
{आपानशालाविचितास्तथा पुष्पगृहाणि च}
{चित्रशालाश्च विचिता भूयः क्रीडागृहाणि च} %||5-10-13||

\twolineshloka
{निष्कुटान्तररथ्याश्च विमानानि च सर्वशः}
{इति संचिन्त्य भूयोऽपि विचेतुमुपचक्रमे} %||5-10-14||

\twolineshloka
{भूमीगृहांश्चैत्यगृहान्गृहातिगृहकानपि}
{उत्पतन्निपतंश्चापि तिष्ठन्गच्छन्पुनः क्वचित्} %||5-10-15||

\threelineshloka
{अपावृण्वंश्च द्वाराणि कपाटान्यवघट्टयन्}
{प्रविशन्निष्पतंश्चापि प्रपतन्नुत्पतन्नपि}
{सर्वमप्यवकाशं स विचचार महाकपिः} %||5-10-16||

\twolineshloka
{चतुरङ्गुलमात्रोऽपि नावकाशः स विद्यते}
{रावणान्तःपुरे तस्मिन्यं कपिर्न जगाम सः} %||5-10-17||

\twolineshloka
{प्राकरान्तररथ्याश्च वेदिकश्चैत्यसंश्रयाः}
{श्वभ्राश्च पुष्करिण्यश्च सर्वं तेनावलोकितम्} %||5-10-18||

\twolineshloka
{राक्षस्यो विविधाकारा विरूपा विकृतास्तथा}
{दृष्टा हनूमता तत्र न तु सा जनकात्मजा} %||5-10-19||

\twolineshloka
{रूपेणाप्रतिमा लोके वरा विद्याधर स्त्रियः}
{दृटा हनूमता तत्र न तु राघवनन्दिनी} %||5-10-20||

\twolineshloka
{नागकन्या वरारोहाः पूर्णचन्द्रनिभाननाः}
{दृष्टा हनूमता तत्र न तु सीता सुमध्यमा} %||5-10-21||

\twolineshloka
{प्रमथ्य राक्षसेन्द्रेण नागकन्या बलाद्धृताः}
{दृष्टा हनूमता तत्र न सा जनकनन्दिनी} %||5-10-22||

\twolineshloka
{सोऽपश्यंस्तां महाबाहुः पश्यंश्चान्या वरस्त्रियः}
{विषसाद महाबाहुर्हनूमान्मारुतात्मजः} %||5-10-23||

\twolineshloka
{उद्योगं वानरेन्द्राणं प्लवनं सागरस्य च}
{व्यर्थं वीक्ष्यानिलसुतश्चिन्तां पुनरुपागमत्} %||5-10-24||

\twolineshloka
{अवतीर्य विमानाच्च हनूमान्मारुतात्मजः}
{चिन्तामुपजगामाथ शोकोपहतचेतनः} %||5-11-25||


{॥इत्यार्षे श्रीमद्रामायणे वाल्मीकीये आदिकाव्ये सुन्दरकाण्डे दशमः सर्गः॥}

\sect{एकादशः सर्गः}

\twolineshloka
{विमानात्तु सुसंक्रम्य प्राकारं हरियूथपः}
{हनूमान्वेगवानासीद्यथा विद्युद्घनान्तरे} %||5-11-1||

\twolineshloka
{संपरिक्रम्य हनुमान्रावणस्य निवेशनान्}
{अदृष्ट्वा जानकीं सीतामब्रवीद्वचनं कपिः} %||5-11-2||

\twolineshloka
{भूयिष्ठं लोडिता लङ्का रामस्य चरता प्रियम्}
{न हि पश्यामि वैदेहीं सीतां सर्वाङ्गशोभनाम्} %||5-11-3||

\threelineshloka
{पल्वलानि तटाकानि सरांसि सरितस्तथा}
{नद्योऽनूपवनान्ताश्च दुर्गाश्च धरणीधराः}
{लोडिता वसुधा सर्वा न च पश्यामि जानकीम्} %||5-11-4||

\twolineshloka
{इह संपातिना सीता रावणस्य निवेशने}
{आख्याता गृध्रराजेन न च पश्यामि तामहम्} %||5-11-5||

\twolineshloka
{किं नु सीताथ वैदेही मैथिली जनकात्मजा}
{उपतिष्ठेत विवशा रावणं दुष्टचारिणम्} %||5-11-6||

\twolineshloka
{क्षिप्रमुत्पततो मन्ये सीतामादाय रक्षसः}
{बिभ्यतो रामबाणानामन्तरा पतिता भवेत्} %||5-11-7||

\twolineshloka
{अथ वा ह्रियमाणायाः पथि सिद्धनिषेविते}
{मन्ये पतितमार्याया हृदयं प्रेक्ष्य सागरम्} %||5-11-8||

\twolineshloka
{रावणस्योरुवेगेन भुजाभ्यां पीडितेन च}
{तया मन्ये विशालाक्ष्या त्यक्तं जीवितमार्यया} %||5-11-9||

\twolineshloka
{उपर्युपरि वा नूनं सागरं क्रमतस्तदा}
{विवेष्टमाना पतिता समुद्रे जनकात्मजा} %||5-11-10||

\twolineshloka
{आहो क्षुद्रेण चानेन रक्षन्ती शीलमात्मनः}
{अबन्धुर्भक्षिता सीता रावणेन तपस्विनी} %||5-11-11||

\twolineshloka
{अथ वा राक्षसेन्द्रस्य पत्नीभिरसितेक्षणा}
{अदुष्टा दुष्टभावाभिर्भक्षिता सा भविष्यति} %||5-11-12||

\twolineshloka
{संपूर्णचन्द्रप्रतिमं पद्मपत्रनिभेक्षणम्}
{रामस्य ध्यायती वक्त्रं पञ्चत्वं कृपणा गता} %||5-11-13||

\twolineshloka
{हा राम लक्ष्मणेत्येव हायोध्येति च मैथिली}
{विलप्य बहु वैदेही न्यस्तदेहा भविष्यति} %||5-11-14||

\twolineshloka
{अथ वा निहिता मन्ये रावणस्य निवेशने}
{नूनं लालप्यते मन्दं पञ्जरस्थेव शारिका} %||5-11-15||

\twolineshloka
{जनकस्य कुले जाता रामपत्नी सुमध्यमा}
{कथमुत्पलपत्राक्षी रावणस्य वशं व्रजेत्} %||5-11-16||

\twolineshloka
{विनष्टा वा प्रनष्टा वा मृता वा जनकात्मजा}
{रामस्य प्रियभार्यस्य न निवेदयितुं क्षमम्} %||5-11-17||

\twolineshloka
{निवेद्यमाने दोषः स्याद्दोषः स्यादनिवेदने}
{कथं नु खलु कर्तव्यं विषमं प्रतिभाति मे} %||5-11-18||

\twolineshloka
{अस्मिन्नेवंगते कर्ये प्राप्तकालं क्षमं च किम्}
{भवेदिति मतिं भूयो हनुमान्प्रविचारयन्} %||5-11-19||

\twolineshloka
{यदि सीतामदृष्ट्वाहं वानरेन्द्रपुरीमितः}
{गमिष्यामि ततः को मे पुरुषार्थो भविष्यति} %||5-11-20||

\twolineshloka
{ममेदं लङ्घनं व्यर्थं सागरस्य भविष्यति}
{प्रवेशश्चिव लङ्काया राक्षसानां च दर्शनम्} %||5-11-21||

\twolineshloka
{किं वा वक्ष्यति सुग्रीवो हरयो व समागताः}
{किष्किन्धां समनुप्राप्तौ तौ वा दशरथात्मजौ} %||5-11-22||

\twolineshloka
{गत्वा तु यदि काकुत्स्थं वक्ष्यामि परमप्रियम्}
{न दृष्टेति मया सीता ततस्त्यक्ष्यन्ति जीवितम्} %||5-11-23||

\twolineshloka
{परुषं दारुणं क्रूरं तीक्ष्णमिन्द्रियतापनम्}
{सीतानिमित्तं दुर्वाक्यं श्रुत्वा स न भविष्यति} %||5-11-24||

\twolineshloka
{तं तु कृच्छ्रगतं दृष्ट्वा पञ्चत्वगतमानसं}
{भृशानुरक्तो मेधावी न भविष्यति लक्ष्मणः} %||5-11-25||

\twolineshloka
{विनष्टौ भ्रातरौ श्रुत्वा भरतोऽपि मरिष्यति}
{भरतं च मृतं दृष्ट्वा शत्रुघ्नो न भविष्यति} %||5-11-26||

\twolineshloka
{पुत्रान्मृतान्समीक्ष्याथ न भविष्यन्ति मातरः}
{कौसल्या च सुमित्रा च कैकेयी च न संशयः} %||5-11-27||

\twolineshloka
{कृतज्ञः सत्यसंधश्च सुग्रीवः प्लवगाधिपः}
{रामं तथा गतं दृष्ट्वा ततस्त्यक्ष्यन्ति जीवितम्} %||5-11-28||

\twolineshloka
{दुर्मना व्यथिता दीना निरानन्दा तपस्विनी}
{पीडिता भर्तृशोकेन रुमा त्यक्ष्यति जीवितम्} %||5-11-29||

\twolineshloka
{वालिजेन तु दुःखेन पीडिता शोककर्शिता}
{पञ्चत्वगमने राज्ञस्तारापि न भविष्यति} %||5-11-30||

\twolineshloka
{मातापित्रोर्विनाशेन सुग्रीव व्यसनेन च}
{कुमारोऽप्यङ्गदः कस्माद्धारयिष्यति जीवितम्} %||5-11-31||

\twolineshloka
{भर्तृजेन तु शोकेन अभिभूता वनौकसः}
{शिरांस्यभिहनिष्यन्ति तलैर्मुष्टिभिरेव च} %||5-11-32||

\twolineshloka
{सान्त्वेनानुप्रदानेन मानेन च यशस्विना}
{लालिताः कपिराजेन प्राणांस्त्यक्ष्यन्ति वानराः} %||5-11-33||

\twolineshloka
{न वनेषु न शैलेषु न निरोधेषु वा पुनः}
{क्रीडामनुभविष्यन्ति समेत्य कपिकुञ्जराः} %||5-11-34||

\twolineshloka
{सपुत्रदाराः सामात्या भर्तृव्यसनपीडिताः}
{शैलाग्रेभ्यः पतिष्यन्ति समेत्य विषमेषु च} %||5-11-35||

\twolineshloka
{विषमुद्बन्धनं वापि प्रवेशं ज्वलनस्य वा}
{उपवासमथो शस्त्रं प्रचरिष्यन्ति वानराः} %||5-11-36||

\twolineshloka
{घोरमारोदनं मन्ये गते मयि भविष्यति}
{इक्ष्वाकुकुलनाशश्च नाशश्चैव वनौकसाम्} %||5-11-37||

\twolineshloka
{सोऽहं नैव गमिष्यामि किष्किन्धां नगरीमितः}
{न हि शक्ष्याम्यहं द्रष्टुं सुग्रीवं मैथिलीं विना} %||5-11-38||

\twolineshloka
{मय्यगच्छति चेहस्थे धर्मात्मानौ महारथौ}
{आशया तौ धरिष्येते वनराश्च मनस्विनः} %||5-11-39||

\twolineshloka
{हस्तादानो मुखादानो नियतो वृक्षमूलिकः}
{वानप्रस्थो भविष्यामि अदृष्ट्वा जनकात्मजाम्} %||5-11-40||

\twolineshloka
{सागरानूपजे देशे बहुमूलफलोदके}
{चितां कृत्वा प्रवेक्ष्यामि समिद्धमरणीसुतम्} %||5-11-41||

\twolineshloka
{उपविष्टस्य वा सम्यग्लिङ्गिनं साधयिष्यतः}
{शरीरं भक्षयिष्यन्ति वायसाः श्वापदानि च} %||5-11-42||

\twolineshloka
{इदमप्यृषिभिर्दृष्टं निर्याणमिति मे मतिः}
{सम्यगापः प्रवेक्ष्यामि न चेत्पश्यामि जानकीम्} %||5-11-43||

\twolineshloka
{सुजातमूला सुभगा कीर्तिमालायशस्विनी}
{प्रभग्ना चिररात्रीयं मम सीतामपश्यतः} %||5-11-44||

\twolineshloka
{तापसो वा भविष्यामि नियतो वृक्षमूलिकः}
{नेतः प्रतिगमिष्यामि तामदृष्ट्वासितेक्षणाम्} %||5-11-45||

\twolineshloka
{यदीतः प्रतिगच्छामि सीतामनधिगम्य ताम्}
{अङ्गदः सहितैः सर्वैर्वानरैर्न भविष्यति} %||5-11-46||

\twolineshloka
{विनाशे बहवो दोषा जीवन्प्राप्नोति भद्रकम्}
{तस्मात्प्राणान्धरिष्यामि ध्रुवो जीवति संगमः} %||5-11-47||

\twolineshloka
{एवं बहुविधं दुःखं मनसा धारयन्मुहुः}
{नाध्यगच्छत्तदा पारं शोकस्य कपिकुञ्जरः} %||5-11-48||

\twolineshloka
{रावणं वा वधिष्यामि दशग्रीवं महाबलम्}
{काममस्तु हृता सीता प्रत्याचीर्णं भविष्यति} %||5-11-49||

\twolineshloka
{अथवैनं समुत्क्षिप्य उपर्युपरि सागरम्}
{रामायोपहरिष्यामि पशुं पशुपतेरिव} %||5-11-50||

\twolineshloka
{इति चिन्ता समापन्नः सीतामनधिगम्य ताम्}
{ध्यानशोका परीतात्मा चिन्तयामास वानरः} %||5-11-51||

\twolineshloka
{यावत्सीतां न पश्यामि रामपत्नीं यशस्विनीम्}
{तावदेतां पुरीं लङ्कां विचिनोमि पुनः पुनः} %||5-11-52||

\twolineshloka
{संपाति वचनाच्चापि रामं यद्यानयाम्यहम्}
{अपश्यन्राघवो भार्यां निर्दहेत्सर्ववानरान्} %||5-11-53||

\twolineshloka
{इहैव नियताहारो वत्स्यामि नियतेन्द्रियः}
{न मत्कृते विनश्येयुः सर्वे ते नरवानराः} %||5-11-54||

\twolineshloka
{अशोकवनिका चापि महतीयं महाद्रुमा}
{इमामभिगमिष्यामि न हीयं विचिता मया} %||5-11-55||

\twolineshloka
{वसून्रुद्रांस्तथादित्यानश्विनौ मरुतोऽपि च}
{नमस्कृत्वा गमिष्यामि रक्षसां शोकवर्धनः} %||5-11-56||

\twolineshloka
{जित्वा तु राक्षसान्देवीमिक्ष्वाकुकुलनन्दिनीम्}
{संप्रदास्यामि रामाया यथासिद्धिं तपस्विने} %||5-11-57||

\twolineshloka
{स मुहूर्तमिव ध्यात्वा चिन्ताविग्रथितेन्द्रियः}
{उदतिष्ठन्महाबाहुर्हनूमान्मारुतात्मजः} %||5-11-58||

\fourlineindentedshloka
{नमोऽस्तु रामाय सलक्ष्मणाय}
{ देव्यै च तस्यै जनकात्मजायै}
{नमोऽस्तु रुद्रेन्द्रयमानिलेभ्यो}
{ नमोऽस्तु चन्द्रार्कमरुद्गणेभ्यः} %||5-11-59||

\twolineshloka
{स तेभ्यस्तु नमस्कृत्वा सुग्रीवाय च मारुतिः}
{दिशः सर्वाः समालोक्य अशोकवनिकां प्रति} %||5-11-60||

\twolineshloka
{स गत्वा मनसा पूर्वमशोकवनिकां शुभाम्}
{उत्तरं चिन्तयामास वानरो मारुतात्मजः} %||5-11-61||

\twolineshloka
{ध्रुवं तु रक्षोबहुला भविष्यति वनाकुला}
{अशोकवनिका चिन्त्या सर्वसंस्कारसंस्कृता} %||5-11-62||

\twolineshloka
{रक्षिणश्चात्र विहिता नूनं रक्षन्ति पादपान्}
{भगवानपि सर्वात्मा नातिक्षोभं प्रवायति} %||5-11-63||

\twolineshloka
{संक्षिप्तोऽयं मयात्मा च रामार्थे रावणस्य च}
{सिद्धिं मे संविधास्यन्ति देवाः सर्षिगणास्त्विह} %||5-11-64||

\twolineshloka
{ब्रह्मा स्वयम्भूर्भगवान्देवाश्चैव दिशन्तु मे}
{सिद्धिमग्निश्च वायुश्च पुरुहूतश्च वज्रधृत्} %||5-11-65||

\twolineshloka
{वरुणः पाशहस्तश्च सोमादित्यै तथैव च}
{अश्विनौ च महात्मानौ मरुतः सर्व एव च} %||5-11-66||

\twolineshloka
{सिद्धिं सर्वाणि भूतानि भूतानां चैव यः प्रभुः}
{दास्यन्ति मम ये चान्ये अदृष्टाः पथि गोचराः} %||5-11-67||

\fourlineindentedshloka
{तदुन्नसं पाण्डुरदन्तमव्रणं}
{ शुचिस्मितं पद्मपलाशलोचनम्}
{द्रक्ष्ये तदार्यावदनं कदा न्वहं}
{ प्रसन्नताराधिपतुल्यदर्शनम्} %||5-11-68||

\fourlineindentedshloka
{क्षुद्रेण पापेन नृशंसकर्मणा}
{ सुदारुणालांकृतवेषधारिणा}
{बलाभिभूता अबला तपस्विनी}
{ कथं नु मे दृष्टपथेऽद्य सा भवेत्} %||5-12-69||


{॥इत्यार्षे श्रीमद्रामायणे वाल्मीकीये आदिकाव्ये सुन्दरकाण्डे एकादशः सर्गः॥}

\sect{द्वादशः सर्गः}

\twolineshloka
{स मुहूर्तमिव ध्यत्वा मनसा चाधिगम्य ताम्}
{अवप्लुतो महातेजाः प्राकारं तस्य वेश्मनः} %||5-12-1||

\twolineshloka
{स तु संहृष्टसर्वाङ्गः प्राकारस्थो महाकपिः}
{पुष्पिताग्रान्वसन्तादौ ददर्श विविधान्द्रुमान्} %||5-12-2||

\twolineshloka
{सालानशोकान्भव्यांश्च चम्पकांश्च सुपुष्पितान्}
{उद्दालकान्नागवृक्षांश्चूतान्कपिमुखानपि} %||5-12-3||

\twolineshloka
{अथाम्रवणसंछन्नां लताशतसमावृताम्}
{ज्यामुक्त इव नाराचः पुप्लुवे वृक्षवाटिकाम्} %||5-12-4||

\twolineshloka
{स प्रविष्य विचित्रां तां विहगैरभिनादिताम्}
{राजतैः काञ्चनैश्चैव पादपैः सर्वतोवृताम्} %||5-12-5||

\twolineshloka
{विहगैर्मृगसंघैश्च विचित्रां चित्रकाननाम्}
{उदितादित्यसंकाशां ददर्श हनुमान्कपिः} %||5-12-6||

\twolineshloka
{वृतां नानाविधैर्वृक्षैः पुष्पोपगफलोपगैः}
{कोकिलैर्भृङ्गराजैश्च मत्तैर्नित्यनिषेविताम्} %||5-12-7||

\twolineshloka
{प्रहृष्टमनुजे कले मृगपक्षिसमाकुले}
{मत्तबर्हिणसंघुष्टां नानाद्विजगणायुताम्} %||5-12-8||

\twolineshloka
{मार्गमाणो वरारोहां राजपुत्रीमनिन्दिताम्}
{सुखप्रसुप्तान्विहगान्बोधयामास वानरः} %||5-12-9||

\twolineshloka
{उत्पतद्भिर्द्विजगणैः पक्षैः सालाः समाहताः}
{अनेकवर्णा विविधा मुमुचुः पुष्पवृष्टयः} %||5-12-10||

\twolineshloka
{पुष्पावकीर्णः शुशुभे हनुमान्मारुतात्मजः}
{अशोकवनिकामध्ये यथा पुष्पमयो गिरिः} %||5-12-11||

\twolineshloka
{दिशः सर्वाभिदावन्तं वृक्षषण्डगतं कपिम्}
{दृष्ट्वा सर्वाणि भूतानि वसन्त इति मेनिरे} %||5-12-12||

\twolineshloka
{वृक्षेभ्यः पतितैः पुष्पैरवकीर्णा पृथग्विधैः}
{रराज वसुधा तत्र प्रमदेव विभूषिता} %||5-12-13||

\twolineshloka
{तरस्विना ते तरवस्तरसाभिप्रकम्पिताः}
{कुसुमानि विचित्राणि ससृजुः कपिना तदा} %||5-12-14||

\twolineshloka
{निर्धूतपत्रशिखराः शीर्णपुष्पफलद्रुमाः}
{निक्षिप्तवस्त्राभरणा धूर्ता इव पराजिताः} %||5-12-15||

\twolineshloka
{हनूमता वेगवता कम्पितास्ते नगोत्तमाः}
{पुष्पपर्णफलान्याशु मुमुचुः पुष्पशालिनः} %||5-12-16||

\twolineshloka
{विहंगसंघैर्हीनास्ते स्कन्धमात्राश्रया द्रुमाः}
{बभूवुरगमाः सर्वे मारुतेनेव निर्धुताः} %||5-12-17||

\twolineshloka
{विधूतकेशी युवतिर्यथा मृदितवर्णिका}
{निष्पीतशुभदन्तौष्ठी नखैर्दन्तैश्च विक्षता} %||5-12-18||

\twolineshloka
{तथा लाङ्गूलहस्तैश्च चरणाभ्यां च मर्दिता}
{बभूवाशोकवनिका प्रभग्नवरपादपा} %||5-12-19||

\twolineshloka
{महालतानां दामानि व्यधमत्तरसा कपिः}
{यथा प्रावृषि विन्ध्यस्य मेघजालानि मारुतः} %||5-12-20||

\twolineshloka
{स तत्र मणिभूमीश्च राजतीश्च मनोरमाः}
{तथा काञ्चनभूमीश्च विचरन्ददृशे कपिः} %||5-12-21||

\twolineshloka
{वापीश्च विविधाकाराः पूर्णाः परमवारिणा}
{महार्हैर्मणिसोपानैरुपपन्नास्ततस्ततः} %||5-12-22||

\twolineshloka
{मुक्ताप्रवालसिकता स्फटिकान्तरकुट्टिमाः}
{काञ्चनैस्तरुभिश्चित्रैस्तीरजैरुपशोभिताः} %||5-12-23||

\twolineshloka
{फुल्लपद्मोत्पलवनाश्चक्रवाकोपकूजिताः}
{नत्यूहरुतसंघुष्टा हंससारसनादिताः} %||5-12-24||

\twolineshloka
{दीर्घाभिर्द्रुमयुक्ताभिः सरिद्भिश्च समन्ततः}
{अमृतोपमतोयाभिः शिवाभिरुपसंस्कृताः} %||5-12-25||

\twolineshloka
{लताशतैरवतताः सन्तानकसमावृताः}
{नानागुल्मावृतवनाः करवीरकृतान्तराः} %||5-12-26||

\twolineshloka
{ततोऽम्बुधरसंकाशं प्रवृद्धशिखरं गिरिम्}
{विचित्रकूटं कूटैश्च सर्वतः परिवारितम्} %||5-12-27||

\twolineshloka
{शिलागृहैरवततं नानावृक्षैः समावृतम्}
{ददर्श कपिशार्दूलो रम्यं जगति पर्वतम्} %||5-12-28||

\twolineshloka
{ददर्श च नगात्तस्मान्नदीं निपतितां कपिः}
{अङ्कादिव समुत्पत्य प्रियस्य पतितां प्रियाम्} %||5-12-29||

\twolineshloka
{जले निपतिताग्रैश्च पादपैरुपशोभिताम्}
{वार्यमाणामिव क्रुद्धां प्रमदां प्रियबन्धुभिः} %||5-12-30||

\twolineshloka
{पुनरावृत्ततोयां च ददर्श स महाकपिः}
{प्रसन्नामिव कान्तस्य कान्तां पुनरुपस्थिताम्} %||5-12-31||

\twolineshloka
{तस्यादूरात्स पद्मिन्यो नानाद्विजगणायुताः}
{ददर्श कपिशार्दूलो हनुमान्मारुतात्मजः} %||5-12-32||

\twolineshloka
{कृत्रिमां दीर्घिकां चापि पूर्णां शीतेन वारिणा}
{मणिप्रवरसोपानां मुक्तासिकतशोभिताम्} %||5-12-33||

\threelineshloka
{विविधैर्मृगसंघैश्च विचित्रां चित्रकाननाम्}
{प्रासादैः सुमहद्भिश्च निर्मितैर्विश्वकर्मणा}
{काननैः कृत्रिमैश्चापि सर्वतः समलंकृताम्} %||5-12-34||

\twolineshloka
{ये केचित्पादपास्तत्र पुष्पोपगफलोपगाः}
{सच्छत्राः सवितर्दीकाः सर्वे सौवर्णवेदिकाः} %||5-12-35||

\twolineshloka
{लताप्रतानैर्बहुभिः पर्णैश्च बहुभिर्वृताम्}
{काञ्चनीं शिंशुपामेकां ददर्श स महाकपिः} %||5-12-36||

\twolineshloka
{सोऽपश्यद्भूमिभागांश्च गर्तप्रस्रवणानि च}
{सुवर्णवृक्षानपरान्ददर्श शिखिसंनिभान्} %||5-12-37||

\twolineshloka
{तेषां द्रुमाणां प्रभया मेरोरिव महाकपिः}
{अमन्यत तदा वीरः काञ्चनोऽस्मीति वानरः} %||5-12-38||

\twolineshloka
{तां काञ्चनैस्तरुगणैर्मारुतेन च वीजिताम्}
{किङ्किणीशतनिर्घोषां दृष्ट्वा विस्मयमागमत्} %||5-12-39||

\twolineshloka
{सुपुष्पिताग्रां रुचिरां तरुणाङ्कुरपल्लवाम्}
{तामारुह्य महावेगः शिंशपां पर्णसंवृताम्} %||5-12-40||

\twolineshloka
{इतो द्रक्ष्यामि वैदेहीं राम दर्शनलालसाम्}
{इतश्चेतश्च दुःखार्तां संपतन्तीं यदृच्छया} %||5-12-41||

\twolineshloka
{अशोकवनिका चेयं दृढं रम्या दुरात्मनः}
{चम्पकैश्चन्दनैश्चापि बकुलैश्च विभूषिता} %||5-12-42||

\twolineshloka
{इयं च नलिनी रम्या द्विजसंघनिषेविता}
{इमां सा राममहिषी नूनमेष्यति जानकी} %||5-12-43||

\twolineshloka
{सा राम राममहिषी राघवस्य प्रिया सदा}
{वनसंचारकुशला नूनमेष्यति जानकी} %||5-12-44||

\twolineshloka
{अथ वा मृगशावाक्षी वनस्यास्य विचक्षणा}
{वनमेष्यति सा चेह रामचिन्तानुकर्शिता} %||5-12-45||

\twolineshloka
{रामशोकाभिसंतप्ता सा देवी वामलोचना}
{वनवासरता नित्यमेष्यते वनचारिणी} %||5-12-46||

\twolineshloka
{वनेचराणां सततं नूनं स्पृहयते पुरा}
{रामस्य दयिता भार्या जनकस्य सुता सती} %||5-12-47||

\twolineshloka
{संध्याकालमनाः श्यामा ध्रुवमेष्यति जानकी}
{नदीं चेमां शिवजलां संध्यार्थे वरवर्णिनी} %||5-12-48||

\twolineshloka
{तस्याश्चाप्यनुरूपेयमशोकवनिका शुभा}
{शुभा या पार्थिवेन्द्रस्य पत्नी रामस्य संमिता} %||5-12-49||

\twolineshloka
{यदि जिवति सा देवी ताराधिपनिभानना}
{आगमिष्यति सावश्यमिमां शिवजलां नदीम्} %||5-12-50||

\fourlineindentedshloka
{एवं तु मत्वा हनुमान्महात्मा}
{ प्रतीक्षमाणो मनुजेन्द्रपत्नीम्}
{अवेक्षमाणश्च ददर्श सर्वं}
{ सुपुष्पिते पर्णघने निलीनः} %||5-13-51||


{॥इत्यार्षे श्रीमद्रामायणे वाल्मीकीये आदिकाव्ये सुन्दरकाण्डे द्वादशः सर्गः॥}

\sect{त्रयोदशः सर्गः}

\twolineshloka
{स वीक्षमाणस्तत्रस्थो मार्गमाणश्च मैथिलीम्}
{अवेक्षमाणश्च महीं सर्वां तामन्ववैक्षत} %||5-13-1||

\twolineshloka
{सन्तान कलताभिश्च पादपैरुपशोभिताम्}
{दिव्यगन्धरसोपेतां सर्वतः समलंकृताम्} %||5-13-2||

\twolineshloka
{तां स नन्दनसंकाशां मृगपक्षिभिरावृताम्}
{हर्म्यप्रासादसंबाधां कोकिलाकुलनिःस्वनाम्} %||5-13-3||

\twolineshloka
{काञ्चनोत्पलपद्माभिर्वापीभिरुपशोभिताम्}
{बह्वासनकुथोपेतां बहुभूमिगृहायुताम्} %||5-13-4||

\twolineshloka
{सर्वर्तुकुसुमै रम्यैः फलवद्भिश्च पादपैः}
{पुष्पितानामशोकानां श्रिया सूर्योदयप्रभाम्} %||5-13-5||

\threelineshloka
{प्रदीप्तामिव तत्रस्थो मारुतिः समुदैक्षत}
{निष्पत्रशाखां विहगैः क्रियमाणामिवासकृत्}
{विनिष्पतद्भिः शतशश्चित्रैः पुष्पावतंसकैः} %||5-13-6||

\twolineshloka
{आमूलपुष्पनिचितैरशोकैः शोकनाशनैः}
{पुष्पभारातिभारैश्च स्पृशद्भिरिव मेदिनीम्} %||5-13-7||

\twolineshloka
{कर्णिकारैः कुसुमितैः किंशुकैश्च सुपुष्पितैः}
{स देशः प्रभया तेषां प्रदीप्त इव सर्वतः} %||5-13-8||

\twolineshloka
{पुंनागाः सप्तपर्णाश्च चम्पकोद्दालकास्तथा}
{विवृद्धमूला बहवः शोभन्ते स्म सुपुष्पिताः} %||5-13-9||

\twolineshloka
{शातकुम्भनिभाः केचित्केचिदग्निशिखोपमाः}
{नीलाञ्जननिभाः केचित्तत्राशोकाः सहस्रशः} %||5-13-10||

\twolineshloka
{नन्दनं विविधोद्यानं चित्रं चैत्ररथं यथा}
{अतिवृत्तमिवाचिन्त्यं दिव्यं रम्यं श्रिया वृतम्} %||5-13-11||

\twolineshloka
{द्वितीयमिव चाकाशं पुष्पज्योतिर्गणायुतम्}
{पुष्परत्नशतैश्चित्रं पञ्चमं सागरं यथा} %||5-13-12||

\twolineshloka
{सर्वर्तुपुष्पैर्निचितं पादपैर्मधुगन्धिभिः}
{नानानिनादैरुद्यानं रम्यं मृगगणैर्द्विजैः} %||5-13-13||

\twolineshloka
{अनेकगन्धप्रवहं पुण्यगन्धं मनोरमम्}
{शैलेन्द्रमिव गन्धाढ्यं द्वितीयं गन्धमादनम्} %||5-13-14||

\twolineshloka
{अशोकवनिकायां तु तस्यां वानरपुंगवः}
{स ददर्शाविदूरस्थं चैत्यप्रासादमूर्जितम्} %||5-13-15||

\twolineshloka
{मध्ये स्तम्भसहस्रेण स्थितं कैलासपाण्डुरम्}
{प्रवालकृतसोपानं तप्तकाञ्चनवेदिकम्} %||5-13-16||

\twolineshloka
{मुष्णन्तमिव चक्षूंषि द्योतमानमिव श्रिया}
{विमलं प्रांशुभावत्वादुल्लिखन्तमिवाम्बरम्} %||5-13-17||

\threelineshloka
{ततो मलिनसंवीतां राक्षसीभिः समावृताम्}
{उपवासकृशां दीनां निःश्वसान्तीं पुनः पुनः}
{ददर्श शुक्लपक्षादौ चन्द्ररेखामिवामलाम्} %||5-13-18||

\twolineshloka
{मन्दप्रख्यायमानेन रूपेण रुचिरप्रभाम्}
{पिनद्धां धूमजालेन शिखामिव विभावसोः} %||5-13-19||

\twolineshloka
{पीतेनैकेन संवीतां क्लिष्टेनोत्तमवाससा}
{सपङ्कामनलंकारां विपद्मामिव पद्मिनीम्} %||5-13-20||

\twolineshloka
{व्रीडितां दुःखसंतप्तां परिम्लानां तपस्विनीम्}
{ग्रहेणाङ्गारकेणैव पीडितामिव रोहिणीम्} %||5-13-21||

\twolineshloka
{अश्रुपूर्णमुखीं दीनां कृशामननशेन च}
{शोकध्यानपरां दीनां नित्यं दुःखपरायणाम्} %||5-13-22||

\twolineshloka
{प्रियं जनमपश्यन्तीं पश्यन्तीं राक्षसीगणम्}
{स्वगणेन मृगीं हीनां श्वगणाभिवृतामिव} %||5-13-23||

\twolineshloka
{नीलनागाभया वेण्या जघनं गतयैकया}
{सुखार्हां दुःखसंतप्तां व्यसनानामकोदिवाम्} %||5-13-24||

\twolineshloka
{तां समीक्ष्य विशालाक्षीमधिकं मलिनां कृशाम्}
{तर्कयामास सीतेति कारणैरुपपादिभिः} %||5-13-25||

\twolineshloka
{ह्रियमाणा तदा तेन रक्षसा कामरूपिणा}
{यथारूपा हि दृष्टा वै तथारूपेयमङ्गना} %||5-13-26||

\twolineshloka
{पूर्णचन्द्राननां सुभ्रूं चारुवृत्तपयोधराम्}
{कुर्वन्तीं प्रभया देवीं सर्वा वितिमिरा दिशः} %||5-13-27||

\twolineshloka
{तां नीलकेशीं बिम्बौष्ठीं सुमध्यां सुप्रतिष्ठिताम्}
{सीतां पद्मपलाशाक्षीं मन्मथस्य रतिं यथा} %||5-13-28||

\twolineshloka
{इष्टां सर्वस्य जगतः पूर्णचन्द्रप्रभामिव}
{भूमौ सुतनुमासीनां नियतामिव तापसीम्} %||5-13-29||

\twolineshloka
{निःश्वासबहुलां भीरुं भुजगेन्द्रवधूमिव}
{शोकजालेन महता विततेन न राजतीम्} %||5-13-30||

\twolineshloka
{संसक्तां धूमजालेन शिखामिव विभावसोः}
{तां स्मृतीमिव संदिघ्दामृद्धिं निपतितामिव} %||5-13-31||

\twolineshloka
{विहतामिव च श्रद्धामाशां प्रतिहतामिव}
{सोपसर्गां यथा सिद्धिं बुद्धिं सकलुषामिव} %||5-13-32||

\twolineshloka
{अभूतेनापवादेन कीर्तिं निपतितामिव}
{रामोपरोधव्यथितां रक्षोहरणकर्शिताम्} %||5-13-33||

\threelineshloka
{अबलां मृगशावाक्षीं वीक्षमाणां ततस्ततः}
{बाष्पाम्बुप्रतिपूर्णेन कृष्णवक्त्राक्षिपक्ष्मणा}
{वदनेनाप्रसन्नेन निःश्वसन्तीं पुनः पुनः} %||5-13-34||

\twolineshloka
{मलपङ्कधरां दीनां मण्डनार्हाममण्डिताम्}
{प्रभां नक्षत्रराजस्य कालमेघैरिवावृताम्} %||5-13-35||

\twolineshloka
{तस्य संदिदिहे बुद्धिर्मुहुः सीतां निरीक्ष्य तु}
{आम्नायानामयोगेन विद्यां प्रशिथिलामिव} %||5-13-36||

\twolineshloka
{दुःखेन बुबुधे सीतां हनुमाननलंकृताम्}
{संस्कारेण यथाहीनां वाचमर्थान्तरं गताम्} %||5-13-37||

\twolineshloka
{तां समीक्ष्य विशालाक्षीं राजपुत्रीमनिन्दिताम्}
{तर्कयामास सीतेति कारणैरुपपादयन्} %||5-13-38||

\twolineshloka
{वैदेह्या यानि चाङ्गेषु तदा रामोऽन्वकीर्तयत्}
{तान्याभरणजालानि गात्रशोभीन्यलक्षयत्} %||5-13-39||

\twolineshloka
{सुकृतौ कर्णवेष्टौ च श्वदंष्ट्रौ च सुसंस्थितौ}
{मणिविद्रुमचित्राणि हस्तेष्वाभरणानि च} %||5-13-40||

\twolineshloka
{श्यामानि चिरयुक्तत्वात्तथा संस्थानवन्ति च}
{तान्येवैतानि मन्येऽहं यानि रामोऽव्नकीर्तयत्} %||5-13-41||

\twolineshloka
{तत्र यान्यवहीनानि तान्यहं नोपलक्षये}
{यान्यस्या नावहीनानि तानीमानि न संशयः} %||5-13-42||

\twolineshloka
{पीतं कनकपट्टाभं स्रस्तं तद्वसनं शुभम्}
{उत्तरीयं नगासक्तं तदा दृष्टं प्लवंगमैः} %||5-13-43||

\twolineshloka
{भूषणानि च मुख्यानि दृष्टानि धरणीतले}
{अनयैवापविद्धानि स्वनवन्ति महान्ति च} %||5-13-44||

\twolineshloka
{इदं चिरगृहीतत्वाद्वसनं क्लिष्टवत्तरम्}
{तथा हि नूनं तद्वर्णं तथा श्रीमद्यथेतरत्} %||5-13-45||

\twolineshloka
{इयं कनकवर्णाङ्गी रामस्य महिषी प्रिया}
{प्रनष्टापि सती यस्य मनसो न प्रणश्यति} %||5-13-46||

\twolineshloka
{इयं सा यत्कृते रामश्चतुर्भिः परितप्यते}
{कारुण्येनानृशंस्येन शोकेन मदनेन च} %||5-13-47||

\twolineshloka
{स्त्री प्रनष्टेति कारुण्यादाश्रितेत्यानृशंस्यतः}
{पत्नी नष्टेति शोकेन प्रियेति मदनेन च} %||5-13-48||

\twolineshloka
{अस्या देव्या यथा रूपमङ्गप्रत्यङ्गसौष्ठवम्}
{रामस्य च यथारूपं तस्येयमसितेक्षणा} %||5-13-49||

\twolineshloka
{अस्या देव्या मनस्तस्मिंस्तस्य चास्यां प्रतिष्ठितम्}
{तेनेयं स च धर्मात्मा मुहूर्तमपि जीवति} %||5-13-50||

\twolineshloka
{दुष्करं कुरुते रामो य इमां मत्तकाशिनीम्}
{सीतां विना महाबाहुर्मुहूर्तमपि जीवति} %||5-13-51||

\twolineshloka
{एवं सीतां तदा दृष्ट्वा हृष्टः पवनसंभवः}
{जगाम मनसा रामं प्रशशंस च तं प्रभुम्} %||5-14-52||


{॥इत्यार्षे श्रीमद्रामायणे वाल्मीकीये आदिकाव्ये सुन्दरकाण्डे त्रयोदशः सर्गः॥}

\sect{चतुर्दशः सर्गः}

\twolineshloka
{प्रशस्य तु प्रशस्तव्यां सीतां तां हरिपुंगवः}
{गुणाभिरामं रामं च पुनश्चिन्तापरोऽभवत्} %||5-14-1||

\twolineshloka
{स मुहूर्तमिव ध्यात्वा बाष्पपर्याकुलेक्षणः}
{सीतामाश्रित्य तेजस्वी हनुमान्विललाप ह} %||5-14-2||

\twolineshloka
{मान्या गुरुविनीतस्य लक्ष्मणस्य गुरुप्रिया}
{यदि सीतापि दुःखार्ता कालो हि दुरतिक्रमः} %||5-14-3||

\twolineshloka
{रामस्य व्यवसायज्ञा लक्ष्मणस्य च धीमतः}
{नात्यर्थं क्षुभ्यते देवी गङ्गेव जलदागमे} %||5-14-4||

\twolineshloka
{तुल्यशीलवयोवृत्तां तुल्याभिजनलक्षणाम्}
{राघवोऽर्हति वैदेहीं तं चेयमसितेक्षणा} %||5-14-5||

\twolineshloka
{तां दृष्ट्वा नवहेमाभां लोककान्तामिव श्रियम्}
{जगाम मनसा रामं वचनं चेदमब्रवीत्} %||5-14-6||

\twolineshloka
{अस्या हेतोर्विशालाक्ष्या हतो वाली महाबलः}
{रावणप्रतिमो वीर्ये कबन्धश्च निपातितः} %||5-14-7||

\twolineshloka
{विराधश्च हतः संख्ये राक्षसो भीमविक्रमः}
{वने रामेण विक्रम्य महेन्द्रेणेव शम्बरः} %||5-14-8||

\twolineshloka
{चतुर्दशसहस्राणि रक्षसां भीमकर्मणाम्}
{निहतानि जनस्थाने शरैरग्निशिखोपमैः} %||5-14-9||

\twolineshloka
{खरश्च निहतः संख्ये त्रिशिराश्च निपातितः}
{दूषणश्च महातेजा रामेण विदितात्मना} %||5-14-10||

\twolineshloka
{ऐश्वर्यं वानराणां च दुर्लभं वालिपालितम्}
{अस्या निमित्ते सुग्रीवः प्राप्तवाँल्लोकसत्कृतम्} %||5-14-11||

\twolineshloka
{सागरश्च मया क्रान्तः श्रीमान्नदनदीपतिः}
{अस्या हेतोर्विशालाक्ष्याः पुरी चेयं निरीक्षिता} %||5-14-12||

\twolineshloka
{यदि रामः समुद्रान्तां मेदिनीं परिवर्तयेत्}
{अस्याः कृते जगच्चापि युक्तमित्येव मे मतिः} %||5-14-13||

\twolineshloka
{राज्यं वा त्रिषु लोकेषु सीता वा जनकात्मजा}
{त्रैलोक्यराज्यं सकलं सीताया नाप्नुयात्कलाम्} %||5-14-14||

\twolineshloka
{इयं सा धर्मशीलस्य मैथिलस्य महात्मनः}
{सुता जनकराजस्य सीता भर्तृदृढव्रता} %||5-14-15||

\twolineshloka
{उत्थिता मेदिनीं भित्त्वा क्षेत्रे हलमुखक्षते}
{पद्मरेणुनिभैः कीर्णा शुभैः केदारपांसुभिः} %||5-14-16||

\twolineshloka
{विक्रान्तस्यार्यशीलस्य संयुगेष्वनिवर्तिनः}
{स्नुषा दशरथस्यैषा ज्येष्ठा राज्ञो यशस्विनी} %||5-14-17||

\twolineshloka
{धर्मज्ञस्य कृतज्ञस्य रामस्य विदितात्मनः}
{इयं सा दयिता भार्या राक्षसी वशमागता} %||5-14-18||

\twolineshloka
{सर्वान्भोगान्परित्यज्य भर्तृस्नेहबलात्कृता}
{अचिन्तयित्वा दुःखानि प्रविष्टा निर्जनं वनम्} %||5-14-19||

\twolineshloka
{संतुष्टा फलमूलेन भर्तृशुश्रूषणे रता}
{या परां भजते प्रीतिं वनेऽपि भवने यथा} %||5-14-20||

\twolineshloka
{सेयं कनकवर्णाङ्गी नित्यं सुस्मितभाषिणी}
{सहते यातनामेतामनर्थानामभागिनी} %||5-14-21||

\twolineshloka
{इमां तु शीलसंपन्नां द्रष्टुमिच्छति राघवः}
{रावणेन प्रमथितां प्रपामिव पिपासितः} %||5-14-22||

\twolineshloka
{अस्या नूनं पुनर्लाभाद्राघवः प्रीतिमेष्यति}
{राजा राज्यपरिभ्रष्टः पुनः प्राप्येव मेदिनीम्} %||5-14-23||

\twolineshloka
{कामभोगैः परित्यक्ता हीना बन्धुजनेन च}
{धारयत्यात्मनो देहं तत्समागमकाङ्क्षिणी} %||5-14-24||

\twolineshloka
{नैषा पश्यति राक्षस्यो नेमान्पुष्पफलद्रुमान्}
{एकस्थहृदया नूनं राममेवानुपश्यति} %||5-14-25||

\twolineshloka
{भर्ता नाम परं नार्या भूषणं भूषणादपि}
{एषा हि रहिता तेन शोभनार्हा न शोभते} %||5-14-26||

\twolineshloka
{दुष्करं कुरुते रामो हीनो यदनया प्रभुः}
{धारयत्यात्मनो देहं न दुःखेनावसीदति} %||5-14-27||

\twolineshloka
{इमामसितकेशान्तां शतपत्रनिभेक्षणाम्}
{सुखार्हां दुःखितां दृष्ट्वा ममापि व्यथितं मनः} %||5-14-28||

\fourlineindentedshloka
{क्षितिक्षमा पुष्करसंनिभाक्षी}
{ या रक्षिता राघवलक्ष्मणाभ्याम्}
{सा राक्षसीभिर्विकृतेक्षणाभिः}
{ संरक्ष्यते संप्रति वृक्षमूले} %||5-14-29||

\fourlineindentedshloka
{हिमहतनलिनीव नष्टशोभा}
{ व्यसनपरम्परया निपीड्यमाना}
{सहचररहितेव चक्रवाकी}
{ जनकसुता कृपणां दशां प्रपन्ना} %||5-14-30||

\fourlineindentedshloka
{अस्या हि पुष्पावनताग्रशाखाः}
{ शोकं दृढं वै जनयत्यशोकाः}
{हिमव्यपायेन च मन्दरश्मि}
{रभ्युत्थितो नैकसहस्ररश्मिः} %||5-14-31||

\fourlineindentedshloka
{इत्येवमर्थं कपिरन्ववेक्ष्य}
{ सीतेयमित्येव निविष्टबुद्धिः}
{संश्रित्य तस्मिन्निषसाद वृक्षे}
{ बली हरीणामृषभस्तरस्वी} %||5-15-32||


{॥इत्यार्षे श्रीमद्रामायणे वाल्मीकीये आदिकाव्ये सुन्दरकाण्डे चतुर्दशः सर्गः॥}

\sect{पञ्चदशः सर्गः}

\twolineshloka
{ततः कुमुदषण्डाभो निर्मलं निर्मलः स्वयम्}
{प्रजगाम नभश्चन्द्रो हंसो नीलमिवोदकम्} %||5-15-1||

\twolineshloka
{साचिव्यमिव कुर्वन्स प्रभया निर्मलप्रभः}
{चन्द्रमा रश्मिभिः शीतैः सिषेवे पवनात्मजम्} %||5-15-2||

\twolineshloka
{स ददर्श ततः सीतां पूर्णचन्द्रनिभाननाम्}
{शोकभारैरिव न्यस्तां भारैर्नावमिवाम्भसि} %||5-15-3||

\twolineshloka
{दिदृक्षमाणो वैदेहीं हनूमान्मारुतात्मजः}
{स ददर्शाविदूरस्था राक्षसीर्घोरदर्शनाः} %||5-15-4||

\twolineshloka
{एकाक्षीमेककर्णां च कर्णप्रावरणां तथा}
{अकर्णां शङ्कुकर्णां च मस्तकोच्छ्वासनासिकाम्} %||5-15-5||

\twolineshloka
{अतिकायोत्तमाङ्गीं च तनुदीर्घशिरोधराम्}
{ध्वस्तकेशीं तथाकेशीं केशकम्बलधारिणीम्} %||5-15-6||

\twolineshloka
{लम्बकर्णललाटां च लम्बोदरपयोधराम्}
{लम्बौष्ठीं चिबुकौष्ठीं च लम्बास्यां लम्बजानुकाम्} %||5-15-7||

\twolineshloka
{ह्रस्वां दीर्घां च कुब्जां च विकटां वामनां तथा}
{करालां भुग्नवस्त्रां च पिङ्गाक्षीं विकृताननाम्} %||5-15-8||

\twolineshloka
{विकृताः पिङ्गलाः कालीः क्रोधनाः कलहप्रियाः}
{कालायसमहाशूलकूटमुद्गरधारिणीः} %||5-15-9||

\twolineshloka
{वराहमृगशार्दूलमहिषाजशिवा मुखाः}
{गजोष्ट्रहयपादाश्च निखातशिरसोऽपराः} %||5-15-10||

\twolineshloka
{एकहस्तैकपादाश्च खरकर्ण्यश्वकर्णिकाः}
{गोकर्णीर्हस्तिकर्णीश्च हरिकर्णीस्तथापराः} %||5-15-11||

\twolineshloka
{अनासा अतिनासाश्च तिर्यन्नासा विनासिकाः}
{गजसंनिभनासाश्च ललाटोच्छ्वासनासिकाः} %||5-15-12||

\twolineshloka
{हस्तिपादा महापादा गोपादाः पादचूलिकाः}
{अतिमात्रशिरोग्रीवा अतिमात्रकुचोदरीः} %||5-15-13||

\twolineshloka
{अतिमात्रास्य नेत्राश्च दीर्घजिह्वानखास्तथा}
{अजामुखीर्हस्तिमुखीर्गोमुखीः सूकरीमुखीः} %||5-15-14||

\twolineshloka
{हयोष्ट्रखरवक्त्राश्च राक्षसीर्घोरदर्शनाः}
{शूलमुद्गरहस्ताश्च क्रोधनाः कलहप्रियाः} %||5-15-15||

\twolineshloka
{कराला धूम्रकेशीश्च रक्षसीर्विकृताननाः}
{पिबन्तीः सततं पानं सदा मांससुराप्रियाः} %||5-15-16||

\twolineshloka
{मांसशोणितदिग्धाङ्गीर्मांसशोणितभोजनाः}
{ता ददर्श कपिश्रेष्ठो रोमहर्षणदर्शनाः} %||5-15-17||

\twolineshloka
{स्कन्धवन्तमुपासीनाः परिवार्य वनस्पतिम्}
{तस्याधस्ताच्च तां देवीं राजपुत्रीमनिन्दिताम्} %||5-15-18||

\twolineshloka
{लक्षयामास लक्ष्मीवान्हनूमाञ्जनकात्मजाम्}
{निष्प्रभां शोकसंतप्तां मलसंकुलमूर्धजाम्} %||5-15-19||

\twolineshloka
{क्षीणपुण्यां च्युतां भूमौ तारां निपतितामिव}
{चारित्र्य व्यपदेशाढ्यां भर्तृदर्शनदुर्गताम्} %||5-15-20||

\twolineshloka
{भूषणैरुत्तमैर्हीनां भर्तृवात्सल्यभूषिताम्}
{राक्षसाधिपसंरुद्धां बन्धुभिश्च विनाकृताम्} %||5-15-21||

\twolineshloka
{वियूथां सिंहसंरुद्धां बद्धां गजवधूमिव}
{चन्द्रलेखां पयोदान्ते शारदाभ्रैरिवावृताम्} %||5-15-22||

\twolineshloka
{क्लिष्टरूपामसंस्पर्शादयुक्तामिव वल्लकीम्}
{सीतां भर्तृहिते युक्तामयुक्तां रक्षसां वशे} %||5-15-23||

\twolineshloka
{अशोकवनिकामध्ये शोकसागरमाप्लुताम्}
{ताभिः परिवृतां तत्र सग्रहामिव रोहिणीम्} %||5-15-24||

\twolineshloka
{ददर्श हनुमान्देवीं लतामकुसुमामिव}
{सा मलेन च दिग्धाङ्गी वपुषा चाप्यलंकृता} %||5-15-25||

\twolineshloka
{मृणाली पङ्कदिघ्देव विभाति च न भाति च}
{मलिनेन तु वस्त्रेण परिक्लिष्टेन भामिनीम्} %||5-15-26||

\twolineshloka
{संवृतां मृगशावाक्षीं ददर्श हनुमान्कपिः}
{तां देवीं दीनवदनामदीनां भर्तृतेजसा} %||5-15-27||

\twolineshloka
{रक्षितां स्वेन शीलेन सीतामसितलोचनाम्}
{तां दृष्ट्वा हनुमान्सीतां मृगशावनिभेक्षणाम्} %||5-15-28||

\twolineshloka
{मृगकन्यामिव त्रस्तां वीक्षमाणां समन्ततः}
{दहन्तीमिव निःश्वासैर्वृक्षान्पल्लवधारिणः} %||5-15-29||

\twolineshloka
{संघातमिव शोकानां दुःखस्योर्मिमिवोत्थिताम्}
{तां क्षामां सुविभक्ताङ्गीं विनाभरणशोभिनीम्} %||5-15-30||

\twolineshloka
{प्रहर्षमतुलं लेभे मारुतिः प्रेक्ष्य मैथिलीम्}
{हर्षजानि च सोऽश्रूणि तां दृष्ट्वा मदिरेक्षणाम्} %||5-15-31||

\twolineshloka
{मुमोच हनुमांस्तत्र नमश्चक्रे च राघवम्}
{नमस्कृत्वा च रामाय लक्ष्मणाय च वीर्यवान्} %||5-15-32||

{सीतादर्शनसंहृष्टो हनूमान्संवृतोऽभवत्॥\devanumber{33}॥} %||5-16-33|| (Check)

\addtocounter{shlokacount}{1}


{॥इत्यार्षे श्रीमद्रामायणे वाल्मीकीये आदिकाव्ये सुन्दरकाण्डे पञ्चदशः सर्गः॥}

\sect{षोडशः सर्गः}

\twolineshloka
{तथा विप्रेक्षमाणस्य वनं पुष्पितपादपम्}
{विचिन्वतश्च वैदेहीं किंचिच्छेषा निशाभवत्} %||5-16-1||

\twolineshloka
{षडङ्गवेदविदुषां क्रतुप्रवरयाजिनाम्}
{शुश्राव ब्रह्मघोषांश्च विरात्रे ब्रह्मरक्षसाम्} %||5-16-2||

\twolineshloka
{अथ मङ्गलवादित्रैः शब्दैः श्रोत्रमनोहरैः}
{प्राबोध्यत महाबाहुर्दशग्रीवो महाबलः} %||5-16-3||

\twolineshloka
{विबुध्य तु यथाकालं राक्षसेन्द्रः प्रतावपान्}
{स्रस्तमाल्याम्बरधरो वैदेहीमन्वचिन्तयत्} %||5-16-4||

\twolineshloka
{भृशं नियुक्तस्तस्यां च मदनेन मदोत्कटः}
{न स तं राक्षसः कामं शशाकात्मनि गूहितुम्} %||5-16-5||

\twolineshloka
{स सर्वाभरणैर्युक्तो बिभ्रच्छ्रियमनुत्तमाम्}
{तां नगैर्विविधैर्जुष्टां सर्वपुष्पफलोपगैः} %||5-16-6||

\twolineshloka
{वृतां पुष्करिणीभिश्च नानापुष्पोपशोभिताम्}
{सदामदैश्च विहगैर्विचित्रां परमाद्भुताम्} %||5-16-7||

\twolineshloka
{ईहामृगैश्च विविधैश्वृतां दृष्टिमनोहरैः}
{वीथीः संप्रेक्षमाणश्च मणिकाञ्चनतोरणाः} %||5-16-8||

\twolineshloka
{नानामृगगणाकीर्णां फलैः प्रपतितैर्वृताम्}
{अशोकवनिकामेव प्राविशत्संततद्रुमाम्} %||5-16-9||

\twolineshloka
{अङ्गनाशतमात्रं तु तं व्रजन्तमनुव्रजत्}
{महेन्द्रमिव पौलस्त्यं देवगन्धर्वयोषितः} %||5-16-10||

\twolineshloka
{दीपिकाः काञ्चनीः काश्चिज्जगृहुस्तत्र योषितः}
{बालव्यजनहस्ताश्च तालवृन्तानि चापराः} %||5-16-11||

\twolineshloka
{काञ्चनैरपि भृङ्गारैर्जह्रुः सलिलमग्रतः}
{मण्डलाग्रानसींश्चैव गृह्यान्याः पृष्ठतो ययुः} %||5-16-12||

\twolineshloka
{काचिद्रत्नमयीं पात्रीं पूर्णां पानस्य भामिनी}
{दक्षिणा दक्षिणेनैव तदा जग्राह पाणिना} %||5-16-13||

\twolineshloka
{राजहंसप्रतीकाशं छत्रं पूर्णशशिप्रभम्}
{सौवर्णदण्डमपरा गृहीत्वा पृष्ठतो ययौ} %||5-16-14||

\twolineshloka
{निद्रामदपरीताक्ष्यो रावणस्योत्तमस्त्रियः}
{अनुजग्मुः पतिं वीरं घनं विद्युल्लता इव} %||5-16-15||

\twolineshloka
{ततः काञ्चीनिनादं च नूपुराणां च निःस्वनम्}
{शुश्राव परमस्त्रीणां स कपिर्मारुतात्मजः} %||5-16-16||

\twolineshloka
{तं चाप्रतिमकर्माणमचिन्त्यबलपौरुषम्}
{द्वारदेशमनुप्राप्तं ददर्श हनुमान्कपिः} %||5-16-17||

\twolineshloka
{दीपिकाभिरनेकाभिः समन्तादवभासितम्}
{गन्धतैलावसिक्ताभिर्ध्रियमाणाभिरग्रतः} %||5-16-18||

\twolineshloka
{कामदर्पमदैर्युक्तं जिह्मताम्रायतेक्षणम्}
{समक्षमिव कन्दर्पमपविद्ध शरासनम्} %||5-16-19||

\twolineshloka
{मथितामृतफेनाभमरजो वस्त्रमुत्तमम्}
{सलीलमनुकर्षन्तं विमुक्तं सक्तमङ्गदे} %||5-16-20||

\twolineshloka
{तं पत्रविटपे लीनः पत्रपुष्पघनावृतः}
{समीपमुपसंक्रान्तं निध्यातुमुपचक्रमे} %||5-16-21||

\twolineshloka
{अवेक्षमाणश्च ततो ददर्श कपिकुञ्जरः}
{रूपयौवनसंपन्ना रावणस्य वरस्त्रियः} %||5-16-22||

\twolineshloka
{ताभिः परिवृतो राजा सुरूपाभिर्महायशाः}
{तन्मृगद्विजसंघुष्टं प्रविष्टः प्रमदावनम्} %||5-16-23||

\twolineshloka
{क्षीबो विचित्राभरणः शङ्कुकर्णो महाबलः}
{तेन विश्रवसः पुत्रः स दृष्टो राक्षसाधिपः} %||5-16-24||

\twolineshloka
{वृतः परमनारीभिस्ताराभिरिव चन्द्रमाः}
{तं ददर्श महातेजास्तेजोवन्तं महाकपिः} %||5-16-25||

\twolineshloka
{रावणोऽयं महाबाहुरिति संचिन्त्य वानरः}
{अवप्लुतो महातेजा हनूमान्मारुतात्मजः} %||5-16-26||

\twolineshloka
{स तथाप्युग्रतेजाः सन्निर्धूतस्तस्य तेजसा}
{पत्रगुह्यान्तरे सक्तो हनूमान्संवृतोऽभवत्} %||5-16-27||

\twolineshloka
{स तामसितकेशान्तां सुश्रोणीं संहतस्तनीम्}
{दिदृक्षुरसितापाङ्गीमुपावर्तत रावणः} %||5-17-28||


{॥इत्यार्षे श्रीमद्रामायणे वाल्मीकीये आदिकाव्ये सुन्दरकाण्डे षोडशः सर्गः॥}

\sect{सप्तदशः सर्गः}

\twolineshloka
{तस्मिन्नेव ततः काले राजपुत्री त्वनिन्दिता}
{रूपयौवनसंपन्नं भूषणोत्तमभूषितम्} %||5-17-1||

\twolineshloka
{ततो दृष्ट्वैव वैदेही रावणं राक्षसाधिपम्}
{प्रावेपत वरारोहा प्रवाते कदली यथा} %||5-17-2||

\twolineshloka
{ऊरुभ्यामुदरं छाद्य बाहुभ्यां च पयोधरौ}
{उपविष्टा विशालाक्षी रुदन्ती वरवर्णिनी} %||5-17-3||

\twolineshloka
{दशग्रीवस्तु वैदेहीं रक्षितां राक्षसीगणैः}
{ददर्श दीनां दुःखार्तं नावं सन्नामिवार्णवे} %||5-17-4||

\threelineshloka
{असंवृतायामासीनां धरण्यां संशितव्रताम्}
{छिन्नां प्रपतितां भूमौ शाखामिव वनस्पतेः}
{मलमण्डनदिग्धाङ्गीं मण्डनार्हाममण्डिताम्} %||5-17-5||

\twolineshloka
{समीपं राजसिंहस्य रामस्य विदितात्मनः}
{संकल्पहयसंयुक्तैर्यान्तीमिव मनोरथैः} %||5-17-6||

\twolineshloka
{शुष्यन्तीं रुदतीमेकां ध्यानशोकपरायणाम्}
{दुःखस्यान्तमपश्यन्तीं रामां राममनुव्रताम्} %||5-17-7||

\twolineshloka
{वेष्टमानामथाविष्टां पन्नगेन्द्रवधूमिव}
{धूप्यमानां ग्रहेणेव रोहिणीं धूमकेतुना} %||5-17-8||

\twolineshloka
{वृत्तशीले कुले जातामाचारवति धार्मिके}
{पुनः संस्कारमापन्नां जातमिव च दुष्कुले} %||5-17-9||

\twolineshloka
{सन्नामिव महाकीर्तिं श्रद्धामिव विमानिताम्}
{प्रज्ञामिव परिक्षीणामाशां प्रतिहतामिव} %||5-17-10||

\twolineshloka
{आयतीमिव विध्वस्तामाज्ञां प्रतिहतामिव}
{दीप्तामिव दिशं काले पूजामपहृतामिव} %||5-17-11||

\twolineshloka
{पद्मिनीमिव विध्वस्तां हतशूरां चमूमिव}
{प्रभामिव तपोध्वस्तामुपक्षीणामिवापगाम्} %||5-17-12||

\twolineshloka
{वेदीमिव परामृष्टां शान्तामग्निशिखामिव}
{पौर्णमासीमिव निशां राहुग्रस्तेन्दुमण्डलाम्} %||5-17-13||

\twolineshloka
{उत्कृष्टपर्णकमलां वित्रासितविहंगमाम्}
{हस्तिहस्तपरामृष्टामाकुलां पद्मिनीमिव} %||5-17-14||

\twolineshloka
{पतिशोकातुरां शुष्कां नदीं विस्रावितामिव}
{परया मृजया हीनां कृष्णपक्षे निशामिव} %||5-17-15||

\twolineshloka
{सुकुमारीं सुजाताङ्गीं रत्नगर्भगृहोचिताम्}
{तप्यमानामिवोष्णेन मृणालीमचिरोद्धृताम्} %||5-17-16||

\twolineshloka
{गृहीतामालितां स्तम्भे यूथपेन विनाकृताम्}
{निःश्वसन्तीं सुदुःखार्तां गजराजवधूमिव} %||5-17-17||

\twolineshloka
{एकया दीर्घया वेण्या शोभमानामयत्नतः}
{नीलया नीरदापाये वनराज्या महीमिव} %||5-17-18||

\twolineshloka
{उपवासेन शोकेन ध्यानेन च भयेन च}
{परिक्षीणां कृशां दीनामल्पाहारां तपोधनाम्} %||5-17-19||

\twolineshloka
{आयाचमानां दुःखार्तां प्राञ्जलिं देवतामिव}
{भावेन रघुमुख्यस्य दशग्रीवपराभवम्} %||5-17-20||

\fourlineindentedshloka
{समीक्षमाणां रुदतीमनिन्दितां}
{ सुपक्ष्मताम्रायतशुक्ललोचनाम्}
{अनुव्रतां राममतीव मैथिलीं}
{ प्रलोभयामास वधाय रावणः} %||5-18-21||


{॥इत्यार्षे श्रीमद्रामायणे वाल्मीकीये आदिकाव्ये सुन्दरकाण्डे सप्तदशः सर्गः॥}

\sect{अष्टादशः सर्गः}

\twolineshloka
{स तां परिवृतां दीनां निरानन्दां तपस्विनीम्}
{साकारैर्मधुरैर्वाक्यैर्न्यदर्शयत रावणः} %||5-18-1||

\twolineshloka
{मां दृष्ट्वा नागनासोरुगूहमाना स्तनोदरम्}
{अदर्शनमिवात्मानं भयान्नेतुं त्वमिच्छसि} %||5-18-2||

\twolineshloka
{कामये त्वां विशालाक्षि बहुमन्यस्व मां प्रिये}
{सर्वाङ्गगुणसंपन्ने सर्वलोकमनोहरे} %||5-18-3||

\twolineshloka
{नेह केचिन्मनुष्या वा राक्षसाः कामरूपिणः}
{व्यपसर्पतु ते सीते भयं मत्तः समुत्थितम्} %||5-18-4||

\twolineshloka
{स्वधर्मे रक्षसां भीरु सर्वथैष न संशयः}
{गमनं वा परस्त्रीणां हरणं संप्रमथ्य वा} %||5-18-5||

\twolineshloka
{एवं चैतदकामां च न त्वां स्प्रक्ष्यामि मैथिलि}
{कामं कामः शरीरे मे यथाकामं प्रवर्तताम्} %||5-18-6||

\twolineshloka
{देवि नेह भयं कार्यं मयि विश्वसिहि प्रिये}
{प्रणयस्व च तत्त्वेन मैवं भूः शोकलालसा} %||5-18-7||

\twolineshloka
{एकवेणी धराशय्या ध्यानं मलिनमम्बरम्}
{अस्थानेऽप्युपवासश्च नैतान्यौपयिकानि ते} %||5-18-8||

\twolineshloka
{विचित्राणि च माल्यानि चन्दनान्यगरूणि च}
{विविधानि च वासांसि दिव्यान्याभरणानि च} %||5-18-9||

\twolineshloka
{महार्हाणि च पानानि यानानि शयनानि च}
{गीतं नृत्तं च वाद्यं च लभ मां प्राप्य मैथिलि} %||5-18-10||

\twolineshloka
{स्त्रीरत्नमसि मैवं भूः कुरु गात्रेषु भूषणम्}
{मां प्राप्य तु कथं हि स्यास्त्वमनर्हा सुविग्रहे} %||5-18-11||

\twolineshloka
{इदं ते चारुसंजातं यौवनं व्यतिवर्तते}
{यदतीतं पुनर्नैति स्रोतः शीघ्रमपामिव} %||5-18-12||

\twolineshloka
{त्वां कृत्वोपरतो मन्ये रूपकर्ता स विश्वकृत्}
{न हि रूपोपमा त्वन्या तवास्ति शुभदर्शने} %||5-18-13||

\twolineshloka
{त्वां समासाद्य वैदेहि रूपयौवनशालिनीम्}
{कः पुमानतिवर्तेत साक्षादपि पितामहः} %||5-18-14||

\twolineshloka
{यद्यत्पश्यामि ते गात्रं शीतांशुसदृशानने}
{तस्मिंस्तस्मिन्पृथुश्रोणि चक्षुर्मम निबध्यते} %||5-18-15||

\twolineshloka
{भव मैथिलि भार्या मे मोहमेनं विसर्जय}
{बह्वीनामुत्तमस्त्रीणां ममाग्रमहिषी भव} %||5-18-16||

\twolineshloka
{लोकेभ्यो यानि रत्नानि संप्रमथ्याहृतानि मे}
{तानि ते भीरु सर्वाणि राज्यं चैतदहं च ते} %||5-18-17||

\twolineshloka
{विजित्य पृथिवीं सर्वां नानानगरमालिनीम्}
{जनकाय प्रदास्यामि तव हेतोर्विलासिनि} %||5-18-18||

\twolineshloka
{नेह पश्यामि लोकेऽन्यं यो मे प्रतिबलो भवेत्}
{पश्य मे सुमहद्वीर्यमप्रतिद्वन्द्वमाहवे} %||5-18-19||

\twolineshloka
{असकृत्संयुगे भग्ना मया विमृदितध्वजाः}
{अशक्ताः प्रत्यनीकेषु स्थातुं मम सुरासुराः} %||5-18-20||

\threelineshloka
{इच्छ मां क्रियतामद्य प्रतिकर्म तवोत्तमम्}
{सप्रभाण्यवसज्जन्तां तवाङ्गे भूषणानि च}
{साधु पश्यामि ते रूपं संयुक्तं प्रतिकर्मणा} %||5-18-21||

\threelineshloka
{प्रतिकर्माभिसंयुक्ता दाक्षिण्येन वरानने}
{भुङ्क्ष्व भोगान्यथाकामं पिब भीरु रमस्व च}
{यथेष्टं च प्रयच्छ त्वं पृथिवीं वा धनानि च} %||5-18-22||

\twolineshloka
{ललस्व मयि विस्रब्धा धृष्टमाज्ञापयस्व च}
{मत्प्रभावाल्ललन्त्याश्च ललन्तां बान्धवास्तव} %||5-18-23||

\twolineshloka
{ऋद्धिं ममानुपश्य त्वं श्रियं भद्रे यशश्च मे}
{किं करिष्यसि रामेण सुभगे चीरवाससा} %||5-18-24||

\twolineshloka
{निक्षिप्तविजयो रामो गतश्रीर्वनगोचरः}
{व्रती स्थण्डिलशायी च शङ्के जीवति वा न वा} %||5-18-25||

\twolineshloka
{न हि वैदेहि रामस्त्वां द्रष्टुं वाप्युपलप्स्यते}
{पुरो बलाकैरसितैर्मेघैर्ज्योत्स्नामिवावृताम्} %||5-18-26||

\twolineshloka
{न चापि मम हस्तात्त्वां प्राप्तुमर्हति राघवः}
{हिरण्यकशिपुः कीर्तिमिन्द्रहस्तगतामिव} %||5-18-27||

\twolineshloka
{चारुस्मिते चारुदति चारुनेत्रे विलासिनि}
{मनो हरसि मे भीरु सुपर्णः पन्नगं यथा} %||5-18-28||

\twolineshloka
{क्लिष्टकौशेयवसनां तन्वीमप्यनलंकृताम्}
{तां दृष्ट्वा स्वेषु दारेषु रतिं नोपलभाम्यहम्} %||5-18-29||

\twolineshloka
{अन्तःपुरनिवासिन्यः स्त्रियः सर्वगुणान्विताः}
{यावन्त्यो मम सर्वासामैश्वर्यं कुरु जानकि} %||5-18-30||

\twolineshloka
{मम ह्यसितकेशान्ते त्रैलोक्यप्रवराः स्त्रियः}
{तास्त्वां परिचरिष्यन्ति श्रियमप्सरसो यथा} %||5-18-31||

\twolineshloka
{यानि वैश्रवणे सुभ्रु रत्नानि च धनानि च}
{तानि लोकांश्च सुश्रोणि मां च भुङ्क्ष्व यथासुखम्} %||5-18-32||

\twolineshloka
{न रामस्तपसा देवि न बलेन न विक्रमैः}
{न धनेन मया तुल्यस्तेजसा यशसापि वा} %||5-18-33||

\fourlineindentedshloka
{पिब विहर रमस्व भुङ्क्ष्व भोगा}
{न्धननिचयं प्रदिशामि मेदिनीं च}
{मयि लल ललने यथासुखं त्वं}
{ त्वयि च समेत्य ललन्तु बान्धवास्ते} %||5-18-34||

\fourlineindentedshloka
{कुसुमिततरुजालसंततानि}
{ भ्रमरयुतानि समुद्रतीरजानि}
{कनकविमलहारभूषिताङ्गी}
{ विहर मया सह भीरु काननानि} %||5-19-35||


{॥इत्यार्षे श्रीमद्रामायणे वाल्मीकीये आदिकाव्ये सुन्दरकाण्डे अष्टादशः सर्गः॥}

\sect{एकोनविंशः सर्गः}

\twolineshloka
{तस्य तद्वचनं श्रुत्वा सीता रौद्रस्य रक्षसः}
{आर्ता दीनस्वरा दीनं प्रत्युवाच शनैर्वचः} %||5-19-1||

\twolineshloka
{दुःखार्ता रुदती सीता वेपमाना तपस्विनी}
{चिन्तयन्ती वरारोहा पतिमेव पतिव्रता} %||5-19-2||

\twolineshloka
{तृणमन्तरतः कृत्वा प्रत्युवाच शुचिस्मिता}
{निवर्तय मनो मत्तः स्वजने क्रियतां मनः} %||5-19-3||

\threelineshloka
{न मां प्रार्थयितुं युक्तस्त्वं सिद्धिमिव पापकृत्}
{अकार्यं न मया कार्यमेकपत्न्या विगर्हितम्}
{कुलं संप्राप्तया पुण्यं कुले महति जातया} %||5-19-4||

\twolineshloka
{एवमुक्त्वा तु वैदेही रावणं तं यशस्विनी}
{राक्षसं पृष्ठतः कृत्वा भूयो वचनमब्रवीत्} %||5-19-5||

\twolineshloka
{नाहमौपयिकी भार्या परभार्या सती तव}
{साधु धर्ममवेक्षस्व साधु साधुव्रतं चर} %||5-19-6||

\twolineshloka
{यथा तव तथान्येषां रक्ष्या दारा निशाचर}
{आत्मानमुपमां कृत्वा स्वेषु दारेषु रम्यताम्} %||5-19-7||

\twolineshloka
{अतुष्टं स्वेषु दारेषु चपलं चलितेन्द्रियम्}
{नयन्ति निकृतिप्रज्ञां परदाराः पराभवम्} %||5-19-8||

\twolineshloka
{इह सन्तो न वा सन्ति सतो वा नानुवर्तसे}
{वचो मिथ्या प्रणीतात्मा पथ्यमुक्तं विचक्षणैः} %||5-19-9||

\twolineshloka
{अकृतात्मानमासाद्य राजानमनये रतम्}
{समृद्धानि विनश्यन्ति राष्ट्राणि नगराणि च} %||5-19-10||

\twolineshloka
{तथेयं त्वां समासाद्य लङ्का रत्नौघ संकुला}
{अपराधात्तवैकस्य नचिराद्विनशिष्यति} %||5-19-11||

\twolineshloka
{स्वकृतैर्हन्यमानस्य रावणादीर्घदर्शिनः}
{अभिनन्दन्ति भूतानि विनाशे पापकर्मणः} %||5-19-12||

\twolineshloka
{एवं त्वां पापकर्माणं वक्ष्यन्ति निकृता जनाः}
{दिष्ट्यैतद्व्यसनं प्राप्तो रौद्र इत्येव हर्षिताः} %||5-19-13||

\twolineshloka
{शक्या लोभयितुं नाहमैश्वर्येण धनेन वा}
{अनन्या राघवेणाहं भास्करेण प्रभा यथा} %||5-19-14||

\twolineshloka
{उपधाय भुजं तस्य लोकनाथस्य सत्कृतम्}
{कथं नामोपधास्यामि भुजमन्यस्य कस्यचित्} %||5-19-15||

\twolineshloka
{अहमौपयिकी भार्या तस्यैव वसुधापतेः}
{व्रतस्नातस्य विप्रस्य विद्येव विदितात्मनः} %||5-19-16||

\twolineshloka
{साधु रावण रामेण मां समानय दुःखिताम्}
{वने वाशितया सार्धं करेण्वेव गजाधिपम्} %||5-19-17||

\twolineshloka
{मित्रमौपयिकं कर्तुं रामः स्थानं परीप्सता}
{वधं चानिच्छता घोरं त्वयासौ पुरुषर्षभः} %||5-19-18||

\twolineshloka
{वर्जयेद्वज्रमुत्सृष्टं वर्जयेदन्तकश्चिरम्}
{त्वद्विधं न तु संक्रुद्धो लोकनाथः स राघवः} %||5-19-19||

\twolineshloka
{रामस्य धनुषः शब्दं श्रोष्यसि त्वं महास्वनम्}
{शतक्रतुविसृष्टस्य निर्घोषमशनेरिव} %||5-19-20||

\twolineshloka
{इह शीघ्रं सुपर्वाणो ज्वलितास्या इवोरगाः}
{इषवो निपतिष्यन्ति रामलक्ष्मणलक्षणाः} %||5-19-21||

\twolineshloka
{रक्षांसि परिनिघ्नन्तः पुर्यामस्यां समन्ततः}
{असंपातं करिष्यन्ति पतन्तः कङ्कवाससः} %||5-19-22||

\twolineshloka
{राक्षसेन्द्रमहासर्पान्स रामगरुडो महान्}
{उद्धरिष्यति वेगेन वैनतेय इवोरगान्} %||5-19-23||

\twolineshloka
{अपनेष्यति मां भर्ता त्वत्तः शीघ्रमरिंदमः}
{असुरेभ्यः श्रियं दीप्तां विष्णुस्त्रिभिरिव क्रमैः} %||5-19-24||

\twolineshloka
{जनस्थाने हतस्थाने निहते रक्षसां बले}
{अशक्तेन त्वया रक्षः कृतमेतदसाधु वै} %||5-19-25||

\twolineshloka
{आश्रमं तु तयोः शून्यं प्रविश्य नरसिंहयोः}
{गोचरं गतयोर्भ्रात्रोरपनीता त्वयाधम} %||5-19-26||

\twolineshloka
{न हि गन्धमुपाघ्राय रामलक्ष्मणयोस्त्वया}
{शक्यं संदर्शने स्थातुं शुना शार्दूलयोरिव} %||5-19-27||

\twolineshloka
{तस्य ते विग्रहे ताभ्यां युगग्रहणमस्थिरम्}
{वृत्रस्येवेन्द्रबाहुभ्यां बाहोरेकस्य निग्रहः} %||5-19-28||

\twolineshloka
{क्षिप्रं तव स नाथो मे रामः सौमित्रिणा सह}
{तोयमल्पमिवादित्यः प्राणानादास्यते शरैः} %||5-19-29||

\fourlineindentedshloka
{गिरिं कुबेरस्य गतोऽथवालयं}
{ सभां गतो वा वरुणस्य राज्ञः}
{असंशयं दाशरथेर्न मोक्ष्यसे}
{ महाद्रुमः कालहतोऽशनेरिव} %||5-20-30||


{॥इत्यार्षे श्रीमद्रामायणे वाल्मीकीये आदिकाव्ये सुन्दरकाण्डे एकोनविंशः सर्गः॥}

\sect{विंशः सर्गः}

\twolineshloka
{सीताया वचनं श्रुत्वा परुषं राक्षसाधिपः}
{प्रत्युवाच ततः सीतां विप्रियं प्रियदर्शनाम्} %||5-20-1||

\twolineshloka
{यथा यथा सान्त्वयिता वश्यः स्त्रीणां तथा तथा}
{यथा यथा प्रियं वक्ता परिभूतस्तथा तथा} %||5-20-2||

\twolineshloka
{संनियच्छति मे क्रोधं त्वयि कामः समुत्थितः}
{द्रवतो मार्गमासाद्य हयानिव सुसारथिः} %||5-20-3||

\twolineshloka
{वामः कामो मनुष्याणां यस्मिन्किल निबध्यते}
{जने तस्मिंस्त्वनुक्रोशः स्नेहश्च किल जायते} %||5-20-4||

\twolineshloka
{एतस्मात्कारणान्न तां घतयामि वरानने}
{वधार्हामवमानार्हां मिथ्याप्रव्रजिते रताम्} %||5-20-5||

\twolineshloka
{परुषाणि हि वाक्यानि यानि यानि ब्रवीषि माम्}
{तेषु तेषु वधो युक्तस्तव मैथिलि दारुणः} %||5-20-6||

\twolineshloka
{एवमुक्त्वा तु वैदेहीं रावणो राक्षसाधिपः}
{क्रोधसंरम्भसंयुक्तः सीतामुत्तरमब्रवीत्} %||5-20-7||

\twolineshloka
{द्वौ मासौ रक्षितव्यौ मे योऽवधिस्ते मया कृतः}
{ततः शयनमारोह मम त्वं वरवर्णिनि} %||5-20-8||

\twolineshloka
{द्वाभ्यामूर्ध्वं तु मासाभ्यां भर्तारं मामनिच्छतीम्}
{मम त्वां प्रातराशार्थमारभन्ते महानसे} %||5-20-9||

\twolineshloka
{तां तर्ज्यमानां संप्रेक्ष्य राक्षसेन्द्रेण जानकीम्}
{देवगन्धर्वकन्यास्ता विषेदुर्विपुलेक्षणाः} %||5-20-10||

\twolineshloka
{ओष्ठप्रकारैरपरा नेत्रवक्त्रैस्तथापराः}
{सीतामाश्वासयामासुस्तर्जितां तेन रक्षसा} %||5-20-11||

\twolineshloka
{ताभिराश्वासिता सीता रावणं राक्षसाधिपम्}
{उवाचात्महितं वाक्यं वृत्तशौण्डीर्यगर्वितम्} %||5-20-12||

\twolineshloka
{नूनं न ते जनः कश्चिदसिन्निःश्रेयसे स्थितः}
{निवारयति यो न त्वां कर्मणोऽस्माद्विगर्हितात्} %||5-20-13||

\twolineshloka
{मां हि धर्मात्मनः पत्नीं शचीमिव शचीपतेः}
{त्वदन्यस्त्रिषु लोकेषु प्रार्थयेन्मनसापि कः} %||5-20-14||

\twolineshloka
{राक्षसाधम रामस्य भार्याममिततेजसः}
{उक्तवानसि यत्पापं क्व गतस्तस्य मोक्ष्यसे} %||5-20-15||

\twolineshloka
{यथा दृप्तश्च मातङ्गः शशश्च सहितौ वने}
{तथा द्विरदवद्रामस्त्वं नीच शशवत्स्मृतः} %||5-20-16||

\twolineshloka
{स त्वमिक्ष्वाकुनाथं वै क्षिपन्निह न लज्जसे}
{चक्षुषो विषयं तस्य न तावदुपगच्छसि} %||5-20-17||

\twolineshloka
{इमे ते नयने क्रूरे विरूपे कृष्णपिङ्गले}
{क्षितौ न पतिते कस्मान्मामनार्यनिरीक्षितः} %||5-20-18||

\twolineshloka
{तस्य धर्मात्मनः पत्नीं स्नुषां दशरथस्य च}
{कथं व्याहरतो मां ते न जिह्वा पाप शीर्यते} %||5-20-19||

\twolineshloka
{असंदेशात्तु रामस्य तपसश्चानुपालनात्}
{न त्वां कुर्मि दशग्रीव भस्म भस्मार्हतेजसा} %||5-20-20||

\twolineshloka
{नापहर्तुमहं शक्या तस्य रामस्य धीमतः}
{विधिस्तव वधार्थाय विहितो नात्र संशयः} %||5-20-21||

\twolineshloka
{शूरेण धनदभ्राता बलैः समुदितेन च}
{अपोह्य रामं कस्माद्धि दारचौर्यं त्वया कृतम्} %||5-20-22||

\twolineshloka
{सीताया वचनं श्रुत्वा रावणो राक्षसाधिपः}
{विवृत्य नयने क्रूरे जानकीमन्ववैक्षत} %||5-20-23||

\twolineshloka
{नीलजीमूतसंकाशो महाभुजशिरोधरः}
{सिंहसत्त्वगतिः श्रीमान्दीप्तजिह्वोग्रलोचनः} %||5-20-24||

\twolineshloka
{चलाग्रमकुटः प्रांशुश्चित्रमाल्यानुलेपनः}
{रक्तमाल्याम्बरधरस्तप्ताङ्गदविभूषणः} %||5-20-25||

\twolineshloka
{श्रोणीसूत्रेण महता मेककेन सुसंवृतः}
{अमृतोत्पादनद्धेन भुजंगेनेव मन्दरः} %||5-20-26||

\twolineshloka
{तरुणादित्यवर्णाभ्यां कुण्डलाभ्यां विभूषितः}
{रक्तपल्लवपुष्पाभ्यामशोकाभ्यामिवाचलः} %||5-20-27||

\twolineshloka
{अवेक्षमाणो वैदेहीं कोपसंरक्तलोचनः}
{उवाच रावणः सीतां भुजंग इव निःश्वसन्} %||5-20-28||

\twolineshloka
{अनयेनाभिसंपन्नमर्थहीनमनुव्रते}
{नाशयाम्यहमद्य त्वां सूर्यः संध्यामिवौजसा} %||5-20-29||

\twolineshloka
{इत्युक्त्वा मैथिलीं राजा रावणः शत्रुरावणः}
{संदिदेश ततः सर्वा राक्षसीर्घोरदर्शनाः} %||5-20-30||

\twolineshloka
{एकाक्षीमेककर्णां च कर्णप्रावरणां तथा}
{गोकर्णीं हस्तिकर्णीं च लम्बकर्णीमकर्णिकाम्} %||5-20-31||

\twolineshloka
{हस्तिपद्य श्वपद्यौ च गोपदीं पादचूलिकाम्}
{एकाक्षीमेकपादीं च पृथुपादीमपादिकाम्} %||5-20-32||

\threelineshloka
{अतिमात्रशिरोग्रीवामतिमात्रकुचोदरीम्}
{अतिमात्रास्यनेत्रां च दीर्घजिह्वामजिह्विकाम्}
{अनासिकां सिंहमुखीं गोमुखीं सूकरीमुखीम्} %||5-20-33||

\twolineshloka
{यथा मद्वशगा सीता क्षिप्रं भवति जानकी}
{तथा कुरुत राक्षस्यः सर्वाः क्षिप्रं समेत्य च} %||5-20-34||

\twolineshloka
{प्रतिलोमानुलोमैश्च सामदानादिभेदनैः}
{आवर्तयत वैदेहीं दण्डस्योद्यमनेन च} %||5-20-35||

\twolineshloka
{इति प्रतिसमादिश्य राक्षसेन्द्रः पुनः पुनः}
{काममन्युपरीतात्मा जानकीं पर्यतर्जयत्} %||5-20-36||

\twolineshloka
{उपगम्य ततः क्षिप्रं राक्षसी धान्यमालिनी}
{परिष्वज्य दशग्रीवमिदं वचनमब्रवीत्} %||5-20-37||

\threelineshloka
{मया क्रीड महाराजसीतया किं तवानया}
{अकामां कामयानस्य शरीरमुपतप्यते}
{इच्छन्तीं कामयानस्य प्रीतिर्भवति शोभना} %||5-20-38||

\twolineshloka
{एवमुक्तस्तु राक्षस्या समुत्क्षिप्तस्ततो बली}
{ज्वलद्भास्करवर्णाभं प्रविवेश निवेशनम्} %||5-20-39||

\twolineshloka
{देवगन्धर्वकन्याश्च नागकन्याश्च तास्ततः}
{परिवार्य दशग्रीवं विविशुस्तद्गृहोत्तमम्} %||5-20-40||

\fourlineindentedshloka
{स मैथिलीं धर्मपरामवस्थितां}
{ प्रवेपमानां परिभर्त्स्य रावणः}
{विहाय सीतां मदनेन मोहितः}
{ स्वमेव वेश्म प्रविवेश भास्वरम्} %||5-21-41||


{॥इत्यार्षे श्रीमद्रामायणे वाल्मीकीये आदिकाव्ये सुन्दरकाण्डे विंशः सर्गः॥}

\sect{एकविंशः सर्गः}

\twolineshloka
{इत्युक्त्वा मैथिलीं राजा रावणः शत्रुरावणः}
{संदिश्य च ततः सर्वा राक्षसीर्निर्जगाम ह} %||5-21-1||

\twolineshloka
{निष्क्रान्ते राक्षसेन्द्रे तु पुनरन्तःपुरं गते}
{राक्षस्यो भीमरूपास्ताः सीतां समभिदुद्रुवुः} %||5-21-2||

\twolineshloka
{ततः सीतामुपागम्य राक्षस्यः क्रोधमूर्छिताः}
{परं परुषया वाचा वैदेहीमिदमब्रुवन्} %||5-21-3||

\twolineshloka
{पौलस्त्यस्य वरिष्ठस्य रावणस्य महात्मनः}
{दशग्रीवस्य भार्यात्वं सीते न बहु मन्यसे} %||5-21-4||

\twolineshloka
{ततस्त्वेकजटा नाम राक्षसी वाक्यमब्रवीत्}
{आमन्त्र्य क्रोधताम्राक्षी सीतां करतलोदरीम्} %||5-21-5||

\twolineshloka
{प्रजापतीनां षण्णां तु चतुर्थो यः प्रजापतिः}
{मानसो ब्रह्मणः पुत्रः पुलस्त्य इति विश्रुतः} %||5-21-6||

\twolineshloka
{पुलस्त्यस्य तु तेजस्वी महर्षिर्मानसः सुतः}
{नाम्ना स विश्रवा नाम प्रजापतिसमप्रभः} %||5-21-7||

\threelineshloka
{तस्य पुत्रो विशालाक्षि रावणः शत्रुरावणः}
{तस्य त्वं राक्षसेन्द्रस्य भार्या भवितुमर्हसि}
{मयोक्तं चारुसर्वाङ्गि वाक्यं किं नानुमन्यसे} %||5-21-8||

\twolineshloka
{ततो हरिजटा नाम राक्षसी वाक्यमब्रवीत्}
{विवृत्य नयने कोपान्मार्जारसदृशेक्षणा} %||5-21-9||

\twolineshloka
{येन देवास्त्रयस्त्रिंशद्देवराजश्च निर्जितः}
{तस्य त्वं राक्षसेन्द्रस्य भार्या भवितुमर्हसि} %||5-21-10||

\twolineshloka
{वीर्योत्सिक्तस्य शूरस्य संग्रामेष्वनिवर्तिनः}
{बलिनो वीर्ययुक्तस्या भार्यात्वं किं न लप्स्यसे} %||5-21-11||

\twolineshloka
{प्रियां बहुमतां भार्यां त्यक्त्वा राजा महाबलः}
{सर्वासां च महाभागां त्वामुपैष्यति रावणः} %||5-21-12||

\twolineshloka
{समृद्धं स्त्रीसहस्रेण नानारत्नोपशोभितम्}
{अन्तःपुरं समुत्सृज्य त्वामुपैष्यति रावणः} %||5-21-13||

\twolineshloka
{असकृद्देवता युद्धे नागगन्धर्वदानवाः}
{निर्जिताः समरे येन स ते पार्श्वमुपागतः} %||5-21-14||

\twolineshloka
{तस्य सर्वसमृद्धस्या रावणस्य महात्मनः}
{किमर्थं राक्षसेन्द्रस्य भार्यात्वं नेच्छसेऽधमे} %||5-21-15||

\twolineshloka
{यस्य सूर्यो न तपति भीतो यस्य च मारुतः}
{न वाति स्मायतापाङ्गे किं त्वं तस्य न तिष्ठसि} %||5-21-16||

\twolineshloka
{पुष्पवृष्टिं च तरवो मुमुचुर्यस्य वै भयात्}
{शैलाश्च सुभ्रु पानीयं जलदाश्च यदेच्छति} %||5-21-17||

\twolineshloka
{तस्य नैरृतराजस्य राजराजस्य भामिनि}
{किं त्वं न कुरुषे बुद्धिं भार्यार्थे रावणस्य हि} %||5-21-18||

\twolineshloka
{साधु ते तत्त्वतो देवि कथितं साधु भामिनि}
{गृहाण सुस्मिते वाक्यमन्यथा न भविष्यसि} %||5-22-19||


{॥इत्यार्षे श्रीमद्रामायणे वाल्मीकीये आदिकाव्ये सुन्दरकाण्डे एकविंशः सर्गः॥}

\sect{द्वाविंशः सर्गः}

\twolineshloka
{ततः सीतामुपागम्य राक्षस्यो विकृताननाः}
{परुषं परुषा नार्य ऊचुस्ता वाक्यमप्रियम्} %||5-22-1||

\twolineshloka
{किं त्वमन्तःपुरे सीते सर्वभूतमनोहरे}
{महार्हशयनोपेते न वासमनुमन्यसे} %||5-22-2||

\twolineshloka
{मानुषी मानुषस्यैव भार्यात्वं बहु मन्यसे}
{प्रत्याहर मनो रामान्न त्वं जातु भविष्यसि} %||5-22-3||

\twolineshloka
{मानुषी मानुषं तं तु राममिच्छसि शोभने}
{राज्याद्भ्रष्टमसिद्धार्थं विक्लवं तमनिन्दिते} %||5-22-4||

\twolineshloka
{राक्षसीनां वचः श्रुत्वा सीता पद्मनिभेक्षणा}
{नेत्राभ्यामश्रुपूर्णाभ्यामिदं वचनमब्रवीत्} %||5-22-5||

\twolineshloka
{यदिदं लोकविद्विष्टमुदाहरथ संगताः}
{नैतन्मनसि वाक्यं मे किल्बिषं प्रतितिष्ठति} %||5-22-6||

\threelineshloka
{न मानुषी राक्षसस्य भार्या भवितुमर्हति}
{कामं खादत मां सर्वा न करिष्यामि वो वचः}
{दीनो वा राज्यहीनो वा यो मे भर्ता स मे गुरुः} %||5-22-7||

\twolineshloka
{सीताया वचनं श्रुत्वा राक्षस्यः क्रोधमूर्छिताः}
{भर्त्सयन्ति स्म परुषैर्वाक्यै रावणचोदिताः} %||5-22-8||

\twolineshloka
{अवलीनः स निर्वाक्यो हनुमाञ्शिंशपाद्रुमे}
{सीतां संतर्जयन्तीस्ता राक्षसीरशृणोत्कपिः} %||5-22-9||

\twolineshloka
{तामभिक्रम्य संरब्धा वेपमानां समन्ततः}
{भृशं संलिलिहुर्दीप्तान्प्रलम्बदशनच्छदान्} %||5-22-10||

\twolineshloka
{ऊचुश्च परमक्रुद्धाः प्रगृह्याशु परश्वधान्}
{नेयमर्हति भर्तारं रावणं राक्षसाधिपम्} %||5-22-11||

\twolineshloka
{सा भर्त्स्यमाना भीमाभी राक्षसीभिर्वरानना}
{सा बाष्पमपमार्जन्ती शिंशपां तामुपागमत्} %||5-22-12||

\twolineshloka
{ततस्तां शिंशपां सीता राक्षसीभिः समावृता}
{अभिगम्य विशालाक्षी तस्थौ शोकपरिप्लुता} %||5-22-13||

\twolineshloka
{तां कृशां दीनवदनां मलिनाम्बरधारिणीम्}
{भर्त्सयां चक्रिरे भीमा राक्षस्यस्ताः समन्ततः} %||5-22-14||

\twolineshloka
{ततस्तां विनता नाम राक्षसी भीमदर्शना}
{अब्रवीत्कुपिताकारा कराला निर्णतोदरी} %||5-22-15||

\twolineshloka
{सीते पर्याप्तमेतावद्भर्तृस्नेहो निदर्शितः}
{सर्वत्रातिकृतं भद्रे व्यसनायोपकल्पते} %||5-22-16||

\twolineshloka
{परितुष्टास्मि भद्रं ते मानुषस्ते कृतो विधिः}
{ममापि तु वचः पथ्यं ब्रुवन्त्याः कुरु मैथिलि} %||5-22-17||

\twolineshloka
{रावणं भज भर्तारं भर्तारं सर्वरक्षसाम्}
{विक्रान्तं रूपवन्तं च सुरेशमिव वासवम्} %||5-22-18||

\twolineshloka
{दक्षिणं त्यागशीलं च सर्वस्य प्रियवादिनम्}
{मानुषं कृपणं रामं त्यक्त्वा रावणमाश्रय} %||5-22-19||

\threelineshloka
{दिव्याङ्गरागा वैदेहि दिव्याभरणभूषिता}
{अद्य प्रभृति सर्वेषां लोकानामीश्वरी भव}
{अग्नेः स्वाहा यथा देवी शचीवेन्द्रस्य शोभने} %||5-22-20||

{किं ते रामेण वैदेहि कृपणेन गतायुषा॥\devanumber{21}॥} %||5-22-21|| (Check)

\addtocounter{shlokacount}{1}

\twolineshloka
{एतदुक्तं च मे वाक्यं यदि त्वं न करिष्यसि}
{अस्मिन्मुहूर्ते सर्वास्त्वां भक्षयिष्यामहे वयम्} %||5-22-22||

\twolineshloka
{अन्या तु विकटा नाम लम्बमानपयोधरा}
{अब्रवीत्कुपिता सीतां मुष्टिमुद्यम्य गर्जती} %||5-22-23||

\threelineshloka
{बहून्यप्रतिरूपाणि वचनानि सुदुर्मते}
{अनुक्रोशान्मृदुत्वाच्च सोढानि तव मैथिलि}
{न च नः कुरुषे वाक्यं हितं कालपुरस्कृतम्} %||5-22-24||

\twolineshloka
{आनीतासि समुद्रस्य पारमन्यैर्दुरासदम्}
{रावणान्तःपुरं घोरं प्रविष्टा चासि मैथिलि} %||5-22-25||

\twolineshloka
{रावणस्य गृहे रुधा अस्माभिस्तु सुरक्षिता}
{न त्वां शक्तः परित्रातुमपि साक्षात्पुरंदरः} %||5-22-26||

\twolineshloka
{कुरुष्व हितवादिन्या वचनं मम मैथिलि}
{अलमश्रुप्रपातेन त्यज शोकमनर्थकम्} %||5-22-27||

\twolineshloka
{भज प्रीतिं प्रहर्षं च त्यजैतां नित्यदैन्यताम्}
{सीते राक्षसराजेन सह क्रीड यथासुखम्} %||5-22-28||

\twolineshloka
{जानासि हि यथा भीरु स्त्रीणां यौवनमध्रुवम्}
{यावन्न ते व्यतिक्रामेत्तावत्सुखमवाप्नुहि} %||5-22-29||

\twolineshloka
{उद्यानानि च रम्याणि पर्वतोपवनानि च}
{सह राक्षसराजेन चर त्वं मदिरेक्षणे} %||5-22-30||

\twolineshloka
{स्त्रीसहस्राणि ते सप्त वशे स्थास्यन्ति सुन्दरि}
{रावणं भज भर्तारं भर्तारं सर्वरक्षसाम्} %||5-22-31||

\twolineshloka
{उत्पाट्य वा ते हृदयं भक्षयिष्यामि मैथिलि}
{यदि मे व्याहृतं वाक्यं न यथावत्करिष्यसि} %||5-22-32||

\twolineshloka
{ततश्चण्डोदरी नाम राक्षसी क्रूरदर्शना}
{भ्रामयन्ती महच्छूलमिदं वचनमब्रवीत्} %||5-22-33||

\twolineshloka
{इमां हरिणलोकाक्षीं त्रासोत्कम्पपयोधराम्}
{रावणेन हृतां दृष्ट्वा दौर्हृदो मे महानभूत्} %||5-22-34||

\twolineshloka
{यकृत्प्लीहमथोत्पीडं हृदयं च सबन्धनम्}
{अन्त्राण्यपि तथा शीर्षं खादेयमिति मे मतिः} %||5-22-35||

\twolineshloka
{ततस्तु प्रघसा नाम राक्षसी वाक्यमब्रवीत्}
{कण्ठमस्या नृशंसायाः पीडयामः किमास्यते} %||5-22-36||

\twolineshloka
{निवेद्यतां ततो राज्ञे मानुषी सा मृतेति ह}
{नात्र कश्चन संदेहः खादतेति स वक्ष्यति} %||5-22-37||

\twolineshloka
{ततस्त्वजामुखी नाम राक्षसी वाक्यमब्रवीत्}
{विशस्येमां ततः सर्वान्समान्कुरुत पीलुकान्} %||5-22-38||

\twolineshloka
{विभजाम ततः सर्वा विवादो मे न रोचते}
{पेयमानीयतां क्षिप्रं माल्यं च विविधं बहु} %||5-22-39||

\twolineshloka
{ततः शूर्पणखा नाम राक्षसी वाक्यमब्रवीत्}
{अजामुखा यदुक्तं हि तदेव मम रोचते} %||5-22-40||

\twolineshloka
{सुरा चानीयतां क्षिप्रं सर्वशोकविनाशिनी}
{मानुषं मांसमासाद्य नृत्यामोऽथ निकुम्भिलाम्} %||5-22-41||

\twolineshloka
{एवं संभर्त्स्यमाना सा सीता सुरसुतोपमा}
{राक्षसीभिः सुघोराभिर्धैर्यमुत्सृज्य रोदिति} %||5-23-42||


{॥इत्यार्षे श्रीमद्रामायणे वाल्मीकीये आदिकाव्ये सुन्दरकाण्डे द्वाविंशः सर्गः॥}

\sect{त्रयोविंशः सर्गः}

\twolineshloka
{तथा तासां वदन्तीनां परुषं दारुणं बहु}
{राक्षसीनामसौम्यानां रुरोद जनकात्मजा} %||5-23-1||

\twolineshloka
{एवमुक्ता तु वैदेही राक्षसीभिर्मनस्विनी}
{उवाच परमत्रस्ता बाष्पगद्गदया गिरा} %||5-23-2||

\twolineshloka
{न मानुषी राक्षसस्य भार्या भवितुमर्हति}
{कामं खादत मां सर्वा न करिष्यामि वो वचः} %||5-23-3||

\twolineshloka
{सा राक्षसी मध्यगता सीता सुरसुतोपमा}
{न शर्म लेभे दुःखार्ता रावणेन च तर्जिता} %||5-23-4||

\twolineshloka
{वेपते स्माधिकं सीता विशन्तीवाङ्गमात्मनः}
{वने यूथपरिभ्रष्टा मृगी कोकैरिवार्दिता} %||5-23-5||

\twolineshloka
{सा त्वशोकस्य विपुलां शाखामालम्ब्य पुष्पिताम्}
{चिन्तयामास शोकेन भर्तारं भग्नमानसा} %||5-23-6||

\twolineshloka
{सा स्नापयन्ती विपुलौ स्तनौ नेत्रजलस्रवैः}
{चिन्तयन्ती न शोकस्य तदान्तमधिगच्छति} %||5-23-7||

\twolineshloka
{सा वेपमाना पतिता प्रवाते कदली यथा}
{राक्षसीनां भयत्रस्ता विवर्णवदनाभवत्} %||5-23-8||

\twolineshloka
{तस्या सा दीर्घविपुला वेपन्त्याः सीतया तदा}
{ददृशे कम्पिनी वेणी व्यालीव परिसर्पती} %||5-23-9||

\twolineshloka
{सा निःश्वसन्ती दुःखार्ता शोकोपहतचेतना}
{आर्ता व्यसृजदश्रूणि मैथिली विललाप ह} %||5-23-10||

\twolineshloka
{हा रामेति च दुःखार्ता पुनर्हा लक्ष्मणेति च}
{हा श्वश्रु मम कौसल्ये हा सुमित्रेति भाविनि} %||5-23-11||

\twolineshloka
{लोकप्रवादः सत्योऽयं पण्डितैः समुदाहृतः}
{अकाले दुर्लभो मृत्युः स्त्रिया वा पुरुषस्य वा} %||5-23-12||

\twolineshloka
{यत्राहमाभिः क्रूराभी राक्षसीभिरिहार्दिता}
{जीवामि हीना रामेण मुहूर्तमपि दुःखिता} %||5-23-13||

\twolineshloka
{एषाल्पपुण्या कृपणा विनशिष्याम्यनाथवत्}
{समुद्रमध्ये नौ पूर्णा वायुवेगैरिवाहता} %||5-23-14||

\twolineshloka
{भर्तारं तमपश्यन्ती राक्षसीवशमागता}
{सीदामि खलु शोकेन कूलं तोयहतं यथा} %||5-23-15||

\twolineshloka
{तं पद्मदलपत्राक्षं सिंहविक्रान्तगामिनम्}
{धन्याः पश्यन्ति मे नाथं कृतज्ञं प्रियवादिनम्} %||5-23-16||

\twolineshloka
{सर्वथा तेन हीनाया रामेण विदितात्मना}
{तीष्क्णं विषमिवास्वाद्य दुर्लभं मम जीवितम्} %||5-23-17||

\twolineshloka
{कीदृशं तु मया पापं पुरा देहान्तरे कृतम्}
{येनेदं प्राप्यते दुःखं मया घोरं सुदारुणम्} %||5-23-18||

\twolineshloka
{जीवितं त्यक्तुमिच्छामि शोकेन महता वृता}
{राक्षसीभिश्च रक्षन्त्या रामो नासाद्यते मया} %||5-23-19||

\twolineshloka
{धिगस्तु खलु मानुष्यं धिगस्तु परवश्यताम्}
{न शक्यं यत्परित्यक्तुमात्मच्छन्देन जीवितम्} %||5-24-20||


{॥इत्यार्षे श्रीमद्रामायणे वाल्मीकीये आदिकाव्ये सुन्दरकाण्डे त्रयोविंशः सर्गः॥}

\sect{चतुर्विंशः सर्गः}

\twolineshloka
{प्रसक्ताश्रुमुखीत्येवं ब्रुवन्ती जनकात्मजा}
{अधोमुखमुखी बाला विलप्तुमुपचक्रमे} %||5-24-1||

\twolineshloka
{उन्मत्तेव प्रमत्तेव भ्रान्तचित्तेव शोचती}
{उपावृत्ता किशोरीव विवेष्टन्ती महीतले} %||5-24-2||

\twolineshloka
{राघवस्याप्रमत्तस्य रक्षसा कामरूपिणा}
{रावणेन प्रमथ्याहमानीता क्रोशती बलात्} %||5-24-3||

\twolineshloka
{राक्षसी वशमापन्ना भर्त्यमाना सुदारुणम्}
{चिन्तयन्ती सुदुःखार्ता नाहं जीवितुमुत्सहे} %||5-24-4||

\twolineshloka
{न हि मे जीवितेनार्थो नैवार्थैर्न च भूषणैः}
{वसन्त्या राक्षसी मध्ये विना रामं महारथम्} %||5-24-5||

\twolineshloka
{धिङ्मामनार्यामसतीं याहं तेन विना कृता}
{मुहूर्तमपि रक्षामि जीवितं पापजीविता} %||5-24-6||

\twolineshloka
{का च मे जीविते श्रद्धा सुखे वा तं प्रियं विना}
{भर्तारं सागरान्ताया वसुधायाः प्रियं वदम्} %||5-24-7||

\twolineshloka
{भिद्यतां भक्ष्यतां वापि शरीरं विसृजाम्यहम्}
{न चाप्यहं चिरं दुःखं सहेयं प्रियवर्जिता} %||5-24-8||

\twolineshloka
{चरणेनापि सव्येन न स्पृशेयं निशाचरम्}
{रावणं किं पुनरहं कामयेयं विगर्हितम्} %||5-24-9||

\twolineshloka
{प्रत्याख्यातं न जानाति नात्मानं नात्मनः कुलम्}
{यो नृशंस स्वभावेन मां प्रार्थयितुमिच्छति} %||5-24-10||

\twolineshloka
{छिन्ना भिन्ना विभक्ता वा दीप्ते वाग्नौ प्रदीपिता}
{रावणं नोपतिष्ठेयं किं प्रलापेन वश्चिरम्} %||5-24-11||

\twolineshloka
{ख्यातः प्राज्ञः कृतज्ञश्च सानुक्रोशश्च राघवः}
{सद्वृत्तो निरनुक्रोशः शङ्के मद्भाग्यसंक्षयात्} %||5-24-12||

\twolineshloka
{राक्षसानां जनस्थाने सहस्राणि चतुर्दश}
{येनैकेन निरस्तानि स मां किं नाभिपद्यते} %||5-24-13||

\twolineshloka
{निरुद्धा रावणेनाहमल्पवीर्येण रक्षसा}
{समर्थः खलु मे भर्ता रावणं हन्तुमाहवे} %||5-24-14||

\twolineshloka
{विराधो दण्डकारण्ये येन राक्षसपुंगवः}
{रणे रामेण निहतः स मां किं नाभिपद्यते} %||5-24-15||

\twolineshloka
{कामं मध्ये समुद्रस्य लङ्केयं दुष्प्रधर्षणा}
{न तु राघवबाणानां गतिरोधी ह विद्यते} %||5-24-16||

\twolineshloka
{किं नु तत्कारणं येन रामो दृढपराक्रमः}
{रक्षसापहृतां भार्यामिष्टां नाभ्यवपद्यते} %||5-24-17||

\twolineshloka
{इहस्थां मां न जानीते शङ्के लक्ष्मणपूर्वजः}
{जानन्नपि हि तेजस्वी धर्षणां मर्षयिष्यति} %||5-24-18||

\twolineshloka
{हृतेति योऽधिगत्वा मां राघवाय निवेदयेत्}
{गृध्रराजोऽपि स रणे रावणेन निपातितः} %||5-24-19||

\twolineshloka
{कृतं कर्म महत्तेन मां तदाभ्यवपद्यता}
{तिष्ठता रावणद्वन्द्वे वृद्धेनापि जटायुषा} %||5-24-20||

\twolineshloka
{यदि मामिह जानीयाद्वर्तमानां स राघवः}
{अद्य बाणैरभिक्रुद्धः कुर्याल्लोकमराक्षसं} %||5-24-21||

\twolineshloka
{विधमेच्च पुरीं लङ्कां शोषयेच्च महोदधिम्}
{रावणस्य च नीचस्य कीर्तिं नाम च नाशयेत्} %||5-24-22||

\threelineshloka
{ततो निहतनथानां राक्षसीनां गृहे गृहे}
{यथाहमेवं रुदती तथा भूयो न संशयः}
{अन्विष्य रक्षसां लङ्कां कुर्याद्रामः सलक्ष्मणः} %||5-24-23||

\threelineshloka
{न हि ताभ्यां रिपुर्दृष्टो मुहूतमपि जीवति}
{चिता धूमाकुलपथा गृध्रमण्डलसंकुला}
{अचिरेण तु लङ्केयं श्मशानसदृशी भवेत्} %||5-24-24||

\twolineshloka
{अचिरेणैव कालेन प्राप्स्याम्येव मनोरथम्}
{दुष्प्रस्थानोऽयमाख्याति सर्वेषां वो विपर्ययः} %||5-24-25||

\twolineshloka
{यादृशानि तु दृश्यन्ते लङ्कायामशुभानि तु}
{अचिरेणैव कालेन भविष्यति हतप्रभा} %||5-24-26||

\twolineshloka
{नूनं लङ्का हते पापे रावणे राक्षसाधिपे}
{शोषं यास्यति दुर्धर्षा प्रमदा विधवा यथा} %||5-24-27||

\twolineshloka
{पुष्योत्सवसमृद्धा च नष्टभर्त्री सराक्षसा}
{भविष्यति पुरी लङ्का नष्टभर्त्री यथाङ्गना} %||5-24-28||

\twolineshloka
{नूनं राक्षसकन्यानां रुदन्तीनां गृहे गृहे}
{श्रोष्यामि नचिरादेव दुःखार्तानामिह ध्वनिम्} %||5-24-29||

\twolineshloka
{सान्धकारा हतद्योता हतराक्षसपुंगवा}
{भविष्यति पुरी लङ्का निर्दग्धा रामसायकैः} %||5-24-30||

\twolineshloka
{यदि नाम स शूरो मां रामो रक्तान्तलोचनः}
{जानीयाद्वर्तमानां हि रावणस्य निवेशने} %||5-24-31||

\twolineshloka
{अनेन तु नृशंसेन रावणेनाधमेन मे}
{समयो यस्तु निर्दिष्टस्तस्य कालोऽयमागतः} %||5-24-32||

\twolineshloka
{अकार्यं ये न जानन्ति नैरृताः पापकारिणः}
{अधर्मात्तु महोत्पातो भविष्यति हि साम्प्रतम्} %||5-24-33||

\twolineshloka
{नैते धर्मं विजानन्ति राक्षसाः पिशिताशनाः}
{ध्रुवं मां प्रातराशार्थे राक्षसः कल्पयिष्यति} %||5-24-34||

\twolineshloka
{साहं कथं करिष्यामि तं विना प्रियदर्शनम्}
{रामं रक्तान्तनयनमपश्यन्ती सुदुःखिता} %||5-24-35||

\twolineshloka
{यदि कश्चित्प्रदाता मे विषस्याद्य भवेदिह}
{क्षिप्रं वैवस्वतं देवं पश्येयं पतिना विना} %||5-24-36||

\twolineshloka
{नाजानाज्जीवतीं रामः स मां लक्ष्मणपूर्वजः}
{जानन्तौ तौ न कुर्यातां नोर्व्यां हि मम मार्गणम्} %||5-24-37||

\twolineshloka
{नूनं ममैव शोकेन स वीरो लक्ष्मणाग्रजः}
{देवलोकमितो यातस्त्यक्त्वा देहं महीतले} %||5-24-38||

\twolineshloka
{धन्या देवाः सगन्धर्वाः सिद्धाश्च परमर्षयः}
{मम पश्यन्ति ये नाथं रामं राजीवलोचनम्} %||5-24-39||

\twolineshloka
{अथ वा न हि तस्यार्थे धर्मकामस्य धीमतः}
{मया रामस्य राजर्षेर्भार्यया परमात्मनः} %||5-24-40||

\twolineshloka
{दृश्यमाने भवेत्प्रीतः सौहृदं नास्त्यपश्यतः}
{नाशयन्ति कृतघ्रास्तु न रामो नाशयिष्यति} %||5-24-41||

\twolineshloka
{किं नु मे न गुणाः केचित्किं वा भाग्य क्षयो हि मे}
{याहं सीता वरार्हेण हीना रामेण भामिनी} %||5-24-42||

\twolineshloka
{श्रेयो मे जीवितान्मर्तुं विहीना या महात्मना}
{रामादक्लिष्टचारित्राच्छूराच्छत्रुनिबर्हणात्} %||5-24-43||

\twolineshloka
{अथ वा न्यस्तशस्त्रौ तौ वने मूलफलाशनौ}
{भ्रातरौ हि नर श्रेष्ठौ चरन्तौ वनगोचरौ} %||5-24-44||

\twolineshloka
{अथ वा राक्षसेन्द्रेण रावणेन दुरात्मना}
{छद्मना घातितौ शूरौ भ्रातरौ रामलक्ष्मणौ} %||5-24-45||

\twolineshloka
{साहमेवंगते काले मर्तुमिच्छामि सर्वथा}
{न च मे विहितो मृत्युरस्मिन्दुःखेऽपि वर्तति} %||5-24-46||

\twolineshloka
{धन्याः खलु महात्मानो मुनयः सत्यसंमताः}
{जितात्मानो महाभागा येषां न स्तः प्रियाप्रिये} %||5-24-47||

\twolineshloka
{प्रियान्न संभवेद्दुःखमप्रियादधिकं भयम्}
{ताभ्यां हि ये वियुज्यन्ते नमस्तेषां महात्मनाम्} %||5-24-48||

\twolineshloka
{साहं त्यक्ता प्रियेणेह रामेण विदितात्मना}
{प्राणांस्त्यक्ष्यामि पापस्य रावणस्य गता वशम्} %||5-25-49||


{॥इत्यार्षे श्रीमद्रामायणे वाल्मीकीये आदिकाव्ये सुन्दरकाण्डे चतुर्विंशः सर्गः॥}

\sect{पञ्चविंशः सर्गः}

\twolineshloka
{इत्युक्ताः सीतया घोरं राक्षस्यः क्रोधमूर्छिताः}
{काश्चिज्जग्मुस्तदाख्यातुं रावणस्य तरस्विनः} %||5-25-1||

\twolineshloka
{ततः सीतामुपागम्य राक्षस्यो घोरदर्शनाः}
{पुनः परुषमेकार्थमनर्थार्थमथाब्रुवन्} %||5-25-2||

\twolineshloka
{हन्तेदानीं तवानार्ये सीते पापविनिश्चये}
{राक्षस्यो भक्षयिष्यन्ति मांसमेतद्यथासुखम्} %||5-25-3||

\twolineshloka
{सीतां ताभिरनार्याभिर्दृष्ट्वा संतर्जितां तदा}
{राक्षसी त्रिजटावृद्धा शयाना वाक्यमब्रवीत्} %||5-25-4||

\twolineshloka
{आत्मानं खादतानार्या न सीतां भक्षयिष्यथ}
{जनकस्य सुतामिष्टां स्नुषां दशरथस्य च} %||5-25-5||

\twolineshloka
{स्वप्नो ह्यद्य मया दृष्टो दारुणो रोमहर्षणः}
{राक्षसानामभावाय भर्तुरस्या भवाय च} %||5-25-6||

\twolineshloka
{एवमुक्तास्त्रिजटया राक्षस्यः क्रोधमूर्छिताः}
{सर्वा एवाब्रुवन्भीतास्त्रिजटां तामिदं वचः} %||5-25-7||

{कथयस्व त्वया दृष्टः स्वप्नेऽयं कीदृशो निशि॥\devanumber{8}॥} %||5-25-8|| (Check)

\addtocounter{shlokacount}{1}

\twolineshloka
{तासां श्रुत्वा तु वचनं राक्षसीनां मुखोद्गतम्}
{उवाच वचनं काले त्रिजटास्वप्नसंश्रितम्} %||5-25-9||

\twolineshloka
{गजदन्तमयीं दिव्यां शिबिकामन्तरिक्षगाम्}
{युक्तां वाजिसहस्रेण स्वयमास्थाय राघवः} %||5-25-10||

\threelineshloka
{स्वप्ने चाद्य मया दृष्टा सीता शुक्लाम्बरावृता}
{सागरेण परिक्षिप्तं श्वेतपर्वतमास्थिता}
{रामेण संगता सीता भास्करेण प्रभा यथा} %||5-25-11||

\twolineshloka
{राघवश्च मया दृष्टश्चतुर्दन्तं महागजम्}
{आरूढः शैलसंकाशं चचार सहलक्ष्मणः} %||5-25-12||

\twolineshloka
{ततस्तौ नरशार्दूलौ दीप्यमानौ स्वतेजसा}
{शुक्लमाल्याम्बरधरौ जानकीं पर्युपस्थितौ} %||5-25-13||

\twolineshloka
{ततस्तस्य नगस्याग्रे आकाशस्थस्य दन्तिनः}
{भर्त्रा परिगृहीतस्य जानकी स्कन्धमाश्रिता} %||5-25-14||

\twolineshloka
{भर्तुरङ्कात्समुत्पत्य ततः कमललोचना}
{चन्द्रसूर्यौ मया दृष्टा पाणिभ्यां परिमार्जती} %||5-25-15||

\twolineshloka
{ततस्ताभ्यां कुमाराभ्यामास्थितः स गजोत्तमः}
{सीतया च विशालाक्ष्या लङ्काया उपरि स्थितः} %||5-25-16||

\threelineshloka
{पाण्डुरर्षभयुक्तेन रथेनाष्टयुजा स्वयम्}
{शुक्लमाल्याम्बरधरो लक्ष्मणेन समागतः}
{लक्ष्मणेन सह भ्रात्रा सीतया सह भार्यया} %||5-25-17||

\twolineshloka
{विमानात्पुष्पकादद्य रावणः पतितो भुवि}
{कृष्यपाणः स्त्रिया दृष्टो मुण्डः कृष्णाम्बरः पुनः} %||5-25-18||

\twolineshloka
{रथेन खरयुक्तेन रक्तमाल्यानुलेपनः}
{प्रयातो दक्षिणामाशां प्रविष्टः कर्दमं ह्रदम्} %||5-25-19||

\twolineshloka
{कण्ठे बद्ध्वा दशग्रीवं प्रमदा रक्तवासिनी}
{काली कर्दमलिप्ताङ्गी दिशं याम्यां प्रकर्षति} %||5-25-20||

\twolineshloka
{वराहेण दशग्रीवः शिंशुमारेण चेन्द्रजित्}
{उष्ट्रेण कुम्भकर्णश्च प्रयातो दक्षिणां दिशम्} %||5-25-21||

\twolineshloka
{समाजश्च महान्वृत्तो गीतवादित्रनिःस्वनः}
{पिबतां रक्तमाल्यानां रक्षसां रक्तवाससाम्} %||5-25-22||

\twolineshloka
{लङ्का चेयं पुरी रम्या सवाजिरथसंकुला}
{सागरे पतिता दृष्टा भग्नगोपुरतोरणा} %||5-25-23||

\twolineshloka
{पीत्व तैलं प्रनृत्ताश्च प्रहसन्त्यो महास्वनाः}
{लङ्कायां भस्मरूक्षायां सर्वा राक्षसयोषितः} %||5-25-24||

\twolineshloka
{कुम्भकर्णादयश्चेमे सर्वे राक्षसपुंगवाः}
{रक्तं निवसनं गृह्य प्रविष्टा गोमयह्रदे} %||5-25-25||

\twolineshloka
{अपगच्छत नश्यध्वं सीतामाप्नोति राघवः}
{घातयेत्परमामर्षी सर्वैः सार्धं हि राक्षसैः} %||5-25-26||

\twolineshloka
{प्रियां बहुमतां भार्यां वनवासमनुव्रताम्}
{भर्त्सितां तर्जितां वापि नानुमंस्यति राघवः} %||5-25-27||

\twolineshloka
{तदलं क्रूरवाक्यैर्वः सान्त्वमेवाभिधीयताम्}
{अभियाचाम वैदेहीमेतद्धि मम रोचते} %||5-25-28||

\twolineshloka
{यस्या ह्येवं विधः स्वप्नो दुःखितायाः प्रदृश्यते}
{सा दुःखैर्बहुभिर्मुक्ता प्रियं प्राप्नोत्यनुत्तमम्} %||5-25-29||

\twolineshloka
{भर्त्सितामपि याचध्वं राक्षस्यः किं विवक्षया}
{राघवाद्धि भयं घोरं राक्षसानामुपस्थितम्} %||5-25-30||

\twolineshloka
{प्रणिपात प्रसन्ना हि मैथिली जनकात्मजा}
{अलमेषा परित्रातुं राक्षस्यो महतो भयात्} %||5-25-31||

\twolineshloka
{अपि चास्या विशालाक्ष्या न किंचिदुपलक्षये}
{विरुद्धमपि चाङ्गेषु सुसूक्ष्ममपि लक्ष्मणम्} %||5-25-32||

\twolineshloka
{छाया वैगुण्य मात्रं तु शङ्के दुःखमुपस्थितम्}
{अदुःखार्हामिमां देवीं वैहायसमुपस्थिताम्} %||5-25-33||

\twolineshloka
{अर्थसिद्धिं तु वैदेह्याः पश्याम्यहमुपस्थिताम्}
{राक्षसेन्द्रविनाशं च विजयं राघवस्य च} %||5-25-34||

\twolineshloka
{निमित्तभूतमेतत्तु श्रोतुमस्या महत्प्रियम्}
{दृश्यते च स्फुरच्चक्षुः पद्मपत्रमिवायतम्} %||5-25-35||

\twolineshloka
{ईषच्च हृषितो वास्या दक्षिणाया ह्यदक्षिणः}
{अकस्मादेव वैदेह्या बाहुरेकः प्रकम्पते} %||5-25-36||

\twolineshloka
{करेणुहस्तप्रतिमः सव्यश्चोरुरनुत्तमः}
{वेपन्सूचयतीवास्या राघवं पुरतः स्थितम्} %||5-25-37||

\fourlineindentedshloka
{पक्षी च शाखा निलयं प्रविष्टः}
{ पुनः पुनश्चोत्तमसान्त्ववादी}
{सुखागतां वाचमुदीरयाणः}
{ पुनः पुनश्चोदयतीव हृष्टः} %||5-26-38||


{॥इत्यार्षे श्रीमद्रामायणे वाल्मीकीये आदिकाव्ये सुन्दरकाण्डे पञ्चविंशः सर्गः॥}

\sect{षड्विंशः सर्गः}

\fourlineindentedshloka
{सा राक्षसेन्द्रस्य वचो निशम्य}
{ तद्रावणस्याप्रियमप्रियार्ता}
{सीता वितत्रास यथा वनान्ते}
{ सिंहाभिपन्ना गजराजकन्या} %||5-26-1||

\fourlineindentedshloka
{सा राक्षसी मध्यगता च भीरु}
{र्वाग्भिर्भृशं रावणतर्जिता च}
{कान्तारमध्ये विजने विसृष्टा}
{ बालेव कन्या विललाप सीता} %||5-26-2||

\fourlineindentedshloka
{सत्यं बतेदं प्रवदन्ति लोके}
{ नाकालमृत्युर्भवतीति सन्तः}
{यत्राहमेवं परिभर्त्स्यमाना}
{ जीवामि किंचित्क्षणमप्यपुण्या} %||5-26-3||

\fourlineindentedshloka
{सुखाद्विहीनं बहुदुःखपूर्ण}
{मिदं तु नूनं हृदयं स्थिरं मे}
{विदीर्यते यन्न सहस्रधाद्य}
{ वज्राहतं शृङ्गमिवाचलस्य} %||5-26-4||

\fourlineindentedshloka
{नैवास्ति नूनं मम दोषमत्र}
{ वध्याहमस्याप्रियदर्शनस्य}
{भावं न चास्याहमनुप्रदातु}
{मलं द्विजो मन्त्रमिवाद्विजाय} %||5-26-5||

\fourlineindentedshloka
{नूनं ममाङ्गान्यचिरादनार्यः}
{ शस्त्रैः शितैश्छेत्स्यति राक्षसेन्द्रः}
{तस्मिन्ननागच्छति लोकनाथे}
{ गर्भस्थजन्तोरिव शल्यकृन्तः} %||5-26-6||

\fourlineindentedshloka
{दुःखं बतेदं मम दुःखिताया}
{ मासौ चिरायाभिगमिष्यतो द्वौ}
{बद्धस्य वध्यस्य यथा निशान्ते}
{ राजापराधादिव तस्करस्य} %||5-26-7||

\fourlineindentedshloka
{हा राम हा लक्ष्मण हा सुमित्रे}
{ हा राम मातः सह मे जनन्या}
{एषा विपद्याम्यहमल्पभाग्या}
{ महार्णवे नौरिव मूढ वाता} %||5-26-8||

\fourlineindentedshloka
{तरस्विनौ धारयता मृगस्य}
{ सत्त्वेन रूपं मनुजेन्द्रपुत्रौ}
{नूनं विशस्तौ मम कारणात्तौ}
{ सिंहर्षभौ द्वाविव वैद्युतेन} %||5-26-9||

\fourlineindentedshloka
{नूनं स कालो मृगरूपधारी}
{ मामल्पभाग्यां लुलुभे तदानीम्}
{यत्रार्यपुत्रं विससर्ज मूढा}
{ रामानुजं लक्ष्मणपूर्वकं च} %||5-26-10||

\fourlineindentedshloka
{हा राम सत्यव्रत दीर्घवाहो}
{ हा पूर्णचन्द्रप्रतिमानवक्त्र}
{हा जीवलोकस्य हितः प्रियश्च}
{ वध्यां न मां वेत्सि हि राक्षसानाम्} %||5-26-11||

\fourlineindentedshloka
{अनन्यदेवत्वमियं क्षमा च}
{ भूमौ च शय्या नियमश्च धर्मे}
{पतिव्रतात्वं विफलं ममेदं}
{ कृतं कृतघ्नेष्विव मानुषाणाम्} %||5-26-12||

\fourlineindentedshloka
{मोघो हि धर्मश्चरितो ममायं}
{ तथैकपत्नीत्वमिदं निरर्थम्}
{या त्वां न पश्यामि कृशा विवर्णा}
{ हीना त्वया संगमने निराशा} %||5-26-13||

\fourlineindentedshloka
{पितुर्निर्देशं नियमेन कृत्वा}
{ वनान्निवृत्तश्चरितव्रतश्च}
{स्त्रीभिस्तु मन्ये विपुलेक्षणाभिः}
{ संरंस्यसे वीतभयः कृतार्थः} %||5-26-14||

\fourlineindentedshloka
{अहं तु राम त्वयि जातकामा}
{ चिरं विनाशाय निबद्धभावा}
{मोघं चरित्वाथ तपोव्रतं च}
{ त्यक्ष्यामि धिग्जीवितमल्पभाग्या} %||5-26-15||

\fourlineindentedshloka
{सा जीवितं क्षिप्रमहं त्यजेयं}
{ विषेण शस्त्रेण शितेन वापि}
{विषस्य दाता न तु मेऽस्ति कश्चि}
{च्छस्त्रस्य वा वेश्मनि राक्षसस्य} %||5-26-16||

\fourlineindentedshloka
{शोकाभितप्ता बहुधा विचिन्त्य}
{ सीताथ वेण्युद्ग्रथनं गृहीत्वा}
{उद्बध्य वेण्युद्ग्रथनेन शीघ्र}
{महं गमिष्यामि यमस्य मूलम्} %||5-26-17||

\fourlineindentedshloka
{इतीव सीता बहुधा विलप्य}
{ सर्वात्मना राममनुस्मरन्ती}
{प्रवेपमाना परिशुष्कवक्त्रा}
{ नगोत्तमं पुष्पितमाससाद} %||5-26-18||

\fourlineindentedshloka
{उपस्थिता सा मृदुर्सर्वगात्री}
{ शाखां गृहीत्वाथ नगस्य तस्य}
{तस्यास्तु रामं प्रविचिन्तयन्त्या}
{ रामानुजं स्वं च कुलं शुभाङ्ग्याः} %||5-26-19||

\fourlineindentedshloka
{शोकानिमित्तानि तदा बहूनि}
{ धैर्यार्जितानि प्रवराणि लोके}
{प्रादुर्निमित्तानि तदा बभूवुः}
{ पुरापि सिद्धान्युपलक्षितानि} %||5-27-20||


{॥इत्यार्षे श्रीमद्रामायणे वाल्मीकीये आदिकाव्ये सुन्दरकाण्डे षड्विंशः सर्गः॥}

\sect{सप्तविंशः सर्गः}

\fourlineindentedshloka
{तथागतां तां व्यथितामनिन्दितां}
{ व्यपेतहर्षां परिदीनमानसाम्}
{शुभां निमित्तानि शुभानि भेजिरे}
{ नरं श्रिया जुष्टमिवोपजीविनः} %||5-27-1||

\fourlineindentedshloka
{तस्याः शुभं वाममरालपक्ष्म}
{ राजीवृतं कृष्णविशालशुक्लम्}
{प्रास्पन्दतैकं नयनं सुकेश्या}
{ मीनाहतं पद्ममिवाभिताम्रम्} %||5-27-2||

\fourlineindentedshloka
{भुजश्च चार्वञ्चितपीनवृत्तः}
{ परार्ध्य कालागुरुचन्दनार्हः}
{अनुत्तमेनाध्युषितः प्रियेण}
{ चिरेण वामः समवेपताशु} %||5-27-3||

\fourlineindentedshloka
{गजेन्द्रहस्तप्रतिमश्च पीन}
{स्तयोर्द्वयोः संहतयोः सुजातः}
{प्रस्पन्दमानः पुनरूरुरस्या}
{ रामं पुरस्तात्स्थितमाचचक्षे} %||5-27-4||

\fourlineindentedshloka
{शुभं पुनर्हेमसमानवर्ण}
{मीषद्रजोध्वस्तमिवामलाक्ष्याः}
{वासः स्थितायाः शिखराग्रदन्त्याः}
{ किंचित्परिस्रंसत चारुगात्र्याः} %||5-27-5||

\fourlineindentedshloka
{एतैर्निमित्तैरपरैश्च सुभ्रूः}
{ संबोधिता प्रागपि साधुसिद्धैः}
{वातातपक्लान्तमिव प्रनष्टं}
{ वर्षेण बीजं प्रतिसंजहर्ष} %||5-27-6||

\fourlineindentedshloka
{तस्याः पुनर्बिम्बफलोपमौष्ठं}
{ स्वक्षिभ्रुकेशान्तमरालपक्ष्म}
{वक्त्रं बभासे सितशुक्लदंष्ट्रं}
{ राहोर्मुखाच्चन्द्र इव प्रमुक्तः} %||5-27-7||

\fourlineindentedshloka
{सा वीतशोका व्यपनीततन्द्री}
{ शान्तज्वरा हर्षविबुद्धसत्त्वा}
{अशोभतार्या वदनेन शुक्ले}
{ शीतान्शुना रात्रिरिवोदितेन} %||5-28-8||


{॥इत्यार्षे श्रीमद्रामायणे वाल्मीकीये आदिकाव्ये सुन्दरकाण्डे सप्तविंशः सर्गः॥}

\sect{अष्टाविंशः सर्गः}

\twolineshloka
{हनुमानपि विक्रान्तः सर्वं शुश्राव तत्त्वतः}
{सीतायास्त्रिजटायाश्च राक्षसीनां च तर्जनम्} %||5-28-1||

\twolineshloka
{अवेक्षमाणस्तां देवीं देवतामिव नन्दने}
{ततो बहुविधां चिन्तां चिन्तयामास वानरः} %||5-28-2||

\twolineshloka
{यां कपीनां सहस्राणि सुबहून्ययुतानि च}
{दिक्षु सर्वासु मार्गन्ते सेयमासादिता मया} %||5-28-3||

\twolineshloka
{चारेण तु सुयुक्तेन शत्रोः शक्तिमवेक्षिता}
{गूढेन चरता तावदवेक्षितमिदं मया} %||5-28-4||

\twolineshloka
{राक्षसानां विशेषश्च पुरी चेयमवेक्षिता}
{राक्षसाधिपतेरस्य प्रभावो रावणस्य च} %||5-28-5||

\twolineshloka
{युक्तं तस्याप्रमेयस्य सर्वसत्त्वदयावतः}
{समाश्वासयितुं भार्यां पतिदर्शनकाङ्क्षिणीम्} %||5-28-6||

\twolineshloka
{अहमाश्वासयाम्येनां पूर्णचन्द्रनिभाननाम्}
{अदृष्टदुःखां दुःखस्य न ह्यन्तमधिगच्छतीम्} %||5-28-7||

\twolineshloka
{यदि ह्यहमिमां देवीं शोकोपहतचेतनाम्}
{अनाश्वास्य गमिष्यामि दोषवद्गमनं भवेत्} %||5-28-8||

\twolineshloka
{गते हि मयि तत्रेयं राजपुत्री यशस्विनी}
{परित्राणमविन्दन्ती जानकी जीवितं त्यजेत्} %||5-28-9||

\twolineshloka
{मया च स महाबाहुः पूर्णचन्द्रनिभाननः}
{समाश्वासयितुं न्याय्यः सीतादर्शनलालसः} %||5-28-10||

\twolineshloka
{निशाचरीणां प्रत्यक्षमक्षमं चाभिभाषणम्}
{कथं नु खलु कर्तव्यमिदं कृच्छ्र गतो ह्यहम्} %||5-28-11||

\twolineshloka
{अनेन रात्रिशेषेण यदि नाश्वास्यते मया}
{सर्वथा नास्ति संदेहः परित्यक्ष्यति जीवितम्} %||5-28-12||

\twolineshloka
{रामश्च यदि पृच्छेन्मां किं मां सीताब्रवीद्वचः}
{किमहं तं प्रतिब्रूयामसंभाष्य सुमध्यमाम्} %||5-28-13||

\twolineshloka
{सीतासंदेशरहितं मामितस्त्वरया गतम्}
{निर्दहेदपि काकुत्स्थः क्रुद्धस्तीव्रेण चक्षुषा} %||5-28-14||

\twolineshloka
{यदि चेद्योजयिष्यामि भर्तारं रामकारणात्}
{व्यर्थमागमनं तस्य ससैन्यस्य भविष्यति} %||5-28-15||

\twolineshloka
{अन्तरं त्वहमासाद्य राक्षसीनामिह स्थितः}
{शनैराश्वासयिष्यामि संतापबहुलामिमाम्} %||5-28-16||

\twolineshloka
{अहं ह्यतितनुश्चैव वनरश्च विशेषतः}
{वाचं चोदाहरिष्यामि मानुषीमिह संस्कृताम्} %||5-28-17||

\twolineshloka
{यदि वाचं प्रदास्यामि द्विजातिरिव संस्कृताम्}
{रावणं मन्यमाना मां सीता भीता भविष्यति} %||5-28-18||

\twolineshloka
{अवश्यमेव वक्तव्यं मानुषं वाक्यमर्थवत्}
{मया सान्त्वयितुं शक्या नान्यथेयमनिन्दिता} %||5-28-19||

\twolineshloka
{सेयमालोक्य मे रूपं जानकी भाषितं तथा}
{रक्षोभिस्त्रासिता पूर्वं भूयस्त्रासं गमिष्यति} %||5-28-20||

\twolineshloka
{ततो जातपरित्रासा शब्दं कुर्यान्मनस्विनी}
{जानमाना विशालाक्षी रावणं कामरूपिणम्} %||5-28-21||

\twolineshloka
{सीतया च कृते शब्दे सहसा राक्षसीगणः}
{नानाप्रहरणो घोरः समेयादन्तकोपमः} %||5-28-22||

\twolineshloka
{ततो मां संपरिक्षिप्य सर्वतो विकृताननाः}
{वधे च ग्रहणे चैव कुर्युर्यत्नं यथाबलम्} %||5-28-23||

\twolineshloka
{तं मां शाखाः प्रशाखाश्च स्कन्धांश्चोत्तमशाखिनाम्}
{दृष्ट्वा विपरिधावन्तं भवेयुर्भयशङ्किताः} %||5-28-24||

\twolineshloka
{मम रूपं च संप्रेक्ष्य वनं विचरतो महत्}
{राक्षस्यो भयवित्रस्ता भवेयुर्विकृताननाः} %||5-28-25||

\twolineshloka
{ततः कुर्युः समाह्वानं राक्षस्यो रक्षसामपि}
{राक्षसेन्द्रनियुक्तानां राक्षसेन्द्रनिवेशने} %||5-28-26||

\twolineshloka
{ते शूलशरनिस्त्रिंश विविधायुधपाणयः}
{आपतेयुर्विमर्देऽस्मिन्वेगेनोद्विग्नकारिणः} %||5-28-27||

\twolineshloka
{संक्रुद्धस्तैस्तु परितो विधमन्रक्षसां बलम्}
{शक्नुयं न तु संप्राप्तुं परं पारं महोदधेः} %||5-28-28||

\twolineshloka
{मां वा गृह्णीयुराप्लुत्य बहवः शीघ्रकारिणः}
{स्यादियं चागृहीतार्था मम च ग्रहणं भवेत्} %||5-28-29||

\twolineshloka
{हिंसाभिरुचयो हिंस्युरिमां वा जनकात्मजाम्}
{विपन्नं स्यात्ततः कार्यं रामसुग्रीवयोरिदम्} %||5-28-30||

\twolineshloka
{उद्देशे नष्टमार्गेऽस्मिन्राक्षसैः परिवारिते}
{सागरेण परिक्षिप्ते गुप्ते वसति जानकी} %||5-28-31||

\twolineshloka
{विशस्ते वा गृहीते वा रक्षोभिर्मयि संयुगे}
{नान्यं पश्यामि रामस्य सहायं कार्यसाधने} %||5-28-32||

\twolineshloka
{विमृशंश्च न पश्यामि यो हते मयि वानरः}
{शतयोजनविस्तीर्णं लङ्घयेत महोदधिम्} %||5-28-33||

\twolineshloka
{कामं हन्तुं समर्थोऽस्मि सहस्राण्यपि रक्षसाम्}
{न तु शक्ष्यामि संप्राप्तुं परं पारं महोदधेः} %||5-28-34||

\twolineshloka
{असत्यानि च युद्धानि संशयो मे न रोचते}
{कश्च निःसंशयं कार्यं कुर्यात्प्राज्ञः ससंशयम्} %||5-28-35||

\twolineshloka
{एष दोषो महान्हि स्यान्मम सीताभिभाषणे}
{प्राणत्यागश्च वैदेह्या भवेदनभिभाषणे} %||5-28-36||

\twolineshloka
{भूताश्चार्था विनश्यन्ति देशकालविरोधिताः}
{विक्लवं दूतमासाद्य तमः सूर्योदये यथा} %||5-28-37||

\twolineshloka
{अर्थानर्थान्तरे बुद्धिर्निश्चितापि न शोभते}
{घातयन्ति हि कार्याणि दूताः पण्डितमानिनः} %||5-28-38||

\twolineshloka
{न विनश्येत्कथं कार्यं वैक्लव्यं न कथं भवेत्}
{लङ्घनं च समुद्रस्य कथं नु न वृथा भवेत्} %||5-28-39||

\twolineshloka
{कथं नु खलु वाक्यं मे शृणुयान्नोद्विजेत च}
{इति संचिन्त्य हनुमांश्चकार मतिमान्मतिम्} %||5-28-40||

\twolineshloka
{राममक्लिष्टकर्माणं स्वबन्धुमनुकीर्तयन्}
{नैनामुद्वेजयिष्यामि तद्बन्धुगतमानसाम्} %||5-28-41||

\twolineshloka
{इक्ष्वाकूणां वरिष्ठस्य रामस्य विदितात्मनः}
{शुभानि धर्मयुक्तानि वचनानि समर्पयन्} %||5-28-42||

\twolineshloka
{श्रावयिष्यामि सर्वाणि मधुरां प्रब्रुवन्गिरम्}
{श्रद्धास्यति यथा हीयं तथा सर्वं समादधे} %||5-28-43||

\fourlineindentedshloka
{इति स बहुविधं महानुभावो}
{ जगतिपतेः प्रमदामवेक्षमाणः}
{मधुरमवितथं जगाद वाक्यं}
{ द्रुमविटपान्तरमास्थितो हनूमान्} %||5-29-44||


{॥इत्यार्षे श्रीमद्रामायणे वाल्मीकीये आदिकाव्ये सुन्दरकाण्डे अष्टाविंशः सर्गः॥}

\sect{एकोनत्रिंशः सर्गः}

\twolineshloka
{एवं बहुविधां चिन्तां चिन्तयित्व महाकपिः}
{संश्रवे मधुरं वाक्यं वैदेह्या व्याजहार ह} %||5-29-1||

\threelineshloka
{राजा दशरथो नाम रथकुञ्जरवाजिनाम्}
{पुण्यशीलो महाकीर्तिरृजुरासीन्महायशाः}
{चक्रवर्तिकुले जातः पुरंदरसमो बले} %||5-29-2||

\twolineshloka
{अहिंसारतिरक्षुद्रो घृणी सत्यपराक्रमः}
{मुख्यश्चेक्ष्वाकुवंशस्य लक्ष्मीवाँल्लक्ष्मिवर्धनः} %||5-29-3||

\twolineshloka
{पार्थिवव्यञ्जनैर्युक्तः पृथुश्रीः पार्थिवर्षभः}
{पृथिव्यां चतुरन्तयां विश्रुतः सुखदः सुखी} %||5-29-4||

\twolineshloka
{तस्य पुत्रः प्रियो ज्येष्ठस्ताराधिपनिभाननः}
{रामो नाम विशेषज्ञः श्रेष्ठः सर्वधनुष्मताम्} %||5-29-5||

\twolineshloka
{रक्षिता स्वस्य वृत्तस्य स्वजनस्यापि रक्षिता}
{रक्षिता जीवलोकस्य धर्मस्य च परंतपः} %||5-29-6||

\twolineshloka
{तस्य सत्याभिसंधस्य वृद्धस्य वचनात्पितुः}
{सभार्यः सह च भ्रात्रा वीरः प्रव्रजितो वनम्} %||5-29-7||

\threelineshloka
{तेन तत्र महारण्ये मृगयां परिधावता}
{जनस्थानवधं श्रुत्वा हतौ च खरदूषणौ}
{ततस्त्वमर्षापहृता जानकी रावणेन तु} %||5-29-8||

\twolineshloka
{यथारूपां यथावर्णां यथालक्ष्मीं विनिश्चिताम्}
{अश्रौषं राघवस्याहं सेयमासादिता मया} %||5-29-9||

\twolineshloka
{विररामैवमुक्त्वासौ वाचं वानरपुंगवः}
{जानकी चापि तच्छ्रुत्वा विस्मयं परमं गता} %||5-29-10||

\twolineshloka
{ततः सा वक्रकेशान्ता सुकेशी केशसंवृतम्}
{उन्नम्य वदनं भीरुः शिंशपावृक्षमैक्षत} %||5-29-11||

\fourlineindentedshloka
{सा तिर्यगूर्ध्वं च तथाप्यधस्ता}
{न्निरीक्षमाणा तमचिन्त्य बुद्धिम्}
{ददर्श पिङ्गाधिपतेरमात्यं}
{ वातात्मजं सूर्यमिवोदयस्थम्} %||5-30-12||


{॥इत्यार्षे श्रीमद्रामायणे वाल्मीकीये आदिकाव्ये सुन्दरकाण्डे एकोनत्रिंशः सर्गः॥}

\sect{त्रिंशः सर्गः}

\twolineshloka
{ततः शाखान्तरे लीनं दृष्ट्वा चलितमानसा}
{सा ददर्श कपिं तत्र प्रश्रितं प्रियवादिनम्} %||5-30-1||

\twolineshloka
{सा तु दृष्ट्वा हरिश्रेष्ठं विनीतवदुपस्थितम्}
{मैथिली चिन्तयामास स्वप्नोऽयमिति भामिनी} %||5-30-2||

\fourlineindentedshloka
{सा तं समीक्ष्यैव भृशं विसंज्ञा}
{ गतासुकल्पेव बभूव सीता}
{चिरेण संज्ञां प्रतिलभ्य चैव}
{ विचिन्तयामास विशालनेत्रा} %||5-30-3||

\fourlineindentedshloka
{स्वप्नो मयायं विकृतोऽद्य दृष्टः}
{ शाखामृगः शास्त्रगणैर्निषिद्धः}
{स्वस्त्यस्तु रामाय सलक्ष्मणाय}
{ तथा पितुर्मे जनकस्य राज्ञः} %||5-30-4||

\fourlineindentedshloka
{स्वप्नोऽपि नायं न हि मेऽस्ति निद्रा}
{ शोकेन दुःखेन च पीडितायाः}
{सुखं हि मे नास्ति यतोऽस्मि हीना}
{ तेनेन्दुपूर्णप्रतिमाननेन} %||5-30-5||

\fourlineindentedshloka
{अहं हि तस्याद्य मनो भवेन}
{ संपीडिता तद्गतसर्वभावा}
{विचिन्तयन्ती सततं तमेव}
{ तथैव पश्यामि तथा शृणोमि} %||5-30-6||

\fourlineindentedshloka
{मनोरथः स्यादिति चिन्तयामि}
{ तथापि बुद्ध्या च वितर्कयामि}
{किं कारणं तस्य हि नास्ति रूपं}
{ सुव्यक्तरूपश्च वदत्ययं माम्} %||5-30-7||

\fourlineindentedshloka
{नमोऽस्तु वाचस्पतये सवज्रिणे}
{ स्वयम्भुवे चैव हुताशनाय}
{अनेन चोक्तं यदिदं ममाग्रतो}
{ वनौकसा तच्च तथास्तु नान्यथा} %||5-31-8||


{॥इत्यार्षे श्रीमद्रामायणे वाल्मीकीये आदिकाव्ये सुन्दरकाण्डे त्रिंशः सर्गः॥}

\sect{एकत्रिंशः सर्गः}

\twolineshloka
{तामब्रवीन्महातेजा हनूमान्मारुतात्मजः}
{शिरस्यञ्जलिमाधाय सीतां मधुरया गिरा} %||5-31-1||

\twolineshloka
{का नु पद्मपलाशाक्षी क्लिष्टकौशेयवासिनी}
{द्रुमस्य शाखामालम्ब्य तिष्ठसि त्वमनिन्दिता} %||5-31-2||

\twolineshloka
{किमर्थं तव नेत्राभ्यां वारि स्रवति शोकजम्}
{पुण्डरीकपलाशाभ्यां विप्रकीर्णमिवोदकम्} %||5-31-3||

\twolineshloka
{सुराणामसुराणां च नागगन्धर्वरक्षसाम्}
{यक्षाणां किंनराणां च का त्वं भवसि शोभने} %||5-31-4||

\twolineshloka
{का त्वं भवसि रुद्राणां मरुतां वा वरानने}
{वसूनां वा वरारोहे देवता प्रतिभासि मे} %||5-31-5||

\twolineshloka
{किं नु चन्द्रमसा हीना पतिता विबुधालयात्}
{रोहिणी ज्योतिषां श्रेष्ठा श्रेष्ठा सर्वगुणान्विता} %||5-31-6||

\twolineshloka
{कोपाद्वा यदि वा मोहाद्भर्तारमसितेक्षणा}
{वसिष्ठं कोपयित्वा त्वं नासि कल्याण्यरुन्धती} %||5-31-7||

\twolineshloka
{को नौ पुत्रः पिता भ्रात भर्ता वा ते सुमध्यमे}
{अस्माल्लोकादमुं लोकं गतं त्वमनुशोचसि} %||5-31-8||

\twolineshloka
{व्यञ्जनानि हि ते यानि लक्षणानि च लक्षये}
{महिषी भूमिपालस्य राजकन्यासि मे मता} %||5-31-9||

\twolineshloka
{रावणेन जनस्थानाद्बलादपहृता यदि}
{सीता त्वमसि भद्रं ते तन्ममाचक्ष्व पृच्छतः} %||5-31-10||

\twolineshloka
{सा तस्य वचनं श्रुत्वा रामकीर्तनहर्षिता}
{उवाच वाक्यं वैदेही हनूमन्तं द्रुमाश्रितम्} %||5-31-11||

\twolineshloka
{दुहिता जनकस्याहं वैदेहस्य महात्मनः}
{सीता च नाम नाम्नाहं भार्या रामस्य धीमतः} %||5-31-12||

\twolineshloka
{समा द्वादश तत्राहं राघवस्य निवेशने}
{भुञ्जाना मानुषान्भोगान्सर्वकामसमृद्धिनी} %||5-31-13||

\twolineshloka
{ततस्त्रयोदशे वर्षे राज्येनेक्ष्वाकुनन्दनम्}
{अभिषेचयितुं राजा सोपाध्यायः प्रचक्रमे} %||5-31-14||

\twolineshloka
{तस्मिन्संभ्रियमाणे तु राघवस्याभिषेचने}
{कैकेयी नाम भर्तारं देवी वचनमब्रवीत्} %||5-31-15||

\twolineshloka
{न पिबेयं न खादेयं प्रत्यहं मम भोजनम्}
{एष मे जीवितस्यान्तो रामो यद्यभिषिच्यते} %||5-31-16||

\twolineshloka
{यत्तदुक्तं त्वया वाक्यं प्रीत्या नृपतिसत्तम}
{तच्चेन्न वितथं कार्यं वनं गच्छतु राघवः} %||5-31-17||

\twolineshloka
{स राजा सत्यवाग्देव्या वरदानमनुस्मरन्}
{मुमोह वचनं श्रुत्वा कैकेय्याः क्रूरमप्रियम्} %||5-31-18||

\twolineshloka
{ततस्तु स्थविरो राजा सत्यधर्मे व्यवस्थितः}
{ज्येष्ठं यशस्विनं पुत्रं रुदन्राज्यमयाचत} %||5-31-19||

\twolineshloka
{स पितुर्वचनं श्रीमानभिषेकात्परं प्रियम्}
{मनसा पूर्वमासाद्य वाचा प्रतिगृहीतवान्} %||5-31-20||

\twolineshloka
{दद्यान्न प्रतिगृह्णीयान्न ब्रूयत्किंचिदप्रियम्}
{अपि जीवितहेतोर्हि रामः सत्यपराक्रमः} %||5-31-21||

\twolineshloka
{स विहायोत्तरीयाणि महार्हाणि महायशाः}
{विसृज्य मनसा राज्यं जनन्यै मां समादिशत्} %||5-31-22||

\twolineshloka
{साहं तस्याग्रतस्तूर्णं प्रस्थिता वनचारिणी}
{न हि मे तेन हीनाया वासः स्वर्गेऽपि रोचते} %||5-31-23||

\twolineshloka
{प्रागेव तु महाभागः सौमित्रिर्मित्रनन्दनः}
{पूर्वजस्यानुयात्रार्थे द्रुमचीरैरलंकृतः} %||5-31-24||

\twolineshloka
{ते वयं भर्तुरादेशं बहु मान्यदृढव्रताः}
{प्रविष्टाः स्म पुराद्दृष्टं वनं गम्भीरदर्शनम्} %||5-31-25||

\twolineshloka
{वसतो दण्डकारण्ये तस्याहममितौजसः}
{रक्षसापहृता भार्या रावणेन दुरात्मना} %||5-31-26||

\twolineshloka
{द्वौ मासौ तेन मे कालो जीवितानुग्रहः कृतः}
{ऊर्ध्वं द्वाभ्यां तु मासाभ्यां ततस्त्यक्ष्यामि जीवितम्} %||5-32-27||


{॥इत्यार्षे श्रीमद्रामायणे वाल्मीकीये आदिकाव्ये सुन्दरकाण्डे एकत्रिंशः सर्गः॥}

\sect{द्वात्रिंशः सर्गः}

\twolineshloka
{तस्यास्तद्वचनं श्रुत्वा हनूमान्हरियूथपः}
{दुःखाद्दुःखाभिभूतायाः सान्तमुत्तरमब्रवीत्} %||5-32-1||

\twolineshloka
{अहं रामस्य संदेशाद्देवि दूतस्तवागतः}
{वैदेहि कुशली रामस्त्वां च कौशलमब्रवीत्} %||5-32-2||

\twolineshloka
{यो ब्राह्ममस्त्रं वेदांश्च वेद वेदविदां वरः}
{स त्वां दाशरथी रामो देवि कौशलमब्रवीत्} %||5-32-3||

\twolineshloka
{लक्ष्मणश्च महातेजा भर्तुस्तेऽनुचरः प्रियः}
{कृतवाञ्शोकसंतप्तः शिरसा तेऽभिवादनम्} %||5-32-4||

\twolineshloka
{सा तयोः कुशलं देवी निशम्य नरसिंहयोः}
{प्रीतिसंहृष्टसर्वाङ्गी हनूमान्तमथाब्रवीत्} %||5-32-5||

\twolineshloka
{कल्याणी बत गथेयं लौकिकी प्रतिभाति मे}
{एहि जीवन्तमानदो नरं वर्षशतादपि} %||5-32-6||

\twolineshloka
{तयोः समागमे तस्मिन्प्रीतिरुत्पादिताद्भुता}
{परस्परेण चालापं विश्वस्तौ तौ प्रचक्रतुः} %||5-32-7||

\twolineshloka
{तस्यास्तद्वचनं श्रुत्वा हनूमान्हरियूथपः}
{सीतायाः शोकदीनायाः समीपमुपचक्रमे} %||5-32-8||

\twolineshloka
{यथा यथा समीपं स हनूमानुपसर्पति}
{तथा तथा रावणं सा तं सीता परिशङ्कते} %||5-32-9||

\twolineshloka
{अहो धिग्धिक्कृतमिदं कथितं हि यदस्य मे}
{रूपान्तरमुपागम्य स एवायं हि रावणः} %||5-32-10||

\twolineshloka
{तामशोकस्य शाखां सा विमुक्त्वा शोककर्शिता}
{तस्यामेवानवद्याङ्गी धरण्यां समुपाविशत्} %||5-32-11||

\twolineshloka
{अवन्दत महाबाहुस्ततस्तां जनकात्मजाम्}
{सा चैनं भयवित्रस्ता भूयो नैवाभ्युदैक्षत} %||5-32-12||

\twolineshloka
{तं दृष्ट्वा वन्दमानं तु सीता शशिनिभानना}
{अब्रवीद्दीर्घमुच्छ्वस्य वानरं मधुरस्वरा} %||5-32-13||

\twolineshloka
{मायां प्रविष्टो मायावी यदि त्वं रावणः स्वयम्}
{उत्पादयसि मे भूयः संतापं तन्न शोभनम्} %||5-32-14||

\twolineshloka
{स्वं परित्यज्य रूपं यः परिव्राजकरूपधृत्}
{जनस्थाने मया दृष्टस्त्वं स एवासि रावणः} %||5-32-15||

\twolineshloka
{उपवासकृशां दीनां कामरूप निशाचर}
{संतापयसि मां भूयः संतापं तन्न शोभनम्} %||5-32-16||

\twolineshloka
{यदि रामस्य दूतस्त्वमागतो भद्रमस्तु ते}
{पृच्छामि त्वां हरिश्रेष्ठ प्रिया राम कथा हि मे} %||5-32-17||

\twolineshloka
{गुणान्रामस्य कथय प्रियस्य मम वानर}
{चित्तं हरसि मे सौम्य नदीकूलं यथा रयः} %||5-32-18||

\twolineshloka
{अहो स्वप्नस्य सुखता याहमेवं चिराहृता}
{प्रेषितं नाम पश्यामि राघवेण वनौकसं} %||5-32-19||

\twolineshloka
{स्वप्नेऽपि यद्यहं वीरं राघवं सहलक्ष्मणम्}
{पश्येयं नावसीदेयं स्वप्नोऽपि मम मत्सरी} %||5-32-20||

\twolineshloka
{नाहं स्वप्नमिमं मन्ये स्वप्ने दृष्ट्वा हि वानरम्}
{न शक्योऽभ्युदयः प्राप्तुं प्राप्तश्चाभ्युदयो मम} %||5-32-21||

\twolineshloka
{किं नु स्याच्चित्तमोहोऽयं भवेद्वातगतिस्त्वियम्}
{उन्मादजो विकारो वा स्यादियं मृगतृष्णिका} %||5-32-22||

\twolineshloka
{अथ वा नायमुन्मादो मोहोऽप्युन्मादलक्ष्मणः}
{संबुध्ये चाहमात्मानमिमं चापि वनौकसं} %||5-32-23||

\twolineshloka
{इत्येवं बहुधा सीता संप्रधार्य बलाबलम्}
{रक्षसां कामरूपत्वान्मेने तं राक्षसाधिपम्} %||5-32-24||

\twolineshloka
{एतां बुद्धिं तदा कृत्वा सीता सा तनुमध्यमा}
{न प्रतिव्याजहाराथ वानरं जनकात्मजा} %||5-32-25||

\twolineshloka
{सीतायाश्चिन्तितं बुद्ध्वा हनूमान्मारुतात्मजः}
{श्रोत्रानुकूलैर्वचनैस्तदा तां संप्रहर्षयत्} %||5-32-26||

\twolineshloka
{आदित्य इव तेजस्वी लोककान्तः शशी यथा}
{राजा सर्वस्य लोकस्य देवो वैश्रवणो यथा} %||5-32-27||

\twolineshloka
{विक्रमेणोपपन्नश्च यथा विष्णुर्महायशाः}
{सत्यवादी मधुरवाग्देवो वाचस्पतिर्यथा} %||5-32-28||

\threelineshloka
{रूपवान्सुभगः श्रीमान्कन्दर्प इव मूर्तिमान्}
{स्थानक्रोधप्रहर्ता च श्रेष्ठो लोके महारथः}
{बाहुच्छायामवष्टब्धो यस्य लोको महात्मनः} %||5-32-29||

\twolineshloka
{अपकृष्याश्रमपदान्मृगरूपेण राघवम्}
{शून्ये येनापनीतासि तस्य द्रक्ष्यसि यत्फलम्} %||5-32-30||

\twolineshloka
{नचिराद्रावणं संख्ये यो वधिष्यति वीर्यवान्}
{रोषप्रमुक्तैरिषुभिर्ज्वलद्भिरिव पावकैः} %||5-32-31||

\twolineshloka
{तेनाहं प्रेषितो दूतस्त्वत्सकाशमिहागतः}
{त्वद्वियोगेन दुःखार्तः स त्वां कौशलमब्रवीत्} %||5-32-32||

\twolineshloka
{लक्ष्मणश्च महातेजाः सुमित्रानन्दवर्धनः}
{अभिवाद्य महाबाहुः सोऽपि कौशलमब्रवीत्} %||5-32-33||

\twolineshloka
{रामस्य च सखा देवि सुग्रीवो नाम वानरः}
{राजा वानरमुख्यानां स त्वां कौशलमब्रवीत्} %||5-32-34||

\twolineshloka
{नित्यं स्मरति रामस्त्वां ससुग्रीवः सलक्ष्मणः}
{दिष्ट्या जीवसि वैदेहि राक्षसी वशमागता} %||5-32-35||

\twolineshloka
{नचिराद्द्रक्ष्यसे रामं लक्ष्मणं च महारथम्}
{मध्ये वानरकोटीनां सुग्रीवं चामितौजसं} %||5-32-36||

\twolineshloka
{अहं सुग्रीवसचिवो हनूमान्नाम वानरः}
{प्रविष्टो नगरीं लङ्कां लङ्घयित्वा महोदधिम्} %||5-32-37||

\twolineshloka
{कृत्वा मूर्ध्नि पदन्यासं रावणस्य दुरात्मनः}
{त्वां द्रष्टुमुपयातोऽहं समाश्रित्य पराक्रमम्} %||5-32-38||

\twolineshloka
{नाहमस्मि तथा देवि यथा मामवगच्छसि}
{विशङ्का त्यज्यतामेषा श्रद्धत्स्व वदतो मम} %||5-33-39||


{॥इत्यार्षे श्रीमद्रामायणे वाल्मीकीये आदिकाव्ये सुन्दरकाण्डे द्वात्रिंशः सर्गः॥}

\sect{त्रयस्त्रिंशः सर्गः}

\twolineshloka
{तां तु राम कथां श्रुत्वा वैदेही वानरर्षभात्}
{उवाच वचनं सान्त्वमिदं मधुरया गिरा} %||5-33-1||

\twolineshloka
{क्व ते रामेण संसर्गः कथं जानासि लक्ष्मणम्}
{वानराणां नराणां च कथमासीत्समागमः} %||5-33-2||

\twolineshloka
{यानि रामस्य लिङ्गानि लक्ष्मणस्य च वानर}
{तानि भूयः समाचक्ष्व न मां शोकः समाविशेत्} %||5-33-3||

\twolineshloka
{कीदृशं तस्य संस्थानं रूपं रामस्य कीदृशम्}
{कथमूरू कथं बाहू लक्ष्मणस्य च शंस मे} %||5-33-4||

\twolineshloka
{एवमुक्तस्तु वैदेह्या हनूमान्मारुतात्मजः}
{ततो रामं यथातत्त्वमाख्यातुमुपचक्रमे} %||5-33-5||

\twolineshloka
{जानन्ती बत दिष्ट्या मां वैदेहि परिपृच्छसि}
{भर्तुः कमलपत्राक्षि संख्यानं लक्ष्मणस्य च} %||5-33-6||

\twolineshloka
{यानि रामस्य चिह्नानि लक्ष्मणस्य च यानि वै}
{लक्षितानि विशालाक्षि वदतः शृणु तानि मे} %||5-33-7||

\twolineshloka
{रामः कमलपत्राक्षः सर्वभूतमनोहरः}
{रूपदाक्षिण्यसंपन्नः प्रसूतो जनकात्मजे} %||5-33-8||

\twolineshloka
{तेजसादित्यसंकाशः क्षमया पृथिवीसमः}
{बृहस्पतिसमो बुद्ध्या यशसा वासवोपमः} %||5-33-9||

\twolineshloka
{रक्षिता जीवलोकस्य स्वजनस्य च रक्षिता}
{रक्षिता स्वस्य वृत्तस्य धर्मस्य च परंतपः} %||5-33-10||

\twolineshloka
{रामो भामिनि लोकस्य चातुर्वर्ण्यस्य रक्षिता}
{मर्यादानां च लोकस्य कर्ता कारयिता च सः} %||5-33-11||

\twolineshloka
{अर्चिष्मानर्चितोऽत्यर्थं ब्रह्मचर्यव्रते स्थितः}
{साधूनामुपकारज्ञः प्रचारज्ञश्च कर्मणाम्} %||5-33-12||

\twolineshloka
{राजविद्याविनीतश्च ब्राह्मणानामुपासिता}
{श्रुतवाञ्शीलसंपन्नो विनीतश्च परंतपः} %||5-33-13||

\twolineshloka
{यजुर्वेदविनीतश्च वेदविद्भिः सुपूजितः}
{धनुर्वेदे च वेदे च वेदाङ्गेषु च निष्ठितः} %||5-33-14||

\twolineshloka
{विपुलांसो महाबाहुः कम्बुग्रीवः शुभाननः}
{गूढजत्रुः सुताम्राक्षो रामो देवि जनैः श्रुतः} %||5-33-15||

\twolineshloka
{दुन्दुभिस्वननिर्घोषः स्निग्धवर्णः प्रतापवान्}
{समः समविभक्ताङ्गो वर्णं श्यामं समाश्रितः} %||5-33-16||

\twolineshloka
{त्रिस्थिरस्त्रिप्रलम्बश्च त्रिसमस्त्रिषु चोन्नतः}
{त्रिवलीवांस्त्र्यवणतश्चतुर्व्यङ्गस्त्रिशीर्षवान्} %||5-33-17||

\twolineshloka
{चतुष्कलश्चतुर्लेखश्चतुष्किष्कुश्चतुःसमः}
{चतुर्दशसमद्वन्द्वश्चतुर्दष्टश्चतुर्गतिः} %||5-33-18||

\threelineshloka
{महौष्ठहनुनासश्च पञ्चस्निग्धोऽष्टवंशवान्}
{दशपद्मो दशबृहत्त्रिभिर्व्याप्तो द्विशुक्लवान्}
{षडुन्नतो नवतनुस्त्रिभिर्व्याप्नोति राघवः} %||5-33-19||

\twolineshloka
{सत्यधर्मपरः श्रीमान्संग्रहानुग्रहे रतः}
{देशकालविभागज्ञः सर्वलोकप्रियंवदः} %||5-33-20||

\twolineshloka
{भ्राता च तस्य द्वैमात्रः सौमित्रिरपराजितः}
{अनुरागेण रूपेण गुणैश्चैव तथाविधः} %||5-33-21||

\twolineshloka
{त्वामेव मार्गमाणो तौ विचरन्तौ वसुंधराम्}
{ददर्शतुर्मृगपतिं पूर्वजेनावरोपितम्} %||5-33-22||

\twolineshloka
{ऋश्यमूकस्य पृष्ठे तु बहुपादपसंकुले}
{भ्रातुर्भार्यार्तमासीनं सुग्रीवं प्रियदर्शनम्} %||5-33-23||

\twolineshloka
{वयं तु हरिराजं तं सुग्रीवं सत्यसंगरम्}
{परिचर्यामहे राज्यात्पूर्वजेनावरोपितम्} %||5-33-24||

\twolineshloka
{ततस्तौ चीरवसनौ धनुःप्रवरपाणिनौ}
{ऋश्यमूकस्य शैलस्य रम्यं देशमुपागतौ} %||5-33-25||

\twolineshloka
{स तौ दृष्ट्वा नरव्याघ्रौ धन्विनौ वानरर्षभः}
{अभिप्लुतो गिरेस्तस्य शिखरं भयमोहितः} %||5-33-26||

\twolineshloka
{ततः स शिखरे तस्मिन्वानरेन्द्रो व्यवस्थितः}
{तयोः समीपं मामेव प्रेषयामास सत्वरः} %||5-33-27||

\twolineshloka
{तावहं पुरुषव्याघ्रौ सुग्रीववचनात्प्रभू}
{रूपलक्षणसंपन्नौ कृताञ्जलिरुपस्थितः} %||5-33-28||

\twolineshloka
{तौ परिज्ञाततत्त्वार्थौ मया प्रीतिसमन्वितौ}
{पृष्ठमारोप्य तं देशं प्रापितौ पुरुषर्षभौ} %||5-33-29||

\twolineshloka
{निवेदितौ च तत्त्वेन सुग्रीवाय महात्मने}
{तयोरन्योन्यसंभाषाद्भृशं प्रीतिरजायत} %||5-33-30||

\twolineshloka
{तत्र तौ कीर्तिसंपन्नौ हरीश्वरनरेश्वरौ}
{परस्परकृताश्वासौ कथया पूर्ववृत्तया} %||5-33-31||

\twolineshloka
{तं ततः सान्त्वयामास सुग्रीवं लक्ष्मणाग्रजः}
{स्त्रीहेतोर्वालिना भ्रात्रा निरस्तमुरु तेजसा} %||5-33-32||

\twolineshloka
{ततस्त्वन्नाशजं शोकं रामस्याक्लिष्टकर्मणः}
{लक्ष्मणो वानरेन्द्राय सुग्रीवाय न्यवेदयत्} %||5-33-33||

\twolineshloka
{स श्रुत्वा वानरेन्द्रस्तु लक्ष्मणेनेरितं वचः}
{तदासीन्निष्प्रभोऽत्यर्थं ग्रहग्रस्त इवांशुमान्} %||5-33-34||

\twolineshloka
{ततस्त्वद्गात्रशोभीनि रक्षसा ह्रियमाणया}
{यान्याभरणजालानि पातितानि महीतले} %||5-33-35||

\twolineshloka
{तानि सर्वाणि रामाय आनीय हरियूथपाः}
{संहृष्टा दर्शयामासुर्गतिं तु न विदुस्तव} %||5-33-36||

\twolineshloka
{तानि रामाय दत्तानि मयैवोपहृतानि च}
{स्वनवन्त्यवकीर्णन्ति तस्मिन्विहतचेतसि} %||5-33-37||

\twolineshloka
{तान्यङ्के दर्शनीयानि कृत्वा बहुविधं ततः}
{तेन देवप्रकाशेन देवेन परिदेवितम्} %||5-33-38||

\twolineshloka
{पश्यतस्तस्या रुदतस्ताम्यतश्च पुनः पुनः}
{प्रादीपयन्दाशरथेस्तानि शोकहुताशनम्} %||5-33-39||

\twolineshloka
{शयितं च चिरं तेन दुःखार्तेन महात्मना}
{मयापि विविधैर्वाक्यैः कृच्छ्रादुत्थापितः पुनः} %||5-33-40||

\twolineshloka
{तानि दृष्ट्वा महार्हाणि दर्शयित्वा मुहुर्मुहुः}
{राघवः सहसौमित्रिः सुग्रीवे स न्यवेदयत्} %||5-33-41||

\twolineshloka
{स तवादर्शनादार्ये राघवः परितप्यते}
{महता ज्वलता नित्यमग्निनेवाग्निपर्वतः} %||5-33-42||

\twolineshloka
{त्वत्कृते तमनिद्रा च शोकश्चिन्ता च राघवम्}
{तापयन्ति महात्मानमग्न्यगारमिवाग्नयः} %||5-33-43||

\twolineshloka
{तवादर्शनशोकेन राघवः प्रविचाल्यते}
{महता भूमिकम्पेन महानिव शिलोच्चयः} %||5-33-44||

\twolineshloka
{कानानानि सुरम्याणि नदीप्रस्रवणानि च}
{चरन्न रतिमाप्नोति त्वमपश्यन्नृपात्मजे} %||5-33-45||

\twolineshloka
{स त्वां मनुजशार्दूलः क्षिप्रं प्राप्स्यति राघवः}
{समित्रबान्धवं हत्वा रावणं जनकात्मजे} %||5-33-46||

\twolineshloka
{सहितौ रामसुग्रीवावुभावकुरुतां तदा}
{समयं वालिनं हन्तुं तव चान्वेषणं तथा} %||5-33-47||

\twolineshloka
{ततो निहत्य तरसा रामो वालिनमाहवे}
{सर्वर्क्षहरिसंघानां सुग्रीवमकरोत्पतिम्} %||5-33-48||

\twolineshloka
{रामसुग्रीवयोरैक्यं देव्येवं समजायत}
{हनूमन्तं च मां विद्धि तयोर्दूतमिहागतम्} %||5-33-49||

\twolineshloka
{स्वराज्यं प्राप्य सुग्रीवः समनीय महाहरीन्}
{त्वदर्थं प्रेषयामास दिशो दश महाबलान्} %||5-33-50||

\twolineshloka
{आदिष्टा वानरेन्द्रेण सुग्रीवेण महौजसः}
{अद्रिराजप्रतीकाशाः सर्वतः प्रस्थिता महीम्} %||5-33-51||

\twolineshloka
{अङ्गदो नाम लक्ष्मीवान्वालिसूनुर्महाबलः}
{प्रस्थितः कपिशार्दूलस्त्रिभागबलसंवृतः} %||5-33-52||

\twolineshloka
{तेषां नो विप्रनष्टानां विन्ध्ये पर्वतसत्तमे}
{भृशं शोकपरीतनामहोरात्रगणा गताः} %||5-33-53||

\twolineshloka
{ते वयं कार्यनैराश्यात्कालस्यातिक्रमेण च}
{भयाच्च कपिराजस्य प्राणांस्त्यक्तुं व्यवस्थिताः} %||5-33-54||

\twolineshloka
{विचित्य वनदुर्गाणि गिरिप्रस्रवणानि च}
{अनासाद्य पदं देव्याः प्राणांस्त्यक्तुं व्यवस्थिताः} %||5-33-55||

\threelineshloka
{भृशं शोकार्णवे मग्नः पर्यदेवयदङ्गदः}
{तव नाशं च वैदेहि वालिनश्च तथा वधम्}
{प्रायोपवेशमस्माकं मरणं च जटायुषः} %||5-33-56||

\twolineshloka
{तेषां नः स्वामिसंदेशान्निराशानां मुमूर्षताम्}
{कार्यहेतोरिवायातः शकुनिर्वीर्यवान्महान्} %||5-33-57||

\twolineshloka
{गृध्रराजस्य सोदर्यः संपातिर्नाम गृध्रराट्}
{श्रुत्वा भ्रातृवधं कोपादिदं वचनमब्रवीत्} %||5-33-58||

\twolineshloka
{यवीयान्केन मे भ्राता हतः क्व च विनाशितः}
{एतदाख्यातुमिच्छामि भवद्भिर्वानरोत्तमाः} %||5-33-59||

\twolineshloka
{अङ्गदोऽकथयत्तस्य जनस्थाने महद्वधम्}
{रक्षसा भीमरूपेण त्वामुद्दिश्य यथातथम्} %||5-33-60||

\twolineshloka
{जटायोस्तु वधं श्रुत्वा दुःखितः सोऽरुणात्मजः}
{त्वामाह स वरारोहे वसन्तीं रावणालये} %||5-33-61||

\threelineshloka
{तस्य तद्वचनं श्रुत्वा संपातेः प्रीतिवर्धनम्}
{अङ्गदप्रमुखाः सर्वे ततः संप्रस्थिता वयम्}
{त्वद्दर्शनकृतोत्साहा हृष्टास्तुष्टाः प्लवंगमाः} %||5-33-62||

\twolineshloka
{अथाहं हरिसैन्यस्य सागरं दृश्य सीदतः}
{व्यवधूय भयं तीव्रं योजनानां शतं प्लुतः} %||5-33-63||

\twolineshloka
{लङ्का चापि मया रात्रौ प्रविष्टा राक्षसाकुला}
{रावणश्च मया दृष्टस्त्वं च शोकनिपीडिता} %||5-33-64||

\twolineshloka
{एतत्ते सर्वमाख्यातं यथावृत्तमनिन्दिते}
{अभिभाषस्व मां देवि दूतो दाशरथेरहम्} %||5-33-65||

\twolineshloka
{त्वं मां रामकृतोद्योगं त्वन्निमित्तमिहागतम्}
{सुग्रीव सचिवं देवि बुध्यस्व पवनात्मजम्} %||5-33-66||

\twolineshloka
{कुशली तव काकुत्स्थः सर्वशस्त्रभृतां वरः}
{गुरोराराधने युक्तो लक्ष्मणश्च सुलक्षणः} %||5-33-67||

\twolineshloka
{तस्य वीर्यवतो देवि भर्तुस्तव हिते रतः}
{अहमेकस्तु संप्राप्तः सुग्रीववचनादिह} %||5-33-68||

\twolineshloka
{मयेयमसहायेन चरता कामरूपिणा}
{दक्षिणा दिगनुक्रान्ता त्वन्मार्गविचयैषिणा} %||5-33-69||

\twolineshloka
{दिष्ट्याहं हरिसैन्यानां त्वन्नाशमनुशोचताम्}
{अपनेष्यामि संतापं तवाभिगमशंसनात्} %||5-33-70||

\twolineshloka
{दिष्ट्या हि न मम व्यर्थं देवि सागरलङ्घनम्}
{प्राप्स्याम्यहमिदं दिष्ट्या त्वद्दर्शनकृतं यशः} %||5-33-71||

\twolineshloka
{राघवश्च महावीर्यः क्षिप्रं त्वामभिपत्स्यते}
{समित्रबान्धवं हत्वा रावणं राक्षसाधिपम्} %||5-33-72||

\twolineshloka
{कौरजो नाम वैदेहि गिरीणामुत्तमो गिरिः}
{ततो गच्छति गोकर्णं पर्वतं केसरी हरिः} %||5-33-73||

\twolineshloka
{स च देवर्षिभिर्दृष्टः पिता मम महाकपिः}
{तीर्थे नदीपतेः पुण्ये शम्बसादनमुद्धरत्} %||5-33-74||

\threelineshloka
{तस्याहं हरिणः क्षेत्रे जातो वातेन मैथिलि}
{हनूमानिति विख्यातो लोके स्वेनैव कर्मणा}
{विश्वासार्थं तु वैदेहि भर्तुरुक्ता मया गुणाः} %||5-33-75||

\twolineshloka
{एवं विश्वासिता सीता हेतुभिः शोककर्शिता}
{उपपन्नैरभिज्ञानैर्दूतं तमवगच्छति} %||5-33-76||

\twolineshloka
{अतुलं च गता हर्षं प्रहर्षेण तु जानकी}
{नेत्राभ्यां वक्रपक्ष्माभ्यां मुमोचानन्दजं जलम्} %||5-33-77||

\threelineshloka
{चारु तच्चाननं तस्यास्ताम्रशुक्लायतेक्षणम्}
{अशोभत विशालाक्ष्या राहुमुक्त इवोडुराट्}
{हनूमन्तं कपिं व्यक्तं मन्यते नान्यथेति सा} %||5-33-78||

{अथोवाच हनूमांस्तामुत्तरं प्रियदर्शनाम्॥\devanumber{79}॥} %||5-33-79|| (Check)

\addtocounter{shlokacount}{1}

\fourlineindentedshloka
{हतेऽसुरे संयति शम्बसादने}
{ कपिप्रवीरेण महर्षिचोदनात्}
{ततोऽस्मि वायुप्रभवो हि मैथिलि}
{ प्रभावतस्तत्प्रतिमश्च वानरः} %||5-34-80||


{॥इत्यार्षे श्रीमद्रामायणे वाल्मीकीये आदिकाव्ये सुन्दरकाण्डे त्रयस्त्रिंशः सर्गः॥}

\sect{चतुस्त्रिंशः सर्गः}

\twolineshloka
{भूय एव महातेजा हनूमान्मारुतात्मजः}
{अब्रवीत्प्रश्रितं वाक्यं सीताप्रत्ययकारणात्} %||5-34-1||

\threelineshloka
{वानरोऽहं महाभागे दूतो रामस्य धीमतः}
{रामनामाङ्कितं चेदं पश्य देव्यङ्गुलीयकम्}
{समाश्वसिहि भद्रं ते क्षीणदुःखफला ह्यसि} %||5-34-2||

\twolineshloka
{गृहीत्वा प्रेक्षमाणा सा भर्तुः करविभूषणम्}
{भर्तारमिव संप्राप्ता जानकी मुदिताभवत्} %||5-34-3||

\twolineshloka
{चारु तद्वदनं तस्यास्ताम्रशुक्लायतेक्षणम्}
{बभूव प्रहर्षोदग्रं राहुमुक्त इवोडुराट्} %||5-34-4||

\twolineshloka
{ततः सा ह्रीमती बाला भर्तुः संदेशहर्षिता}
{परितुट्षा प्रियं श्रुत्वा प्राशंसत महाकपिम्} %||5-34-5||

\twolineshloka
{विक्रान्तस्त्वं समर्थस्त्वं प्राज्ञस्त्वं वानरोत्तम}
{येनेदं राक्षसपदं त्वयैकेन प्रधर्षितम्} %||5-34-6||

\twolineshloka
{शतयोजनविस्तीर्णः सागरो मकरालयः}
{विक्रमश्लाघनीयेन क्रमता गोष्पदीकृतः} %||5-34-7||

\twolineshloka
{न हि त्वां प्राकृतं मन्ये वनरं वनरर्षभ}
{यस्य ते नास्ति संत्रासो रावणान्नापि संभ्रमः} %||5-34-8||

\twolineshloka
{अर्हसे च कपिश्रेष्ठ मया समभिभाषितुम्}
{यद्यसि प्रेषितस्तेन रामेण विदितात्मना} %||5-34-9||

\twolineshloka
{प्रेषयिष्यति दुर्धर्षो रामो न ह्यपरीक्षितम्}
{पराक्रममविज्ञाय मत्सकाशं विशेषतः} %||5-34-10||

\twolineshloka
{दिष्ट्या च कुशली रामो धर्मात्मा धर्मवत्सलः}
{लक्ष्मणश्च महातेजाः सुमित्रानन्दवर्धनः} %||5-34-11||

\twolineshloka
{कुशली यदि काकुत्स्थः किं नु सागरमेखलाम्}
{महीं दहति कोपेन युगान्ताग्निरिवोत्थितः} %||5-34-12||

\twolineshloka
{अथ वा शक्तिमन्तौ तौ सुराणामपि निग्रहे}
{ममैव तु न दुःखानामस्ति मन्ये विपर्ययः} %||5-34-13||

\twolineshloka
{कच्चिच्च व्यथते रामः कच्चिन्न परिपत्यते}
{उत्तराणि च कार्याणि कुरुते पुरुषोत्तमः} %||5-34-14||

\twolineshloka
{कच्चिन्न दीनः संभ्रान्तः कार्येषु च न मुह्यति}
{कच्चिन्पुरुषकार्याणि कुरुते नृपतेः सुतः} %||5-34-15||

\twolineshloka
{द्विविधं त्रिविधोपायमुपायमपि सेवते}
{विजिगीषुः सुहृत्कच्चिन्मित्रेषु च परंतपः} %||5-34-16||

\twolineshloka
{कच्चिन्मित्राणि लभते मित्रैश्चाप्यभिगम्यते}
{कच्चित्कल्याणमित्रश्च मित्रैश्चापि पुरस्कृतः} %||5-34-17||

\twolineshloka
{कच्चिदाशास्ति देवानां प्रसादं पार्थिवात्मजः}
{कच्चित्पुरुषकारं च दैवं च प्रतिपद्यते} %||5-34-18||

\twolineshloka
{कच्चिन्न विगतस्नेहो विवासान्मयि राघवः}
{कच्चिन्मां व्यसनादस्मान्मोक्षयिष्यति वानरः} %||5-34-19||

\twolineshloka
{सुखानामुचितो नित्यमसुखानामनूचितः}
{दुःखमुत्तरमासाद्य कच्चिद्रामो न सीदति} %||5-34-20||

\twolineshloka
{कौसल्यायास्तथा कच्चित्सुमित्रायास्तथैव च}
{अभीक्ष्णं श्रूयते कच्चित्कुशलं भरतस्य च} %||5-34-21||

\twolineshloka
{मन्निमित्तेन मानार्हः कच्चिच्छोकेन राघवः}
{कच्चिन्नान्यमना रामः कच्चिन्मां तारयिष्यति} %||5-34-22||

\twolineshloka
{कच्चिदक्षौहिणीं भीमां भरतो भ्रातृवत्सलः}
{ध्वजिनीं मन्त्रिभिर्गुप्तां प्रेषयिष्यति मत्कृते} %||5-34-23||

\twolineshloka
{वानराधिपतिः श्रीमान्सुग्रीवः कच्चिदेष्यति}
{मत्कृते हरिभिर्वीरैर्वृतो दन्तनखायुधैः} %||5-34-24||

\twolineshloka
{कच्चिच्च लक्ष्मणः शूरः सुमित्रानन्दवर्धनः}
{अस्त्रविच्छरजालेन राक्षसान्विधमिष्यति} %||5-34-25||

\twolineshloka
{रौद्रेण कच्चिदस्त्रेण रामेण निहतं रणे}
{द्रक्ष्याम्यल्पेन कालेन रावणं ससुहृज्जनम्} %||5-34-26||

\fourlineindentedshloka
{कच्चिन्न तद्धेमसमानवर्णं}
{ तस्याननं पद्मसमानगन्धि}
{मया विना शुष्यति शोकदीनं}
{ जलक्षये पद्ममिवातपेन} %||5-34-27||

\fourlineindentedshloka
{धर्मापदेशात्त्यजतश्च राज्यां}
{ मां चाप्यरण्यं नयतः पदातिम्}
{नासीद्व्यथा यस्य न भीर्न शोकः}
{ कच्चित्स धैर्यं हृदये करोति} %||5-34-28||

\fourlineindentedshloka
{न चास्य माता न पिता न चान्यः}
{ स्नेहाद्विशिष्टोऽस्ति मया समो वा}
{तावद्ध्यहं दूतजिजीविषेयं}
{ यावत्प्रवृत्तिं शृणुयां प्रियस्य} %||5-34-29||

\fourlineindentedshloka
{इतीव देवी वचनं महार्थं}
{ तं वानरेन्द्रं मधुरार्थमुक्त्वा}
{श्रोतुं पुनस्तस्य वचोऽभिरामं}
{ रामार्थयुक्तं विरराम रामा} %||5-34-30||

\twolineshloka
{सीताया वचनं श्रुत्वा मारुतिर्भीमविक्रमः}
{शिरस्यञ्जलिमाधाय वाक्यमुत्तरमब्रवीत्} %||5-34-31||

\twolineshloka
{न त्वामिहस्थां जानीते रामः कमललोचनः}
{श्रुत्वैव तु वचो मह्यं क्षिप्रमेष्यति राघवः} %||5-34-32||

\threelineshloka
{चमूं प्रकर्षन्महतीं हर्यृष्कगणसंकुलाम्}
{विष्टम्भयित्वा बाणौघैरक्षोभ्यं वरुणालयम्}
{करिष्यति पुरीं लङ्कां काकुत्स्थः शान्तराक्षसाम्} %||5-34-33||

\twolineshloka
{तत्र यद्यन्तरा मृत्युर्यदि देवाः सहासुराः}
{स्थास्यन्ति पथि रामस्य स तानपि वधिष्यति} %||5-34-34||

\twolineshloka
{तवादर्शनजेनार्ये शोकेन स परिप्लुतः}
{न शर्म लभते रामः सिंहार्दित इव द्विपः} %||5-34-35||

\twolineshloka
{दर्दरेण च ते देवि शपे मूलफलेन च}
{मलयेन च विन्ध्येन मेरुणा मन्दरेण च} %||5-34-36||

\twolineshloka
{यथा सुनयनं वल्गु बिम्बौष्ठं चारुकुण्डलम्}
{मुखं द्रक्ष्यसि रामस्य पूर्णचन्द्रमिवोदितम्} %||5-34-37||

\twolineshloka
{क्षिप्रं द्रक्ष्यसि वैदेहि रामं प्रस्रवणे गिरौ}
{शतक्रतुमिवासीनं नाकपृष्ठस्य मूर्धनि} %||5-34-38||

\twolineshloka
{न मांसं राघवो भुङ्क्ते न चापि मधुसेवते}
{वन्यं सुविहितं नित्यं भक्तमश्नाति पञ्चमम्} %||5-34-39||

\twolineshloka
{नैव दंशान्न मशकान्न कीटान्न सरीसृपान्}
{राघवोऽपनयेद्गत्रात्त्वद्गतेनान्तरात्मना} %||5-34-40||

\twolineshloka
{नित्यं ध्यानपरो रामो नित्यं शोकपरायणः}
{नान्यच्चिन्तयते किंचित्स तु कामवशं गतः} %||5-34-41||

\twolineshloka
{अनिद्रः सततं रामः सुप्तोऽपि च नरोत्तमः}
{सीतेति मधुरां वाणीं व्याहरन्प्रतिबुध्यते} %||5-34-42||

\twolineshloka
{दृष्ट्वा फलं वा पुष्पं वा यच्चान्यत्स्त्रीमनोहरम्}
{बहुशो हा प्रियेत्येवं श्वसंस्त्वामभिभाषते} %||5-34-43||

\fourlineindentedshloka
{स देवि नित्यं परितप्यमान}
{स्त्वामेव सीतेत्यभिभाषमाणः}
{धृतव्रतो राजसुतो महात्मा}
{ तवैव लाभाय कृतप्रयत्नः} %||5-34-44||

\fourlineindentedshloka
{सा रामसंकीर्तनवीतशोका}
{ रामस्य शोकेन समानशोका}
{शरन्मुखेनाम्बुदशेषचन्द्रा}
{ निशेव वैदेहसुता बभूव} %||5-35-45||


{॥इत्यार्षे श्रीमद्रामायणे वाल्मीकीये आदिकाव्ये सुन्दरकाण्डे चतुस्त्रिंशः सर्गः॥}

\sect{पञ्चत्रिंशः सर्गः}

\twolineshloka
{सीता तद्वचनं श्रुत्वा पूर्णचन्द्रनिभानना}
{हनूमन्तमुवाचेदं धर्मार्थसहितं वचः} %||5-35-1||

\twolineshloka
{अमृतं विषसंसृष्टं त्वया वानरभाषितम्}
{यच्च नान्यमना रामो यच्च शोकपरायणः} %||5-35-2||

\twolineshloka
{ऐश्वर्ये वा सुविस्तीर्णे व्यसने वा सुदारुणे}
{रज्ज्वेव पुरुषं बद्ध्वा कृतान्तः परिकर्षति} %||5-35-3||

\twolineshloka
{विधिर्नूनमसंहार्यः प्राणिनां प्लवगोत्तम}
{सौमित्रिं मां च रामं च व्यसनैः पश्य मोहितान्} %||5-35-4||

\twolineshloka
{शोकस्यास्य कदा पारं राघवोऽधिगमिष्यति}
{प्लवमानः परिश्रान्तो हतनौः सागरे यथा} %||5-35-5||

\twolineshloka
{राक्षसानां क्षयं कृत्वा सूदयित्वा च रावणम्}
{लङ्कामुन्मूलितां कृत्वा कदा द्रक्ष्यति मां पतिः} %||5-35-6||

\twolineshloka
{स वाच्यः संत्वरस्वेति यावदेव न पूर्यते}
{अयं संवत्सरः कालस्तावद्धि मम जीवितम्} %||5-35-7||

\twolineshloka
{वर्तते दशमो मासो द्वौ तु शेषौ प्लवंगम}
{रावणेन नृशंसेन समयो यः कृतो मम} %||5-35-8||

\twolineshloka
{विभीषणेन च भ्रात्रा मम निर्यातनं प्रति}
{अनुनीतः प्रयत्नेन न च तत्कुरुते मतिम्} %||5-35-9||

\twolineshloka
{मम प्रतिप्रदानं हि रावणस्य न रोचते}
{रावणं मार्गते संख्ये मृत्युः कालवशं गतम्} %||5-35-10||

\twolineshloka
{ज्येष्ठा कन्यानला नम विभीषणसुता कपे}
{तया ममैतदाख्यातं मात्रा प्रहितया स्वयम्} %||5-35-11||

\twolineshloka
{अविन्ध्यो नाम मेधावी विद्वान्राक्षसपुंगवः}
{धृतिमाञ्शीलवान्वृद्धो रावणस्य सुसंमतः} %||5-35-12||

\twolineshloka
{रामात्क्षयमनुप्राप्तं रक्षसां प्रत्यचोदयत्}
{न च तस्यापि दुष्टात्मा शृणोति वचनं हितम्} %||5-35-13||

\twolineshloka
{आशंसेति हरिश्रेष्ठ क्षिप्रं मां प्राप्स्यते पतिः}
{अन्तरात्मा हि मे शुद्धस्तस्मिंश्च बहवो गुणाः} %||5-35-14||

\twolineshloka
{उत्साहः पौरुषं सत्त्वमानृशंस्यं कृतज्ञता}
{विक्रमश्च प्रभावश्च सन्ति वानरराघवे} %||5-35-15||

\twolineshloka
{चतुर्दशसहस्राणि राक्षसानां जघान यः}
{जनस्थाने विना भ्रात्रा शत्रुः कस्तस्य नोद्विजेत्} %||5-35-16||

\twolineshloka
{न स शक्यस्तुलयितुं व्यसनैः पुरुषर्षभः}
{अहं तस्यानुभावज्ञा शक्रस्येव पुलोमजा} %||5-35-17||

\twolineshloka
{शरजालांशुमाञ्शूरः कपे रामदिवाकरः}
{शत्रुरक्षोमयं तोयमुपशोषं नयिष्यति} %||5-35-18||

\twolineshloka
{इति संजल्पमानां तां रामार्थे शोककर्शिताम्}
{अश्रुसंपूर्णवदनामुवाच हनुमान्कपिः} %||5-35-19||

\twolineshloka
{श्रुत्वैव तु वचो मह्यं क्षिप्रमेष्यति राघवः}
{चमूं प्रकर्षन्महतीं हर्यृक्षगणसंकुलाम्} %||5-35-20||

\twolineshloka
{अथ वा मोचयिष्यामि तामद्यैव हि राक्षसात्}
{अस्माद्दुःखादुपारोह मम पृष्ठमनिन्दिते} %||5-35-21||

\twolineshloka
{त्वं हि पृष्ठगतां कृत्वा संतरिष्यामि सागरम्}
{शक्तिरस्ति हि मे वोढुं लङ्कामपि सरावणाम्} %||5-35-22||

\twolineshloka
{अहं प्रस्रवणस्थाय राघवायाद्य मैथिलि}
{प्रापयिष्यामि शक्राय हव्यं हुतमिवानलः} %||5-35-23||

\twolineshloka
{द्रक्ष्यस्यद्यैव वैदेहि राघवं सहलक्ष्मणम्}
{व्यवसाय समायुक्तं विष्णुं दैत्यवधे यथा} %||5-35-24||

\twolineshloka
{त्वद्दर्शनकृतोत्साहमाश्रमस्थं महाबलम्}
{पुरंदरमिवासीनं नागराजस्य मूर्धनि} %||5-35-25||

\twolineshloka
{पृष्ठमारोह मे देवि मा विकाङ्क्षस्व शोभने}
{योगमन्विच्छ रामेण शशाङ्केनेव रोहिणी} %||5-35-26||

\twolineshloka
{कथयन्तीव चन्द्रेण सूर्येणेव सुवर्चला}
{मत्पृष्ठमधिरुह्य त्वं तराकाशमहार्णवम्} %||5-35-27||

\twolineshloka
{न हि मे संप्रयातस्य त्वामितो नयतोऽङ्गने}
{अनुगन्तुं गतिं शक्ताः सर्वे लङ्कानिवासिनः} %||5-35-28||

\twolineshloka
{यथैवाहमिह प्राप्तस्तथैवाहमसंशयम्}
{यास्यामि पश्य वैदेहि त्वामुद्यम्य विहायसं} %||5-35-29||

\twolineshloka
{मैथिली तु हरिश्रेष्ठाच्छ्रुत्वा वचनमद्भुतम्}
{हर्षविस्मितसर्वाङ्गी हनूमन्तमथाब्रवीत्} %||5-35-30||

\twolineshloka
{हनूमन्दूरमध्वनं कथं मां वोढुमिच्छसि}
{तदेव खलु ते मन्ये कपित्वं हरियूथप} %||5-35-31||

\twolineshloka
{कथं वाल्पशरीरस्त्वं मामितो नेतुमिच्छसि}
{सकाशं मानवेन्द्रस्य भर्तुर्मे प्लवगर्षभ} %||5-35-32||

\twolineshloka
{सीताया वचनं श्रुत्वा हनूमान्मारुतात्मजः}
{चिन्तयामास लक्ष्मीवान्नवं परिभवं कृतम्} %||5-35-33||

\twolineshloka
{न मे जानाति सत्त्वं वा प्रभावं वासितेक्षणा}
{तस्मात्पश्यतु वैदेही यद्रूपं मम कामतः} %||5-35-34||

\twolineshloka
{इति संचिन्त्य हनुमांस्तदा प्लवगसत्तमः}
{दर्शयामास वैदेह्याः स्वरूपमरिमर्दनः} %||5-35-35||

\twolineshloka
{स तस्मात्पादपाद्धीमानाप्लुत्य प्लवगर्षभः}
{ततो वर्धितुमारेभे सीताप्रत्ययकारणात्} %||5-35-36||

\twolineshloka
{मेरुमन्दारसंकाशो बभौ दीप्तानलप्रभः}
{अग्रतो व्यवतस्थे च सीताया वानरर्षभः} %||5-35-37||

\twolineshloka
{हरिः पर्वतसंकाशस्ताम्रवक्त्रो महाबलः}
{वज्रदंष्ट्रनखो भीमो वैदेहीमिदमब्रवीत्} %||5-35-38||

\twolineshloka
{सपर्वतवनोद्देशां साट्टप्राकारतोरणाम्}
{लङ्कामिमां सनथां वा नयितुं शक्तिरस्ति मे} %||5-35-39||

\twolineshloka
{तदवस्थाप्य तां बुद्धिरलं देवि विकाङ्क्षया}
{विशोकं कुरु वैदेहि राघवं सहलक्ष्मणम्} %||5-35-40||

\twolineshloka
{तं दृष्ट्वाचलसंकाशमुवाच जनकात्मजा}
{पद्मपत्रविशालाक्षी मारुतस्यौरसं सुतम्} %||5-35-41||

\twolineshloka
{तव सत्त्वं बलं चैव विजानामि महाकपे}
{वायोरिव गतिं चापि तेजश्चाग्निरिवाद्भुतम्} %||5-35-42||

\twolineshloka
{प्राकृतोऽन्यः कथं चेमां भूमिमागन्तुमर्हति}
{उदधेरप्रमेयस्य पारं वानरपुंगव} %||5-35-43||

\twolineshloka
{जानामि गमने शक्तिं नयने चापि ते मम}
{अवश्यं साम्प्रधार्याशु कार्यसिद्धिरिहात्मनः} %||5-35-44||

\twolineshloka
{अयुक्तं तु कपिश्रेष्ठ मया गन्तुं त्वया सह}
{वायुवेगसवेगस्य वेगो मां मोहयेत्तव} %||5-35-45||

\twolineshloka
{अहमाकाशमासक्ता उपर्युपरि सागरम्}
{प्रपतेयं हि ते पृष्ठाद्भयाद्वेगेन गच्छतः} %||5-35-46||

\twolineshloka
{पतिता सागरे चाहं तिमिनक्रझषाकुले}
{भयेयमाशु विवशा यादसामन्नमुत्तमम्} %||5-35-47||

\twolineshloka
{न च शक्ष्ये त्वया सार्धं गन्तुं शत्रुविनाशन}
{कलत्रवति संदेहस्त्वय्यपि स्यादसंशयम्} %||5-35-48||

\twolineshloka
{ह्रियमाणां तु मां दृष्ट्वा राक्षसा भीमविक्रमाः}
{अनुगच्छेयुरादिष्टा रावणेन दुरात्मना} %||5-35-49||

\twolineshloka
{तैस्त्वं परिवृतः शूरैः शूलमुद्गर पाणिभिः}
{भवेस्त्वं संशयं प्राप्तो मया वीर कलत्रवान्} %||5-35-50||

\twolineshloka
{सायुधा बहवो व्योम्नि राक्षसास्त्वं निरायुधः}
{कथं शक्ष्यसि संयातुं मां चैव परिरक्षितुम्} %||5-35-51||

\twolineshloka
{युध्यमानस्य रक्षोभिस्ततस्तैः क्रूरकर्मभिः}
{प्रपतेयं हि ते पृष्ठद्भयार्ता कपिसत्तम} %||5-35-52||

\twolineshloka
{अथ रक्षांसि भीमानि महान्ति बलवन्ति च}
{कथंचित्साम्पराये त्वां जयेयुः कपिसत्तम} %||5-35-53||

\twolineshloka
{अथ वा युध्यमानस्य पतेयं विमुखस्य ते}
{पतितां च गृहीत्वा मां नयेयुः पापराक्षसाः} %||5-35-54||

\twolineshloka
{मां वा हरेयुस्त्वद्धस्ताद्विशसेयुरथापि वा}
{अव्यवस्थौ हि दृश्येते युद्धे जयपराजयौ} %||5-35-55||

\twolineshloka
{अहं वापि विपद्येयं रक्षोभिरभितर्जिता}
{त्वत्प्रयत्नो हरिश्रेष्ठ भवेन्निष्फल एव तु} %||5-35-56||

\twolineshloka
{कामं त्वमपि पर्याप्तो निहन्तुं सर्वराक्षसान्}
{राघवस्य यशो हीयेत्त्वया शस्तैस्तु राक्षसैः} %||5-35-57||

\twolineshloka
{अथ वादाय रक्षांसि न्यस्येयुः संवृते हि माम्}
{यत्र ते नाभिजानीयुर्हरयो नापि राघवः} %||5-35-58||

\twolineshloka
{आरम्भस्तु मदर्थोऽयं ततस्तव निरर्थकः}
{त्वया हि सह रामस्य महानागमने गुणः} %||5-35-59||

\twolineshloka
{मयि जीवितमायत्तं राघवस्य महात्मनः}
{भ्रातॄणां च महाबाहो तव राजकुलस्य च} %||5-35-60||

\twolineshloka
{तौ निराशौ मदर्थे तु शोकसंतापकर्शितौ}
{सह सर्वर्क्षहरिभिस्त्यक्ष्यतः प्राणसंग्रहम्} %||5-35-61||

\twolineshloka
{भर्तुर्भक्तिं पुरस्कृत्य रामादन्यस्य वानर}
{नाहं स्प्रष्टुं पदा गात्रमिच्छेयं वानरोत्तम} %||5-35-62||

\twolineshloka
{यदहं गात्रसंस्पर्शं रावणस्य गता बलात्}
{अनीशा किं करिष्यामि विनाथा विवशा सती} %||5-35-63||

\twolineshloka
{यदि रामो दशग्रीवमिह हत्वा सराक्षसं}
{मामितो गृह्य गच्छेत तत्तस्य सदृशं भवेत्} %||5-35-64||

\fourlineindentedshloka
{श्रुता हि दृष्टाश्च मया पराक्रमा}
{ महात्मनस्तस्य रणावमर्दिनः}
{न देवगन्धर्वभुजंगराक्षसा}
{ भवन्ति रामेण समा हि संयुगे} %||5-35-65||

\fourlineindentedshloka
{समीक्ष्य तं संयति चित्रकार्मुकं}
{ महाबलं वासवतुल्यविक्रमम्}
{सलक्ष्मणं को विषहेत राघवं}
{ हुताशनं दीप्तमिवानिलेरितम्} %||5-35-66||

\fourlineindentedshloka
{सलक्ष्मणं राघवमाजिमर्दनं}
{ दिशागजं मत्तमिव व्यवस्थितम्}
{सहेत को वानरमुख्य संयुगे}
{ युगान्तसूर्यप्रतिमं शरार्चिषम्} %||5-35-67||

\fourlineindentedshloka
{स मे हरिश्रेष्ठ सलक्ष्मणं पतिं}
{ सयूथपं क्षिप्रमिहोपपादय}
{चिराय रामं प्रति शोककर्शितां}
{ कुरुष्व मां वानरमुख्य हर्षिताम्} %||5-36-68||


{॥इत्यार्षे श्रीमद्रामायणे वाल्मीकीये आदिकाव्ये सुन्दरकाण्डे पञ्चत्रिंशः सर्गः॥}

\sect{षट्त्रिंशः सर्गः}

\twolineshloka
{ततः स कपिशार्दूलस्तेन वाक्येन हर्षितः}
{सीतामुवाच तच्छ्रुत्वा वाक्यं वाक्यविशारदः} %||5-36-1||

\twolineshloka
{युक्तरूपं त्वया देवि भाषितं शुभदर्शने}
{सदृशं स्त्रीस्वभावस्य साध्वीनां विनयस्य च} %||5-36-2||

\twolineshloka
{स्त्रीत्वं न तु समर्थं हि सागरं व्यतिवर्तितुम्}
{मामधिष्ठाय विस्तीर्णं शतयोजनमायतम्} %||5-36-3||

\twolineshloka
{द्वितीयं कारणं यच्च ब्रवीषि विनयान्विते}
{रामादन्यस्य नार्हामि संस्पर्शमिति जानकि} %||5-36-4||

\twolineshloka
{एतत्ते देवि सदृशं पत्न्यास्तस्य महात्मनः}
{का ह्यन्या त्वामृते देवि ब्रूयाद्वचनमीदृशम्} %||5-36-5||

\twolineshloka
{श्रोष्यते चैव काकुत्स्थः सर्वं निरवशेषतः}
{चेष्टितं यत्त्वया देवि भाषितं मम चाग्रतः} %||5-36-6||

\twolineshloka
{कारणैर्बहुभिर्देवि राम प्रियचिकीर्षया}
{स्नेहप्रस्कन्नमनसा मयैतत्समुदीरितम्} %||5-36-7||

\twolineshloka
{लङ्काया दुष्प्रवेशत्वाद्दुस्तरत्वान्महोदधेः}
{सामर्थ्यादात्मनश्चैव मयैतत्समुदाहृतम्} %||5-36-8||

\twolineshloka
{इच्छामि त्वां समानेतुमद्यैव रघुबन्धुना}
{गुरुस्नेहेन भक्त्या च नान्यथा तदुदाहृतम्} %||5-36-9||

\twolineshloka
{यदि नोत्सहसे यातुं मया सार्धमनिन्दिते}
{अभिज्ञानं प्रयच्छ त्वं जानीयाद्राघवो हि यत्} %||5-36-10||

\twolineshloka
{एवमुक्ता हनुमता सीता सुरसुतोपमा}
{उवाच वचनं मन्दं बाष्पप्रग्रथिताक्षरम्} %||5-36-11||

\twolineshloka
{इदं श्रेष्ठमभिज्ञानं ब्रूयास्त्वं तु मम प्रियम्}
{शैलस्य चित्रकूटस्य पादे पूर्वोत्तरे तदा} %||5-36-12||

\twolineshloka
{तापसाश्रमवासिन्याः प्राज्यमूलफलोदके}
{तस्मिन्सिद्धाश्रमे देशे मन्दाकिन्या अदूरतः} %||5-36-13||

\twolineshloka
{तस्योपवनषण्डेषु नानापुष्पसुगन्धिषु}
{विहृत्य सलिलक्लिन्ना तवाङ्के समुपाविशम्} %||5-36-14||

{पर्यायेण प्रसुप्तश्च ममाङ्के भरताग्रजः॥\devanumber{15}॥} %||5-36-15|| (Check)

\addtocounter{shlokacount}{1}

\twolineshloka
{ततो मांससमायुक्तो वायसः पर्यतुण्डयत्}
{तमहं लोष्टमुद्यम्य वारयामि स्म वायसं} %||5-36-16||

\twolineshloka
{दारयन्स च मां काकस्तत्रैव परिलीयते}
{न चाप्युपरमन्मांसाद्भक्षार्थी बलिभोजनः} %||5-36-17||

\twolineshloka
{उत्कर्षन्त्यां च रशनां क्रुद्धायां मयि पक्षिणे}
{स्रंसमाने च वसने ततो दृष्टा त्वया ह्यहम्} %||5-36-18||

\twolineshloka
{त्वया विहसिता चाहं क्रुद्धा संलज्जिता तदा}
{भक्ष्य गृद्धेन कालेन दारिता त्वामुपागता} %||5-36-19||

\twolineshloka
{आसीनस्य च ते श्रान्ता पुनरुत्सङ्गमाविशम्}
{क्रुध्यन्ती च प्रहृष्टेन त्वयाहं परिसान्त्विता} %||5-36-20||

\twolineshloka
{बाष्पपूर्णमुखी मन्दं चक्षुषी परिमार्जती}
{लक्षिताहं त्वया नाथ वायसेन प्रकोपिता} %||5-36-21||

\threelineshloka
{आशीविष इव क्रुद्धः श्वसान्वाक्यमभाषथाः}
{केन ते नागनासोरु विक्षतं वै स्तनान्तरम्}
{कः क्रीडति सरोषेण पञ्चवक्त्रेण भोगिना} %||5-36-22||

\twolineshloka
{वीक्षमाणस्ततस्तं वै वायसं समवैक्षथाः}
{नखैः सरुधिरैस्तीक्ष्णैर्मामेवाभिमुखं स्थितम्} %||5-36-23||

\twolineshloka
{पुत्रः किल स शक्रस्य वायसः पततां वरः}
{धरान्तरचरः शीघ्रं पवनस्य गतौ समः} %||5-36-24||

\twolineshloka
{ततस्तस्मिन्महाबाहुः कोपसंवर्तितेक्षणः}
{वायसे कृतवान्क्रूरां मतिं मतिमतां वर} %||5-36-25||

\twolineshloka
{स दर्भसंस्तराद्गृह्य ब्रह्मणोऽस्त्रेण योजयः}
{स दीप्त इव कालाग्निर्जज्वालाभिमुखो द्विजम्} %||5-36-26||

\threelineshloka
{चिक्षेपिथ प्रदीप्तां तामिषीकां वायसं प्रति}
{अनुसृष्टस्तदा कालो जगाम विविधां गतिम्}
{त्राणकाम इमं लोकं सर्वं वै विचचार ह} %||5-36-27||

\twolineshloka
{स पित्रा च परित्यक्तः सुरैः सर्वैर्महर्षिभिः}
{त्रीँल्लोकान्संपरिक्रम्य त्वामेव शरणं गतः} %||5-36-28||

\threelineshloka
{तं त्वं निपतितं भूमौ शरण्यः शरणागतम्}
{वधार्हमपि काकुत्स्थ कृपया पर्यपालयः}
{न शर्म लब्ध्वा लोकेषु त्वामेव शरणं गतः} %||5-36-29||

\twolineshloka
{परिद्यूनं विषण्णं च स त्वमायान्तमुक्तवान्}
{मोघं कर्तुं न शक्यं तु ब्राह्ममस्त्रं तदुच्यताम्} %||5-36-30||

{ततस्तस्याक्षि काकस्य हिनस्ति स्म स दक्षिणम्॥\devanumber{31}॥} %||5-36-31|| (Check)

\addtocounter{shlokacount}{1}

\twolineshloka
{स ते तदा नमस्कृत्वा राज्ञे दशरथाय च}
{त्वया वीर विसृष्टस्तु प्रतिपेदे स्वमालयम्} %||5-36-32||

\twolineshloka
{मत्कृते काकमात्रेऽपि ब्रह्मास्त्रं समुदीरितम्}
{कस्माद्यो मां हरत्त्वत्तः क्षमसे तं महीपते} %||5-36-33||

\twolineshloka
{स कुरुष्व महोत्साहं कृपां मयि नरर्षभ}
{आनृशंस्यं परो धर्मस्त्वत्त एव मया श्रुतः} %||5-36-34||

\threelineshloka
{जानामि त्वां महावीर्यं महोत्साहं महाबलम्}
{अपारपारमक्षोभ्यं गाम्भीर्यात्सागरोपमम्}
{भर्तारं ससमुद्राया धरण्या वासवोपमम्} %||5-36-35||

\twolineshloka
{एवमस्त्रविदां श्रेष्ठः सत्त्ववान्बलवानपि}
{किमर्थमस्त्रं रक्षःसु न योजयसि राघव} %||5-36-36||

\twolineshloka
{न नागा नापि गन्धर्वा नासुरा न मरुद्गणाः}
{रामस्य समरे वेगं शक्ताः प्रति समाधितुम्} %||5-36-37||

\twolineshloka
{तस्या वीर्यवतः कश्चिद्यद्यस्ति मयि संभ्रमः}
{किमर्थं न शरैस्तीक्ष्णैः क्षयं नयति राक्षसान्} %||5-36-38||

\twolineshloka
{भ्रातुरादेशमादाय लक्ष्मणो वा परंतपः}
{कस्य हेतोर्न मां वीरः परित्राति महाबलः} %||5-36-39||

\twolineshloka
{यदि तौ पुरुषव्याघ्रौ वाय्विन्द्रसमतेजसौ}
{सुराणामपि दुर्धर्षो किमर्थं मामुपेक्षतः} %||5-36-40||

\twolineshloka
{ममैव दुष्कृतं किंचिन्महदस्ति न संशयः}
{समर्थावपि तौ यन्मां नावेक्षेते परंतपौ} %||5-36-41||

\twolineshloka
{कौसल्या लोकभर्तारं सुषुवे यं मनस्विनी}
{तं ममार्थे सुखं पृच्छ शिरसा चाभिवादय} %||5-36-42||

\twolineshloka
{स्रजश्च सर्वरत्नानि प्रिया याश्च वराङ्गनाः}
{ऐश्वर्यं च विशालायां पृथिव्यामपि दुर्लभम्} %||5-36-43||

\threelineshloka
{पितरं मातरं चैव संमान्याभिप्रसाद्य च}
{अनुप्रव्रजितो रामं सुमित्रा येन सुप्रजाः}
{आनुकूल्येन धर्मात्मा त्यक्त्वा सुखमनुत्तमम्} %||5-36-44||

\twolineshloka
{अनुगच्छति काकुत्स्थं भ्रातरं पालयन्वने}
{सिंहस्कन्धो महाबाहुर्मनस्वी प्रियदर्शनः} %||5-36-45||

\twolineshloka
{पितृवद्वर्तते रामे मातृवन्मां समाचरन्}
{ह्रियमाणां तदा वीरो न तु मां वेद लक्ष्मणः} %||5-36-46||

\twolineshloka
{वृद्धोपसेवी लक्ष्मीवाञ्शक्तो न बहुभाषिता}
{राजपुत्रः प्रियश्रेष्ठः सदृशः श्वशुरस्य मे} %||5-36-47||

\twolineshloka
{मत्तः प्रियतरो नित्यं भ्राता रामस्य लक्ष्मणः}
{नियुक्तो धुरि यस्यां तु तामुद्वहति वीर्यवान्} %||5-36-48||

\threelineshloka
{यं दृष्ट्वा राघवो नैव वृद्धमार्यमनुस्मरत्}
{स ममार्थाय कुशलं वक्तव्यो वचनान्मम}
{मृदुर्नित्यं शुचिर्दक्षः प्रियो रामस्य लक्ष्मणः} %||5-36-49||

\threelineshloka
{इदं ब्रूयाश्च मे नाथं शूरं रामं पुनः पुनः}
{जीवितं धारयिष्यामि मासं दशरथात्मज}
{ऊर्ध्वं मासान्न जीवेयं सत्येनाहं ब्रवीमि ते} %||5-36-50||

\twolineshloka
{रावणेनोपरुद्धां मां निकृत्या पापकर्मणा}
{त्रातुमर्हसि वीर त्वं पातालादिव कौशिकीम्} %||5-36-51||

\twolineshloka
{ततो वस्त्रगतं मुक्त्वा दिव्यं चूडामणिं शुभम्}
{प्रदेयो राघवायेति सीता हनुमते ददौ} %||5-36-52||

\twolineshloka
{प्रतिगृह्य ततो वीरो मणिरत्नमनुत्तमम्}
{अङ्गुल्या योजयामास न ह्यस्या प्राभवद्भुजः} %||5-36-53||

\twolineshloka
{मणिरत्नं कपिवरः प्रतिगृह्याभिवाद्य च}
{सीतां प्रदक्षिणं कृत्वा प्रणतः पार्श्वतः स्थितः} %||5-36-54||

\twolineshloka
{हर्षेण महता युक्तः सीतादर्शनजेन सः}
{हृदयेन गतो रामं शरीरेण तु विष्ठितः} %||5-36-55||

\fourlineindentedshloka
{मणिवरमुपगृह्य तं महार्हं}
{ जनकनृपात्मजया धृतं प्रभावात्}
{गिरिवरपवनावधूतमुक्तः}
{ सुखितमनाः प्रतिसंक्रमं प्रपेदे} %||5-37-56||


{॥इत्यार्षे श्रीमद्रामायणे वाल्मीकीये आदिकाव्ये सुन्दरकाण्डे षट्त्रिंशः सर्गः॥}

\sect{सप्तत्रिंशः सर्गः}

\twolineshloka
{मणिं दत्त्वा ततः सीता हनूमन्तमथाब्रवीत्}
{अभिज्ञानमभिज्ञातमेतद्रामस्य तत्त्वतः} %||5-37-1||

\twolineshloka
{मणिं तु दृष्ट्वा रामो वै त्रयाणां संस्मरिष्यति}
{वीरो जनन्या मम च राज्ञो दशरथस्य च} %||5-37-2||

\twolineshloka
{स भूयस्त्वं समुत्साहे चोदितो हरिसत्तम}
{अस्मिन्कार्यसमारम्भे प्रचिन्तय यदुत्तरम्} %||5-37-3||

\twolineshloka
{त्वमस्मिन्कार्यनिर्योगे प्रमाणं हरिसत्तम}
{तस्य चिन्तय यो यत्नो दुःखक्षयकरो भवेत्} %||5-37-4||

\twolineshloka
{स तथेति प्रतिज्ञाय मारुतिर्भीमविक्रमः}
{शिरसावन्द्य वैदेहीं गमनायोपचक्रमे} %||5-37-5||

\twolineshloka
{ज्ञात्वा संप्रस्थितं देवी वानरं मारुतात्मजम्}
{बाष्पगद्गदया वाचा मैथिली वाक्यमब्रवीत्} %||5-37-6||

\twolineshloka
{कुशलं हनुमन्ब्रूयाः सहितौ रामलक्ष्मणौ}
{सुग्रीवं च सहामात्यं वृद्धान्सर्वांश्च वानरान्} %||5-37-7||

\twolineshloka
{यथा च स महाबाहुर्मां तारयति राघवः}
{अस्माद्दुःखाम्बुसंरोधात्त्वं समाधातुमर्हसि} %||5-37-8||

\twolineshloka
{जीवन्तीं मां यथा रामः संभावयति कीर्तिमान्}
{तत्त्वया हनुमन्वाच्यं वाचा धर्ममवाप्नुहि} %||5-37-9||

\twolineshloka
{नित्यमुत्साहयुक्ताश्च वाचः श्रुत्वा मयेरिताः}
{वर्धिष्यते दाशरथेः पौरुषं मदवाप्तये} %||5-37-10||

\twolineshloka
{मत्संदेशयुता वाचस्त्वत्तः श्रुत्वैव राघवः}
{पराक्रमविधिं वीरो विधिवत्संविधास्यति} %||5-37-11||

\twolineshloka
{सीतायास्तद्वचः श्रुत्वा हनुमान्मारुतात्मजः}
{शिरस्यञ्जलिमाधाय वाक्यमुत्तरमब्रवीत्} %||5-37-12||

\twolineshloka
{क्षिप्रमेष्यति काकुत्स्थो हर्यृक्षप्रवरैर्वृतः}
{यस्ते युधि विजित्यारीञ्शोकं व्यपनयिष्यति} %||5-37-13||

\twolineshloka
{न हि पश्यामि मर्त्येषु नामरेष्वसुरेषु वा}
{यस्तस्य वमतो बाणान्स्थातुमुत्सहतेऽग्रतः} %||5-37-14||

\twolineshloka
{अप्यर्कमपि पर्जन्यमपि वैवस्वतं यमम्}
{स हि सोढुं रणे शक्तस्तवहेतोर्विशेषतः} %||5-37-15||

\twolineshloka
{स हि सागरपर्यन्तां महीं शासितुमीहते}
{त्वन्निमित्तो हि रामस्य जयो जनकनन्दिनि} %||5-37-16||

\twolineshloka
{तस्य तद्वचनं श्रुत्वा सम्यक्सत्यं सुभाषितम्}
{जानकी बहु मेनेऽथ वचनं चेदमब्रवीत्} %||5-37-17||

\twolineshloka
{ततस्तं प्रस्थितं सीता वीक्षमाणा पुनः पुनः}
{भर्तुः स्नेहान्वितं वाक्यं सौहार्दादनुमानयत्} %||5-37-18||

\twolineshloka
{यदि वा मन्यसे वीर वसैकाहमरिंदम}
{कस्मिंश्चित्संवृते देशे विश्रान्तः श्वो गमिष्यसि} %||5-37-19||

\twolineshloka
{मम चेदल्पभाग्यायाः साम्निध्यात्तव वीर्यवान्}
{अस्य शोकस्य महतो मुहूर्तं मोक्षणं भवेत्} %||5-37-20||

\twolineshloka
{गते हि हरिशार्दूल पुनरागमनाय तु}
{प्राणानामपि संदेहो मम स्यान्नात्र संशयः} %||5-37-21||

\twolineshloka
{तवादर्शनजः शोको भूयो मां परितापयेत्}
{दुःखाद्दुःखपरामृष्टां दीपयन्निव वानर} %||5-37-22||

\twolineshloka
{अयं च वीर संदेहस्तिष्ठतीव ममाग्रतः}
{सुमहांस्त्वत्सहायेषु हर्यृक्षेषु हरीश्वर} %||5-37-23||

\twolineshloka
{कथं नु खलु दुष्पारं तरिष्यन्ति महोदधिम्}
{तानि हर्यृक्षसैन्यानि तौ वा नरवरात्मजौ} %||5-37-24||

\twolineshloka
{त्रयाणामेव भूतानां सागरस्येह लङ्घने}
{शक्तिः स्याद्वैनतेयस्य तव वा मारुतस्य वा} %||5-37-25||

\twolineshloka
{तदस्मिन्कार्यनिर्योगे वीरैवं दुरतिक्रमे}
{किं पश्यसि समाधानं त्वं हि कार्यविदां वरः} %||5-37-26||

\twolineshloka
{काममस्य त्वमेवैकः कार्यस्य परिसाधने}
{पर्याप्तः परवीरघ्न यशस्यस्ते बलोदयः} %||5-37-27||

\twolineshloka
{बलैः समग्रैर्यदि मां रावणं जित्य संयुगे}
{विजयी स्वपुरं यायात्तत्तु मे स्याद्यशस्करम्} %||5-37-28||

\twolineshloka
{बलैस्तु संकुलां कृत्वा लङ्कां परबलार्दनः}
{मां नयेद्यदि काकुत्स्थस्तत्तस्य सदृशं भवेत्} %||5-37-29||

\twolineshloka
{तद्यथा तस्य विक्रान्तमनुरूपं महात्मनः}
{भवेदाहव शूरस्य तथा त्वमुपपादय} %||5-37-30||

\twolineshloka
{तदर्थोपहितं वाक्यं सहितं हेतुसंहितम्}
{निशम्य हनुमाञ्शेषं वाक्यमुत्तरमब्रवीत्} %||5-37-31||

\twolineshloka
{देवि हर्यृक्षसैन्यानामीश्वरः प्लवतां वरः}
{सुग्रीवः सत्त्वसंपन्नस्तवार्थे कृतनिश्चयः} %||5-37-32||

\twolineshloka
{स वानरसहस्राणां कोटीभिरभिसंवृतः}
{क्षिप्रमेष्यति वैदेहि राक्षसानां निबर्हणः} %||5-37-33||

\twolineshloka
{तस्य विक्रमसंपन्नाः सत्त्ववन्तो महाबलाः}
{मनःसंकल्पसंपाता निदेशे हरयः स्थिताः} %||5-37-34||

\twolineshloka
{येषां नोपरि नाधस्तान्न तिर्यक्सज्जते गतिः}
{न च कर्मसु सीदन्ति महत्स्वमिततेजसः} %||5-37-35||

\twolineshloka
{असकृत्तैर्महोत्सहैः ससागरधराधरा}
{प्रदक्षिणीकृता भूमिर्वायुमार्गानुसारिभिः} %||5-37-36||

\twolineshloka
{मद्विशिष्टाश्च तुल्याश्च सन्ति तत्र वनौकसः}
{मत्तः प्रत्यवरः कश्चिन्नास्ति सुग्रीवसंनिधौ} %||5-37-37||

\twolineshloka
{अहं तावदिह प्राप्तः किं पुनस्ते महाबलाः}
{न हि प्रकृष्टाः प्रेष्यन्ते प्रेष्यन्ते हीतरे जनाः} %||5-37-38||

\twolineshloka
{तदलं परितापेन देवि शोको व्यपैतु ते}
{एकोत्पातेन ते लङ्कामेष्यन्ति हरियूथपाः} %||5-37-39||

\twolineshloka
{मम पृष्ठगतौ तौ च चन्द्रसूर्याविवोदितौ}
{त्वत्सकाशं महासत्त्वौ नृसिंहावागमिष्यतः} %||5-37-40||

\twolineshloka
{तौ हि वीरौ नरवरौ सहितौ रामलक्ष्मणौ}
{आगम्य नगरीं लङ्कां सायकैर्विधमिष्यतः} %||5-37-41||

\twolineshloka
{सगणं रावणं हत्वा राघवो रघुनन्दनः}
{त्वामादाय वरारोहे स्वपुरं प्रतियास्यति} %||5-37-42||

\twolineshloka
{तदाश्वसिहि भद्रं ते भव त्वं कालकाङ्क्षिणी}
{नचिराद्द्रक्ष्यसे रामं प्रज्वजन्तमिवानिलम्} %||5-37-43||

\twolineshloka
{निहते राक्षसेन्द्रे च सपुत्रामात्यबान्धवे}
{त्वं समेष्यसि रामेण शशाङ्केनेव रोहिणी} %||5-37-44||

\twolineshloka
{क्षिप्रं त्वं देवि शोकस्य पारं यास्यसि मैथिलि}
{रावणं चैव रामेण निहतं द्रक्ष्यसेऽचिरात्} %||5-37-45||

\twolineshloka
{एवमाश्वस्य वैदेहीं हनूमान्मारुतात्मजः}
{गमनाय मतिं कृत्वा वैदेहीं पुनरब्रवीत्} %||5-37-46||

\twolineshloka
{तमरिघ्नं कृतात्मानं क्षिप्रं द्रक्ष्यसि राघवम्}
{लक्ष्मणं च धनुष्पाणिं लङ्काद्वारमुपस्थितम्} %||5-37-47||

\twolineshloka
{नखदंष्ट्रायुधान्वीरान्सिंहशार्दूलविक्रमान्}
{वानरान्वारणेन्द्राभान्क्षिप्रं द्रक्ष्यसि संगतान्} %||5-37-48||

\twolineshloka
{शैलाम्बुदनिकाशानां लङ्कामलयसानुषु}
{नर्दतां कपिमुख्यानामार्ये यूथान्यनेकशः} %||5-37-49||

\twolineshloka
{स तु मर्मणि घोरेण ताडितो मन्मथेषुणा}
{न शर्म लभते रामः सिंहार्दित इव द्विपः} %||5-37-50||

\twolineshloka
{मा रुदो देवि शोकेन मा भूत्ते मनसोऽप्रियम्}
{शचीव पथ्या शक्रेण भर्त्रा नाथवती ह्यसि} %||5-37-51||

\twolineshloka
{रामाद्विशिष्टः कोऽन्योऽस्ति कश्चित्सौमित्रिणा समः}
{अग्निमारुतकल्पौ तौ भ्रातरौ तव संश्रयौ} %||5-37-52||

\fourlineindentedshloka
{नास्मिंश्चिरं वत्स्यसि देवि देशे}
{ रक्षोगणैरध्युषितोऽतिरौद्रे}
{न ते चिरादागमनं प्रियस्य}
{ क्षमस्व मत्संगमकालमात्रम्} %||5-38-53||


{॥इत्यार्षे श्रीमद्रामायणे वाल्मीकीये आदिकाव्ये सुन्दरकाण्डे सप्तत्रिंशः सर्गः॥}

\sect{अष्टात्रिंशः सर्गः}

\twolineshloka
{श्रुत्वा तु वचनं तस्य वायुसूनोर्महात्मनः}
{उवाचात्महितं वाक्यं सीता सुरसुतोपमा} %||5-38-1||

\twolineshloka
{त्वां दृष्ट्वा प्रियवक्तारं संप्रहृष्यामि वानर}
{अर्धसंजातसस्येव वृष्टिं प्राप्य वसुंधरा} %||5-38-2||

\twolineshloka
{यथा तं पुरुषव्याघ्रं गात्रैः शोकाभिकर्शितैः}
{संस्पृशेयं सकामाहं तथा कुरु दयां मयि} %||5-38-3||

\twolineshloka
{अभिज्ञानं च रामस्य दत्तं हरिगणोत्तम}
{क्षिप्तामीषिकां काकस्य कोपादेकाक्षिशातनीम्} %||5-38-4||

\twolineshloka
{मनःशिलायास्तिकलो गण्डपार्श्वे निवेशितः}
{त्वया प्रनष्टे तिलके तं किल स्मर्तुमर्हसि} %||5-38-5||

\twolineshloka
{स वीर्यवान्कथं सीतां हृतां समनुमन्यसे}
{वसन्तीं रक्षसां मध्ये महेन्द्रवरुणोपम} %||5-38-6||

\twolineshloka
{एष चूडामणिर्दिव्यो मया सुपरिरक्षितः}
{एतं दृष्ट्वा प्रहृष्यामि व्यसने त्वामिवानघ} %||5-38-7||

\twolineshloka
{एष निर्यातितः श्रीमान्मया ते वारिसंभवः}
{अतः परं न शक्ष्यामि जीवितुं शोकलालसा} %||5-38-8||

\twolineshloka
{असह्यानि च दुःखानि वाचश्च हृदयच्छिदः}
{राक्षसीनां सुघोराणां त्वत्कृते मर्षयाम्यहम्} %||5-38-9||

\twolineshloka
{धारयिष्यामि मासं तु जीवितं शत्रुसूदन}
{मासादूर्ध्वं न जीविष्ये त्वया हीना नृपात्मज} %||5-38-10||

\twolineshloka
{घोरो राक्षसराजोऽयं दृष्टिश्च न सुखा मयि}
{त्वां च श्रुत्वा विपद्यन्तं न जीवेयमहं क्षणम्} %||5-38-11||

\twolineshloka
{वैदेह्या वचनं श्रुत्वा करुणं साश्रुभाषितम्}
{अथाब्रवीन्महातेजा हनुमान्मारुतात्मजः} %||5-38-12||

\twolineshloka
{त्वच्छोकविमुखो रामो देवि सत्येन ते शपे}
{रामे शोकाभिभूते तु लक्ष्मणः परितप्यते} %||5-38-13||

\twolineshloka
{दृष्टा कथंचिद्भवती न कालः परिशोचितुम्}
{इमं मुहूर्तं दुःखानामन्तं द्रक्ष्यसि भामिनि} %||5-38-14||

\twolineshloka
{तावुभौ पुरुषव्याघ्रौ राजपुत्रावनिन्दितौ}
{त्वद्दर्शनकृतोत्साहौ लङ्कां भस्मीकरिष्यतः} %||5-38-15||

\twolineshloka
{हत्वा तु समरे क्रूरं रावणं सह बान्धवम्}
{राघवौ त्वां विशालाक्षि स्वां पुरीं प्रापयिष्यतः} %||5-38-16||

\twolineshloka
{यत्तु रामो विजानीयादभिज्ञानमनिन्दिते}
{प्रीतिसंजननं तस्य भूयस्त्वं दातुमर्हसि} %||5-38-17||

\threelineshloka
{साब्रवीद्दत्तमेवेह मयाभिज्ञानमुत्तमम्}
{एतदेव हि रामस्य दृष्ट्वा मत्केशभूषणम्}
{श्रद्धेयं हनुमन्वाक्यं तव वीर भविष्यति} %||5-38-18||

\twolineshloka
{स तं मणिवरं गृह्य श्रीमान्प्लवगसत्तमः}
{प्रणम्य शिरसा देवीं गमनायोपचक्रमे} %||5-38-19||

\threelineshloka
{तमुत्पातकृतोत्साहमवेक्ष्य हरिपुंगवम्}
{वर्धमानं महावेगमुवाच जनकात्मजा}
{अश्रुपूर्णमुखी दीना बाष्पगद्गदया गिरा} %||5-38-20||

\twolineshloka
{हनूमन्सिंहसंकाशौ भ्रातरौ रामलक्ष्मणौ}
{सुग्रीवं च सहामात्यं सर्वान्ब्रूया अनामयम्} %||5-38-21||

\twolineshloka
{यथा च स महाबाहुर्मां तारयति राघवः}
{अस्माद्दुःखाम्बुसंरोधात्तत्समाधातुमर्हसि} %||5-38-22||

\fourlineindentedshloka
{इमं च तीव्रं मम शोकवेगं}
{ रक्षोभिरेभिः परिभर्त्सनं च}
{ब्रूयास्तु रामस्य गतः समीपं}
{ शिवश्च तेऽध्वास्तु हरिप्रवीर} %||5-38-23||

\fourlineindentedshloka
{स राजपुत्र्या प्रतिवेदितार्थः}
{ कपिः कृतार्थः परिहृष्टचेताः}
{तदल्पशेषं प्रसमीक्ष्य कार्यं}
{ दिशं ह्युदीचीं मनसा जगाम} %||5-39-24||


{॥इत्यार्षे श्रीमद्रामायणे वाल्मीकीये आदिकाव्ये सुन्दरकाण्डे अष्टात्रिंशः सर्गः॥}

\sect{एकोनचत्वारिंशः सर्गः}

\twolineshloka
{स च वाग्भिः प्रशस्ताभिर्गमिष्यन्पूजितस्तया}
{तस्माद्देशादपक्रम्य चिन्तयामास वानरः} %||5-39-1||

\twolineshloka
{अल्पशेषमिदं कार्यं दृष्टेयमसितेक्षणा}
{त्रीनुपायानतिक्रम्य चतुर्थ इह दृश्यते} %||5-39-2||

\fourlineindentedshloka
{न साम रक्षःसु गुणाय कल्पते}
{ न दनमर्थोपचितेषु वर्तते}
{न भेदसाध्या बलदर्पिता जनाः}
{ पराक्रमस्त्वेष ममेह रोचते} %||5-39-3||

\fourlineindentedshloka
{न चास्य कार्यस्य पराक्रमादृते}
{ विनिश्चयः कश्चिदिहोपपद्यते}
{हृतप्रवीरास्तु रणे हि राक्षसाः}
{ कथंचिदीयुर्यदिहाद्य मार्दवम्} %||5-39-4||

\twolineshloka
{कार्ये कर्मणि निर्दिष्टो यो बहून्यपि साधयेत्}
{पूर्वकार्यविरोधेन स कार्यं कर्तुमर्हति} %||5-39-5||

\twolineshloka
{न ह्येकः साधको हेतुः स्वल्पस्यापीह कर्मणः}
{यो ह्यर्थं बहुधा वेद स समर्थोऽर्थसाधने} %||5-39-6||

\fourlineindentedshloka
{इहैव तावत्कृतनिश्चयो ह्यहं}
{ यदि व्रजेयं प्लवगेश्वरालयम्}
{परात्मसंमर्द विशेषतत्त्ववि}
{त्ततः कृतं स्यान्मम भर्तृशासनम्} %||5-39-7||

\fourlineindentedshloka
{कथं नु खल्वद्य भवेत्सुखागतं}
{ प्रसह्य युद्धं मम राक्षसैः सह}
{तथैव खल्वात्मबलं च सारव}
{त्समानयेन्मां च रणे दशाननः} %||5-39-8||

\twolineshloka
{इदमस्य नृशंसस्य नन्दनोपममुत्तमम्}
{वनं नेत्रमनःकान्तं नानाद्रुमलतायुतम्} %||5-39-9||

\twolineshloka
{इदं विध्वंसयिष्यामि शुष्कं वनमिवानलः}
{अस्मिन्भग्ने ततः कोपं करिष्यति स रावणः} %||5-39-10||

\fourlineindentedshloka
{ततो महत्साश्वमहारथद्विपं}
{ बलं समानेष्वपि राक्षसाधिपः}
{त्रिशूलकालायसपट्टिशायुधं}
{ ततो महद्युद्धमिदं भविष्यति} %||5-39-11||

\fourlineindentedshloka
{अहं तु तैः संयति चण्डविक्रमैः}
{ समेत्य रक्षोभिरसंगविक्रमः}
{निहत्य तद्रावणचोदितं बलं}
{ सुखं गमिष्यामि कपीश्वरालयम्} %||5-39-12||

\twolineshloka
{ततो मारुतवत्क्रुद्धो मारुतिर्भीमविक्रमः}
{ऊरुवेगेन महता द्रुमान्क्षेप्तुमथारभत्} %||5-39-13||

\twolineshloka
{ततस्तद्धनुमान्वीरो बभञ्ज प्रमदावनम्}
{मत्तद्विजसमाघुष्टं नानाद्रुमलतायुतम्} %||5-39-14||

\twolineshloka
{तद्वनं मथितैर्वृक्षैर्भिन्नैश्च सलिलाशयैः}
{चूर्णितैः पर्वताग्रैश्च बभूवाप्रियदर्शनम्} %||5-39-15||

\fourlineindentedshloka
{लतागृहैश्चित्रगृहैश्च नाशितै}
{र्महोरगैर्व्यालमृगैश्च निर्धुतैः}
{शिलागृहैरुन्मथितैस्तथा गृहैः}
{ प्रनष्टरूपं तदभून्महद्वनम्} %||5-39-16||

\fourlineindentedshloka
{स तस्य कृत्वार्थपतेर्महाकपि}
{र्महद्व्यलीकं मनसो महात्मनः}
{युयुत्सुरेको बहुभिर्महाबलैः}
{ श्रिया ज्वलंस्तोरणमाश्रितः कपिः} %||5-40-17||


{॥इत्यार्षे श्रीमद्रामायणे वाल्मीकीये आदिकाव्ये सुन्दरकाण्डे एकोनचत्वारिंशः सर्गः॥}

\sect{चत्वारिंशः सर्गः}

\twolineshloka
{ततः पक्षिनिनादेन वृक्षभङ्गस्वनेन च}
{बभूवुस्त्राससंभ्रान्ताः सर्वे लङ्कानिवासिनः} %||5-40-1||

\twolineshloka
{विद्रुताश्च भयत्रस्ता विनेदुर्मृगपक्षुणः}
{रक्षसां च निमित्तानि क्रूराणि प्रतिपेदिरे} %||5-40-2||

\twolineshloka
{ततो गतायां निद्रायां राक्षस्यो विकृताननाः}
{तद्वनं ददृशुर्भग्नं तं च वीरं महाकपिम्} %||5-40-3||

\twolineshloka
{स ता दृष्ट्व महाबाहुर्महासत्त्वो महाबलः}
{चकार सुमहद्रूपं राक्षसीनां भयावहम्} %||5-40-4||

\twolineshloka
{ततस्तं गिरिसंकाशमतिकायं महाबलम्}
{राक्षस्यो वानरं दृष्ट्वा पप्रच्छुर्जनकात्मजाम्} %||5-40-5||

\twolineshloka
{कोऽयं कस्य कुतो वायं किंनिमित्तमिहागतः}
{कथं त्वया सहानेन संवादः कृत इत्युत} %||5-40-6||

\twolineshloka
{आचक्ष्व नो विशालाक्षि मा भूत्ते सुभगे भयम्}
{संवादमसितापाङ्गे त्वया किं कृतवानयम्} %||5-40-7||

\twolineshloka
{अथाब्रवीत्तदा साध्वी सीता सर्वाङ्गशोभना}
{रक्षसां कामरूपाणां विज्ञाने मम का गतिः} %||5-40-8||

\twolineshloka
{यूयमेवास्य जानीत योऽयं यद्वा करिष्यति}
{अहिरेव अहेः पादान्विजानाति न संशयः} %||5-40-9||

\twolineshloka
{अहमप्यस्य भीतास्मि नैनं जानामि कोऽन्वयम्}
{वेद्मि राक्षसमेवैनं कामरूपिणमागतम्} %||5-40-10||

\twolineshloka
{वैदेह्या वचनं श्रुत्वा राक्षस्यो विद्रुता द्रुतम्}
{स्थिताः काश्चिद्गताः काश्चिद्रावणाय निवेदितुम्} %||5-40-11||

\twolineshloka
{रावणस्य समीपे तु राक्षस्यो विकृताननाः}
{विरूपं वानरं भीममाख्यातुमुपचक्रमुः} %||5-40-12||

\twolineshloka
{अशोकवनिका मध्ये राजन्भीमवपुः कपिः}
{सीतया कृतसंवादस्तिष्ठत्यमितविक्रमः} %||5-40-13||

\twolineshloka
{न च तं जानकी सीता हरिं हरिणलोचणा}
{अस्माभिर्बहुधा पृष्टा निवेदयितुमिच्छति} %||5-40-14||

\twolineshloka
{वासवस्य भवेद्दूतो दूतो वैश्रवणस्य वा}
{प्रेषितो वापि रामेण सीतान्वेषणकाङ्क्षया} %||5-40-15||

\twolineshloka
{तेन त्वद्भूतरूपेण यत्तत्तव मनोहरम्}
{नानामृगगणाकीर्णं प्रमृष्टं प्रमदावनम्} %||5-40-16||

\twolineshloka
{न तत्र कश्चिदुद्देशो यस्तेन न विनाशितः}
{यत्र सा जानकी सीता स तेन न विनाशितः} %||5-40-17||

\twolineshloka
{जानकीरक्षणार्थं वा श्रमाद्वा नोपलभ्यते}
{अथ वा कः श्रमस्तस्य सैव तेनाभिरक्षिता} %||5-40-18||

\twolineshloka
{चारुपल्लवपत्राढ्यं यं सीता स्वयमास्थिता}
{प्रवृद्धः शिंशपावृक्षः स च तेनाभिरक्षितः} %||5-40-19||

\twolineshloka
{तस्योग्ररूपस्योग्रं त्वं दण्डमाज्ञातुमर्हसि}
{सीता संभाषिता येन तद्वनं च विनाशितम्} %||5-40-20||

\twolineshloka
{मनःपरिगृहीतां तां तव रक्षोगणेश्वर}
{कः सीतामभिभाषेत यो न स्यात्त्यक्तजीवितः} %||5-40-21||

\twolineshloka
{राक्षसीनां वचः श्रुत्वा रावणो राक्षसेश्वरः}
{हुतागिरिव जज्वाल कोपसंवर्तितेक्षणः} %||5-40-22||

\twolineshloka
{आत्मनः सदृशाञ्शूरान्किंकरान्नाम राक्षसान्}
{व्यादिदेश महातेजा निग्रहार्थं हनूमतः} %||5-40-23||

\twolineshloka
{तेषामशीतिसाहस्रं किंकराणां तरस्विनाम्}
{निर्ययुर्भवनात्तस्मात्कूटमुद्गरपाणयः} %||5-40-24||

\twolineshloka
{महोदरा महादंष्ट्रा घोररूपा महाबलाः}
{युद्धाभिमनसः सर्वे हनूमद्ग्रहणोन्मुखाः} %||5-40-25||

\twolineshloka
{ते कपिं तं समासाद्य तोरणस्थमवस्थितम्}
{अभिपेतुर्महावेगाः पतङ्गा इव पावकम्} %||5-40-26||

\twolineshloka
{ते गदाभिर्विचित्राभिः परिघैः काञ्चनाङ्गदैः}
{आजघ्नुर्वानरश्रेष्ठं शरैरादित्यसंनिभैः} %||5-40-27||

\twolineshloka
{हनूमानपि तेजस्वी श्रीमान्पर्वतसंनिभः}
{क्षितावाविध्य लाङ्गूलं ननाद च महास्वनम्} %||5-40-28||

\twolineshloka
{तस्य संनादशब्देन तेऽभवन्भयशङ्किताः}
{ददृशुश्च हनूमन्तं संध्यामेघमिवोन्नतम्} %||5-40-29||

\twolineshloka
{स्वामिसंदेशनिःशङ्कास्ततस्ते राक्षसाः कपिम्}
{चित्रैः प्रहरणैर्भीमैरभिपेतुस्ततस्ततः} %||5-40-30||

\twolineshloka
{स तैः परिवृतः शूरैः सर्वतः स महाबलः}
{आससादायसं भीमं परिघं तोरणाश्रितम्} %||5-40-31||

{स तं परिघमादाय जघान रजनीचरान्॥\devanumber{32}॥} %||5-40-32|| (Check)

\addtocounter{shlokacount}{1}

\twolineshloka
{स पन्नगमिवादाय स्फुरन्तं विनतासुतः}
{विचचाराम्बरे वीरः परिगृह्य च मारुतिः} %||5-40-33||

\twolineshloka
{स हत्वा राक्षसान्वीरः किंकरान्मारुतात्मजः}
{युद्धाकाङ्क्षी पुनर्वीरस्तोरणं समुपस्थितः} %||5-40-34||

\twolineshloka
{ततस्तस्माद्भयान्मुक्ताः कतिचित्तत्र राक्षसाः}
{निहतान्किंकरान्सर्वान्रावणाय न्यवेदयन्} %||5-40-35||

\fourlineindentedshloka
{स राक्षसानां निहतं महाबलं}
{ निशम्य राजा परिवृत्तलोचनः}
{समादिदेशाप्रतिमं पराक्रमे}
{ प्रहस्तपुत्रं समरे सुदुर्जयम्} %||5-41-36||


{॥इत्यार्षे श्रीमद्रामायणे वाल्मीकीये आदिकाव्ये सुन्दरकाण्डे चत्वारिंशः सर्गः॥}

\sect{एकचत्वारिंशः सर्गः}

\threelineshloka
{ततः स किंकरान्हत्वा हनूमान्ध्यानमास्थितः}
{वनं भग्नं मया चैत्यप्रासादो न विनाशितः}
{तस्मात्प्रासादमप्येवमिमं विध्वंसयाम्यहम्} %||5-41-1||

\threelineshloka
{इति संचिन्त्य हनुमान्मनसा दर्शयन्बलम्}
{चैत्यप्रासादमाप्लुत्य मेरुशृङ्गमिवोन्नतम्}
{आरुरोह हरिश्रेष्ठो हनूमान्मारुतात्मजः} %||5-41-2||

\twolineshloka
{संप्रधृष्य च दुर्धर्षश्चैत्यप्रासादमुन्नतम्}
{हनूमान्प्रज्वलँल्लक्ष्म्या पारियात्रोपमोऽभवत्} %||5-41-3||

\twolineshloka
{स भूत्वा तु महाकायो हनूमान्मारुतात्मजः}
{धृष्टमास्फोटयामास लङ्कां शब्देन पूरयन्} %||5-41-4||

\twolineshloka
{तस्यास्फोटितशब्देन महता श्रोत्रघातिना}
{पेतुर्विहंगा गगनादुच्चैश्चेदमघोषयत्} %||5-41-5||

\twolineshloka
{जयत्यतिबलो रामो लक्ष्मणश्च महाबलः}
{राजा जयति सुग्रीवो राघवेणाभिपालितः} %||5-41-6||

\twolineshloka
{दासोऽहं कोसलेन्द्रस्य रामस्याक्लिष्टकर्मणः}
{हनुमाञ्शत्रुसैन्यानां निहन्ता मारुतात्मजः} %||5-41-7||

\twolineshloka
{न रावणसहस्रं मे युद्धे प्रतिबलं भवेत्}
{शिलाभिस्तु प्रहरतः पादपैश्च सहस्रशः} %||5-41-8||

\twolineshloka
{अर्दयित्वा पुरीं लङ्कामभिवाद्य च मैथिलीम्}
{समृद्धार्थो गमिष्यामि मिषतां सर्वरक्षसाम्} %||5-41-9||

\twolineshloka
{एवमुक्त्वा विमानस्थश्चैत्यस्थान्हरिपुंगवः}
{ननाद भीमनिर्ह्रादो रक्षसां जनयन्भयम्} %||5-41-10||

\threelineshloka
{तेन शब्देन महता चैत्यपालाः शतं ययुः}
{गृहीत्वा विविधानस्त्रान्प्रासान्खड्गान्परश्वधान्}
{विसृजन्तो महाक्षया मारुतिं पर्यवारयन्} %||5-41-11||

\twolineshloka
{आवर्त इव गङ्गायास्तोयस्य विपुलो महान्}
{परिक्षिप्य हरिश्रेष्ठं स बभौ रक्षसां गणः} %||5-41-12||

{ततो वातात्मजः क्रुद्धो भीमरूपं समास्थितः॥\devanumber{13}॥} %||5-41-13|| (Check)

\addtocounter{shlokacount}{1}

\threelineshloka
{प्रासादस्य महांस्तस्य स्तम्भं हेमपरिष्कृतम्}
{उत्पाटयित्वा वेगेन हनूमान्मारुतात्मजः}
{ततस्तं भ्रामयामास शतधारं महाबलः} %||5-41-14||

\twolineshloka
{स राक्षसशतं हत्वा वज्रेणेन्द्र इवासुरान्}
{अन्तरिक्षस्थितः श्रीमानिदं वचनमब्रवीत्} %||5-41-15||

\twolineshloka
{मादृशानां सहस्राणि विसृष्टानि महात्मनाम्}
{बलिनां वानरेन्द्राणां सुग्रीववशवर्तिनाम्} %||5-41-16||

\twolineshloka
{शतैः शतसहस्रैश्च कोटीभिरयुतैरपि}
{आगमिष्यति सुग्रीवः सर्वेषां वो निषूदनः} %||5-41-17||

\twolineshloka
{नेयमस्ति पुरी लङ्का न यूयं न च रावणः}
{यस्मादिक्ष्वाकुनाथेन बद्धं वैरं महात्मना} %||5-42-18||


{॥इत्यार्षे श्रीमद्रामायणे वाल्मीकीये आदिकाव्ये सुन्दरकाण्डे एकचत्वारिंशः सर्गः॥}

\sect{द्विचत्वारिंशः सर्गः}

\twolineshloka
{संदिष्टो राक्षसेन्द्रेण प्रहस्तस्य सुतो बली}
{जम्बुमाली महादंष्ट्रो निर्जगाम धनुर्धरः} %||5-42-1||

\twolineshloka
{रक्तमाल्याम्बरधरः स्रग्वी रुचिरकुण्डलः}
{महान्विवृत्तनयनश्चण्डः समरदुर्जयः} %||5-42-2||

\twolineshloka
{धनुः शक्रधनुः प्रख्यं महद्रुचिरसायकम्}
{विस्फारयाणो वेगेन वज्राशनिसमस्वनम्} %||5-42-3||

\twolineshloka
{तस्य विस्फारघोषेण धनुषो महता दिशः}
{प्रदिशश्च नभश्चैव सहसा समपूर्यत} %||5-42-4||

\twolineshloka
{रथेन खरयुक्तेन तमागतमुदीक्ष्य सः}
{हनूमान्वेगसंपन्नो जहर्ष च ननाद च} %||5-42-5||

\twolineshloka
{तं तोरणविटङ्कस्थं हनूमन्तं महाकपिम्}
{जम्बुमाली महाबाहुर्विव्याध निशितैः शरैः} %||5-42-6||

\twolineshloka
{अर्धचन्द्रेण वदने शिरस्येकेन कर्णिना}
{बाह्वोर्विव्याध नाराचैर्दशभिस्तं कपीश्वरम्} %||5-42-7||

\twolineshloka
{तस्य तच्छुशुभे ताम्रं शरेणाभिहतं मुखम्}
{शरदीवाम्बुजं फुल्लं विद्धं भास्कररश्मिना} %||5-42-8||

\twolineshloka
{चुकोप बाणाभिहतो राक्षसस्य महाकपिः}
{ततः पार्श्वेऽतिविपुलां ददर्श महतीं शिलाम्} %||5-42-9||

\twolineshloka
{तरसा तां समुत्पाट्य चिक्षेप बलवद्बली}
{तां शरैर्दशभिः क्रुद्धस्ताडयामास राक्षसः} %||5-42-10||

\twolineshloka
{विपन्नं कर्म तद्दृष्ट्वा हनूमांश्चण्डविक्रमः}
{सालं विपुलमुत्पाट्य भ्रामयामास वीर्यवान्} %||5-42-11||

\twolineshloka
{भ्रामयन्तं कपिं दृष्ट्वा सालवृक्षं महाबलम्}
{चिक्षेप सुबहून्बाणाञ्जम्बुमाली महाबलः} %||5-42-12||

\twolineshloka
{सालं चतुर्भिर्चिच्छेद वानरं पञ्चभिर्भुजे}
{उरस्येकेन बाणेन दशभिस्तु स्तनान्तरे} %||5-42-13||

\twolineshloka
{स शरैः पूरिततनुः क्रोधेन महता वृतः}
{तमेव परिघं गृह्य भ्रामयामास वेगितः} %||5-42-14||

\twolineshloka
{अतिवेगोऽतिवेगेन भ्रामयित्वा बलोत्कटः}
{परिघं पातयामास जम्बुमालेर्महोरसि} %||5-42-15||

\twolineshloka
{तस्य चैव शिरो नास्ति न बाहू न च जानुनी}
{न धनुर्न रथो नाश्वास्तत्रादृश्यन्त नेषवः} %||5-42-16||

\twolineshloka
{स हतस्तरसा तेन जम्बुमाली महारथः}
{पपात निहतो भूमौ चूर्णिताङ्गविभूषणः} %||5-42-17||

\twolineshloka
{जम्बुमालिं च निहतं किंकरांश्च महाबलान्}
{चुक्रोध रावणः श्रुत्वा कोपसंरक्तलोचनः} %||5-42-18||

\fourlineindentedshloka
{स रोषसंवर्तितताम्रलोचनः}
{ प्रहस्तपुत्रे निहते महाबले}
{अमात्यपुत्रानतिवीर्यविक्रमा}
{न्समादिदेशाशु निशाचरेश्वरः} %||5-43-19||


{॥इत्यार्षे श्रीमद्रामायणे वाल्मीकीये आदिकाव्ये सुन्दरकाण्डे द्विचत्वारिंशः सर्गः॥}

\sect{त्रिचत्वारिंशः सर्गः}

\twolineshloka
{ततस्ते राक्षसेन्द्रेण चोदिता मन्त्रिणः सुताः}
{निर्ययुर्भवनात्तस्मात्सप्त सप्तार्चिवर्चसः} %||5-43-1||

\twolineshloka
{महाबलपरीवारा धनुष्मन्तो महाबलाः}
{कृतास्त्रास्त्रविदां श्रेष्ठाः परस्परजयैषिणः} %||5-43-2||

\twolineshloka
{हेमजालपरिक्षिप्तैर्ध्वजवद्भिः पताकिभिः}
{तोयदस्वननिर्घोषैर्वाजियुक्तैर्महारथैः} %||5-43-3||

\twolineshloka
{तप्तकाञ्चनचित्राणि चापान्यमितविक्रमाः}
{विस्फारयन्तः संहृष्टास्तडिद्वन्त इवाम्बुदाः} %||5-43-4||

\twolineshloka
{जनन्यस्तास्ततस्तेषां विदित्वा किंकरान्हतान्}
{बभूवुः शोकसंभ्रान्ताः सबान्धवसुहृज्जनाः} %||5-43-5||

\twolineshloka
{ते परस्परसंघर्षास्तप्तकाञ्चनभूषणाः}
{अभिपेतुर्हनूमन्तं तोरणस्थमवस्थितम्} %||5-43-6||

\twolineshloka
{सृजन्तो बाणवृष्टिं ते रथगर्जितनिःस्वनाः}
{वृष्टिमन्त इवाम्भोदा विचेरुर्नैरृतर्षभाः} %||5-43-7||

\twolineshloka
{अवकीर्णस्ततस्ताभिर्हनूमाञ्शरवृष्टिभिः}
{अभवत्संवृताकारः शैलराडिव वृष्टिभिः} %||5-43-8||

\twolineshloka
{स शरान्वञ्चयामास तेषामाशुचरः कपिः}
{रथवेगांश्च वीराणां विचरन्विमलेऽम्बरे} %||5-43-9||

\twolineshloka
{स तैः क्रीडन्धनुष्मद्भिर्व्योम्नि वीरः प्रकाशते}
{धनुष्मद्भिर्यथा मेघैर्मारुतः प्रभुरम्बरे} %||5-43-10||

\twolineshloka
{स कृत्वा निनदं घोरं त्रासयंस्तां महाचमूम्}
{चकार हनुमान्वेगं तेषु रक्षःसु वीर्यवान्} %||5-43-11||

\twolineshloka
{तलेनाभिहनत्कांश्चित्पादैः कांश्चित्परंतपः}
{मुष्टिनाभ्यहनत्कांश्चिन्नखैः कांश्चिद्व्यदारयत्} %||5-43-12||

\twolineshloka
{प्रममाथोरसा कांश्चिदूरुभ्यामपरान्कपिः}
{केचित्तस्यैव नादेन तत्रैव पतिता भुवि} %||5-43-13||

\twolineshloka
{ततस्तेष्ववपन्नेषु भूमौ निपतितेषु च}
{तत्सैन्यमगमत्सर्वं दिशो दशभयार्दितम्} %||5-43-14||

\twolineshloka
{विनेदुर्विस्वरं नागा निपेतुर्भुवि वाजिनः}
{भग्ननीडध्वजच्छत्रैर्भूश्च कीर्णाभवद्रथैः} %||5-43-15||

\fourlineindentedshloka
{स तान्प्रवृद्धान्विनिहत्य राक्षसा}
{न्महाबलश्चण्डपराक्रमः कपिः}
{युयुत्सुरन्यैः पुनरेव राक्षसै}
{स्तदेव वीरोऽभिजगाम तोरणम्} %||5-44-16||


{॥इत्यार्षे श्रीमद्रामायणे वाल्मीकीये आदिकाव्ये सुन्दरकाण्डे त्रिचत्वारिंशः सर्गः॥}

\sect{चतुश्चत्वारिंशः सर्गः}

\twolineshloka
{हतान्मन्त्रिसुतान्बुद्ध्वा वानरेण महात्मना}
{रावणः संवृताकारश्चकार मतिमुत्तमाम्} %||5-44-1||

\twolineshloka
{स विरूपाक्षयूपाक्षौ दुर्धरं चैव राक्षसं}
{प्रघसं भासकर्णं च पञ्चसेनाग्रनायकान्} %||5-44-2||

\twolineshloka
{संदिदेश दशग्रीवो वीरान्नयविशारदान्}
{हनूमद्ग्रहणे व्यग्रान्वायुवेगसमान्युधि} %||5-44-3||

\twolineshloka
{यात सेनाग्रगाः सर्वे महाबलपरिग्रहाः}
{सवाजिरथमातङ्गाः स कपिः शास्यतामिति} %||5-44-4||

\twolineshloka
{यत्तैश्च खलु भाव्यं स्यात्तमासाद्य वनालयम्}
{कर्म चापि समाधेयं देशकालविरोधितम्} %||5-44-5||

\threelineshloka
{न ह्यहं तं कपिं मन्ये कर्मणा प्रतितर्कयन्}
{सर्वथा तन्महद्भूतं महाबलपरिग्रहम्}
{भवेदिन्द्रेण वा सृष्टमस्मदर्थं तपोबलात्} %||5-44-6||

\twolineshloka
{सनागयक्षगन्धर्वा देवासुरमहर्षयः}
{युष्माभिः सहितैः सर्वैर्मया सह विनिर्जिताः} %||5-44-7||

\twolineshloka
{तैरवश्यं विधातव्यं व्यलीकं किंचिदेव नः}
{तदेव नात्र संदेहः प्रसह्य परिगृह्यताम्} %||5-44-8||

\twolineshloka
{नावमन्यो भवद्भिश्च हरिः क्रूरपराक्रमः}
{दृष्टा हि हरयः शीघ्रा मया विपुलविक्रमाः} %||5-44-9||

\twolineshloka
{वाली च सह सुग्रीवो जाम्बवांश्च महाबलः}
{नीलः सेनापतिश्चैव ये चान्ये द्विविदादयः} %||5-44-10||

\twolineshloka
{नैव तेषां गतिर्भीमा न तेजो न पराक्रमः}
{न मतिर्न बलोत्साहो न रूपपरिकल्पनम्} %||5-44-11||

\twolineshloka
{महत्सत्त्वमिदं ज्ञेयं कपिरूपं व्यवस्थितम्}
{प्रयत्नं महदास्थाय क्रियतामस्य निग्रहः} %||5-44-12||

\twolineshloka
{कामं लोकास्त्रयः सेन्द्राः ससुरासुरमानवाः}
{भवतामग्रतः स्थातुं न पर्याप्ता रणाजिरे} %||5-44-13||

\twolineshloka
{तथापि तु नयज्ञेन जयमाकाङ्क्षता रणे}
{आत्मा रक्ष्यः प्रयत्नेन युद्धसिद्धिर्हि चञ्चला} %||5-44-14||

\twolineshloka
{ते स्वामिवचनं सर्वे प्रतिगृह्य महौजसः}
{समुत्पेतुर्महावेगा हुताशसमतेजसः} %||5-44-15||

\twolineshloka
{रथैश्च मत्तैर्नागैश्च वाजिभिश्च महाजवैः}
{शस्त्रैश्च विविधैस्तीक्ष्णैः सर्वैश्चोपचिता बलैः} %||5-44-16||

\twolineshloka
{ततस्तं ददृशुर्वीरा दीप्यमानं महाकपिम्}
{रश्मिमन्तमिवोद्यन्तं स्वतेजोरश्मिमालिनम्} %||5-44-17||

\twolineshloka
{तोरणस्थं महावेगं महासत्त्वं महाबलम्}
{महामतिं महोत्साहं महाकायं महाबलम्} %||5-44-18||

\twolineshloka
{तं समीक्ष्यैव ते सर्वे दिक्षु सर्वास्ववस्थिताः}
{तैस्तैः प्रहरणैर्भीमैरभिपेतुस्ततस्ततः} %||5-44-19||

\twolineshloka
{तस्य पञ्चायसास्तीक्ष्णाः सिताः पीतमुखाः शराः}
{शिरस्त्युत्पलपत्राभा दुर्धरेण निपातिताः} %||5-44-20||

\twolineshloka
{स तैः पञ्चभिराविद्धः शरैः शिरसि वानरः}
{उत्पपात नदन्व्योम्नि दिशो दश विनादयन्} %||5-44-21||

\twolineshloka
{ततस्तु दुर्धरो वीरः सरथः सज्जकार्मुकः}
{किरञ्शरशतैर्नैकैरभिपेदे महाबलः} %||5-44-22||

\twolineshloka
{स कपिर्वारयामास तं व्योम्नि शरवर्षिणम्}
{वृष्टिमन्तं पयोदान्ते पयोदमिव मारुतः} %||5-44-23||

\twolineshloka
{अर्द्यमानस्ततस्तेन दुर्धरेणानिलात्मजः}
{चकार निनदं भूयो व्यवर्धत च वेगवान्} %||5-44-24||

\twolineshloka
{स दूरं सहसोत्पत्य दुर्धरस्य रथे हरिः}
{निपपात महावेगो विद्युद्राशिर्गिराविव} %||5-44-25||

\twolineshloka
{ततस्तं मथिताष्टाश्वं रथं भग्नाक्षकूवरम्}
{विहाय न्यपतद्भूमौ दुर्धरस्त्यक्तजीवितः} %||5-44-26||

\twolineshloka
{तं विरूपाक्षयूपाक्षौ दृष्ट्वा निपतितं भुवि}
{संजातरोषौ दुर्धर्षावुत्पेततुररिंदमौ} %||5-44-27||

\twolineshloka
{स ताभ्यां सहसोत्पत्य विष्ठितो विमलेऽम्बरे}
{मुद्गराभ्यां महाबाहुर्वक्षस्यभिहतः कपिः} %||5-44-28||

\twolineshloka
{तयोर्वेगवतोर्वेगं विनिहत्य महाबलः}
{निपपात पुनर्भूमौ सुपर्णसमविक्रमः} %||5-44-29||

\twolineshloka
{स सालवृक्षमासाद्य समुत्पाट्य च वानरः}
{तावुभौ राक्षसौ वीरौ जघान पवनात्मजः} %||5-44-30||

\twolineshloka
{ततस्तांस्त्रीन्हताञ्ज्ञात्वा वानरेण तरस्विना}
{अभिपेदे महावेगः प्रसह्य प्रघसो हरिम्} %||5-44-31||

\twolineshloka
{भासकर्णश्च संक्रुद्धः शूलमादाय वीर्यवान्}
{एकतः कपिशार्दूलं यशस्विनमवस्थितौ} %||5-44-32||

\twolineshloka
{पट्टिशेन शिताग्रेण प्रघसः प्रत्यपोथयत्}
{भासकर्णश्च शूलेन राक्षसः कपिसत्तमम्} %||5-44-33||

\twolineshloka
{स ताभ्यां विक्षतैर्गात्रैरसृग्दिग्धतनूरुहः}
{अभवद्वानरः क्रुद्धो बालसूर्यसमप्रभः} %||5-44-34||

\twolineshloka
{समुत्पाट्य गिरेः शृङ्गं समृगव्यालपादपम्}
{जघान हनुमान्वीरो राक्षसौ कपिकुञ्जरः} %||5-44-35||

\twolineshloka
{ततस्तेष्ववसन्नेषु सेनापतिषु पञ्चसु}
{बलं तदवशेषं तु नाशयामास वानरः} %||5-44-36||

\twolineshloka
{अश्वैरश्वान्गजैर्नागान्योधैर्योधान्रथै रथान्}
{स कपिर्नाशयामास सहस्राक्ष इवासुरान्} %||5-44-37||

\twolineshloka
{हतैर्नागैश्च तुरगैर्भग्नाक्षैश्च महारथैः}
{हतैश्च राक्षसैर्भूमी रुद्धमार्गा समन्ततः} %||5-44-38||

\fourlineindentedshloka
{ततः कपिस्तान्ध्वजिनीपतीन्रणे}
{ निहत्य वीरान्सबलान्सवाहनान्}
{तदेव वीरः परिगृह्य तोरणं}
{ कृतक्षणः काल इव प्रजाक्षये} %||5-45-39||


{॥इत्यार्षे श्रीमद्रामायणे वाल्मीकीये आदिकाव्ये सुन्दरकाण्डे चतुश्चत्वारिंशः सर्गः॥}

\sect{पञ्चचत्वारिंशः सर्गः}

\fourlineindentedshloka
{सेनापतीन्पञ्च स तु प्रमापिता}
{न्हनूमता सानुचरान्सवाहनान्}
{समीक्ष्य राजा समरोद्धतोन्मुखं}
{ कुमारमक्षं प्रसमैक्षताक्षतम्} %||5-45-1||

\fourlineindentedshloka
{स तस्य दृष्ट्यर्पणसंप्रचोदितः}
{ प्रतापवान्काञ्चनचित्रकार्मुकः}
{समुत्पपाताथ सदस्युदीरितो}
{ द्विजातिमुख्यैर्हविषेव पावकः} %||5-45-2||

\fourlineindentedshloka
{ततो महद्बालदिवाकरप्रभं}
{ प्रतप्तजाम्बूनदजालसंततम्}
{रथां समास्थाय ययौ स वीर्यवा}
{न्महाहरिं तं प्रति नैरृतर्षभः} %||5-45-3||

\fourlineindentedshloka
{ततस्तपःसंग्रहसंचयार्जितं}
{ प्रतप्तजाम्बूनदजालशोभितम्}
{पताकिनं रत्नविभूषितध्वजं}
{ मनोजवाष्टाश्ववरैः सुयोजितम्} %||5-45-4||

\fourlineindentedshloka
{सुरासुराधृष्यमसंगचारिणं}
{ रविप्रभं व्योमचरं समाहितम्}
{सतूणमष्टासिनिबद्धबन्धुरं}
{ यथाक्रमावेशितशक्तितोमरम्} %||5-45-5||

\fourlineindentedshloka
{विराजमानं प्रतिपूर्णवस्तुना}
{ सहेमदाम्ना शशिसूर्यवर्वसा}
{दिवाकराभं रथमास्थितस्ततः}
{ स निर्जगामामरतुल्यविक्रमः} %||5-45-6||

\fourlineindentedshloka
{स पूरयन्खं च महीं च साचलां}
{ तुरंगमतङ्गमहारथस्वनैः}
{बलैः समेतैः स हि तोरणस्थितं}
{ समर्थमासीनमुपागमत्कपिम्} %||5-45-7||

\fourlineindentedshloka
{स तं समासाद्य हरिं हरीक्षणो}
{ युगान्तकालाग्निमिव प्रजाक्षये}
{अवस्थितं विस्मितजातसंभ्रमः}
{ समैक्षताक्षो बहुमानचक्षुषा} %||5-45-8||

\fourlineindentedshloka
{स तस्य वेगं च कपेर्महात्मनः}
{ पराक्रमं चारिषु पार्थिवात्मजः}
{विचारयन्खं च बलं महाबलो}
{ हिमक्षये सूर्य इवाभिवर्धते} %||5-45-9||

\fourlineindentedshloka
{स जातमन्युः प्रसमीक्ष्य विक्रमं}
{ स्थिरः स्थितः संयति दुर्निवारणम्}
{समाहितात्मा हनुमन्तमाहवे}
{ प्रचोदयामास शरैस्त्रिभिः शितैः} %||5-45-10||

\fourlineindentedshloka
{ततः कपिं तं प्रसमीक्ष्य गर्वितं}
{ जितश्रमं शत्रुपराजयोर्जितम्}
{अवैक्षताक्षः समुदीर्णमानसः}
{ सबाणपाणिः प्रगृहीतकार्मुकः} %||5-45-11||

\fourlineindentedshloka
{स हेमनिष्काङ्गदचारुकुण्डलः}
{ समाससादाशु पराक्रमः कपिम्}
{तयोर्बभूवाप्रतिमः समागमः}
{ सुरासुराणामपि संभ्रमप्रदः} %||5-45-12||

\fourlineindentedshloka
{ररास भूमिर्न तताप भानुमा}
{न्ववौ न वायुः प्रचचाल चाचलः}
{कपेः कुमारस्य च वीक्ष्य संयुगं}
{ ननाद च द्यौरुदधिश्च चुक्षुभे} %||5-45-13||

\fourlineindentedshloka
{ततः स वीरः सुमुखान्पतत्रिणः}
{ सुवर्णपुङ्खान्सविषानिवोरगान्}
{समाधिसंयोगविमोक्षतत्त्ववि}
{च्छरानथ त्रीन्कपिमूर्ध्न्यपातयत्} %||5-45-14||

\fourlineindentedshloka
{स तैः शरैर्मूर्ध्नि समं निपातितैः}
{ क्षरन्नसृग्दिग्धविवृत्तलोचनः}
{नवोदितादित्यनिभः शरांशुमा}
{न्व्यराजतादित्य इवांशुमालिकः} %||5-45-15||

\fourlineindentedshloka
{ततः स पिङ्गाधिपमन्त्रिसत्तमः}
{ समीक्ष्य तं राजवरात्मजं रणे}
{उदग्रचित्रायुधचित्रकार्मुकं}
{ जहर्ष चापूर्यत चाहवोन्मुखः} %||5-45-16||

\fourlineindentedshloka
{स मन्दराग्रस्थ इवांशुमाली}
{ विवृद्धकोपो बलवीर्यसंयुतः}
{कुमारमक्षं सबलं सवाहनं}
{ ददाह नेत्राग्निमरीचिभिस्तदा} %||5-45-17||

\fourlineindentedshloka
{ततः स बाणासनशक्रकार्मुकः}
{ शरप्रवर्षो युधि राक्षसाम्बुदः}
{शरान्मुमोचाशु हरीश्वराचले}
{ बलाहको वृष्टिमिवाचलोत्तमे} %||5-45-18||

\fourlineindentedshloka
{ततः कपिस्तं रणचण्डविक्रमं}
{ विवृद्धतेजोबलवीर्यसायकम्}
{कुमारमक्षं प्रसमीक्ष्य संयुगे}
{ ननाद हर्षाद्घनतुल्यविक्रमः} %||5-45-19||

\fourlineindentedshloka
{स बालभावाद्युधि वीर्यदर्पितः}
{ प्रवृद्धमन्युः क्षतजोपमेक्षणः}
{समाससादाप्रतिमं रणे कपिं}
{ गजो महाकूपमिवावृतं तृणैः} %||5-45-20||

\fourlineindentedshloka
{स तेन बाणैः प्रसभं निपातितै}
{श्चकार नादं घननादनिःस्वनः}
{समुत्पपाताशु नभः स मारुति}
{र्भुजोरुविक्षेपण घोरदर्शनः} %||5-45-21||

\fourlineindentedshloka
{समुत्पतन्तं समभिद्रवद्बली}
{ स राक्षसानां प्रवरः प्रतापवान्}
{रथी रथश्रेष्ठतमः किरञ्शरैः}
{ पयोधरः शैलमिवाश्मवृष्टिभिः} %||5-45-22||

\fourlineindentedshloka
{स ताञ्शरांस्तस्य विमोक्षयन्कपि}
{श्चचार वीरः पथि वायुसेविते}
{शरान्तरे मारुतवद्विनिष्पत}
{न्मनोजवः संयति चण्डविक्रमः} %||5-45-23||

\fourlineindentedshloka
{तमात्तबाणासनमाहवोन्मुखं}
{ खमास्तृणन्तं विविधैः शरोत्तमैः}
{अवैक्षताक्षं बहुमानचक्षुषा}
{ जगाम चिन्तां च स मारुतात्मजः} %||5-45-24||

\fourlineindentedshloka
{ततः शरैर्भिन्नभुजान्तरः कपिः}
{ कुमारवर्येण महात्मना नदन्}
{महाभुजः कर्मविशेषतत्त्ववि}
{द्विचिन्तयामास रणे पराक्रमम्} %||5-45-25||

\fourlineindentedshloka
{अबालवद्बालदिवाकरप्रभः}
{ करोत्ययं कर्म महन्महाबलः}
{न चास्य सर्वाहवकर्मशोभिनः}
{ प्रमापणे मे मतिरत्र जायते} %||5-45-26||

\fourlineindentedshloka
{अयं महात्मा च महांश्च वीर्यतः}
{ समाहितश्चातिसहश्च संयुगे}
{असंशयं कर्मगुणोदयादयं}
{ सनागयक्षैर्मुनिभिश्च पूजितः} %||5-45-27||

\fourlineindentedshloka
{पराक्रमोत्साहविवृद्धमानसः}
{ समीक्षते मां प्रमुखागतः स्थितः}
{पराक्रमो ह्यस्य मनांसि कम्पये}
{त्सुरासुराणामपि शीघ्रकारिणः} %||5-45-28||

\fourlineindentedshloka
{न खल्वयं नाभिभवेदुपेक्षितः}
{ पराक्रमो ह्यस्य रणे विवर्धते}
{प्रमापणं त्वेव ममास्य रोचते}
{ न वर्धमानोऽग्निरुपेक्षितुं क्षमः} %||5-45-29||

\fourlineindentedshloka
{इति प्रवेगं तु परस्य तर्कय}
{न्स्वकर्मयोगं च विधाय वीर्यवान्}
{चकार वेगं तु महाबलस्तदा}
{ मतिं च चक्रेऽस्य वधे महाकपिः} %||5-45-30||

\fourlineindentedshloka
{स तस्य तानष्टहयान्महाजवा}
{न्समाहितान्भारसहान्विवर्तने}
{जघान वीरः पथि वायुसेविते}
{ तलप्रहालैः पवनात्मजः कपिः} %||5-45-31||

\fourlineindentedshloka
{ततस्तलेनाभिहतो महारथः}
{ स तस्य पिङ्गाधिपमन्त्रिनिर्जितः}
{स भग्ननीडः परिमुक्तकूबरः}
{ पपात भूमौ हतवाजिरम्बरात्} %||5-45-32||

\fourlineindentedshloka
{स तं परित्यज्य महारथो रथं}
{ सकार्मुकः खड्गधरः खमुत्पतत्}
{तपोऽभियोगादृषिरुग्रवीर्यवा}
{न्विहाय देहं मरुतामिवालयम्} %||5-45-33||

\fourlineindentedshloka
{ततः कपिस्तं विचरन्तमम्बरे}
{ पतत्रिराजानिलसिद्धसेविते}
{समेत्य तं मारुतवेगविक्रमः}
{ क्रमेण जग्राह च पादयोर्दृढम्} %||5-45-34||

\fourlineindentedshloka
{स तं समाविध्य सहस्रशः कपि}
{र्महोरगं गृह्य इवाण्डजेश्वरः}
{मुमोच वेगात्पितृतुल्यविक्रमो}
{ महीतले संयति वानरोत्तमः} %||5-45-35||

\fourlineindentedshloka
{स भग्नबाहूरुकटीशिरो धरः}
{ क्षरन्नसृन्निर्मथितास्थिलोचनः}
{स भिन्नसंधिः प्रविकीर्णबन्धनो}
{ हतः क्षितौ वायुसुतेन राक्षसः} %||5-45-36||

{महाकपिर्भूमितले निपीड्य तं; चकार रक्षोऽधिपतेर्महद्भयम्॥\devanumber{37}॥} %||5-45-37|| (Check)

\addtocounter{shlokacount}{1}

\fourlineindentedshloka
{महर्षिभिश्चक्रचरैर्महाव्रतैः}
{ समेत्य भूतैश्च सयक्षपन्नगैः}
{सुरैश्च सेन्द्रैर्भृशजातविस्मयै}
{र्हते कुमारे स कपिर्निरीक्षितः} %||5-45-38||

\fourlineindentedshloka
{निहत्य तं वज्रसुतोपमप्रभं}
{ कुमारमक्षं क्षतजोपमेक्षणम्}
{तदेव वीरोऽभिजगाम तोरणं}
{ कृतक्षणः काल इव प्रजाक्षये} %||5-46-39||


{॥इत्यार्षे श्रीमद्रामायणे वाल्मीकीये आदिकाव्ये सुन्दरकाण्डे पञ्चचत्वारिंशः सर्गः॥}

\sect{षट्चत्वारिंशः सर्गः}

\fourlineindentedshloka
{ततस्तु रक्षोऽधिपतिर्महात्मा}
{ हनूमताक्षे निहते कुमारे}
{मनः समाधाय तदेन्द्रकल्पं}
{ समादिदेशेन्द्रजितं स रोषात्} %||5-46-1||

\fourlineindentedshloka
{त्वमस्त्रविच्छस्त्रभृतां वरिष्ठः}
{ सुरासुराणामपि शोकदाता}
{सुरेषु सेन्द्रेषु च दृष्टकर्मा}
{ पितामहाराधनसंचितास्त्रः} %||5-46-2||

\twolineshloka
{तवास्त्रबलमासाद्य नासुरा न मरुद्गणाः}
{न कश्चित्त्रिषु लोकेषु संयुगे न गतश्रमः} %||5-46-3||

\twolineshloka
{भुजवीर्याभिगुप्तश्च तपसा चाभिरक्षितः}
{देशकालविभागज्ञस्त्वमेव मतिसत्तमः} %||5-46-4||

\fourlineindentedshloka
{न तेऽस्त्यशक्यं समरेषु कर्मणा}
{ न तेऽस्त्यकार्यं मतिपूर्वमन्त्रणे}
{न सोऽस्ति कश्चित्त्रिषु संग्रहेषु वै}
{ न वेद यस्तेऽस्त्रबलं बलं च ते} %||5-46-5||

\fourlineindentedshloka
{ममानुरूपं तपसो बलं च ते}
{ पराक्रमश्चास्त्रबलं च संयुगे}
{न त्वां समासाद्य रणावमर्दे}
{ मनः श्रमं गच्छति निश्चितार्थम्} %||5-46-6||

\twolineshloka
{निहता इंकराः सर्वे जम्बुमाली च राक्षसः}
{अमात्यपुत्रा वीराश्च पञ्च सेनाग्रयायिनः} %||5-46-7||

\twolineshloka
{सहोदरस्ते दयितः कुमारोऽक्षश्च सूदितः}
{न तु तेष्वेव मे सारो यस्त्वय्यरिनिषूदन} %||5-46-8||

\fourlineindentedshloka
{इदं हि दृष्ट्वा मतिमन्महद्बलं}
{ कपेः प्रभावं च पराक्रमं च}
{त्वमात्मनश्चापि समीक्ष्य सारं}
{ कुरुष्व वेगं स्वबलानुरूपम्} %||5-46-9||

\fourlineindentedshloka
{बलावमर्दस्त्वयि संनिकृष्टे}
{ यथा गते शाम्यति शान्तशत्रौ}
{तथा समीक्ष्यात्मबलं परं च}
{ समारभस्वास्त्रविदां वरिष्ठ} %||5-46-10||

\twolineshloka
{न खल्वियं मतिः श्रेष्ठा यत्त्वां संप्रेषयाम्यहम्}
{इयं च राजधर्माणां क्षत्रस्य च मतिर्मता} %||5-46-11||

\twolineshloka
{नानाशस्त्रैश्च संग्रामे वैशारद्यमरिंदम}
{अवश्यमेव बोद्धव्यं काम्यश्च विजयो रणे} %||5-46-12||

\fourlineindentedshloka
{ततः पितुस्तद्वचनं निशम्य}
{ प्रदक्षिणं दक्षसुतप्रभावः}
{चकार भर्तारमदीनसत्त्वो}
{ रणाय वीरः प्रतिपन्नबुद्धिः} %||5-46-13||

\twolineshloka
{ततस्तैः स्वगणैरिष्टैरिन्द्रजित्प्रतिपूजितः}
{युद्धोद्धतकृतोत्साहः संग्रामं प्रतिपद्यत} %||5-46-14||

\twolineshloka
{श्रीमान्पद्मपलाशाक्षो राक्षसाधिपतेः सुतः}
{निर्जगाम महातेजाः समुद्र इव पर्वसु} %||5-46-15||

\fourlineindentedshloka
{स पक्षि राजोपमतुल्यवेगै}
{र्व्यालैश्चतुर्भिः सिततीक्ष्णदंष्ट्रैः}
{रथं समायुक्तमसंगवेगं}
{ समारुरोहेन्द्रजिदिन्द्रकल्पः} %||5-46-16||

\twolineshloka
{स रथी धन्विनां श्रेष्ठः शस्त्रज्ञोऽस्त्रविदां वरः}
{रथेनाभिययौ क्षिप्रं हनूमान्यत्र सोऽभवत्} %||5-46-17||

\twolineshloka
{स तस्य रथनिर्घोषं ज्यास्वनं कार्मुकस्य च}
{निशम्य हरिवीरोऽसौ संप्रहृष्टतरोऽभवत्} %||5-46-18||

\twolineshloka
{सुमहच्चापमादाय शितशल्यांश्च सायकान्}
{हनूमन्तमभिप्रेत्य जगाम रणपण्डितः} %||5-46-19||

\fourlineindentedshloka
{तस्मिंस्ततः संयति जातहर्षे}
{ रणाय निर्गच्छति बाणपाणौ}
{दिशश्च सर्वाः कलुषा बभूवु}
{र्मृगाश्च रौद्रा बहुधा विनेदुः} %||5-46-20||

\fourlineindentedshloka
{समागतास्तत्र तु नागयक्षा}
{ महर्षयश्चक्रचराश्च सिद्धाः}
{नभः समावृत्य च पक्षिसंघा}
{ विनेदुरुच्चैः परमप्रहृष्टाः} %||5-46-21||

\twolineshloka
{आयन्तं सरथं दृष्ट्वा तूर्णमिन्द्रजितं कपिः}
{विननाद महानादं व्यवर्धत च वेगवान्} %||5-46-22||

\twolineshloka
{इन्द्रजित्तु रथं दिव्यमास्थितश्चित्रकार्मुकः}
{धनुर्विस्फारयामास तडिदूर्जितनिःस्वनम्} %||5-46-23||

\fourlineindentedshloka
{ततः समेतावतितीक्ष्णवेगौ}
{ महाबलौ तौ रणनिर्विशङ्कौ}
{कपिश्च रक्षोऽधिपतेश्च पुत्रः}
{ सुरासुरेन्द्राविव बद्धवैरौ} %||5-46-24||

\fourlineindentedshloka
{स तस्य वीरस्य महारथस्या}
{ धनुष्मतः संयति संमतस्य}
{शरप्रवेगं व्यहनत्प्रवृद्ध}
{श्चचार मार्गे पितुरप्रमेयः} %||5-46-25||

\fourlineindentedshloka
{ततः शरानायततीक्ष्णशल्या}
{न्सुपत्रिणः काञ्चनचित्रपुङ्खान्}
{मुमोच वीरः परवीरहन्ता}
{ सुसंततान्वज्रनिपातवेगान्} %||5-46-26||

\fourlineindentedshloka
{स तस्य तत्स्यन्दननिःस्वनं च}
{ मृदङ्गभेरीपटहस्वनं च}
{विकृष्यमाणस्य च कार्मुकस्य}
{ निशम्य घोषं पुनरुत्पपात} %||5-46-27||

\twolineshloka
{शराणामन्तरेष्वाशु व्यवर्तत महाकपिः}
{हरिस्तस्याभिलक्षस्य मोक्षयँल्लक्ष्यसंग्रहम्} %||5-46-28||

\twolineshloka
{शराणामग्रतस्तस्य पुनः समभिवर्तत}
{प्रसार्य हस्तौ हनुमानुत्पपातानिलात्मजः} %||5-46-29||

\twolineshloka
{तावुभौ वेगसंपन्नौ रणकर्मविशारदौ}
{सर्वभूतमनोग्राहि चक्रतुर्युद्धमुत्तमम्} %||5-46-30||

\fourlineindentedshloka
{हनूमतो वेद न राक्षसोऽन्तरं}
{ न मारुतिस्तस्य महात्मनोऽन्तरम्}
{परस्परं निर्विषहौ बभूवतुः}
{ समेत्य तौ देवसमानविक्रमौ} %||5-46-31||

\fourlineindentedshloka
{ततस्तु लक्ष्ये स विहन्यमाने}
{ शरेषु मोघेषु च संपतत्सु}
{जगाम चिन्तां महतीं महात्मा}
{ समाधिसंयोगसमाहितात्मा} %||5-46-32||

\fourlineindentedshloka
{ततो मतिं राक्षसराजसूनु}
{श्चकार तस्मिन्हरिवीरमुख्ये}
{अवध्यतां तस्य कपेः समीक्ष्य}
{ कथं निगच्छेदिति निग्रहार्थम्} %||5-46-33||

\twolineshloka
{ततः पैतामहां वीरः सोऽस्त्रमस्त्रविदां वरः}
{संदधे सुमहातेजास्तं हरिप्रवरं प्रति} %||5-46-34||

\twolineshloka
{अवध्योऽयमिति ज्ञात्वा तमस्त्रेणास्त्रतत्त्ववित्}
{निजग्राह महाबाहुर्मारुतात्मजमिन्द्रजित्} %||5-46-35||

\twolineshloka
{तेन बद्धस्ततोऽस्त्रेण राक्षसेन स वानरः}
{अभवन्निर्विचेष्टश्च पपात च महीतले} %||5-46-36||

\fourlineindentedshloka
{ततोऽथ बुद्ध्वा स तदास्त्रबन्धं}
{ प्रभोः प्रभावाद्विगताल्पवेगः}
{पितामहानुग्रहमात्मनश्च}
{ विचिन्तयामास हरिप्रवीरः} %||5-46-37||

\twolineshloka
{ततः स्वायम्भुवैर्मन्त्रैर्ब्रह्मास्त्रमभिमन्त्रितम्}
{हनूमांश्चिन्तयामास वरदानं पितामहात्} %||5-46-38||

\fourlineindentedshloka
{न मेऽस्त्रबन्धस्य च शक्तिरस्ति}
{ विमोक्षणे लोकगुरोः प्रभावात्}
{इत्येवमेवंविहितोऽस्त्रबन्धो}
{ मयात्मयोनेरनुवर्तितव्यः} %||5-46-39||

\fourlineindentedshloka
{स वीर्यमस्त्रस्य कपिर्विचार्य}
{ पितामहानुग्रहमात्मनश्च}
{विमोक्षशक्तिं परिचिन्तयित्वा}
{ पितामहाज्ञामनुवर्तते स्म} %||5-46-40||

\twolineshloka
{अस्त्रेणापि हि बद्धस्य भयं मम न जायते}
{पितामहमहेन्द्राभ्यां रक्षितस्यानिलेन च} %||5-46-41||

\twolineshloka
{ग्रहणे चापि रक्षोभिर्महन्मे गुणदर्शनम्}
{राक्षसेन्द्रेण संवादस्तस्माद्गृह्णन्तु मां परे} %||5-46-42||

\fourlineindentedshloka
{स निश्चितार्थः परवीरहन्ता}
{ समीक्ष्य करी विनिवृत्तचेष्टः}
{परैः प्रसह्याभिगतैर्निगृह्य}
{ ननाद तैस्तैः परिभर्त्स्यमानः} %||5-46-43||

\twolineshloka
{ततस्तं राक्षसा दृष्ट्वा निर्विचेष्टमरिंदमम्}
{बबन्धुः शणवल्कैश्च द्रुमचीरैश्च संहतैः} %||5-46-44||

\fourlineindentedshloka
{स रोचयामास परैश्च बन्धनं}
{ प्रसह्य वीरैरभिनिग्रहं च}
{कौतूहलान्मां यदि राक्षसेन्द्रो}
{ द्रष्टुं व्यवस्येदिति निश्चितार्थः} %||5-46-45||

\twolineshloka
{स बद्धस्तेन वल्केन विमुक्तोऽस्त्रेण वीर्यवान्}
{अस्त्रबन्धः स चान्यं हि न बन्धमनुवर्तते} %||5-46-46||

\fourlineindentedshloka
{अथेन्द्रजित्तं द्रुमचीरबन्धं}
{ विचार्य वीरः कपिसत्तमं तम्}
{विमुक्तमस्त्रेण जगाम चिन्ता}
{मन्येन बद्धो ह्यनुवर्ततेऽस्त्रम्} %||5-46-47||

\fourlineindentedshloka
{अहो महत्कर्म कृतं निरर्थकं}
{ न राक्षसैर्मन्त्रगतिर्विमृष्टा}
{पुनश्च नास्त्रे विहतेऽस्त्रमन्य}
{त्प्रवर्तते संशयिताः स्म सर्वे} %||5-46-48||

\twolineshloka
{अस्त्रेण हनुमान्मुक्तो नात्मानमवबुध्यते}
{कृष्यमाणस्तु रक्षोभिस्तैश्च बन्धैर्निपीडितः} %||5-46-49||

\twolineshloka
{हन्यमानस्ततः क्रूरै राक्षसैः काष्ठमुष्टिभिः}
{समीपं राक्षसेन्द्रस्य प्राकृष्यत स वानरः} %||5-46-50||

\fourlineindentedshloka
{अथेन्द्रजित्तं प्रसमीक्ष्य मुक्त}
{मस्त्रेण बद्धं द्रुमचीरसूत्रैः}
{व्यदर्शयत्तत्र महाबलं तं}
{ हरिप्रवीरं सगणाय राज्ञे} %||5-46-51||

\twolineshloka
{तं मत्तमिव मातङ्गं बद्धं कपिवरोत्तमम्}
{राक्षसा राक्षसेन्द्राय रावणाय न्यवेदयन्} %||5-46-52||

\twolineshloka
{कोऽयं कस्य कुतो वापि किं कार्यं को व्यपाश्रयः}
{इति राक्षसवीराणां तत्र संजज्ञिरे कथाः} %||5-46-53||

\twolineshloka
{हन्यतां दह्यतां वापि भक्ष्यतामिति चापरे}
{राक्षसास्तत्र संक्रुद्धाः परस्परमथाब्रुवन्} %||5-46-54||

\fourlineindentedshloka
{अतीत्य मार्गं सहसा महात्मा}
{ स तत्र रक्षोऽधिपपादमूले}
{ददर्श राज्ञः परिचारवृद्धा}
{न्गृहं महारत्नविभूषितं च} %||5-46-55||

\twolineshloka
{स ददर्श महातेजा रावणः कपिसत्तमम्}
{रक्षोभिर्विकृताकारैः कृष्यमाणमितस्ततः} %||5-46-56||

\twolineshloka
{राक्षसाधिपतिं चापि ददर्श कपिसत्तमः}
{तेजोबलसमायुक्तं तपन्तमिव भास्करम्} %||5-46-57||

\fourlineindentedshloka
{स रोषसंवर्तितताम्रदृष्टि}
{र्दशाननस्तं कपिमन्ववेक्ष्य}
{अथोपविष्टान्कुलशीलवृद्धा}
{न्समादिशत्तं प्रति मन्त्रमुख्यान्} %||5-46-58||

\fourlineindentedshloka
{यथाक्रमं तैः स कपिश्च पृष्टः}
{ कार्यार्थमर्थस्य च मूलमादौ}
{निवेदयामास हरीश्वरस्य}
{ दूतः सकाशादहमागतोऽस्मि} %||5-47-59||


{॥इत्यार्षे श्रीमद्रामायणे वाल्मीकीये आदिकाव्ये सुन्दरकाण्डे षट्चत्वारिंशः सर्गः॥}

\sect{सप्तचत्वारिंशः सर्गः}

\twolineshloka
{ततः स कर्मणा तस्य विस्मितो भीमविक्रमः}
{हनुमान्रोषताम्राक्षो रक्षोऽधिपमवैक्षत} %||5-47-1||

\twolineshloka
{भाजमानं महार्हेण काञ्चनेन विराजता}
{मुक्ताजालावृतेनाथ मुकुटेन महाद्युतिम्} %||5-47-2||

\twolineshloka
{वज्रसंयोगसंयुक्तैर्महार्हमणिविग्रहैः}
{हैमैराभरणैश्चित्रैर्मनसेव प्रकल्पितैः} %||5-47-3||

\twolineshloka
{महार्हक्षौमसंवीतं रक्तचन्दनरूषितम्}
{स्वनुलिप्तं विचित्राभिर्विविधभिश्च भक्तिभिः} %||5-47-4||

\twolineshloka
{विपुलैर्दर्शनीयैश्च रक्षाक्षैर्भीमदर्शनैः}
{दीप्ततीक्ष्णमहादंष्ट्रैः प्रलम्बदशनच्छदैः} %||5-47-5||

\twolineshloka
{शिरोभिर्दशभिर्वीरं भ्राजमानं महौजसं}
{नानाव्यालसमाकीर्णैः शिखरैरिव मन्दरम्} %||5-47-6||

\twolineshloka
{नीलाञ्जनचय प्रख्यं हारेणोरसि राजता}
{पूर्णचन्द्राभवक्त्रेण सबलाकमिवाम्बुदम्} %||5-47-7||

\twolineshloka
{बाहुभिर्बद्धकेयूरैश्चन्दनोत्तमरूषितैः}
{भ्राजमानाङ्गदैः पीनैः पञ्चशीर्षैरिवोरगैः} %||5-47-8||

\twolineshloka
{महति स्फाटिके चित्रे रत्नसंयोगसंस्कृते}
{उत्तमास्तरणास्तीर्णे उपविष्टं वरासने} %||5-47-9||

\twolineshloka
{अलंकृताभिरत्यर्थं प्रमदाभिः समन्ततः}
{वालव्यजनहस्ताभिरारात्समुपसेवितम्} %||5-47-10||

\twolineshloka
{दुर्धरेण प्रहस्तेन महापार्श्वेन रक्षसा}
{मन्त्रिभिर्मन्त्रतत्त्वज्ञैर्निकुम्भेन च मन्त्रिणा} %||5-47-11||

\twolineshloka
{उपोपविष्टं रक्षोभिश्चतुर्भिर्बलदर्पितैः}
{कृत्स्नैः परिवृतं लोकं चतुर्भिरिव सागरैः} %||5-47-12||

\twolineshloka
{मन्त्रिभिर्मन्त्रतत्त्वज्ञैरन्यैश्च शुभबुद्धिभिः}
{अन्वास्यमानं सचिवैः सुरैरिव सुरेश्वरम्} %||5-47-13||

\twolineshloka
{अपश्यद्राक्षसपतिं हनूमानतितेजसं}
{विष्ठितं मेरुशिखरे सतोयमिव तोयदम्} %||5-47-14||

\twolineshloka
{स तैः संपीड्यमानोऽपि रक्षोभिर्भीमविक्रमैः}
{विस्मयं परमं गत्वा रक्षोऽधिपमवैक्षत} %||5-47-15||

\twolineshloka
{भ्राजमानं ततो दृष्ट्वा हनुमान्राक्षसेश्वरम्}
{मनसा चिन्तयामास तेजसा तस्य मोहितः} %||5-47-16||

\twolineshloka
{अहो रूपमहो धैर्यमहो सत्त्वमहो द्युतिः}
{अहो राक्षसराजस्य सर्वलक्षणयुक्तता} %||5-47-17||

\twolineshloka
{यद्यधर्मो न बलवान्स्यादयं राक्षसेश्वरः}
{स्यादयं सुरलोकस्य सशक्रस्यापि रक्षिता} %||5-47-18||

\twolineshloka
{तेन बिभ्यति खल्वस्माल्लोकाः सामरदानवाः}
{अयं ह्युत्सहते क्रुद्धः कर्तुमेकार्णवं जगत्} %||5-47-19||

\twolineshloka
{इति चिन्तां बहुविधामकरोन्मतिमान्कपिः}
{दृष्ट्वा राक्षसराजस्य प्रभावममितौजसः} %||5-48-20||


{॥इत्यार्षे श्रीमद्रामायणे वाल्मीकीये आदिकाव्ये सुन्दरकाण्डे सप्तचत्वारिंशः सर्गः॥}

\sect{अष्टचत्वारिंशः सर्गः}

\twolineshloka
{तमुद्वीक्ष्य महाबाहुः पिङ्गाक्षं पुरतः स्थितम्}
{रोषेण महताविष्टो रावणो लोकरावणः} %||5-48-1||

\twolineshloka
{स राजा रोषताम्राक्षः प्रहस्तं मन्त्रिसत्तमम्}
{कालयुक्तमुवाचेदं वचो विपुलमर्थवत्} %||5-48-2||

\twolineshloka
{दुरात्मा पृच्छ्यतामेष कुतः किं वास्य कारणम्}
{वनभङ्गे च कोऽस्यार्थो राक्षसीनां च तर्जने} %||5-48-3||

\twolineshloka
{रावणस्य वचः श्रुत्वा प्रहस्तो वाक्यमब्रवीत्}
{समाश्वसिहि भद्रं ते न भीः कार्या त्वया कपे} %||5-48-4||

\twolineshloka
{यदि तावत्त्वमिन्द्रेण प्रेषितो रावणालयम्}
{तत्त्वमाख्याहि मा ते भूद्भयं वानर मोक्ष्यसे} %||5-48-5||

\twolineshloka
{यदि वैश्रवणस्य त्वं यमस्य वरुणस्य च}
{चारुरूपमिदं कृत्वा यमस्य वरुणस्य च} %||5-48-6||

\twolineshloka
{विष्णुना प्रेषितो वापि दूतो विजयकाङ्क्षिणा}
{न हि ते वानरं तेजो रूपमात्रं तु वानरम्} %||5-48-7||

\twolineshloka
{तत्त्वतः कथयस्वाद्य ततो वानर मोक्ष्यसे}
{अनृतं वदतश्चापि दुर्लभं तव जीवितम्} %||5-48-8||

{अथ वा यन्निमित्तस्ते प्रवेशो रावणालये॥\devanumber{9}॥} %||5-48-9|| (Check)

\addtocounter{shlokacount}{1}

\twolineshloka
{एवमुक्तो हरिवरस्तदा रक्षोगणेश्वरम्}
{अब्रवीन्नास्मि शक्रस्य यमस्य वरुणस्य वा} %||5-48-10||

\twolineshloka
{धनदेन न मे सख्यं विष्णुना नास्मि चोदितः}
{जातिरेव मम त्वेषा वानरोऽहमिहागतः} %||5-48-11||

\twolineshloka
{दर्शने राक्षसेन्द्रस्य दुर्लभे तदिदं मया}
{वनं राक्षसराजस्य दर्शनार्थे विनाशितम्} %||5-48-12||

\twolineshloka
{ततस्ते राक्षसाः प्राप्ता बलिनो युद्धकाङ्क्षिणः}
{रक्षणार्थं च देहस्य प्रतियुद्धा मया रणे} %||5-48-13||

\twolineshloka
{अस्त्रपाशैर्न शक्योऽहं बद्धुं देवासुरैरपि}
{पितामहादेव वरो ममाप्येषोऽभ्युपागतः} %||5-48-14||

\twolineshloka
{राजानं द्रष्टुकामेन मयास्त्रमनुवर्तितम्}
{विमुक्तो अहमस्त्रेण राक्षसैस्त्वतिपीडितः} %||5-48-15||

\twolineshloka
{दूतोऽहमिति विज्ञेयो राघवस्यामितौजसः}
{श्रूयतां चापि वचनं मम पथ्यमिदं प्रभो} %||5-49-16||


{॥इत्यार्षे श्रीमद्रामायणे वाल्मीकीये आदिकाव्ये सुन्दरकाण्डे अष्टचत्वारिंशः सर्गः॥}

\sect{एकोनपञ्चाशत्तमः सर्गः}

\twolineshloka
{तं समीक्ष्य महासत्त्वं सत्त्ववान्हरिसत्तमः}
{वाक्यमर्थवदव्यग्रस्तमुवाच दशाननम्} %||5-49-1||

\twolineshloka
{अहं सुग्रीवसंदेशादिह प्राप्तस्तवालयम्}
{राक्षसेन्द्र हरीशस्त्वां भ्राता कुशलमब्रवीत्} %||5-49-2||

\twolineshloka
{भ्रातुः शृणु समादेशं सुग्रीवस्य महात्मनः}
{धर्मार्थोपहितं वाक्यमिह चामुत्र च क्षमम्} %||5-49-3||

\twolineshloka
{राजा दशरथो नाम रथकुञ्जरवाजिमान्}
{पितेव बन्धुर्लोकस्य सुरेश्वरसमद्युतिः} %||5-49-4||

\twolineshloka
{ज्येष्ठस्तस्य महाबाहुः पुत्रः प्रियकरः प्रभुः}
{पितुर्निदेशान्निष्क्रान्तः प्रविष्टो दण्डकावनम्} %||5-49-5||

\twolineshloka
{लक्ष्मणेन सह भ्रात्रा सीतया चापि भार्यया}
{रामो नाम महातेजा धर्म्यं पन्थानमाश्रितः} %||5-49-6||

\twolineshloka
{तस्य भार्या वने नष्टा सीता पतिमनुव्रता}
{वैदेहस्य सुता राज्ञो जनकस्य महात्मनः} %||5-49-7||

\twolineshloka
{स मार्गमाणस्तां देवीं राजपुत्रः सहानुजः}
{ऋश्यमूकमनुप्राप्तः सुग्रीवेण च संगतः} %||5-49-8||

\twolineshloka
{तस्य तेन प्रतिज्ञातं सीतायाः परिमार्गणम्}
{सुग्रीवस्यापि रामेण हरिराज्यं निवेदितम्} %||5-49-9||

\twolineshloka
{ततस्तेन मृधे हत्वा राजपुत्रेण वालिनम्}
{सुग्रीवः स्थापितो राज्ये हर्यृक्षाणां गणेश्वरः} %||5-49-10||

\twolineshloka
{स सीतामार्गणे व्यग्रः सुग्रीवः सत्यसंगरः}
{हरीन्संप्रेषयामास दिशः सर्वा हरीश्वरः} %||5-49-11||

\twolineshloka
{तां हरीणां सहस्राणि शतानि नियुतानि च}
{दिक्षु सर्वासु मार्गन्ते अधश्चोपरि चाम्बरे} %||5-49-12||

\twolineshloka
{वैनतेय समाः केचित्केचित्तत्रानिलोपमाः}
{असंगगतयः शीघ्रा हरिवीरा महाबलाः} %||5-49-13||

\threelineshloka
{अहं तु हनुमान्नाम मारुतस्यौरसः सुतः}
{सीतायास्तु कृते तूर्णं शतयोजनमायतम्}
{समुद्रं लङ्घयित्वैव तां दिदृक्षुरिहागतः} %||5-49-14||

\twolineshloka
{तद्भवान्दृष्टधर्मार्थस्तपः कृतपरिग्रहः}
{परदारान्महाप्राज्ञ नोपरोद्धुं त्वमर्हसि} %||5-49-15||

\twolineshloka
{न हि धर्मविरुद्धेषु बह्वपायेषु कर्मसु}
{मूलघातिषु सज्जन्ते बुद्धिमन्तो भवद्विधाः} %||5-49-16||

\twolineshloka
{कश्च लक्ष्मणमुक्तानां रामकोपानुवर्तिनाम्}
{शराणामग्रतः स्थातुं शक्तो देवासुरेष्वपि} %||5-49-17||

\twolineshloka
{न चापि त्रिषु लोकेषु राजन्विद्येत कश्चन}
{राघवस्य व्यलीकं यः कृत्वा सुखमवाप्नुयात्} %||5-49-18||

\twolineshloka
{तत्त्रिकालहितं वाक्यं धर्म्यमर्थानुबन्धि च}
{मन्यस्व नरदेवाय जानकी प्रतिदीयताम्} %||5-49-19||

\twolineshloka
{दृष्टा हीयं मया देवी लब्धं यदिह दुर्लभम्}
{उत्तरं कर्म यच्छेषं निमित्तं तत्र राघवः} %||5-49-20||

\twolineshloka
{लक्षितेयं मया सीता तथा शोकपरायणा}
{गृह्य यां नाभिजानासि पञ्चास्यामिव पन्नगीम्} %||5-49-21||

\twolineshloka
{नेयं जरयितुं शक्या सासुरैरमरैरपि}
{विषसंसृष्टमत्यर्थं भुक्तमन्नमिवौजसा} %||5-49-22||

\twolineshloka
{तपःसंतापलब्धस्ते योऽयं धर्मपरिग्रहः}
{न स नाशयितुं न्याय्य आत्मप्राणपरिग्रहः} %||5-49-23||

\twolineshloka
{अवध्यतां तपोभिर्यां भवान्समनुपश्यति}
{आत्मनः सासुरैर्देवैर्हेतुस्तत्राप्ययं महान्} %||5-49-24||

\twolineshloka
{सुग्रीवो न हि देवोऽयं नासुरो न च मानुषः}
{न राक्षसो न गन्धर्वो न यक्षो न च पन्नगः} %||5-49-25||

\twolineshloka
{मानुषो राघवो राजन्सुग्रीवश्च हरीश्वरः}
{तस्मात्प्राणपरित्राणं कथं राजन्करिष्यसि} %||5-49-26||

\twolineshloka
{न तु धर्मोपसंहारमधर्मफलसंहितम्}
{तदेव फलमन्वेति धर्मश्चाधर्मनाशनः} %||5-49-27||

\twolineshloka
{प्राप्तं धर्मफलं तावद्भवता नात्र संशयः}
{फलमस्याप्यधर्मस्य क्षिप्रमेव प्रपत्स्यसे} %||5-49-28||

\twolineshloka
{जनस्थानवधं बुद्ध्वा बुद्ध्वा वालिवधं तथा}
{रामसुग्रीवसख्यं च बुध्यस्व हितमात्मनः} %||5-49-29||

\twolineshloka
{कामं खल्वहमप्येकः सवाजिरथकुञ्जराम्}
{लङ्कां नाशयितुं शक्तस्तस्यैष तु विनिश्चयः} %||5-49-30||

\twolineshloka
{रामेण हि प्रतिज्ञातं हर्यृक्षगणसंनिधौ}
{उत्सादनममित्राणां सीता यैस्तु प्रधर्षिता} %||5-49-31||

\twolineshloka
{अपकुर्वन्हि रामस्य साक्षादपि पुरंदरः}
{न सुखं प्राप्नुयादन्यः किं पुनस्त्वद्विधो जनः} %||5-49-32||

\twolineshloka
{यां सीतेत्यभिजानासि येयं तिष्ठति ते वशे}
{कालरात्रीति तां विद्धि सर्वलङ्काविनाशिनीम्} %||5-49-33||

\twolineshloka
{तदलं कालपाशेन सीता विग्रहरूपिणा}
{स्वयं स्कन्धावसक्तेन क्षममात्मनि चिन्त्यताम्} %||5-49-34||

\twolineshloka
{सीतायास्तेजसा दग्धां रामकोपप्रपीडिताम्}
{दह्यमनामिमां पश्य पुरीं साट्टप्रतोलिकाम्} %||5-49-35||

\fourlineindentedshloka
{स सौष्ठवोपेतमदीनवादिनः}
{ कपेर्निशम्याप्रतिमोऽप्रियं वचः}
{दशाननः कोपविवृत्तलोचनः}
{ समादिशत्तस्य वधं महाकपेः} %||5-50-36||


{॥इत्यार्षे श्रीमद्रामायणे वाल्मीकीये आदिकाव्ये सुन्दरकाण्डे एकोनपञ्चाशत्तमः सर्गः॥}

\sect{पञ्चाशत्तमः सर्गः}

\twolineshloka
{तस्य तद्वचनं श्रुत्वा वानरस्य महात्मनः}
{आज्ञापयद्वधं तस्य रावणः क्रोधमूर्छितः} %||5-50-1||

\twolineshloka
{वधे तस्य समाज्ञप्ते रावणेन दुरात्मना}
{निवेदितवतो दौत्यं नानुमेने विभीषणः} %||5-50-2||

\twolineshloka
{तं रक्षोऽधिपतिं क्रुद्धं तच्च कार्यमुपस्थितम्}
{विदित्वा चिन्तयामास कार्यं कार्यविधौ स्थितः} %||5-50-3||

\twolineshloka
{निश्चितार्थस्ततः साम्नापूज्य शत्रुजिदग्रजम्}
{उवाच हितमत्यर्थं वाक्यं वाक्यविशारदः} %||5-50-4||

\twolineshloka
{राजन्धर्मविरुद्धं च लोकवृत्तेश्च गर्हितम्}
{तव चासदृशं वीर कपेरस्य प्रमापणम्} %||5-50-5||

\fourlineindentedshloka
{असंशयं शत्रुरयं प्रवृद्धः}
{ कृतं ह्यनेनाप्रियमप्रमेयम्}
{न दूतवध्यां प्रवदन्ति सन्तो}
{ दूतस्य दृष्टा बहवो हि दण्डाः} %||5-50-6||

\fourlineindentedshloka
{वैरूप्यामङ्गेषु कशाभिघातो}
{ मौण्ड्यं तथा लक्ष्मणसंनिपातः}
{एतान्हि दूते प्रवदन्ति दण्डा}
{न्वधस्तु दूतस्य न नः श्रुतोऽपि} %||5-50-7||

\fourlineindentedshloka
{कथं च धर्मार्थविनीतबुद्धिः}
{ परावरप्रत्ययनिश्चितार्थः}
{भवद्विधः कोपवशे हि तिष्ठे}
{त्कोपं नियच्छन्ति हि सत्त्ववन्तः} %||5-50-8||

\fourlineindentedshloka
{न धर्मवादे न च लोकवृत्ते}
{ न शास्त्रबुद्धिग्रहणेषु वापि}
{विद्येत कश्चित्तव वीरतुल्य}
{स्त्वं ह्युत्तमः सर्वसुरासुराणाम्} %||5-50-9||

\twolineshloka
{न चाप्यस्य कपेर्घाते कंचित्पश्याम्यहं गुणम्}
{तेष्वयं पात्यतां दण्डो यैरयं प्रेषितः कपिः} %||5-50-10||

\twolineshloka
{साधुर्वा यदि वासाधुर्परैरेष समर्पितः}
{ब्रुवन्परार्थं परवान्न दूतो वधमर्हति} %||5-50-11||

\twolineshloka
{अपि चास्मिन्हते राजन्नान्यं पश्यामि खेचरम्}
{इह यः पुनरागच्छेत्परं पारं महोदधिः} %||5-50-12||

\twolineshloka
{तस्मान्नास्य वधे यत्नः कार्यः परपुरंजय}
{भवान्सेन्द्रेषु देवेषु यत्नमास्थातुमर्हति} %||5-50-13||

\fourlineindentedshloka
{अस्मिन्विनष्टे न हि दूतमन्यं}
{ पश्यामि यस्तौ नरराजपुत्रौ}
{युद्धाय युद्धप्रियदुर्विनीता}
{वुद्योजयेद्दीर्घपथावरुद्धौ} %||5-50-14||

\fourlineindentedshloka
{पराक्रमोत्साहमनस्विनां च}
{ सुरासुराणामपि दुर्जयेन}
{त्वया मनोनन्दन नैरृतानां}
{ युद्धायतिर्नाशयितुं न युक्ता} %||5-50-15||

\fourlineindentedshloka
{हिताश्च शूराश्च समाहिताश्च}
{ कुलेषु जाताश्च महागुणेषु}
{मनस्विनः शस्त्रभृतां वरिष्ठाः}
{ कोट्यग्रशस्ते सुभृताश्च योधाः} %||5-50-16||

\fourlineindentedshloka
{तदेकदेशेन बलस्य ताव}
{त्केचित्तवादेशकृतोऽपयान्तु}
{तौ राजपुत्रौ विनिगृह्य मूढौ}
{ परेषु ते भावयितुं प्रभावम्} %||5-51-17||


{॥इत्यार्षे श्रीमद्रामायणे वाल्मीकीये आदिकाव्ये सुन्दरकाण्डे पञ्चाशत्तमः सर्गः॥}

\sect{एकपञ्चाशत्तमः सर्गः}

\twolineshloka
{तस्य तद्वचनं श्रुत्वा दशग्रीवो महाबलः}
{देशकालहितं वाक्यं भ्रातुरुत्तममब्रवीत्} %||5-51-1||

\twolineshloka
{सम्यगुक्तं हि भवता दूतवध्या विगर्हिता}
{अवश्यं तु वधादन्यः क्रियतामस्य निग्रहः} %||5-51-2||

\twolineshloka
{कपीनां किल लाङ्गूलमिष्टं भवति भूषणम्}
{तदस्य दीप्यतां शीघ्रं तेन दग्धेन गच्छतु} %||5-51-3||

\twolineshloka
{ततः पश्यन्त्विमं दीनमङ्गवैरूप्यकर्शितम्}
{समित्रा ज्ञातयः सर्वे बान्धवाः ससुहृज्जनाः} %||5-51-4||

\twolineshloka
{आज्ञापयद्राक्षसेन्द्रः पुरं सर्वं सचत्वरम्}
{लाङ्गूलेन प्रदीप्तेन रक्षोभिः परिणीयताम्} %||5-51-5||

\twolineshloka
{तस्य तद्वचनं श्रुत्वा राक्षसाः कोपकर्कशाः}
{वेष्टन्ते तस्य लाङ्गूलं जीर्णैः कार्पासिकैः पटैः} %||5-51-6||

\twolineshloka
{संवेष्ट्यमाने लाङ्गूले व्यवर्धत महाकपिः}
{शुष्कमिन्धनमासाद्य वनेष्विव हुताशनः} %||5-51-7||

{तैलेन परिषिच्याथ तेऽग्निं तत्रावपातयन्॥\devanumber{8}॥} %||5-51-8|| (Check)

\addtocounter{shlokacount}{1}

\twolineshloka
{लाङ्गूलेन प्रदीप्तेन राक्षसांस्तानपातयत्}
{रोषामर्षपरीतात्मा बालसूर्यसमाननः} %||5-51-9||

\twolineshloka
{स भूयः संगतैः क्रूरै राकसैर्हरिसत्तमः}
{निबद्धः कृतवान्वीरस्तत्कालसदृशीं मतिम्} %||5-51-10||

\twolineshloka
{कामं खलु न मे शक्ता निबधस्यापि राक्षसाः}
{छित्त्वा पाशान्समुत्पत्य हन्यामहमिमान्पुनः} %||5-51-11||

\twolineshloka
{सर्वेषामेव पर्याप्तो राक्षसानामहं युधि}
{किं तु रामस्य प्रीत्यर्थं विषहिष्येऽहमीदृशम्} %||5-51-12||

\threelineshloka
{लङ्का चरयितव्या मे पुनरेव भवेदिति}
{रात्रौ न हि सुदृष्टा मे दुर्गकर्मविधानतः}
{अवश्यमेव द्रष्टव्या मया लङ्का निशाक्षये} %||5-51-13||

\twolineshloka
{कामं बन्धैश्च मे भूयः पुच्छस्योद्दीपनेन च}
{पीडां कुर्वन्तु रक्षांसि न मेऽस्ति मनसः श्रमः} %||5-51-14||

\twolineshloka
{ततस्ते संवृताकारं सत्त्ववन्तं महाकपिम्}
{परिगृह्य ययुर्हृष्टा राक्षसाः कपिकुञ्जरम्} %||5-51-15||

\twolineshloka
{शङ्खभेरीनिनादैस्तैर्घोषयन्तः स्वकर्मभिः}
{राक्षसाः क्रूरकर्माणश्चारयन्ति स्म तां पुरीम्} %||5-51-16||

\twolineshloka
{हनुमांश्चारयामास राक्षसानां महापुरीम्}
{अथापश्यद्विमानानि विचित्राणि महाकपिः} %||5-51-17||

\twolineshloka
{संवृतान्भूमिभागांश्च सुविभक्तांश्च चत्वरान्}
{रथ्याश्च गृहसंबाधाः कपिः शृङ्गाटकानि च} %||5-51-18||

\twolineshloka
{चत्वरेषु चतुष्केषु राजमार्गे तथैव च}
{घोषयन्ति कपिं सर्वे चारीक इति राक्षसाः} %||5-51-19||

\twolineshloka
{दीप्यमाने ततस्तस्य लाङ्गूलाग्रे हनूमतः}
{राक्षस्यस्ता विरूपाक्ष्यः शंसुर्देव्यास्तदप्रियम्} %||5-51-20||

\twolineshloka
{यस्त्वया कृतसंवादः सीते ताम्रमुखः कपिः}
{लाङ्गूलेन प्रदीप्तेन स एष परिणीयते} %||5-51-21||

\twolineshloka
{श्रुत्वा तद्वचनं क्रूरमात्मापहरणोपमम्}
{वैदेही शोकसंतप्ता हुताशनमुपागमत्} %||5-51-22||

\twolineshloka
{मङ्गलाभिमुखी तस्य सा तदासीन्महाकपेः}
{उपतस्थे विशालाक्षी प्रयता हव्यवाहनम्} %||5-51-23||

\twolineshloka
{यद्यस्ति पतिशुश्रूषा यद्यस्ति चरितं तपः}
{यदि चास्त्येकपत्नीत्वं शीतो भव हनूमतः} %||5-51-24||

\twolineshloka
{यदि कश्चिदनुक्रोशस्तस्य मय्यस्ति धीमतः}
{यदि वा भाग्यशेषं मे शीतो भव हनूमतः} %||5-51-25||

\twolineshloka
{यदि मां वृत्तसंपन्नां तत्समागमलालसाम्}
{स विजानाति धर्मात्मा शीतो भव हनूमतः} %||5-51-26||

\twolineshloka
{यदि मां तारयत्यार्यः सुग्रीवः सत्यसंगरः}
{अस्माद्दुःखान्महाबाहुः शीतो भव हनूमतः} %||5-51-27||

\twolineshloka
{ततस्तीक्ष्णार्चिरव्यग्रः प्रदक्षिणशिखोऽनलः}
{जज्वाल मृगशावाक्ष्याः शंसन्निव शिवं कपेः} %||5-51-28||

\twolineshloka
{दह्यमाने च लाङ्गूले चिन्तयामास वानरः}
{प्रदीप्तोऽग्निरयं कस्मान्न मां दहति सर्वतः} %||5-51-29||

\twolineshloka
{दृश्यते च महाज्वालः करोति च न मे रुजम्}
{शिशिरस्येव संपातो लाङ्गूलाग्रे प्रतिष्ठितः} %||5-51-30||

\twolineshloka
{अथ वा तदिदं व्यक्तं यद्दृष्टं प्लवता मया}
{रामप्रभावादाश्चर्यं पर्वतः सरितां पतौ} %||5-51-31||

\twolineshloka
{यदि तावत्समुद्रस्य मैनाकस्य च धीमथ}
{रामार्थं संभ्रमस्तादृक्किमग्निर्न करिष्यति} %||5-51-32||

\twolineshloka
{सीतायाश्चानृशंस्येन तेजसा राघवस्य च}
{पितुश्च मम सख्येन न मां दहति पावकः} %||5-51-33||

\twolineshloka
{भूयः स चिन्तयामास मुहूर्तं कपिकुञ्जरः}
{उत्पपाताथ वेगेन ननाद च महाकपिः} %||5-51-34||

\twolineshloka
{पुरद्वारं ततः श्रीमाञ्शैलशृङ्गमिवोन्नतम्}
{विभक्तरक्षःसंबाधमाससादानिलात्मजः} %||5-51-35||

\twolineshloka
{स भूत्वा शैलसंकाशः क्षणेन पुनरात्मवान्}
{ह्रस्वतां परमां प्राप्तो बन्धनान्यवशातयत्} %||5-51-36||

\twolineshloka
{विमुक्तश्चाभवच्छ्रीमान्पुनः पर्वतसंनिभः}
{वीक्षमाणश्च ददृशे परिघं तोरणाश्रितम्} %||5-51-37||

\twolineshloka
{स तं गृह्य महाबाहुः कालायसपरिष्कृतम्}
{रक्षिणस्तान्पुनः सर्वान्सूदयामास मारुतिः} %||5-51-38||

\fourlineindentedshloka
{स तान्निहत्वा रणचण्डविक्रमः}
{ समीक्षमाणः पुनरेव लङ्काम्}
{प्रदीप्तलाङ्गूलकृतार्चिमाली}
{ प्रकाशतादित्य इवांशुमाली} %||5-52-39||


{॥इत्यार्षे श्रीमद्रामायणे वाल्मीकीये आदिकाव्ये सुन्दरकाण्डे एकपञ्चाशत्तमः सर्गः॥}

\sect{द्विपञ्चाशत्तमः सर्गः}

\twolineshloka
{वीक्षमाणस्ततो लङ्कां कपिः कृतमनोरथः}
{वर्धमानसमुत्साहः कार्यशेषमचिन्तयत्} %||5-52-1||

\twolineshloka
{किं नु खल्वविशिष्टं मे कर्तव्यमिह साम्प्रतम्}
{यदेषां रक्षसां भूयः संतापजननं भवेत्} %||5-52-2||

\twolineshloka
{वनं तावत्प्रमथितं प्रकृष्टा राक्षसा हताः}
{बलैकदेशः क्षपितः शेषं दुर्गविनाशनम्} %||5-52-3||

\twolineshloka
{दुर्गे विनाशिते कर्म भवेत्सुखपरिश्रमम्}
{अल्पयत्नेन कार्येऽस्मिन्मम स्यात्सफलः श्रमः} %||5-52-4||

\twolineshloka
{यो ह्ययं मम लाङ्गूले दीप्यते हव्यवाहनः}
{अस्य संतर्पणं न्याय्यं कर्तुमेभिर्गृहोत्तमैः} %||5-52-5||

\twolineshloka
{ततः प्रदीप्तलाङ्गूलः सविद्युदिव तोयदः}
{भवनाग्रेषु लङ्काया विचचार महाकपिः} %||5-52-6||

{मुमोच हनुमानग्निं कालानलशिखोपमम्॥\devanumber{7}॥} %||5-52-7|| (Check)

\addtocounter{shlokacount}{1}

\twolineshloka
{श्वसनेन च संयोगादतिवेगो महाबलः}
{कालाग्निरिव जज्वाल प्रावर्धत हुताशनः} %||5-52-8||

{प्रदीप्तमग्निं पवनस्तेषु वेश्मसु चारयत्॥\devanumber{9}॥} %||5-52-9|| (Check)

\addtocounter{shlokacount}{1}

\twolineshloka
{तानि काञ्चनजालानि मुक्तामणिमयानि च}
{भवनान्यवशीर्यन्त रत्नवन्ति महान्ति च} %||5-52-10||

\twolineshloka
{तानि भग्नविमानानि निपेतुर्वसुधातले}
{भवनानीव सिद्धानामम्बरात्पुण्यसंक्षये} %||5-52-11||

\twolineshloka
{वज्रविद्रुमवैदूर्यमुक्तारजतसंहितान्}
{विचित्रान्भवनाद्धातून्स्यन्दमानान्ददर्श सः} %||5-52-12||

\twolineshloka
{नाग्निस्तृप्यति काष्ठानां तृणानां च यथा तथा}
{हनूमान्राक्षसेन्द्राणां वधे किंचिन्न तृप्यति} %||5-52-13||

\fourlineindentedshloka
{हुताशनज्वालसमावृता सा}
{ हतप्रवीरा परिवृत्तयोधा}
{हनूमातः क्रोधबलाभिभूता}
{ बभूव शापोपहतेव लङ्का} %||5-52-14||

\fourlineindentedshloka
{ससंभ्रमं त्रस्तविषण्णराक्षसां}
{ समुज्ज्वलज्ज्वालहुताशनाङ्किताम्}
{ददर्श लङ्कां हनुमान्महामनाः}
{ स्वयम्भुकोपोपहतामिवावनिम्} %||5-52-15||

\fourlineindentedshloka
{स राक्षसांस्तान्सुबहूंश्च हत्वा}
{ वनं च भङ्क्त्वा बहुपादपं तत्}
{विसृज्य रक्षो भवनेषु चाग्निं}
{ जगाम रामं मनसा महात्मा} %||5-52-16||

\twolineshloka
{लङ्कां समस्तां संदीप्य लाङ्गूलाग्निं महाकपिः}
{निर्वापयामास तदा समुद्रे हरिसत्तमः} %||5-53-17||


{॥इत्यार्षे श्रीमद्रामायणे वाल्मीकीये आदिकाव्ये सुन्दरकाण्डे द्विपञ्चाशत्तमः सर्गः॥}

\sect{त्रिपञ्चाशत्तमः सर्गः}

\twolineshloka
{संदीप्यमानां विध्वस्तां त्रस्तरक्षो गणां पुरीम्}
{अवेक्ष्य हानुमाँल्लङ्कां चिन्तयामास वानरः} %||5-53-1||

\twolineshloka
{तस्याभूत्सुमहांस्त्रासः कुत्सा चात्मन्यजायत}
{लङ्कां प्रदहता कर्म किंस्वित्कृतमिदं मया} %||5-53-2||

\twolineshloka
{धन्यास्ते पुरुषश्रेष्ठ ये बुद्ध्या कोपमुत्थितम्}
{निरुन्धन्ति महात्मानो दीप्तमग्निमिवाम्भसा} %||5-53-3||

\twolineshloka
{यदि दग्धा त्वियं लङ्का नूनमार्यापि जानकी}
{दग्धा तेन मया भर्तुर्हतं कार्यमजानता} %||5-53-4||

\twolineshloka
{यदर्थमयमारम्भस्तत्कार्यमवसादितम्}
{मया हि दहता लङ्कां न सीता परिरक्षिता} %||5-53-5||

\twolineshloka
{ईषत्कार्यमिदं कार्यं कृतमासीन्न संशयः}
{तस्य क्रोधाभिभूतेन मया मूलक्षयः कृतः} %||5-53-6||

\twolineshloka
{विनष्टा जानकी व्यक्तं न ह्यदग्धः प्रदृश्यते}
{लङ्कायाः कश्चिदुद्देशः सर्वा भस्मीकृता पुरी} %||5-53-7||

\twolineshloka
{यदि तद्विहतं कार्यं मया प्रज्ञाविपर्ययात्}
{इहैव प्राणसंन्यासो ममापि ह्यतिरोचते} %||5-53-8||

\twolineshloka
{किमग्नौ निपताम्यद्य आहोस्विद्वडवामुखे}
{शरीरमाहो सत्त्वानां दद्मि सागरवासिनाम्} %||5-53-9||

\twolineshloka
{कथं हि जीवता शक्यो मया द्रष्टुं हरीश्वरः}
{तौ वा पुरुषशार्दूलौ कार्यसर्वस्वघातिना} %||5-53-10||

\twolineshloka
{मया खलु तदेवेदं रोषदोषात्प्रदर्शितम्}
{प्रथितं त्रिषु लोकेषु कपितमनवस्थितम्} %||5-53-11||

\twolineshloka
{धिगस्तु राजसं भावमनीशमनवस्थितम्}
{ईश्वरेणापि यद्रागान्मया सीता न रक्षिता} %||5-53-12||

\twolineshloka
{विनष्टायां तु सीतायां तावुभौ विनशिष्यतः}
{तयोर्विनाशे सुग्रीवः सबन्धुर्विनशिष्यति} %||5-53-13||

\twolineshloka
{एतदेव वचः श्रुत्वा भरतो भ्रातृवत्सलः}
{धर्मात्मा सहशत्रुघ्नः कथं शक्ष्यति जीवितुम्} %||5-53-14||

\twolineshloka
{इक्ष्वाकुवंशे धर्मिष्ठे गते नाशमसंशयम्}
{भविष्यन्ति प्रजाः सर्वाः शोकसंतापपीडिताः} %||5-53-15||

\twolineshloka
{तदहं भाग्यरहितो लुप्तधर्मार्थसंग्रहः}
{रोषदोषपरीतात्मा व्यक्तं लोकविनाशनः} %||5-53-16||

\twolineshloka
{इति चिन्तयतस्तस्य निमित्तान्युपपेदिरे}
{पूरमप्युपलब्धानि साक्षात्पुनरचिन्तयत्} %||5-53-17||

\twolineshloka
{अथ वा चारुसर्वाङ्गी रक्षिता स्वेन तेजसा}
{न नशिष्यति कल्याणी नाग्निरग्नौ प्रवर्तते} %||5-53-18||

\twolineshloka
{न हि धर्मान्मनस्तस्य भार्याममिततेजसः}
{स्वचारित्राभिगुप्तां तां स्प्रष्टुमर्हति पावकः} %||5-53-19||

\twolineshloka
{नूनं रामप्रभावेन वैदेह्याः सुकृतेन च}
{यन्मां दहनकर्मायं नादहद्धव्यवाहनः} %||5-53-20||

\twolineshloka
{त्रयाणां भरतादीनां भ्रातॄणां देवता च या}
{रामस्य च मनःकान्ता सा कथं विनशिष्यति} %||5-53-21||

\twolineshloka
{यद्वा दहनकर्मायं सर्वत्र प्रभुरव्ययः}
{न मे दहति लाङ्गूलं कथमार्यां प्रधक्ष्यति} %||5-53-22||

\twolineshloka
{तपसा सत्यवाक्येन अनन्यत्वाच्च भर्तरि}
{अपि सा निर्दहेदग्निं न तामग्निः प्रधक्ष्यति} %||5-53-23||

\twolineshloka
{स तथा चिन्तयंस्तत्र देव्या धर्मपरिग्रहम्}
{शुश्राव हनुमान्वाक्यं चारणानां महात्मनाम्} %||5-53-24||

\twolineshloka
{अहो खलु कृतं कर्म दुर्विषह्यं हनूमता}
{अग्निं विसृजताभीक्ष्णं भीमं राक्षससद्मनि} %||5-53-25||

\twolineshloka
{दग्धेयं नगरी लङ्का साट्टप्राकारतोरणा}
{जानकी न च दग्धेति विस्मयोऽद्भुत एव नः} %||5-53-26||

\twolineshloka
{स निमित्तैश्च दृष्टार्थैः कारणैश्च महागुणैः}
{ऋषिवाक्यैश्च हनुमानभवत्प्रीतमानसः} %||5-53-27||

\fourlineindentedshloka
{ततः कपिः प्राप्तमनोरथार्थ}
{स्तामक्षतां राजसुतां विदित्वा}
{प्रत्यक्षतस्तां पुनरेव दृष्ट्वा}
{ प्रतिप्रयाणाय मतिं चकार} %||5-54-28||


{॥इत्यार्षे श्रीमद्रामायणे वाल्मीकीये आदिकाव्ये सुन्दरकाण्डे त्रिपञ्चाशत्तमः सर्गः॥}

\sect{चतुःपञ्चाशत्तमः सर्गः}

\twolineshloka
{ततस्तु शिंशपामूले जानकीं पर्यवस्थिताम्}
{अभिवाद्याब्रवीद्दिष्ट्या पश्यामि त्वामिहाक्षताम्} %||5-54-1||

\twolineshloka
{ततस्तं प्रस्थितं सीता वीक्षमाणा पुनः पुनः}
{भर्तृस्नेहान्वितं वाक्यं हनूमन्तमभाषत} %||5-54-2||

\twolineshloka
{काममस्य त्वमेवैकः कार्यस्य परिसाधने}
{पर्याप्तः परवीरघ्न यशस्यस्ते बलोदयः} %||5-54-3||

\twolineshloka
{बलैस्तु संकुलां कृत्वा लङ्कां परबलार्दनः}
{मां नयेद्यदि काकुत्स्थस्तस्य तत्सादृशं भवेत्} %||5-54-4||

\twolineshloka
{तद्यथा तस्य विक्रान्तमनुरूपं महात्मनः}
{भवत्याहवशूरस्य तत्त्वमेवोपपादय} %||5-54-5||

\twolineshloka
{तदर्थोपहितं वाक्यं प्रश्रितं हेतुसंहितम्}
{निशम्य हनुमांस्तस्या वाक्यमुत्तरमब्रवीत्} %||5-54-6||

\twolineshloka
{क्षिप्रमेष्यति काकुत्स्थो हर्यृक्षप्रवरैर्वृतः}
{यस्ते युधि विजित्यारीञ्शोकं व्यपनयिष्यति} %||5-54-7||

\twolineshloka
{एवमाश्वास्य वैदेहीं हनूमान्मारुतात्मजः}
{गमनाय मतिं कृत्वा वैदेहीमभ्यवादयत्} %||5-54-8||

\twolineshloka
{ततः स कपिशार्दूलः स्वामिसंदर्शनोत्सुकः}
{आरुरोह गिरिश्रेष्ठमरिष्टमरिमर्दनः} %||5-54-9||

\twolineshloka
{तुङ्गपद्मकजुष्टाभिर्नीलाभिर्वनराजिभिः}
{सालतालाश्वकर्णैश्च वंशैश्च बहुभिर्वृतम्} %||5-54-10||

\twolineshloka
{लतावितानैर्विततैः पुष्पवद्भिरलंकृतम्}
{नानामृगगणाकीर्णं धातुनिष्यन्दभूषितम्} %||5-54-11||

\twolineshloka
{बहुप्रस्रवणोपेतं शिलासंचयसंकटम्}
{महर्षियक्षगन्धर्वकिंनरोरगसेवितम्} %||5-54-12||

\twolineshloka
{लतापादपसंबाधं सिंहाकुलितकन्दरम्}
{व्याघ्रसंघसमाकीर्णं स्वादुमूलफलद्रुमम्} %||5-54-13||

\twolineshloka
{तमारुरोहातिबलः पर्वतं प्लवगोत्तमः}
{रामदर्शनशीघ्रेण प्रहर्षेणाभिचोदितः} %||5-54-14||

\twolineshloka
{तेन पादतलाक्रान्ता रम्येषु गिरिसानुषु}
{सघोषाः समशीर्यन्त शिलाश्चूर्णीकृतास्ततः} %||5-54-15||

\twolineshloka
{स तमारुह्य शैलेन्द्रं व्यवर्धत महाकपिः}
{दक्षिणादुत्तरं पारं प्रार्थयँल्लवणाम्भसः} %||5-54-16||

\twolineshloka
{अधिरुह्य ततो वीरः पर्वतं पवनात्मजः}
{ददर्श सागरं भीमं मीनोरगनिषेवितम्} %||5-54-17||

\twolineshloka
{स मारुत इवाकाशं मारुतस्यात्मसंभवः}
{प्रपेदे हरिशार्दूलो दक्षिणादुत्तरां दिशम्} %||5-54-18||

\threelineshloka
{स तदा पीडितस्तेन कपिना पर्वतोत्तमः}
{ररास सह तैर्भूतैः प्राविशद्वसुधातलम्}
{कम्पमानैश्च शिखरैः पतद्भिरपि च द्रुमैः} %||5-54-19||

\twolineshloka
{तस्योरुवेगान्मथिताः पादपाः पुष्पशालिनः}
{निपेतुर्भूतले रुग्णाः शक्रायुधहता इव} %||5-54-20||

\twolineshloka
{कन्दरोदरसंस्थानां पीडितानां महौजसाम्}
{सिंहानां निनदो भीमो नभो भिन्दन्स शुश्रुवे} %||5-54-21||

\twolineshloka
{स्रस्तव्याविद्धवसना व्याकुलीकृतभूषणा}
{विद्याधर्यः समुत्पेतुः सहसा धरणीधरात्} %||5-54-22||

\twolineshloka
{अतिप्रमाणा बलिनो दीप्तजिह्वा महाविषाः}
{निपीडितशिरोग्रीवा व्यवेष्टन्त महाहयः} %||5-54-23||

\twolineshloka
{किंनरोरगगन्धर्वयक्षविद्याधरास्तथा}
{पीडितं तं नगवरं त्यक्त्वा गगनमास्थिताः} %||5-54-24||

\twolineshloka
{स च भूमिधरः श्रीमान्बलिना तेन पीडितः}
{सवृक्षशिखरोदग्राः प्रविवेश रसातलम्} %||5-54-25||

\twolineshloka
{दशयोजनविस्तारस्त्रिंशद्योजनमुच्छ्रितः}
{धरण्यां समतां यातः स बभूव धराधरः} %||5-55-26||


{॥इत्यार्षे श्रीमद्रामायणे वाल्मीकीये आदिकाव्ये सुन्दरकाण्डे चतुःपञ्चाशत्तमः सर्गः॥}

\sect{पञ्चपञ्चाशत्तमः सर्गः}

\twolineshloka
{सचन्द्रकुमुदं रम्यं सार्ककारण्डवं शुभम्}
{तिष्यश्रवणकदम्बमभ्रशैवलशाद्वलम्} %||5-55-1||

\twolineshloka
{पुनर्वसु महामीनं लोहिताङ्गमहाग्रहम्}
{ऐरावतमहाद्वीपं स्वातीहंसविलोडितम्} %||5-55-2||

\twolineshloka
{वातसंघातजातोर्मिं चन्द्रांशुशिशिराम्बुमत्}
{भुजंगयक्षगन्धर्वप्रबुद्धकमलोत्पलम्} %||5-55-3||

\twolineshloka
{ग्रसमान इवाकाशं ताराधिपमिवालिखन्}
{हरन्निव सनक्षत्रं गगनं सार्कमण्डलम्} %||5-55-4||

\twolineshloka
{मारुतस्यालयं श्रीमान्कपिर्व्योमचरो महान्}
{हनूमान्मेघजालानि विकर्षन्निव गच्छति} %||5-55-5||

\twolineshloka
{पाण्डुरारुणवर्णानि नीलमाञ्जिष्ठकानि च}
{हरितारुणवर्णानि महाभ्राणि चकाशिरे} %||5-55-6||

\twolineshloka
{प्रविशन्नभ्रजालानि निष्क्रमंश्च पुनः पुनः}
{प्रच्छन्नश्च प्रकाशश्च चन्द्रमा इव लक्ष्यते} %||5-55-7||

\twolineshloka
{नदन्नादेन महता मेघस्वनमहास्वनः}
{आजगाम महातेजाः पुनर्मध्येन सागरम्} %||5-55-8||

\twolineshloka
{पर्वतेन्द्रं सुनाभं च समुपस्पृश्य वीर्यवान्}
{ज्यामुक्त इव नाराचो महावेगोऽभ्युपागतः} %||5-55-9||

\twolineshloka
{स किंचिदनुसंप्राप्तः समालोक्य महागिरिम्}
{महेन्द्रमेघसंकाशं ननाद हरिपुंगवः} %||5-55-10||

\twolineshloka
{निशम्य नदतो नादं वानरास्ते समन्ततः}
{बभूवुरुत्सुकाः सर्वे सुहृद्दर्शनकाङ्क्षिणः} %||5-55-11||

\twolineshloka
{जाम्बवान्स हरिश्रेष्ठः प्रीतिसंहृष्टमानसः}
{उपामन्त्र्य हरीन्सर्वानिदं वचनमब्रवीत्} %||5-55-12||

\twolineshloka
{सर्वथा कृतकार्योऽसौ हनूमान्नात्र संशयः}
{न ह्यस्याकृतकार्यस्य नाद एवंविधो भवेत्} %||5-55-13||

\twolineshloka
{तस्या बाहूरुवेगं च निनादं च महात्मनः}
{निशम्य हरयो हृष्टाः समुत्पेतुस्ततस्ततः} %||5-55-14||

\twolineshloka
{ते नगाग्रान्नगाग्राणि शिखराच्छिखराणि च}
{प्रहृष्टाः समपद्यन्त हनूमन्तं दिदृक्षवः} %||5-55-15||

\twolineshloka
{ते प्रीताः पादपाग्रेषु गृह्य शाखाः सुपुष्पिताः}
{वासांसीव प्रकाशानि समाविध्यन्त वानराः} %||5-55-16||

\twolineshloka
{तमभ्रघनसंकाशमापतन्तं महाकपिम्}
{दृष्ट्वा ते वानराः सर्वे तस्थुः प्राञ्जलयस्तदा} %||5-55-17||

\twolineshloka
{ततस्तु वेगवांस्तस्य गिरेर्गिरिनिभः कपिः}
{निपपात महेन्द्रस्य शिखरे पादपाकुले} %||5-55-18||

\twolineshloka
{ततस्ते प्रीतमनसः सर्वे वानरपुंगवाः}
{हनूमन्तं महात्मानं परिवार्योपतस्थिरे} %||5-55-19||

\twolineshloka
{परिवार्य च ते सर्वे परां प्रीतिमुपागताः}
{प्रहृष्टवदनाः सर्वे तमरोगमुपागतम्} %||5-55-20||

\twolineshloka
{उपायनानि चादाय मूलानि च फलानि च}
{प्रत्यर्चयन्हरिश्रेष्ठं हरयो मारुतात्मजम्} %||5-55-21||

\twolineshloka
{विनेदुर्मुदिताः केचिच्चक्रुः किल किलां तथा}
{हृष्टाः पादपशाखाश्च आनिन्युर्वानरर्षभाः} %||5-55-22||

\twolineshloka
{हनूमांस्तु गुरून्वृद्धाञ्जाम्बवत्प्रमुखांस्तदा}
{कुमारमङ्गदं चैव सोऽवन्दत महाकपिः} %||5-55-23||

\twolineshloka
{स ताभ्यां पूजितः पूज्यः कपिभिश्च प्रसादितः}
{दृष्टा देवीति विक्रान्तः संक्षेपेण न्यवेदयत्} %||5-55-24||

\twolineshloka
{निषसाद च हस्तेन गृहीत्वा वालिनः सुतम्}
{रमणीये वनोद्देशे महेन्द्रस्य गिरेस्तदा} %||5-55-25||

\twolineshloka
{हनूमानब्रवीद्धृष्टस्तदा तान्वानरर्षभान्}
{अशोकवनिकासंस्था दृष्टा सा जनकात्मजा} %||5-55-26||

\threelineshloka
{रक्ष्यमाणा सुघोराभी राक्षसीभिरनिन्दिता}
{एकवेणीधरा बाला रामदर्शनलालसा}
{उपवासपरिश्रान्ता मलिना जटिला कृशा} %||5-55-27||

\twolineshloka
{ततो दृष्टेति वचनं महार्थममृतोपमम्}
{निशम्य मारुतेः सर्वे मुदिता वानरा भवन्} %||5-55-28||

\twolineshloka
{क्ष्वेडन्त्यन्ये नदन्त्यन्ये गर्जन्त्यन्ये महाबलाः}
{चक्रुः किल किलामन्ये प्रतिगर्जन्ति चापरे} %||5-55-29||

\twolineshloka
{केचिदुच्छ्रितलाङ्गूलाः प्रहृष्टाः कपिकुञ्जराः}
{अञ्चितायतदीर्घाणि लाङ्गूलानि प्रविव्यधुः} %||5-55-30||

\twolineshloka
{अपरे तु हनूमन्तं वानरा वारणोपमम्}
{आप्लुत्य गिरिशृङ्गेभ्यः संस्पृशन्ति स्म हर्षिताः} %||5-55-31||

\twolineshloka
{उक्तवाक्यं हनूमन्तमङ्गदस्तु तदाब्रवीत्}
{सर्वेषां हरिवीराणां मध्ये वाचमनुत्तमाम्} %||5-55-32||

\twolineshloka
{सत्त्वे वीर्ये न ते कश्चित्समो वानरविद्यते}
{यदवप्लुत्य विस्तीर्णं सागरं पुनरागतः} %||5-55-33||

\twolineshloka
{दिष्ट्या दृष्टा त्वया देवी रामपत्नी यशस्विनी}
{दिष्ट्या त्यक्ष्यति काकुत्स्थः शोकं सीता वियोगजम्} %||5-55-34||

\twolineshloka
{ततोऽङ्गदं हनूमन्तं जाम्बवन्तं च वानराः}
{परिवार्य प्रमुदिता भेजिरे विपुलाः शिलाः} %||5-55-35||

\threelineshloka
{श्रोतुकामाः समुद्रस्य लङ्घनं वानरोत्तमाः}
{दर्शनं चापि लङ्कायाः सीताया रावणस्य च}
{तस्थुः प्राञ्जलयः सर्वे हनूमद्वदनोन्मुखाः} %||5-55-36||

\twolineshloka
{तस्थौ तत्राङ्गदः श्रीमान्वानरैर्बहुभिर्वृतः}
{उपास्यमानो विबुधैर्दिवि देवपतिर्यथा} %||5-55-37||

\fourlineindentedshloka
{हनूमता कीर्तिमता यशस्विना}
{ तथाङ्गदेनाङ्गदबद्धबाहुना}
{मुदा तदाध्यासितमुन्नतं मह}
{न्महीधराग्रं ज्वलितं श्रियाभवत्} %||5-56-38||


{॥इत्यार्षे श्रीमद्रामायणे वाल्मीकीये आदिकाव्ये सुन्दरकाण्डे पञ्चपञ्चाशत्तमः सर्गः॥}

\sect{षट्पञ्चाशत्तमः सर्गः}

\twolineshloka
{ततस्तस्य गिरेः शृङ्गे महेन्द्रस्य महाबलाः}
{हनुमत्प्रमुखाः प्रीतिं हरयो जग्मुरुत्तमाम्} %||5-56-1||

\twolineshloka
{तं ततः प्रतिसंहृष्टः प्रीतिमन्तं महाकपिम्}
{जाम्बवान्कार्यवृत्तान्तमपृच्छदनिलात्मजम्} %||5-56-2||

\twolineshloka
{कथं दृष्टा त्वया देवी कथं वा तत्र वर्तते}
{तस्यां वा स कथं वृत्तः क्रूरकर्मा दशाननः} %||5-56-3||

\twolineshloka
{तत्त्वतः सर्वमेतन्नः प्रब्रूहि त्वं महाकपे}
{श्रुतार्थाश्चिन्तयिष्यामो भूयः कार्यविनिश्चयम्} %||5-56-4||

\twolineshloka
{यश्चार्थस्तत्र वक्तव्यो गतैरस्माभिरात्मवान्}
{रक्षितव्यं च यत्तत्र तद्भवान्व्याकरोतु नः} %||5-56-5||

\twolineshloka
{स नियुक्तस्ततस्तेन संप्रहृष्टतनूरुहः}
{नमस्यञ्शिरसा देव्यै सीतायै प्रत्यभाषत} %||5-56-6||

\twolineshloka
{प्रत्यक्षमेव भवतां महेन्द्राग्रात्खमाप्लुतः}
{उदधेर्दक्षिणं पारं काङ्क्षमाणः समाहितः} %||5-56-7||

\twolineshloka
{गच्छतश्च हि मे घोरं विघ्नरूपमिवाभवत्}
{काञ्चनं शिखरं दिव्यं पश्यामि सुमनोहरम्} %||5-56-8||

{स्थितं पन्थानमावृत्य मेने विघ्नं च तं नगम्॥\devanumber{9}॥} %||5-56-9|| (Check)

\addtocounter{shlokacount}{1}

\twolineshloka
{उपसंगम्य तं दिव्यं काञ्चनं नगसत्तमम्}
{कृता मे मनसा बुद्धिर्भेत्तव्योऽयं मयेति च} %||5-56-10||

\twolineshloka
{प्रहतं च मया तस्य लाङ्गूलेन महागिरेः}
{शिखरं सूर्यसंकाशं व्यशीर्यत सहस्रधा} %||5-56-11||

\twolineshloka
{व्यवसायं च मे बुद्ध्वा स होवाच महागिरिः}
{पुत्रेति मधुरां बाणीं मनःप्रह्लादयन्निव} %||5-56-12||

\twolineshloka
{पितृव्यं चापि मां विद्धि सखायं मातरिश्वनः}
{मैनाकमिति विख्यातं निवसन्तं महोदधौ} %||5-56-13||

\twolineshloka
{पक्ष्ववन्तः पुरा पुत्र बभूवुः पर्वतोत्तमाः}
{छन्दतः पृथिवीं चेरुर्बाधमानाः समन्ततः} %||5-56-14||

\twolineshloka
{श्रुत्वा नगानां चरितं महेन्द्रः पाकशासनः}
{चिच्छेद भगवान्पक्षान्वज्रेणैषां सहस्रशः} %||5-56-15||

\twolineshloka
{अहं तु मोक्षितस्तस्मात्तव पित्रा महात्मना}
{मारुतेन तदा वत्स प्रक्षिप्तोऽस्मि महार्णवे} %||5-56-16||

\twolineshloka
{रामस्य च मया साह्ये वर्तितव्यमरिंदम}
{रामो धर्मभृतां श्रेष्ठो महेन्द्रसमविक्रमः} %||5-56-17||

\twolineshloka
{एतच्छ्रुत्वा मया तस्य मैनाकस्य महात्मनः}
{कार्यमावेद्य तु गिरेरुद्धतं च मनो मम} %||5-56-18||

\twolineshloka
{तेन चाहमनुज्ञातो मैनाकेन महात्मना}
{उत्तमं जवमास्थाय शेषमध्वानमास्थितः} %||5-56-19||

\twolineshloka
{ततोऽहं सुचिरं कालं वेगेनाभ्यगमं पथि}
{ततः पश्याम्यहं देवीं सुरसां नागमातरम्} %||5-56-20||

\threelineshloka
{समुद्रमध्ये सा देवी वचनं मामभाषत}
{मम भक्ष्यः प्रदिष्टस्त्वममारैर्हरिसत्तमम्}
{ततस्त्वां भक्षयिष्यामि विहितस्त्वं चिरस्य मे} %||5-56-21||

\twolineshloka
{एवमुक्तः सुरसया प्राञ्जलिः प्रणतः स्थितः}
{विवर्णवदनो भूत्वा वाक्यं चेदमुदीरयम्} %||5-56-22||

\twolineshloka
{रामो दाशरथिः श्रीमान्प्रविष्टो दण्डकावनम्}
{लक्ष्मणेन सह भ्रात्रा सीतया च परंतपः} %||5-56-23||

\twolineshloka
{तस्य सीता हृता भार्या रावणेन दुरात्मना}
{तस्याः सकाशं दूतोऽहं गमिष्ये रामशासनात्} %||5-56-24||

{कर्तुमर्हसि रामस्य साह्यं विषयवासिनि॥\devanumber{25}॥} %||5-56-25|| (Check)

\addtocounter{shlokacount}{1}

\twolineshloka
{अथ वा मैथिलीं दृष्ट्वा रामं चाक्लिष्टकारिणम्}
{आगमिष्यामि ते वक्त्रं सत्यं प्रतिशृणोति मे} %||5-56-26||

\twolineshloka
{एवमुक्ता मया सा तु सुरसा कामरूपिणी}
{अब्रवीन्नातिवर्तेत कश्चिदेष वरो मम} %||5-56-27||

\twolineshloka
{एवमुक्तः सुरसया दशयोजनमायतः}
{ततोऽर्धगुणविस्तारो बभूवाहं क्षणेन तु} %||5-56-28||

\twolineshloka
{मत्प्रमाणानुरूपं च व्यादितं तन्मुखं तया}
{तद्दृष्ट्वा व्यादितं त्वास्यं ह्रस्वं ह्यकरवं वपुः} %||5-56-29||

\twolineshloka
{तस्मिन्मुहूर्ते च पुनर्बभूवाङ्गुष्ठसंमितः}
{अभिपत्याशु तद्वक्त्रं निर्गतोऽहं ततः क्षणात्} %||5-56-30||

\twolineshloka
{अब्रवीत्सुरसा देवी स्वेन रूपेण मां पुनः}
{अर्थसिद्ध्यै हरिश्रेष्ठ गच्छ सौम्य यथासुखम्} %||5-56-31||

\twolineshloka
{समानय च वैदेहीं राघवेण महात्मना}
{सुखी भव महाबाहो प्रीतास्मि तव वानर} %||5-56-32||

\twolineshloka
{ततोऽहं साधु साध्वीति सर्वभूतैः प्रशंसितः}
{ततोऽन्तरिक्षं विपुलं प्लुतोऽहं गरुडो यथा} %||5-56-33||

\threelineshloka
{छाया मे निगृहीता च न च पश्यामि किंचन}
{सोऽहं विगतवेगस्तु दिशो दश विलोकयन्}
{न किंचित्तत्र पश्यामि येन मेऽपहृता गतिः} %||5-56-34||

\twolineshloka
{ततो मे बुद्धिरुत्पन्ना किं नाम गमने मम}
{ईदृशो विघ्न उत्पन्नो रूपं यत्र न दृश्यते} %||5-56-35||

\twolineshloka
{अधो भागेन मे दृष्टिः शोचता पातिता मया}
{ततोऽद्राक्षमहं भीमां राक्षसीं सलिले शयाम्} %||5-56-36||

\twolineshloka
{प्रहस्य च महानादमुक्तोऽहं भीमया तया}
{अवस्थितमसंभ्रान्तमिदं वाक्यमशोभनम्} %||5-56-37||

\twolineshloka
{क्वासि गन्ता महाकाय क्षुधिताया ममेप्सितः}
{भक्षः प्रीणय मे देहं चिरमाहारवर्जितम्} %||5-56-38||

\twolineshloka
{बाढमित्येव तां वाणीं प्रत्यगृह्णामहं ततः}
{आस्य प्रमाणादधिकं तस्याः कायमपूरयम्} %||5-56-39||

\twolineshloka
{तस्याश्चास्यं महद्भीमं वर्धते मम भक्षणे}
{न च मां सा तु बुबुधे मम वा विकृतं कृतम्} %||5-56-40||

\twolineshloka
{ततोऽहं विपुलं रूपं संक्षिप्य निमिषान्तरात्}
{तस्या हृदयमादाय प्रपतामि नभस्तलम्} %||5-56-41||

\twolineshloka
{सा विसृष्टभुजा भीमा पपात लवणाम्भसि}
{मया पर्वतसंकाशा निकृत्तहृदया सती} %||5-56-42||

\twolineshloka
{शृणोमि खगतानां च सिद्धानां चारणैः सह}
{राक्षसी सिंहिका भीमा क्षिप्रं हनुमता हृता} %||5-56-43||

\threelineshloka
{तां हत्वा पुनरेवाहं कृत्यमात्ययिकं स्मरन्}
{गत्वा च महदध्वानं पश्यामि नगमण्डितम्}
{दक्षिणं तीरमुदधेर्लङ्का यत्र च सा पुरी} %||5-56-44||

\twolineshloka
{अस्तं दिनकरे याते रक्षसां निलयं पुरीम्}
{प्रविष्टोऽहमविज्ञातो रक्षोभिर्भीमविक्रमैः} %||5-56-45||

\twolineshloka
{तत्राहं सर्वरात्रं तु विचिन्वञ्जनकात्मजाम्}
{रावणान्तःपुरगतो न चापश्यं सुमध्यमाम्} %||5-56-46||

\twolineshloka
{ततः सीतामपश्यंस्तु रावणस्य निवेशने}
{शोकसागरमासाद्य न पारमुपलक्षये} %||5-56-47||

\twolineshloka
{शोचता च मया दृष्टं प्राकारेण समावृतम्}
{काञ्चनेन विकृष्टेन गृहोपवनमुत्तमम्} %||5-56-48||

{स प्राकारमवप्लुत्य पश्यामि बहुपादपम्॥\devanumber{49}॥} %||5-56-49|| (Check)

\addtocounter{shlokacount}{1}

\twolineshloka
{अशोकवनिकामध्ये शिंशपापादपो महान्}
{तमारुह्य च पश्यामि काञ्चनं कदली वनम्} %||5-56-50||

\twolineshloka
{अदूराच्छिंशपावृक्षात्पश्यामि वनवर्णिनीम्}
{श्यामां कमलपत्राक्षीमुपवासकृशाननाम्} %||5-56-51||

\twolineshloka
{राक्षसीभिर्विरूपाभिः क्रूराभिरभिसंवृताम्}
{मांसशोणितभक्ष्याभिर्व्याघ्रीभिर्हरिणीं यथा} %||5-56-52||

\twolineshloka
{तां दृष्ट्वा तादृशीं नारीं रामपत्नीमनिन्दिताम्}
{तत्रैव शिंशपावृक्षे पश्यन्नहमवस्थितः} %||5-56-53||

\twolineshloka
{ततो हलहलाशब्दं काञ्चीनूपुरमिश्रितम्}
{शृणोम्यधिकगम्भीरं रावणस्य निवेशने} %||5-56-54||

\twolineshloka
{ततोऽहं परमोद्विग्नः स्वरूपं प्रत्यसंहरम्}
{अहं च शिंशपावृक्षे पक्षीव गहने स्थितः} %||5-56-55||

\twolineshloka
{ततो रावणदाराश्च रावणश्च महाबलः}
{तं देशं समनुप्राप्ता यत्र सीताभवत्स्थिता} %||5-56-56||

\twolineshloka
{तं दृष्ट्वाथ वरारोहा सीता रक्षोगणेश्वरम्}
{संकुच्योरू स्तनौ पीनौ बाहुभ्यां परिरभ्य च} %||5-56-57||

\twolineshloka
{तामुवाच दशग्रीवः सीतां परमदुःखिताम्}
{अवाक्शिराः प्रपतितो बहु मन्यस्व मामिति} %||5-56-58||

\twolineshloka
{यदि चेत्त्वं तु मां दर्पान्नाभिनन्दसि गर्विते}
{द्विमासानन्तरं सीते पास्यामि रुधिरं तव} %||5-56-59||

\twolineshloka
{एतच्छ्रुत्वा वचस्तस्य रावणस्य दुरात्मनः}
{उवाच परमक्रुद्धा सीता वचनमुत्तमम्} %||5-56-60||

\threelineshloka
{राक्षसाधम रामस्य भार्याममिततेजसः}
{इक्ष्वाकुकुलनाथस्य स्नुषां दशरथस्य च}
{अवाच्यं वदतो जिह्वा कथं न पतिता तव} %||5-56-61||

\twolineshloka
{किंस्विद्वीर्यं तवानार्य यो मां भर्तुरसंनिधौ}
{अपहृत्यागतः पाप तेनादृष्टो महात्मना} %||5-56-62||

\twolineshloka
{न त्वं रामस्य सदृशो दास्येऽप्यस्या न युज्यसे}
{यज्ञीयः सत्यवाक्चैव रणश्लाघी च राघवः} %||5-56-63||

\twolineshloka
{जानक्या परुषं वाक्यमेवमुक्तो दशाननः}
{जज्वाल सहसा कोपाच्चितास्थ इव पावकः} %||5-56-64||

\twolineshloka
{विवृत्य नयने क्रूरे मुष्टिमुद्यम्य दक्षिणम्}
{मैथिलीं हन्तुमारब्धः स्त्रीभिर्हाहाकृतं तदा} %||5-56-65||

\twolineshloka
{स्त्रीणां मध्यात्समुत्पत्य तस्य भार्या दुरात्मनः}
{वरा मन्दोदरी नाम तया स प्रतिषेधितः} %||5-56-66||

\threelineshloka
{उक्तश्च मधुरां वाणीं तया स मदनार्दितः}
{सीतया तव किं कार्यं महेन्द्रसमविक्रम}
{मया सह रमस्वाद्य मद्विशिष्टा न जानकी} %||5-56-67||

\twolineshloka
{देवगन्धर्वकन्याभिर्यक्षकन्याभिरेव च}
{सार्धं प्रभो रमस्वेह सीतया किं करिष्यसि} %||5-56-68||

\twolineshloka
{ततस्ताभिः समेताभिर्नारीभिः स महाबलः}
{उत्थाप्य सहसा नीतो भवनं स्वं निशाचरः} %||5-56-69||

\twolineshloka
{याते तस्मिन्दशग्रीवे राक्षस्यो विकृताननाः}
{सीतां निर्भर्त्सयामासुर्वाक्यैः क्रूरैः सुदारुणैः} %||5-56-70||

\twolineshloka
{तृणवद्भाषितं तासां गणयामास जानकी}
{तर्जितं च तदा तासां सीतां प्राप्य निरर्थकम्} %||5-56-71||

\twolineshloka
{वृथागर्जितनिश्चेष्टा राक्षस्यः पिशिताशनाः}
{रावणाय शशंसुस्ताः सीताव्यवसितं महत्} %||5-56-72||

\twolineshloka
{ततस्ताः सहिताः सर्वा विहताशा निरुद्यमाः}
{परिक्षिप्य समन्तात्तां निद्रावशमुपागताः} %||5-56-73||

\twolineshloka
{तासु चैव प्रसुप्तासु सीता भर्तृहिते रता}
{विलप्य करुणं दीना प्रशुशोच सुदुःखिता} %||5-56-74||

\twolineshloka
{तां चाहं तादृशीं दृष्ट्वा सीताया दारुणां दशाम्}
{चिन्तयामास विश्रान्तो न च मे निर्वृतं मनः} %||5-56-75||

\twolineshloka
{संभाषणार्थे च मया जानक्याश्चिन्तितो विधिः}
{इक्ष्वाकुकुलवंशस्तु ततो मम पुरस्कृतः} %||5-56-76||

\twolineshloka
{श्रुत्वा तु गदितां वाचं राजर्षिगणपूजिताम्}
{प्रत्यभाषत मां देवी बाष्पैः पिहितलोचना} %||5-56-77||

\twolineshloka
{कस्त्वं केन कथं चेह प्राप्तो वानरपुंगव}
{का च रामेण ते प्रीतिस्तन्मे शंसितुमर्हसि} %||5-56-78||

\threelineshloka
{तस्यास्तद्वचनं श्रुत्वा अहमप्यब्रुवं वचः}
{देवि रामस्य भर्तुस्ते सहायो भीमविक्रमः}
{सुग्रीवो नाम विक्रान्तो वानरेन्दो महाबलः} %||5-56-79||

\twolineshloka
{तस्य मां विद्धि भृत्यं त्वं हनूमन्तमिहागतम्}
{भर्त्राहं प्रहितस्तुभ्यं रामेणाक्लिष्टकर्मणा} %||5-56-80||

\twolineshloka
{इदं च पुरुषव्याघ्रः श्रीमान्दाशरथिः स्वयम्}
{अङ्गुलीयमभिज्ञानमदात्तुभ्यं यशस्विनि} %||5-56-81||

\twolineshloka
{तदिच्छामि त्वयाज्ञप्तं देवि किं करवाण्यहम्}
{रामलक्ष्मणयोः पार्श्वं नयामि त्वां किमुत्तरम्} %||5-56-82||

\twolineshloka
{एतच्छ्रुत्वा विदित्वा च सीता जनकनन्दिनी}
{आह रावणमुत्साद्य राघवो मां नयत्विति} %||5-56-83||

\twolineshloka
{प्रणम्य शिरसा देवीमहमार्यामनिन्दिताम्}
{राघवस्य मनोह्लादमभिज्ञानमयाचिषम्} %||5-56-84||

\twolineshloka
{एवमुक्ता वरारोहा मणिप्रवरमुत्तमम्}
{प्रायच्छत्परमोद्विग्ना वाचा मां संदिदेश ह} %||5-56-85||

\twolineshloka
{ततस्तस्यै प्रणम्याहं राजपुत्र्यै समाहितः}
{प्रदक्षिणं परिक्राममिहाभ्युद्गतमानसः} %||5-56-86||

\twolineshloka
{उत्तरं पुनरेवाह निश्चित्य मनसा तदा}
{हनूमन्मम वृत्तान्तं वक्तुमर्हसि राघवे} %||5-56-87||

\twolineshloka
{यथा श्रुत्वैव नचिरात्तावुभौ रामलक्ष्मणौ}
{सुग्रीवसहितौ वीरावुपेयातां तथा कुरु} %||5-56-88||

\twolineshloka
{यद्यन्यथा भवेदेतद्द्वौ मासौ जीवितं मम}
{न मां द्रक्ष्यति काकुत्स्थो म्रिये साहमनाथवत्} %||5-56-89||

\twolineshloka
{तच्छ्रुत्वा करुणं वाक्यं क्रोधो मामभ्यवर्तत}
{उत्तरं च मया दृष्टं कार्यशेषमनन्तरम्} %||5-56-90||

\twolineshloka
{ततोऽवर्धत मे कायस्तदा पर्वतसंनिभः}
{युद्धकाङ्क्षी वनं तच्च विनाशयितुमारभे} %||5-56-91||

\twolineshloka
{तद्भग्नं वनषण्डं तु भ्रान्तत्रस्तमृगद्विजम्}
{प्रतिबुद्धा निरीक्षन्ते राक्षस्यो विकृताननाः} %||5-56-92||

\twolineshloka
{मां च दृष्ट्वा वने तस्मिन्समागम्य ततस्ततः}
{ताः समभ्यागताः क्षिप्रं रावणायाचचक्षिरे} %||5-56-93||

\twolineshloka
{राजन्वनमिदं दुर्गं तव भग्नं दुरात्मना}
{वानरेण ह्यविज्ञाय तव वीर्यं महाबल} %||5-56-94||

\twolineshloka
{दुर्बुद्धेस्तस्य राजेन्द्र तव विप्रियकारिणः}
{वधमाज्ञापय क्षिप्रं यथासौ विलयं व्रजेत्} %||5-56-95||

\twolineshloka
{तच्छ्रुत्वा राक्षसेन्द्रेण विसृष्टा भृशदुर्जयाः}
{राक्षसाः किंकरा नाम रावणस्य मनोऽनुगाः} %||5-56-96||

\twolineshloka
{तेषामशीतिसाहस्रं शूलमुद्गरपाणिनाम्}
{मया तस्मिन्वनोद्देशे परिघेण निषूदितम्} %||5-56-97||

\twolineshloka
{तेषां तु हतशेषा ये ते गता लघुविक्रमाः}
{निहतं च मया सैन्यं रावणायाचचक्षिरे} %||5-56-98||

{ततो मे बुद्धिरुत्पन्ना चैत्यप्रासादमाक्रमम्॥\devanumber{99}॥} %||5-56-99|| (Check)

\addtocounter{shlokacount}{1}

\twolineshloka
{तत्रस्थान्राक्षसान्हत्वा शतं स्तम्भेन वै पुनः}
{ललाम भूतो लङ्काया मया विध्वंसितो रुषा} %||5-56-100||

{ततः प्रहस्तस्य सुतं जम्बुमालिनमादिशत्॥\devanumber{101}॥} %||5-56-101|| (Check)

\addtocounter{shlokacount}{1}

\twolineshloka
{तमहं बलसंपन्नं राक्षसं रणकोविदम्}
{परिघेणातिघोरेण सूदयामि सहानुगम्} %||5-56-102||

\threelineshloka
{तच्छ्रुत्वा राक्षसेन्द्रस्तु मन्त्रिपुत्रान्महाबलान्}
{पदातिबलसंपन्नान्प्रेषयामास रावणः}
{परिघेणैव तान्सर्वान्नयामि यमसादनम्} %||5-56-103||

\threelineshloka
{मन्त्रिपुत्रान्हताञ्श्रुत्वा समरे लघुविक्रमान्}
{पञ्चसेनाग्रगाञ्शूरान्प्रेषयामास रावणः}
{तानहं सह सैन्यान्वै सर्वानेवाभ्यसूदयम्} %||5-56-104||

\twolineshloka
{ततः पुनर्दशग्रीवः पुत्रमक्षं महाबलम्}
{बहुभी राकसैः सार्धं प्रेषयामास संयुगे} %||5-56-105||

\threelineshloka
{तं तु मन्दोदरी पुत्रं कुमारं रणपण्डितम्}
{सहसा खं समुत्क्रान्तं पादयोश्च गृहीतवान्}
{चर्मासिनं शतगुणं भ्रामयित्वा व्यपेषयम्} %||5-56-106||

\threelineshloka
{तमक्षमागतं भग्नं निशम्य स दशाननः}
{तत इन्द्रजितं नाम द्वितीयं रावणः सुतम्}
{व्यादिदेश सुसंक्रुद्धो बलिनं युद्धदुर्मदम्} %||5-56-107||

\twolineshloka
{तस्याप्यहं बलं सर्वं तं च राक्षसपुंगवम्}
{नष्टौजसं रणे कृत्वा परं हर्षमुपागमम्} %||5-56-108||

\twolineshloka
{महता हि महाबाहुः प्रत्ययेन महाबलः}
{प्रेषितो रावणेनैष सह वीरैर्मदोत्कटैः} %||5-56-109||

\twolineshloka
{ब्राह्मेणास्त्रेण स तु मां प्रबध्नाच्चातिवेगतः}
{रज्जूभिरभिबध्नन्ति ततो मां तत्र राक्षसाः} %||5-56-110||

\twolineshloka
{रावणस्य समीपं च गृहीत्वा मामुपानयन्}
{दृष्ट्वा संभाषितश्चाहं रावणेन दुरात्मना} %||5-56-111||

\twolineshloka
{पृष्टश्च लङ्कागमनं राक्षसानां च तद्वधम्}
{तत्सर्वं च मया तत्र सीतार्थमिति जल्पितम्} %||5-56-112||

\twolineshloka
{अस्याहं दर्शनाकाङ्क्षी प्राप्तस्त्वद्भवनं विभो}
{मारुतस्यौरसः पुत्रो वानरो हनुमानहम्} %||5-56-113||

\twolineshloka
{रामदूतं च मां विद्धि सुग्रीवसचिवं कपिम्}
{सोऽहं दौत्येन रामस्य त्वत्समीपमिहागतः} %||5-56-114||

\threelineshloka
{शृणु चापि समादेशं यदहं प्रब्रवीमि ते}
{राक्षसेश हरीशस्त्वां वाक्यमाह समाहितम्}
{धर्मार्थकामसहितं हितं पथ्यमिवाशनम्} %||5-56-115||

\twolineshloka
{वसतो ऋष्यमूके मे पर्वते विपुलद्रुमे}
{राघवो रणविक्रान्तो मित्रत्वं समुपागतः} %||5-56-116||

\twolineshloka
{तेन मे कथितं राजन्भार्या मे रक्षसा हृता}
{तत्र साहाय्यहेतोर्मे समयं कर्तुमर्हसि} %||5-56-117||

\twolineshloka
{वालिना हृतराज्येन सुग्रीवेण सह प्रभुः}
{चक्रेऽग्निसाक्षिकं सक्यं राघवः सहलक्ष्मणः} %||5-56-118||

\twolineshloka
{तेन वालिनमुत्साद्य शरेणैकेन संयुगे}
{वानराणां महाराजः कृतः संप्लवतां प्रभुः} %||5-56-119||

\twolineshloka
{तस्य साहाय्यमस्माभिः कार्यं सर्वात्मना त्विह}
{तेन प्रस्थापितस्तुभ्यं समीपमिह धर्मतः} %||5-56-120||

\twolineshloka
{क्षिप्रमानीयतां सीता दीयतां राघवस्य च}
{यावन्न हरयो वीरा विधमन्ति बलं तव} %||5-56-121||

\twolineshloka
{वानराणां प्रभवो हि न केन विदितः पुरा}
{देवतानां सकाशं च ये गच्छन्ति निमन्त्रिताः} %||5-56-122||

\twolineshloka
{इति वानरराजस्त्वामाहेत्यभिहितो मया}
{मामैक्षत ततो रुष्टश्चक्षुषा प्रदहन्निव} %||5-56-123||

{तेन वध्योऽहमाज्ञप्तो रक्षसा रौद्रकर्मणा॥\devanumber{124}॥} %||5-56-124|| (Check)

\addtocounter{shlokacount}{1}

\twolineshloka
{ततो विभीषणो नाम तस्य भ्राता महामतिः}
{तेन राक्षसराजोऽसौ याचितो मम कारणात्} %||5-56-125||

\twolineshloka
{दूतवध्या न दृष्टा हि राजशास्त्रेषु राक्षस}
{दूतेन वेदितव्यं च यथार्थं हितवादिना} %||5-56-126||

\twolineshloka
{सुमहत्यपराधेऽपि दूतस्यातुलविक्रमः}
{विरूपकरणं दृष्टं न वधोऽस्तीह शास्त्रतः} %||5-56-127||

\twolineshloka
{विभीषणेनैवमुक्तो रावणः संदिदेश तान्}
{राक्षसानेतदेवाद्य लाङ्गूलं दह्यतामिति} %||5-56-128||

\twolineshloka
{ततस्तस्य वचः श्रुत्वा मम पुच्छं समन्ततः}
{वेष्टितं शणवल्कैश्च पटैः कार्पासकैस्तथा} %||5-56-129||

\twolineshloka
{राक्षसाः सिद्धसंनाहास्ततस्ते चण्डविक्रमाः}
{तदादीप्यन्त मे पुच्छं हनन्तः काष्ठमुष्टिभिः} %||5-56-130||

\twolineshloka
{बद्धस्य बहुभिः पाशैर्यन्त्रितस्य च राक्षसैः}
{न मे पीडा भवेत्काचिद्दिदृक्षोर्नगरीं दिवा} %||5-56-131||

\twolineshloka
{ततस्ते राक्षसाः शूरा बद्धं मामग्निसंवृतम्}
{अघोषयन्राजमार्गे नगरद्वारमागताः} %||5-56-132||

\twolineshloka
{ततोऽहं सुमहद्रूपं संक्षिप्य पुनरात्मनः}
{विमोचयित्वा तं बन्धं प्रकृतिष्ठः स्थितः पुनः} %||5-56-133||

\twolineshloka
{आयसं परिघं गृह्य तानि रक्षांस्यसूदयम्}
{ततस्तन्नगरद्वारं वेगेनाप्लुतवानहम्} %||5-56-134||

\twolineshloka
{पुच्छेन च प्रदीप्तेन तां पुरीं साट्टगोपुराम्}
{दहाम्यहमसंभ्रान्तो युगान्ताग्निरिव प्रजाः} %||5-56-135||

\twolineshloka
{दग्ध्वा लङ्कां पुनश्चैव शङ्का मामभ्यवर्तत}
{दहता च मया लङ्कां दघ्दा सीता न संशयः} %||5-56-136||

\twolineshloka
{अथाहं वाचमश्रौषं चारणानां शुभाक्षराम्}
{जानकी न च दग्धेति विस्मयोदन्तभाषिणाम्} %||5-56-137||

\twolineshloka
{ततो मे बुद्धिरुत्पन्ना श्रुत्वा तामद्भुतां गिरम्}
{पुनर्दृष्टा च वैदेही विसृष्टश्च तया पुनः} %||5-56-138||

\twolineshloka
{राघवस्य प्रभावेन भवतां चैव तेजसा}
{सुग्रीवस्य च कार्यार्थं मया सर्वमनुष्ठितम्} %||5-56-139||

\twolineshloka
{एतत्सर्वं मया तत्र यथावदुपपादितम्}
{अत्र यन्न कृतं शेषं तत्सर्वं क्रियतामिति} %||5-57-140||


{॥इत्यार्षे श्रीमद्रामायणे वाल्मीकीये आदिकाव्ये सुन्दरकाण्डे षट्पञ्चाशत्तमः सर्गः॥}

\sect{सप्तपञ्चाशत्तमः सर्गः}

\twolineshloka
{एतदाख्यानं तत्सर्वं हनूमान्मारुतात्मजः}
{भूयः समुपचक्राम वचनं वक्तुमुत्तरम्} %||5-57-1||

\twolineshloka
{सफलो राघवोद्योगः सुग्रीवस्य च संभ्रमः}
{शीलमासाद्य सीताया मम च प्लवनं महत्} %||5-57-2||

\twolineshloka
{आर्यायाः सदृशं शीलं सीतायाः प्लवगर्षभाः}
{तपसा धारयेल्लोकान्क्रुद्धा वा निर्दहेदपि} %||5-57-3||

\twolineshloka
{सर्वथातिप्रवृद्धोऽसौ रावणो राक्षसाधिपः}
{यस्य तां स्पृशतो गात्रं तपसा न विनाशितम्} %||5-57-4||

\twolineshloka
{न तदग्निशिखा कुर्यात्संस्पृष्टा पाणिना सती}
{जनकस्यात्मजा कुर्यादुत्क्रोधकलुषीकृता} %||5-57-5||

\twolineshloka
{अशोकवनिकामध्ये रावणस्य दुरात्मनः}
{अधस्ताच्छिंशपावृक्षे साध्वी करुणमास्थिता} %||5-57-6||

\twolineshloka
{राक्षसीभिः परिवृता शोकसंतापकर्शिता}
{मेघलेखापरिवृता चन्द्रलेखेव निष्प्रभा} %||5-57-7||

\twolineshloka
{अचिन्तयन्ती वैदेही रावणं बलदर्पितम्}
{पतिव्रता च सुश्रोणी अवष्टब्धा च जानकी} %||5-57-8||

\twolineshloka
{अनुरक्ता हि वैदेही रामं सर्वात्मना शुभा}
{अनन्यचित्ता रामे च पौलोमीव पुरंदरे} %||5-57-9||

\twolineshloka
{तदेकवासःसंवीता रजोध्वस्ता तथैव च}
{शोकसंतापदीनाङ्गी सीता भर्तृहिते रता} %||5-57-10||

\twolineshloka
{सा मया राक्षसी मध्ये तर्ज्यमाना मुहुर्मुहुः}
{राक्षसीभिर्विरूपाभिर्दृष्टा हि प्रमदा वने} %||5-57-11||

\twolineshloka
{एकवेणीधरा दीना भर्तृचिन्तापरायणा}
{अधःशय्या विवर्णाङ्गी पद्मिनीव हिमागमे} %||5-57-12||

\twolineshloka
{रावणाद्विनिवृत्तार्था मर्तव्यकृतनिश्चया}
{कथंचिन्मृगशावाक्षी विश्वासमुपपादिता} %||5-57-13||

\twolineshloka
{ततः संभाषिता चैव सर्वमर्थं च दर्शिता}
{रामसुग्रीवसख्यं च श्रुत्वा प्रीतिमुपागता} %||5-57-14||

{नियतः समुदाचारो भक्तिर्भर्तरि चोत्तमा॥\devanumber{15}॥} %||5-57-15|| (Check)

\addtocounter{shlokacount}{1}

\twolineshloka
{यन्न हन्ति दशग्रीवं स महात्मा दशाननः}
{निमित्तमात्रं रामस्तु वधे तस्य भविष्यति} %||5-57-16||

\twolineshloka
{एवमास्ते महाभागा सीता शोकपरायणा}
{यदत्र प्रतिकर्तव्यं तत्सर्वमुपपाद्यताम्} %||5-58-17||


{॥इत्यार्षे श्रीमद्रामायणे वाल्मीकीये आदिकाव्ये सुन्दरकाण्डे सप्तपञ्चाशत्तमः सर्गः॥}

\sect{अष्टपञ्चाशत्तमः सर्गः}

\twolineshloka
{तस्य तद्वचनं श्रुत्वा वालिसूनुरभाषत}
{जाम्बवत्प्रमुखान्सर्वाननुज्ञाप्य महाकपीन्} %||5-58-1||

\twolineshloka
{अस्मिन्नेवंगते कार्ये भवतां च निवेदिते}
{न्याय्यं स्म सह वैदेह्या द्रष्टुं तौ पार्थिवात्मजौ} %||5-58-2||

\twolineshloka
{अहमेकोऽपि पर्याप्तः सराक्षसगणां पुरीम्}
{तां लङ्कां तरसा हन्तुं रावणं च महाबलम्} %||5-58-3||

\twolineshloka
{किं पुनः सहितो वीरैर्बलवद्भिः कृतात्मभिः}
{कृतास्त्रैः प्लवगैः शक्तैर्भवद्भिर्विजयैषिभिः} %||5-58-4||

\twolineshloka
{अहं तु रावणं युद्धे ससैन्यं सपुरःसरम्}
{सपुत्रं विधमिष्यामि सहोदरयुतं युधि} %||5-58-5||

\threelineshloka
{ब्राह्ममैन्द्रं च रौद्रं च वायव्यं वारुणं तथा}
{यदि शक्रजितोऽस्त्राणि दुर्निरीक्ष्याणि संयुगे}
{तान्यहं विधमिष्यामि निहनिष्यामि राक्षसान्} %||5-58-6||

{भवतामभ्यनुज्ञातो विक्रमो मे रुणद्धि तम्॥\devanumber{7}॥} %||5-58-7|| (Check)

\addtocounter{shlokacount}{1}

\twolineshloka
{मयातुला विसृष्टा हि शैलवृष्टिर्निरन्तरा}
{देवानपि रणे हन्यात्किं पुनस्तान्निशाचरान्} %||5-58-8||

\twolineshloka
{सागरोऽप्यतियाद्वेलां मन्दरः प्रचलेदपि}
{न जाम्बवन्तं समरे कम्पयेदरिवाहिनी} %||5-58-9||

\twolineshloka
{सर्वराक्षससंघानां राक्षसा ये च पूर्वकाः}
{अलमेको विनाशाय वीरो वायुसुतः कपिः} %||5-58-10||

\twolineshloka
{पनसस्योरुवेगेन नीलस्य च महात्मनः}
{मन्दरोऽप्यवशीर्येत किं पुनर्युधि राक्षसाः} %||5-58-11||

\twolineshloka
{सदेवासुरयुद्धेषु गन्धर्वोरगपक्षिषु}
{मैन्दस्य प्रतियोद्धारं शंसत द्विविदस्य वा} %||5-58-12||

\twolineshloka
{अश्विपुत्रौ महावेगावेतौ प्लवगसत्तमौ}
{पितामहवरोत्सेकात्परमं दर्पमास्थितौ} %||5-58-13||

\twolineshloka
{अश्विनोर्माननार्थं हि सर्वलोकपितामहः}
{सर्वावध्यत्वमतुलमनयोर्दत्तवान्पुरा} %||5-58-14||

\twolineshloka
{वरोत्सेकेन मत्तौ च प्रमथ्य महतीं चमूम्}
{सुराणाममृतं वीरौ पीतवन्तौ प्लवंगमौ} %||5-58-15||

\twolineshloka
{एतावेव हि संक्रुद्धौ सवाजिरथकुञ्जराम्}
{लङ्कां नाशयितुं शक्तौ सर्वे तिष्ठन्तु वानराः} %||5-58-16||

\twolineshloka
{अयुक्तं तु विना देवीं दृष्टबद्भिः प्लवंगमाः}
{समीपं गन्तुमस्माभी राघवस्य महात्मनः} %||5-58-17||

\twolineshloka
{दृष्टा देवी न चानीता इति तत्र निवेदनम्}
{अयुक्तमिव पश्यामि भवद्भिः ख्यातविक्रमैः} %||5-58-18||

\twolineshloka
{न हि वः प्लवते कश्चिन्नापि कश्चित्पराक्रमे}
{तुल्यः सामरदैत्येषु लोकेषु हरिसत्तमाः} %||5-58-19||

\twolineshloka
{तेष्वेवं हतवीरेषु राक्षसेषु हनूमता}
{किमन्यदत्र कर्तव्यं गृहीत्वा याम जानकीम्} %||5-58-20||

\twolineshloka
{तमेवं कृतसंकल्पं जाम्बवान्हरिसत्तमः}
{उवाच परमप्रीतो वाक्यमर्थवदर्थवित्} %||5-58-21||

\fourlineindentedshloka
{न तावदेषा मतिरक्षमा नो}
{ यथा भवान्पश्यति राजपुत्र}
{यथा तु रामस्य मतिर्निविष्टा}
{ तथा भवान्पश्यतु कार्यसिद्धिम्} %||5-59-22||


{॥इत्यार्षे श्रीमद्रामायणे वाल्मीकीये आदिकाव्ये सुन्दरकाण्डे अष्टपञ्चाशत्तमः सर्गः॥}

\sect{एकोनषष्टितमः सर्गः}

\twolineshloka
{ततो जाम्बवतो वाक्यमगृह्णन्त वनौकसः}
{अङ्गदप्रमुखा वीरा हनूमांश्च महाकपिः} %||5-59-1||

\twolineshloka
{प्रीतिमन्तस्ततः सर्वे वायुपुत्रपुरःसराः}
{महेन्द्राग्रं परित्यज्य पुप्लुवुः प्लवगर्षभाः} %||5-59-2||

\twolineshloka
{मेरुमन्दरसंकाशा मत्ता इव महागजाः}
{छादयन्त इवाकाशं महाकाया महाबलाः} %||5-59-3||

\twolineshloka
{सभाज्यमानं भूतैस्तमात्मवन्तं महाबलम्}
{हनूमन्तं महावेगं वहन्त इव दृष्टिभिः} %||5-59-4||

\twolineshloka
{राघवे चार्थनिर्वृत्तिं भर्तुश्च परमं यशः}
{समाधाय समृद्धार्थाः कर्मसिद्धिभिरुन्नताः} %||5-59-5||

\twolineshloka
{प्रियाख्यानोन्मुखाः सर्वे सर्वे युद्धाभिनन्दिनः}
{सर्वे रामप्रतीकारे निश्चितार्था मनस्विनः} %||5-59-6||

\twolineshloka
{प्लवमानाः खमाप्लुत्य ततस्ते काननौक्षकः}
{नन्दनोपममासेदुर्वनं द्रुमलतायुतम्} %||5-59-7||

\twolineshloka
{यत्तन्मधुवनं नाम सुग्रीवस्याभिरक्षितम्}
{अधृष्यं सर्वभूतानां सर्वभूतमनोहरम्} %||5-59-8||

\twolineshloka
{यद्रक्षति महावीर्यः सदा दधिमुखः कपिः}
{मातुलः कपिमुख्यस्य सुग्रीवस्य महात्मनः} %||5-59-9||

\twolineshloka
{ते तद्वनमुपागम्य बभूवुः परमोत्कटाः}
{वानरा वानरेन्द्रस्य मनःकान्ततमं महत्} %||5-59-10||

\twolineshloka
{ततस्ते वानरा हृष्टा दृष्ट्वा मधुवनं महत्}
{कुमारमभ्ययाचन्त मधूनि मधुपिङ्गलाः} %||5-59-11||

\twolineshloka
{ततः कुमारस्तान्वृद्धाञ्जाम्बवत्प्रमुखान्कपीन्}
{अनुमान्य ददौ तेषां निसर्गं मधुभक्षणे} %||5-59-12||

\twolineshloka
{ततश्चानुमताः सर्वे संप्रहृष्टा वनौकसः}
{मुदिताश्च ततस्ते च प्रनृत्यन्ति ततस्ततः} %||5-59-13||

\fourlineindentedshloka
{गायन्ति केचित्प्रणमन्ति के चि}
{न्नृत्यन्ति केचित्प्रहसन्ति केचित्}
{पतन्ति केचिद्विचरन्ति के चि}
{त्प्लवन्ति केचित्प्रलपन्ति केचित्} %||5-59-14||

\fourlineindentedshloka
{परस्परं केचिदुपाश्रयन्ते}
{ परस्परं केचिदतिब्रुवन्ते}
{द्रुमाद्द्रुमं केचिदभिप्लवन्ते}
{ क्षितौ नगाग्रान्निपतन्ति केचित्} %||5-59-15||

\fourlineindentedshloka
{महीतलात्केचिदुदीर्णवेगा}
{ महाद्रुमाग्राण्यभिसंपतन्ते}
{गायन्तमन्यः प्रहसन्नुपैति}
{ हसन्तमन्यः प्रहसन्नुपैति} %||5-59-16||

\threelineshloka
{रुदन्तमन्यः प्ररुदन्नुपैति; नुदन्तमन्यः प्रणुदन्नुपैति}
{समाकुलं तत्कपिसैन्यमासी;न्मधुप्रपानोत्कट सत्त्वचेष्टम्}
{न चात्र कश्चिन्न बभूव मत्तो; न चात्र कश्चिन्न बभूव तृप्तो} %||5-59-17||

\fourlineindentedshloka
{ततो वनं तत्परिभक्ष्यमाणं}
{ द्रुमांश्च विध्वंसितपत्रपुष्पान्}
{समीक्ष्य कोपाद्दधिवक्त्रनामा}
{ निवारयामास कपिः कपींस्तान्} %||5-59-18||

\fourlineindentedshloka
{स तैः प्रवृद्धैः परिभर्त्स्यमानो}
{ वनस्य गोप्ता हरिवीरवृद्धः}
{चकार भूयो मतिमुग्रतेजा}
{ वनस्य रक्षां प्रति वानरेभ्यः} %||5-59-19||

\fourlineindentedshloka
{उवाच कांश्चित्परुषाणि धृष्ट}
{मसक्तमन्यांश्च तलैर्जघान}
{समेत्य कैश्चित्कलहं चकार}
{ तथैव साम्नोपजगाम कांश्चित्} %||5-59-20||

\fourlineindentedshloka
{स तैर्मदाच्चाप्रतिवार्य वेगै}
{र्बलाच्च तेनाप्रतिवार्यमाणैः}
{प्रधर्षितस्त्यक्तभयैः समेत्य}
{ प्रकृष्यते चाप्यनवेक्ष्य दोषम्} %||5-59-21||

\fourlineindentedshloka
{नखैस्तुदन्तो दशनैर्दशन्त}
{स्तलैश्च पादैश्च समाप्नुवन्तः}
{मदात्कपिं तं कपयः समग्रा}
{ महावनं निर्विषयं च चक्रुः} %||5-60-22||


{॥इत्यार्षे श्रीमद्रामायणे वाल्मीकीये आदिकाव्ये सुन्दरकाण्डे एकोनषष्टितमः सर्गः॥}

\sect{षष्टितमः सर्गः}

\twolineshloka
{तानुवाच हरिश्रेष्ठो हनूमान्वानरर्षभः}
{अव्यग्रमनसो यूयं मधु सेवत वानराः} %||5-60-1||

\twolineshloka
{श्रुत्वा हनुमतो वाक्यं हरीणां प्रवरोऽङ्गदः}
{प्रत्युवाच प्रसन्नात्मा पिबन्तु हरयो मधु} %||5-60-2||

\twolineshloka
{अवश्यं कृतकार्यस्य वाक्यं हनुमतो मया}
{अकार्यमपि कर्तव्यं किमङ्ग पुनरीदृशम्} %||5-60-3||

\twolineshloka
{अन्दगस्य मुखाच्छ्रुत्वा वचनं वानरर्षभाः}
{साधु साध्विति संहृष्टा वानराः प्रत्यपूजयन्} %||5-60-4||

\twolineshloka
{पूजयित्वाङ्गदं सर्वे वानरा वानरर्षभम्}
{जग्मुर्मधुवनं यत्र नदीवेग इव द्रुतम्} %||5-60-5||

\twolineshloka
{ते प्रहृष्टा मधुवनं पालानाक्रम्य वीर्यतः}
{अतिसर्गाच्च पटवो दृष्ट्वा श्रुत्वा च मैथिलीम्} %||5-60-6||

\twolineshloka
{उत्पत्य च ततः सर्वे वनपालान्समागताः}
{ताडयन्ति स्म शतशः सक्तान्मधुवने तदा} %||5-60-7||

\twolineshloka
{मधूनि द्रोणमात्राणि बहुभिः परिगृह्य ते}
{घ्नन्ति स्म सहिताः सर्वे भक्षयन्ति तथापरे} %||5-60-8||

\twolineshloka
{केचित्पीत्वापविध्यन्ति मधूनि मधुपिङ्गलाः}
{मधूच्चिष्टेन केचिच्च जघ्नुरन्योन्यमुत्कटाः} %||5-60-9||

\twolineshloka
{अपरे वृक्षमूलेषु शाखां गृह्य व्यवस्थितः}
{अत्यर्थं च मदग्लानाः पर्णान्यास्तीर्य शेरते} %||5-60-10||

\twolineshloka
{उन्मत्तभूताः प्लवगा मधुमत्ताश्च हृष्टवत्}
{क्षिपन्त्यपि तथान्योन्यं स्खलन्त्यपि तथापरे} %||5-60-11||

\twolineshloka
{केचित्क्ष्वेडान्प्रकुर्वन्ति केचित्कूजन्ति हृष्टवत्}
{हरयो मधुना मत्ताः केचित्सुप्ता महीतले} %||5-60-12||

\twolineshloka
{येऽप्यत्र मधुपालाः स्युः प्रेष्या दधिमुखस्य तु}
{तेऽपि तैर्वानरैर्भीमैः प्रतिषिद्धा दिशो गताः} %||5-60-13||

\twolineshloka
{जानुभिश्च प्रकृष्टाश्च देवमार्गं च दर्शिताः}
{अब्रुवन्परमोद्विग्ना गत्वा दधिमुखं वचः} %||5-60-14||

\twolineshloka
{हनूमता दत्तवरैर्हतं मधुवनं बलात्}
{वयं च जानुभिः कृष्टा देवमार्गं च दर्शिताः} %||5-60-15||

\twolineshloka
{ततो दधिमुखः क्रुद्धो वनपस्तत्र वानरः}
{हतं मधुवनं श्रुत्वा सान्त्वयामास तान्हरीन्} %||5-60-16||

\twolineshloka
{एतागच्छत गच्छामो वानरानतिदर्पितान्}
{बलेनावारयिष्यामो मधु भक्षयतो वयम्} %||5-60-17||

\twolineshloka
{श्रुत्वा दधिमुखस्येदं वचनं वानरर्षभाः}
{पुनर्वीरा मधुवनं तेनैव सहिता ययुः} %||5-60-18||

\twolineshloka
{मध्ये चैषां दधिमुखः प्रगृह्य सुमहातरुम्}
{समभ्यधावद्वेगेना ते च सर्वे प्लवंगमाः} %||5-60-19||

\twolineshloka
{ते शिलाः पादपांश्चापि पाषाणांश्चापि वानराः}
{गृहीत्वाभ्यागमन्क्रुद्धा यत्र ते कपिकुञ्जराः} %||5-60-20||

\twolineshloka
{ते स्वामिवचनं वीरा हृदयेष्ववसज्य तत्}
{त्वरया ह्यभ्यधावन्त सालतालशिलायुधाः} %||5-60-21||

\twolineshloka
{वृक्षस्थांश्च तलस्थांश्च वानरान्बलदर्पितान्}
{अभ्यक्रामन्त ते वीराः पालास्तत्र सहस्रशः} %||5-60-22||

\twolineshloka
{अथ दृष्ट्वा दधिमुखं क्रुद्धं वानरपुंगवाः}
{अभ्यधावन्त वेगेन हनूमत्प्रमुखास्तदा} %||5-60-23||

\twolineshloka
{तं सवृक्षं महाबाहुमापतन्तं महाबलम्}
{आर्यकं प्राहरत्तत्र बाहुभ्यां कुपितोऽङ्गदः} %||5-60-24||

\twolineshloka
{मदान्धश न वेदैनमार्यकोऽयं ममेति सः}
{अथैनं निष्पिपेषाशु वेगवद्वसुधातले} %||5-60-25||

\twolineshloka
{स भग्नबाहुर्विमुखो विह्वलः शोणितोक्षितः}
{मुमोह सहसा वीरो मुहूर्तं कपिकुञ्जरः} %||5-60-26||

\twolineshloka
{स कथंचिद्विमुक्तस्तैर्वानरैर्वानरर्षभः}
{उवाचैकान्तमागम्य भृत्यांस्तान्समुपागतान्} %||5-60-27||

\twolineshloka
{एते तिष्ठन्तु गच्छामो भर्ता नो यत्र वानरः}
{सुग्रीवो विपुलग्रीवः सह रामेण तिष्ठति} %||5-60-28||

\twolineshloka
{सर्वं चैवाङ्गदे दोषं श्रावयिष्यामि पार्थिव}
{अमर्षी वचनं श्रुत्वा घातयिष्यति वानरान्} %||5-60-29||

\twolineshloka
{इष्टं मधुवनं ह्येतत्सुग्रीवस्य महात्मनः}
{पितृपैतामहं दिव्यं देवैरपि दुरासदम्} %||5-60-30||

\twolineshloka
{स वानरानिमान्सर्वान्मधुलुब्धान्गतायुषः}
{घातयिष्यति दण्डेन सुग्रीवः ससुहृज्जनान्} %||5-60-31||

\twolineshloka
{वध्या ह्येते दुरात्मानो नृपाज्ञा परिभाविनः}
{अमर्षप्रभवो रोषः सफलो नो भविष्यति} %||5-60-32||

\twolineshloka
{एवमुक्त्वा दधिमुखो वनपालान्महाबलः}
{जगाम सहसोत्पत्य वनपालैः समन्वितः} %||5-60-33||

\twolineshloka
{निमेषान्तरमात्रेण स हि प्राप्तो वनालयः}
{सहस्रांशुसुतो धीमान्सुग्रीवो यत्र वानरः} %||5-60-34||

\twolineshloka
{रामं च लक्ष्मणं चैव दृष्ट्वा सुग्रीवमेव च}
{समप्रतिष्ठां जगतीमाकाशान्निपपात ह} %||5-60-35||

\twolineshloka
{स निपत्य महावीर्यः सर्वैस्तैः परिवारितः}
{हरिर्दधिमुखः पालैः पालानां परमेश्वरः} %||5-60-36||

\twolineshloka
{स दीनवदनो भूत्वा कृत्वा शिरसि चाञ्जलिम्}
{सुग्रीवस्य शुभौ मूर्ध्ना चरणौ प्रत्यपीडयत्} %||5-61-37||


{॥इत्यार्षे श्रीमद्रामायणे वाल्मीकीये आदिकाव्ये सुन्दरकाण्डे षष्टितमः सर्गः॥}

\sect{एकषष्टितमः सर्गः}

\twolineshloka
{ततो मूर्ध्ना निपतितं वानरं वानरर्षभः}
{दृष्ट्वैवोद्विग्नहृदयो वाक्यमेतदुवाच ह} %||5-61-1||

\twolineshloka
{उत्तिष्ठोत्तिष्ठ कस्मात्त्वं पादयोः पतितो मम}
{अभयं ते भवेद्वीर सत्यमेवाभिधीयताम्} %||5-61-2||

\twolineshloka
{स तु विश्वासितस्तेन सुग्रीवेण महात्मना}
{उत्थाय च महाप्राज्ञो वाक्यं दधिमुखोऽब्रवीत्} %||5-61-3||

\twolineshloka
{नैवर्क्षरजसा राजन्न त्वया नापि वालिना}
{वनं निसृष्टपूर्वं हि भक्षितं तत्तु वानरैः} %||5-61-4||

\twolineshloka
{एभिः प्रधर्षिताश्चैव वारिता वनरक्षिभिः}
{मधून्यचिन्तयित्वेमान्भक्षयन्ति पिबन्ति च} %||5-61-5||

\twolineshloka
{शिष्टमत्रापविध्यन्ति भक्षयन्ति तथापरे}
{निवार्यमाणास्ते सर्वे भ्रुवौ वै दर्शयन्ति हि} %||5-61-6||

\twolineshloka
{इमे हि संरब्धतरास्तथा तैः संप्रधर्षिताः}
{वारयन्तो वनात्तस्मात्क्रुद्धैर्वानरपुंगवैः} %||5-61-7||

\twolineshloka
{ततस्तैर्बहुभिर्वीरैर्वानरैर्वानरर्षभाः}
{संरक्तनयनैः क्रोधाद्धरयः संप्रचालिताः} %||5-61-8||

\twolineshloka
{पाणिभिर्निहताः केचित्केचिज्जानुभिराहताः}
{प्रकृष्टाश्च यथाकामं देवमार्गं च दर्शिताः} %||5-61-9||

\twolineshloka
{एवमेते हताः शूरास्त्वयि तिष्ठति भर्तरि}
{कृत्स्नं मधुवनं चैव प्रकामं तैः प्रभक्ष्यते} %||5-61-10||

\twolineshloka
{एवं विज्ञाप्यमानं तु सुग्रीवं वानरर्षभम्}
{अपृच्छत्तं महाप्राज्ञो लक्ष्मणः परवीरहा} %||5-61-11||

\twolineshloka
{किमयं वानरो राजन्वनपः प्रत्युपस्थितः}
{कं चार्थमभिनिर्दिश्य दुःखितो वाक्यमब्रवीत्} %||5-61-12||

\twolineshloka
{एवमुक्तस्तु सुग्रीवो लक्ष्मणेन महात्मना}
{लक्ष्मणं प्रत्युवाचेदं वाक्यं वाक्यविशारदः} %||5-61-13||

\twolineshloka
{आर्य लक्ष्मण संप्राह वीरो दधिमुखः कपिः}
{अङ्गदप्रमुखैर्वीरैर्भक्षितं मधुवानरैः} %||5-61-14||

\twolineshloka
{नैषामकृतकृत्यानामीदृशः स्यादुपक्रमः}
{वनं यथाभिपन्नं तैः साधितं कर्म वानरैः} %||5-61-15||

\twolineshloka
{दृष्टा देवी न संदेहो न चान्येन हनूमता}
{न ह्यन्यः साधने हेतुः कर्मणोऽस्य हनूमतः} %||5-61-16||

\twolineshloka
{कार्यसिद्धिर्हनुमति मतिश्च हरिपुंगव}
{व्यवसायश्च वीर्यं च श्रुतं चापि प्रतिष्ठितम्} %||5-61-17||

\twolineshloka
{जाम्बवान्यत्र नेता स्यादङ्गदस्य बलेश्वरः}
{हनूमांश्चाप्यधिष्ठाता न तस्य गतिरन्यथा} %||5-61-18||

\twolineshloka
{अङ्गदप्रमुखैर्वीरैर्हतं मधुवनं किल}
{विचिन्त्य दक्षिणामाशामागतैर्हरिपुंगवैः} %||5-61-19||

\threelineshloka
{आगतैश्च प्रविष्टं तद्यथा मधुवनं हि तैः}
{धर्षितं च वनं कृत्स्नमुपयुक्तं च वानरैः}
{वारिताः सहिताः पालास्तथा जानुभिराहताः} %||5-61-20||

\twolineshloka
{एतदर्थमयं प्राप्तो वक्तुं मधुरवागिह}
{नाम्ना दधिमुखो नाम हरिः प्रख्यातविक्रमः} %||5-61-21||

\twolineshloka
{दृष्टा सीता महाबाहो सौमित्रे पश्य तत्त्वतः}
{अभिगम्य यथा सर्वे पिबन्ति मधु वानराः} %||5-61-22||

\twolineshloka
{न चाप्यदृष्ट्वा वैदेहीं विश्रुताः पुरुषर्षभ}
{वनं दात्त वरं दिव्यं धर्षयेयुर्वनौकसः} %||5-61-23||

\twolineshloka
{ततः प्रहृष्टो धर्मात्मा लक्ष्मणः सहराघवः}
{श्रुत्वा कर्णसुखां वाणीं सुग्रीववदनाच्च्युताम्} %||5-61-24||

\threelineshloka
{प्राहृष्यत भृशं रामो लक्ष्मणश्च महायशाः}
{श्रुत्वा दधिमुखस्येदं सुग्रीवस्तु प्रहृष्य च}
{वनपालं पुनर्वाक्यं सुग्रीवः प्रत्यभाषत} %||5-61-25||

\twolineshloka
{प्रीतोऽस्मि सौम्य यद्भुक्तं वनं तैः कृतकर्मभिः}
{मर्षितं मर्षणीयं च चेष्टितं कृतकर्मणाम्} %||5-61-26||

\fourlineindentedshloka
{इच्छामि शीघ्रं हनुमत्प्रधाना}
{न्शाखामृगांस्तान्मृगराजदर्पान्}
{द्रष्टुं कृतार्थान्सह राघवाभ्यां}
{ श्रोतुं च सीताधिगमे प्रयत्नम्} %||5-62-27||


{॥इत्यार्षे श्रीमद्रामायणे वाल्मीकीये आदिकाव्ये सुन्दरकाण्डे एकषष्टितमः सर्गः॥}

\sect{द्विषष्टितमः सर्गः}

\twolineshloka
{सुग्रीवेणैवमुक्तस्तु हृष्टो दधिमुखः कपिः}
{राघवं लक्ष्मणं चैव सुग्रीवं चाभ्यवादयत्} %||5-62-1||

\twolineshloka
{स प्रणम्य च सुग्रीवं राघवौ च महाबलौ}
{वानरैः सहितैः शूरैर्दिवमेवोत्पपात ह} %||5-62-2||

\twolineshloka
{स यथैवागतः पूर्वं तथैव त्वरितो गतः}
{निपत्य गगनाद्भूमौ तद्वनं प्रविवेश ह} %||5-62-3||

\twolineshloka
{स प्रविष्टो मधुवनं ददर्श हरियूथपान्}
{विमदानुद्धतान्सर्वान्मेहमानान्मधूदकम्} %||5-62-4||

\twolineshloka
{स तानुपागमद्वीरो बद्ध्वा करपुटाञ्जलिम्}
{उवाच वचनं श्लक्ष्णमिदं हृष्टवदङ्गदम्} %||5-62-5||

\twolineshloka
{सौम्य रोषो न कर्तव्यो यदेभिरभिवारितः}
{अज्ञानाद्रक्षिभिः क्रोधाद्भवन्तः प्रतिषेधिताः} %||5-62-6||

\twolineshloka
{युवराजस्त्वमीशश्च वनस्यास्य महाबल}
{मौर्ख्यात्पूर्वं कृतो दोषस्तद्भवान्क्षन्तुमर्हति} %||5-62-7||

\twolineshloka
{यथैव हि पिता तेऽभूत्पूर्वं हरिगणेश्वरः}
{तथा त्वमपि सुग्रीवो नान्यस्तु हरिसत्तम} %||5-62-8||

\twolineshloka
{आख्यातं हि मया गत्वा पितृव्यस्य तवानघ}
{इहोपयानं सर्वेषामेतेषां वनचारिणाम्} %||5-62-9||

\twolineshloka
{स त्वदागमनं श्रुत्वा सहैभिर्हरियूथपैः}
{प्रहृष्टो न तु रुष्टोऽसौ वनं श्रुत्वा प्रधर्षितम्} %||5-62-10||

\twolineshloka
{प्रहृष्टो मां पितृव्यस्ते सुग्रीवो वानरेश्वरः}
{शीघ्रं प्रेषय सर्वांस्तानिति होवाच पार्थिवः} %||5-62-11||

\twolineshloka
{श्रुत्वा दधिमुखस्यैतद्वचनं श्लक्ष्णमङ्गदः}
{अब्रवीत्तान्हरिश्रेष्ठो वाक्यं वाक्यविशारदः} %||5-62-12||

\twolineshloka
{शङ्के श्रुतोऽयं वृत्तान्तो रामेण हरियूथपाः}
{तत्क्षमं नेह नः स्थातुं कृते कार्ये परंतपाः} %||5-62-13||

\twolineshloka
{पीत्वा मधु यथाकामं विश्रान्ता वनचारिणः}
{किं शेषं गमनं तत्र सुग्रीवो यत्र मे गुरुः} %||5-62-14||

\twolineshloka
{सर्वे यथा मां वक्ष्यन्ति समेत्य हरियूथपाः}
{तथास्मि कर्ता कर्तव्ये भवद्भिः परवानहम्} %||5-62-15||

\twolineshloka
{नाज्ञापयितुमीशोऽहं युवराजोऽस्मि यद्यपि}
{अयुक्तं कृतकर्माणो यूयं धर्षयितुं मया} %||5-62-16||

\twolineshloka
{ब्रुवतश्चाङ्गदश्चैवं श्रुत्वा वचनमव्ययम्}
{प्रहृष्टमनसो वाक्यमिदमूचुर्वनौकसः} %||5-62-17||

\twolineshloka
{एवं वक्ष्यति को राजन्प्रभुः सन्वानरर्षभ}
{ऐश्वर्यमदमत्तो हि सर्वोऽहमिति मन्यते} %||5-62-18||

\twolineshloka
{तव चेदं सुसदृशं वाक्यं नान्यस्य कस्यचित्}
{संनतिर्हि तवाख्याति भविष्यच्छुभभाग्यताम्} %||5-62-19||

\twolineshloka
{सर्वे वयमपि प्राप्तास्तत्र गन्तुं कृतक्षणाः}
{स यत्र हरिवीराणां सुग्रीवः पतिरव्ययः} %||5-62-20||

\twolineshloka
{त्वया ह्यनुक्तैर्हरिभिर्नैव शक्यं पदात्पदम्}
{क्वचिद्गन्तुं हरिश्रेष्ठ ब्रूमः सत्यमिदं तु ते} %||5-62-21||

\twolineshloka
{एवं तु वदतां तेषामङ्गदः प्रत्यभाषत}
{बाढं गच्छाम इत्युक्त्वा उत्पपात महीतलात्} %||5-62-22||

\twolineshloka
{उत्पतन्तमनूत्पेतुः सर्वे ते हरियूथपाः}
{कृत्वाकाशं निराकाशं यज्ञोत्क्षिप्ता इवानलाः} %||5-62-23||

\twolineshloka
{तेऽम्बरं सहसोत्पत्य वेगवन्तः प्लवंगमाः}
{विनदन्तो महानादं घना वातेरिता यथा} %||5-62-24||

\twolineshloka
{अङ्गदे ह्यननुप्राप्ते सुग्रीवो वानराधिपः}
{उवाच शोकोपहतं रामं कमललोचनम्} %||5-62-25||

\twolineshloka
{समाश्वसिहि भद्रं ते दृष्टा देवी न संशयः}
{नागन्तुमिह शक्यं तैरतीते समये हि नः} %||5-62-26||

\twolineshloka
{न मत्सकाशमागच्छेत्कृत्ये हि विनिपातिते}
{युवराजो महाबाहुः प्लवतां प्रवरोऽङ्गदः} %||5-62-27||

\twolineshloka
{यद्यप्यकृतकृत्यानामीदृशः स्यादुपक्रमः}
{भवेत्तु दीनवदनो भ्रान्तविप्लुतमानसः} %||5-62-28||

\twolineshloka
{पितृपैतामहं चैतत्पूर्वकैरभिरक्षितम्}
{न मे मधुवनं हन्यादहृष्टः प्लवगेश्वरः} %||5-62-29||

\threelineshloka
{कौसल्या सुप्रजा राम समाश्वसिहि सुव्रत}
{दृष्टा देवी न संदेहो न चान्येन हनूमता}
{न ह्यन्यः कर्मणो हेतुः साधने तद्विधो भवेत्} %||5-62-30||

\twolineshloka
{हनूमति हि सिद्धिश्च मतिश्च मतिसत्तम}
{व्यवसायश्च वीर्यं च सूर्ये तेज इव ध्रुवम्} %||5-62-31||

\twolineshloka
{जाम्बवान्यत्र नेता स्यादङ्गदश्च बलेश्वरः}
{हनूमांश्चाप्यधिष्ठाता न तस्य गतिरन्यथा} %||5-62-32||

{मा भूश्चिन्ता समायुक्तः संप्रत्यमितविक्रम॥\devanumber{33}॥} %||5-62-33|| (Check)

\addtocounter{shlokacount}{1}

\threelineshloka
{ततः किल किला शब्दं शुश्रावासन्नमम्बरे}
{हनूमत्कर्मदृप्तानां नर्दतां काननौकसाम्}
{किष्किन्धामुपयातानां सिद्धिं कथयतामिव} %||5-62-34||

\twolineshloka
{ततः श्रुत्वा निनादं तं कपीनां कपिसत्तमः}
{आयताञ्चितलाङ्गूलः सोऽभवद्धृष्टमानसः} %||5-62-35||

\twolineshloka
{आजग्मुस्तेऽपि हरयो रामदर्शनकाङ्क्षिणः}
{अङ्गदं पुरतः कृत्वा हनूमन्तं च वानरम्} %||5-62-36||

\twolineshloka
{तेऽङ्गदप्रमुखा वीराः प्रहृष्टाश्च मुदान्विताः}
{निपेतुर्हरिराजस्य समीपे राघवस्य च} %||5-62-37||

\twolineshloka
{हनूमांश्च महाबहुः प्रणम्य शिरसा ततः}
{नियतामक्षतां देवीं राघवाय न्यवेदयत्} %||5-62-38||

\twolineshloka
{निश्चितार्थं ततस्तस्मिन्सुग्रीवं पवनात्मजे}
{लक्ष्मणः प्रीतिमान्प्रीतं बहुमानादवैक्षत} %||5-62-39||

\twolineshloka
{प्रीत्या च रममाणोऽथ राघवः परवीरहा}
{बहु मानेन महता हनूमन्तमवैक्षत} %||5-63-40||


{॥इत्यार्षे श्रीमद्रामायणे वाल्मीकीये आदिकाव्ये सुन्दरकाण्डे द्विषष्टितमः सर्गः॥}

\sect{त्रिषष्टितमः सर्गः}

\twolineshloka
{ततः प्रस्रवणं शैलं ते गत्वा चित्रकाननम्}
{प्रणम्य शिरसा रामं लक्ष्मणं च महाबलम्} %||5-63-1||

\twolineshloka
{युवराजं पुरस्कृत्य सुग्रीवमभिवाद्य च}
{प्रवृत्तमथ सीतायाः प्रवक्तुमुपचक्रमुः} %||5-63-2||

\twolineshloka
{रावणान्तःपुरे रोधं राक्षसीभिश्च तर्जनम्}
{रामे समनुरागं च यश्चापि समयः कृतः} %||5-63-3||

\twolineshloka
{एतदाख्यान्ति ते सर्वे हरयो राम संनिधौ}
{वैदेहीमक्षतां श्रुत्वा रामस्तूत्तरमब्रवीत्} %||5-63-4||

\twolineshloka
{क्व सीता वर्तते देवी कथं च मयि वर्तते}
{एतन्मे सर्वमाख्यात वैदेहीं प्रति वानराः} %||5-63-5||

\twolineshloka
{रामस्य गदितं श्रुत्व हरयो रामसंनिधौ}
{चोदयन्ति हनूमन्तं सीतावृत्तान्तकोविदम्} %||5-63-6||

\twolineshloka
{श्रुत्वा तु वचनं तेषां हनूमान्मारुतात्मजः}
{उवाच वाक्यं वाक्यज्ञः सीताया दर्शनं यथा} %||5-63-7||

\twolineshloka
{समुद्रं लङ्घयित्वाहं शतयोजनमायतम्}
{अगच्छं जानकीं सीतां मार्गमाणो दिदृक्षया} %||5-63-8||

\twolineshloka
{तत्र लङ्केति नगरी रावणस्य दुरात्मनः}
{दक्षिणस्य समुद्रस्य तीरे वसति दक्षिणे} %||5-63-9||

\twolineshloka
{तत्र दृष्टा मया सीता रावणान्तःपुरे सती}
{संन्यस्य त्वयि जीवन्ती रामा राम मनोरथम्} %||5-63-10||

\twolineshloka
{दृष्टा मे राक्षसी मध्ये तर्ज्यमाना मुहुर्मुहुः}
{राक्षसीभिर्विरूपाभी रक्षिता प्रमदावने} %||5-63-11||

\twolineshloka
{दुःखमापद्यते देवी तवादुःखोचिता सती}
{रावणान्तःपुरे रुद्ध्वा राक्षसीभिः सुरक्षिता} %||5-63-12||

\twolineshloka
{एकवेणीधरा दीना त्वयि चिन्तापरायणा}
{अधःशय्या विवर्णाङ्गी पद्मिनीव हिमागमे} %||5-63-13||

\twolineshloka
{रावणाद्विनिवृत्तार्था मर्तव्यकृतनिश्चया}
{देवी कथंचित्काकुत्स्थ त्वन्मना मार्गिता मया} %||5-63-14||

\twolineshloka
{इक्ष्वाकुवंशविख्यातिं शनैः कीर्तयतानघ}
{स मया नरशार्दूल विश्वासमुपपादिता} %||5-63-15||

\twolineshloka
{ततः संभाषिता देवी सर्वमर्थं च दर्शिता}
{रामसुग्रीवसख्यं च श्रुत्वा प्रीतिमुपागता} %||5-63-16||

\threelineshloka
{नियतः समुदाचारो भक्तिश्चास्यास्तथा त्वयि}
{एवं मया महाभागा दृष्टा जनकनन्दिनी}
{उग्रेण तपसा युक्ता त्वद्भक्त्या पुरुषर्षभ} %||5-63-17||

\twolineshloka
{अभिज्ञानं च मे दत्तं यथावृत्तं तवान्तिके}
{चित्रकूटे महाप्राज्ञ वायसं प्रति राघव} %||5-63-18||

\twolineshloka
{विज्ञाप्यश्च नर व्याघ्रो रामो वायुसुत त्वया}
{अखिलेनेह यद्दृष्टमिति मामाह जानकी} %||5-63-19||

\twolineshloka
{इदं चास्मै प्रदातव्यं यत्नात्सुपरिरक्षितम्}
{ब्रुवता वचनान्येवं सुग्रीवस्योपशृण्वतः} %||5-63-20||

\twolineshloka
{एष चूडामणिः श्रीमान्मया ते यत्नरक्षितः}
{मनःशिलायास्तिकलस्तं स्मरस्वेति चाब्रवीत्} %||5-63-21||

\twolineshloka
{एष निर्यातितः श्रीमान्मया ते वारिसंभवः}
{एतं दृष्ट्वा प्रमोदिष्ये व्यसने त्वामिवानघ} %||5-63-22||

\twolineshloka
{जीवितं धारयिष्यामि मासं दशरथात्मज}
{ऊर्ध्वं मासान्न जीवेयं रक्षसां वशमागता} %||5-63-23||

\twolineshloka
{इति मामब्रवीत्सीता कृशाङ्गी धर्म चारिणी}
{रावणान्तःपुरे रुद्धा मृगीवोत्फुल्ललोचना} %||5-63-24||

\twolineshloka
{एतदेव मयाख्यातं सर्वं राघव यद्यथा}
{सर्वथा सागरजले संतारः प्रविधीयताम्} %||5-63-25||

\fourlineindentedshloka
{तौ जाताश्वासौ राजपुत्रौ विदित्वा}
{ तच्चाभिज्ञानं राघवाय प्रदाय}
{देव्या चाख्यातं सर्वमेवानुपूर्व्या}
{द्वाचा संपूर्णं वायुपुत्रः शशंस} %||5-64-26||


{॥इत्यार्षे श्रीमद्रामायणे वाल्मीकीये आदिकाव्ये सुन्दरकाण्डे त्रिषष्टितमः सर्गः॥}

\sect{चतुःषष्टितमः सर्गः}

\twolineshloka
{एवमुक्तो हनुमता रामो दशरथात्मजः}
{तं मणिं हृदये कृत्वा प्ररुरोद सलक्ष्मणः} %||5-64-1||

\twolineshloka
{तं तु दृष्ट्वा मणिश्रेष्ठं राघवः शोककर्शितः}
{नेत्राभ्यामश्रुपूर्णाभ्यां सुग्रीवमिदमब्रवीत्} %||5-64-2||

\twolineshloka
{यथैव धेनुः स्रवति स्नेहाद्वत्सस्य वत्सला}
{तथा ममापि हृदयं मणिरत्नस्य दर्शनात्} %||5-64-3||

\twolineshloka
{मणिरत्नमिदं दत्तं वैदेह्याः श्वशुरेण मे}
{वधूकाले यथा बद्धमधिकं मूर्ध्नि शोभते} %||5-64-4||

\twolineshloka
{अयं हि जलसंभूतो मणिः प्रवरपूजितः}
{यज्ञे परमतुष्टेन दत्तः शक्रेण धीमता} %||5-64-5||

\twolineshloka
{इमं दृष्ट्वा मणिश्रेष्ठं तथा तातस्य दर्शनम्}
{अद्यास्म्यवगतः सौम्य वैदेहस्य तथा विभोः} %||5-64-6||

\twolineshloka
{अयं हि शोभते तस्याः प्रियाया मूर्ध्नि मे मणिः}
{अद्यास्य दर्शनेनाहं प्राप्तां तामिव चिन्तये} %||5-64-7||

\twolineshloka
{किमाह सीता वैदेही ब्रूहि सौम्य पुनः पुनः}
{परासुमिव तोयेन सिञ्चन्ती वाक्यवारिणा} %||5-64-8||

\twolineshloka
{इतस्तु किं दुःखतरं यदिमं वारिसंभवम्}
{मणिं पश्यामि सौमित्रे वैदेहीमागतं विना} %||5-64-9||

\twolineshloka
{चिरं जीवति वैदेही यदि मासं धरिष्यति}
{क्षणं सौम्य न जीवेयं विना तामसितेक्षणाम्} %||5-64-10||

\twolineshloka
{नय मामपि तं देशं यत्र दृष्टा मम प्रिया}
{न तिष्ठेयं क्षणमपि प्रवृत्तिमुपलभ्य च} %||5-64-11||

\twolineshloka
{कथं सा मम सुश्रोणि भीरु भीरुः सती तदा}
{भयावहानां घोराणां मध्ये तिष्ठति रक्षसाम्} %||5-64-12||

\twolineshloka
{शारदस्तिमिरोन्मुखो नूनं चन्द्र इवाम्बुदैः}
{आवृतं वदनं तस्या न विराजति राक्षसैः} %||5-64-13||

\twolineshloka
{किमाह सीता हनुमंस्तत्त्वतः कथयस्व मे}
{एतेन खलु जीविष्ये भेषजेनातुरो यथा} %||5-64-14||

\threelineshloka
{मधुरा मधुरालापा किमाह मम भामिनी}
{मद्विहीना वरारोहा हनुमन्कथयस्व मे}
{दुःखाद्दुःखतरं प्राप्य कथं जीवति जानकी} %||5-65-15||


{॥इत्यार्षे श्रीमद्रामायणे वाल्मीकीये आदिकाव्ये सुन्दरकाण्डे चतुःषष्टितमः सर्गः॥}

\sect{पञ्चषष्टितमः सर्गः}

\twolineshloka
{एवमुक्तस्तु हनुमान्राघवेण महात्मना}
{सीताया भाषितं सर्वं न्यवेदयत राघवे} %||5-65-1||

\twolineshloka
{इदमुक्तवती देवी जानकी पुरुषर्षभ}
{पूर्ववृत्तमभिज्ञानं चित्रकूटे यथा तथम्} %||5-65-2||

\twolineshloka
{सुखसुप्ता त्वया सार्धं जानकी पूर्वमुत्थिता}
{वायसः सहसोत्पत्य विरराद स्तनान्तरे} %||5-65-3||

\twolineshloka
{पर्यायेण च सुप्तस्त्वं देव्यङ्के भरताग्रज}
{पुनश्च किल पक्षी स देव्या जनयति व्यथाम्} %||5-65-4||

\twolineshloka
{ततः पुनरुपागम्य विरराद भृशं किल}
{ततस्त्वं बोधितस्तस्याः शोणितेन समुक्षितः} %||5-65-5||

\twolineshloka
{वायसेन च तेनैव सततं बाध्यमानया}
{बोधितः किल देव्यास्त्वं सुखसुप्तः परंतप} %||5-65-6||

\twolineshloka
{तां तु दृष्ट्वा महाबाहो रादितां च स्तनान्तरे}
{आशीविष इव क्रुद्धो निःश्वसन्नभ्यभाषथाः} %||5-65-7||

\twolineshloka
{नखाग्रैः केन ते भीरु दारितं तु स्तनान्तरम्}
{कः क्रीडति सरोषेण पञ्चवक्त्रेण भोगिना} %||5-65-8||

\twolineshloka
{निरीक्षमाणः सहसा वायसं समवैक्षताः}
{नखैः सरुधिरैस्तीक्ष्णैर्मामेवाभिमुखं स्थितम्} %||5-65-9||

\twolineshloka
{सुतः किल स शक्रस्य वायसः पततां वरः}
{धरान्तरचरः शीघ्रं पवनस्य गतौ समः} %||5-65-10||

\twolineshloka
{ततस्तस्मिन्महाबाहो कोपसंवर्तितेक्षणः}
{वायसे त्वं कृत्वाः क्रूरां मतिं मतिमतां वर} %||5-65-11||

\twolineshloka
{स दर्भं संस्तराद्गृह्य ब्रह्मास्त्रेण न्ययोजयः}
{स दीप्त इव कालाग्निर्जज्वालाभिमुखः खगम्} %||5-65-12||

\twolineshloka
{स त्वं प्रदीप्तं चिक्षेप दर्भं तं वायसं प्रति}
{ततस्तु वायसं दीप्तः स दर्भोऽनुजगाम ह} %||5-65-13||

\twolineshloka
{स पित्रा च परित्यक्तः सुरैः सर्वैर्महर्षिभिः}
{त्रीँल्लोकान्संपरिक्रम्य त्रातारं नाधिगच्छति} %||5-65-14||

\twolineshloka
{तं त्वं निपतितं भूमौ शरण्यः शरणागतम्}
{वधार्हमपि काकुत्स्थ कृपया परिपालयः} %||5-65-15||

\twolineshloka
{मोघमस्त्रं न शक्यं तु कर्तुमित्येव राघव}
{ततस्तस्याक्षिकाकस्य हिनस्ति स्म स दक्षिणम्} %||5-65-16||

\twolineshloka
{राम त्वां स नमस्कृत्वा राज्ञो दशरथस्य च}
{विसृष्टस्तु तदा काकः प्रतिपेदे खमालयम्} %||5-65-17||

\twolineshloka
{एवमस्त्रविदां श्रेष्ठः सत्त्ववाञ्शीलवानपि}
{किमर्थमस्त्रं रक्षःसु न योजयसि राघव} %||5-65-18||

\twolineshloka
{न नागा नापि गन्धर्वा नासुरा न मरुद्गणाः}
{तव राम मुखे स्थातुं शक्ताः प्रतिसमाधितुम्} %||5-65-19||

\twolineshloka
{तव वीर्यवतः कच्चिन्मयि यद्यस्ति संभ्रमः}
{क्षिप्रं सुनिशितैर्बाणैर्हन्यतां युधि रावणः} %||5-65-20||

\twolineshloka
{भ्रातुरादेशमादाय लक्ष्मणो वा परंतपः}
{स किमर्थं नरवरो न मां रक्षति राघवः} %||5-65-21||

\twolineshloka
{शक्तौ तौ पुरुषव्याघ्रौ वाय्वग्निसमतेजसौ}
{सुराणामपि दुर्धर्षौ किमर्थं मामुपेक्षतः} %||5-65-22||

\twolineshloka
{ममैव दुष्कृतं किंचिन्महदस्ति न संशयः}
{समर्थौ सहितौ यन्मां नापेक्षेते परंतपौ} %||5-65-23||

\twolineshloka
{वैदेह्या वचनं श्रुत्वा करुणं साश्रुभाषितम्}
{पुनरप्यहमार्यां तामिदं वचनमब्रुवम्} %||5-65-24||

\twolineshloka
{त्वच्छोकविमुखो रामो देवि सत्येन ते शपे}
{रामे दुःखाभिभूते च लक्ष्मणः परितप्यते} %||5-65-25||

\twolineshloka
{कथंचिद्भवती दृष्टा न कालः परिशोचितुम्}
{इमं मुहूर्तं दुःखानामन्तं द्रक्ष्यसि भामिनि} %||5-65-26||

\twolineshloka
{तावुभौ नरशार्दूलौ राजपुत्रावरिंदमौ}
{त्वद्दर्शनकृतोत्साहौ लङ्कां भस्मीकरिष्यतः} %||5-65-27||

\twolineshloka
{हत्वा च समरे रौद्रं रावणं सह बान्धवम्}
{राघवस्त्वां महाबाहुः स्वां पुरीं नयते ध्रुवम्} %||5-65-28||

\twolineshloka
{यत्तु रामो विजानीयादभिज्ञानमनिन्दिते}
{प्रीतिसंजननं तस्य प्रदातुं तत्त्वमर्हसि} %||5-65-29||

\twolineshloka
{साभिवीक्ष्य दिशः सर्वा वेण्युद्ग्रथनमुत्तमम्}
{मुक्त्वा वस्त्राद्ददौ मह्यं मणिमेतं महाबल} %||5-65-30||

\twolineshloka
{प्रतिगृह्य मणिं दिव्यं तव हेतो रघूत्तम}
{शिरसा संप्रणम्यैनामहमागमने त्वरे} %||5-65-31||

\threelineshloka
{गमने च कृतोत्साहमवेक्ष्य वरवर्णिनी}
{विवर्धमानं च हि मामुवाच जनकात्मजा}
{अश्रुपूर्णमुखी दीना बाष्पसंदिग्धभाषिणी} %||5-65-32||

\twolineshloka
{हनुमन्सिंहसंकाशौ तावुभौ रामलक्ष्मणौ}
{सुग्रीवं च सहामात्यं सर्वान्ब्रूया अनामयम्} %||5-65-33||

\twolineshloka
{यथा च स महाबाहुर्मां तारयति राघवः}
{अस्माद्दुःखाम्बुसंरोधात्तत्समाधातुमर्हसि} %||5-65-34||

\fourlineindentedshloka
{इमं च तीव्रं मम शोकवेगं}
{ रक्षोभिरेभिः परिभर्त्सनं च}
{ब्रूयास्तु रामस्य गतः समीपं}
{ शिवश्च तेऽध्वास्तु हरिप्रवीर} %||5-65-35||

\fourlineindentedshloka
{एतत्तवार्या नृपराजसिंह}
{ सीता वचः प्राह विषादपूर्वम्}
{एतच्च बुद्ध्वा गदितं मया त्वं}
{ श्रद्धत्स्व सीतां कुशलां समग्राम्} %||5-66-36||


{॥इत्यार्षे श्रीमद्रामायणे वाल्मीकीये आदिकाव्ये सुन्दरकाण्डे पञ्चषष्टितमः सर्गः॥}

\sect{षट्षष्टितमः सर्गः}

\twolineshloka
{अथाहमुत्तरं देव्या पुनरुक्तः ससंभ्रमम्}
{तव स्नेहान्नरव्याघ्र सौहार्यादनुमान्य च} %||5-66-1||

\twolineshloka
{एवं बहुविधं वाच्यो रामो दाशरथिस्त्वया}
{यथा मामाप्नुयाच्छीघ्रं हत्वा रावणमाहवे} %||5-66-2||

\twolineshloka
{यदि वा मन्यसे वीर वसैकाहमरिंदम}
{कस्मिंश्चित्संवृते देशे विश्रान्तः श्वो गमिष्यसि} %||5-66-3||

\twolineshloka
{मम चाप्यल्पभाग्यायाः साम्निध्यात्तव वानर}
{अस्य शोकविपाकस्य मुहूर्तं स्याद्विमोक्षणम्} %||5-66-4||

\twolineshloka
{गते हि त्वयि विक्रान्ते पुनरागमनाय वै}
{प्राणानामपि संदेहो मम स्यान्नात्र संशयः} %||5-66-5||

\twolineshloka
{तवादर्शनजः शोको भूयो मां परितापयेत्}
{दुःखाद्दुःखपराभूतां दुर्गतां दुःखभागिनीम्} %||5-66-6||

\twolineshloka
{अयं तु वीरसंदेहस्तिष्ठतीव ममाग्रतः}
{सुमहांस्त्वत्सहायेषु हर्यृक्षेषु असंशयः} %||5-66-7||

\twolineshloka
{कथं नु खलु दुष्पारं तरिष्यन्ति महोदधिम्}
{तानि हर्यृक्षसैन्यानि तौ वा नरवरात्मजौ} %||5-66-8||

\twolineshloka
{त्रयाणामेव भूतानां सागरस्यास्य लङ्घने}
{शक्तिः स्याद्वैनतेयस्य वायोर्वा तव वानघ} %||5-66-9||

\twolineshloka
{तदस्मिन्कार्यनियोगे वीरैवं दुरतिक्रमे}
{किं पश्यसि समाधानं ब्रूहि कार्यविदां वर} %||5-66-10||

\twolineshloka
{काममस्य त्वमेवैकः कार्यस्य परिसाधने}
{पर्याप्तः परवीरघ्न यशस्यस्ते बलोदयः} %||5-66-11||

\twolineshloka
{बलैः समग्रैर्यदि मां हत्वा रावणमाहवे}
{विजयी स्वां पुरीं रामो नयेत्तत्स्याद्यशस्करम्} %||5-66-12||

\twolineshloka
{यथाहं तस्य वीरस्य वनादुपधिना हृता}
{रक्षसा तद्भयादेव तथा नार्हति राघवः} %||5-66-13||

\twolineshloka
{बलैस्तु संकुलां कृत्वा लङ्कां परबलार्दनः}
{मां नयेद्यदि काकुत्स्थस्तत्तस्य सदृशं भवेत्} %||5-66-14||

\twolineshloka
{तद्यथा तस्य विक्रान्तमनुरूपं महात्मनः}
{भवत्याहवशूरस्य तथा त्वमुपपादय} %||5-66-15||

\twolineshloka
{तदर्थोपहितं वाक्यं प्रश्रितं हेतुसंहितम्}
{निशम्याहं ततः शेषं वाक्यमुत्तरमब्रुवम्} %||5-66-16||

\twolineshloka
{देवि हर्यृक्षसैन्यानामीश्वरः प्लवतां वरः}
{सुग्रीवः सत्त्वसंपन्नस्तवार्थे कृतनिश्चयः} %||5-66-17||

\twolineshloka
{तस्य विक्रमसंपन्नाः सत्त्ववन्तो महाबलाः}
{मनःसंकल्पसंपाता निदेशे हरयः स्थिताः} %||5-66-18||

\twolineshloka
{येषां नोपरि नाधस्तान्न तिर्यक्सज्जते गतिः}
{न च कर्मसु सीदन्ति महत्स्वमिततेजसः} %||5-66-19||

\twolineshloka
{असकृत्तैर्महाभागैर्वानरैर्बलसंयुतैः}
{प्रदक्षिणीकृता भूमिर्वायुमार्गानुसारिभिः} %||5-66-20||

\twolineshloka
{मद्विशिष्टाश्च तुल्याश्च सन्ति तत्र वनौकसः}
{मत्तः प्रत्यवरः कश्चिन्नास्ति सुग्रीवसंनिधौ} %||5-66-21||

\twolineshloka
{अहं तावदिह प्राप्तः किं पुनस्ते महाबलाः}
{न हि प्रकृष्टाः प्रेष्यन्ते प्रेष्यन्ते हीतरे जनाः} %||5-66-22||

\twolineshloka
{तदलं परितापेन देवि मन्युर्व्यपैतु ते}
{एकोत्पातेन ते लङ्कामेष्यन्ति हरियूथपाः} %||5-66-23||

\twolineshloka
{मम पृष्ठगतौ तौ च चन्द्रसूर्याविवोदितौ}
{त्वत्सकाशं महाभागे नृसिंहावागमिष्यतः} %||5-66-24||

\twolineshloka
{अरिघ्नं सिंहसंकाशं क्षिप्रं द्रक्ष्यसि राघवम्}
{लक्ष्मणं च धनुष्पाणिं लङ्का द्वारमुपस्थितम्} %||5-66-25||

\twolineshloka
{नखदंष्ट्रायुधान्वीरान्सिंहशार्दूलविक्रमान्}
{वानरान्वानरेन्द्राभान्क्षिप्रं द्रक्ष्यसि संगतान्} %||5-66-26||

\twolineshloka
{शैलाम्बुदन्निकाशानां लङ्कामलयसानुषु}
{नर्दतां कपिमुख्यानामचिराच्छोष्यसे स्वनम्} %||5-66-27||

\twolineshloka
{निवृत्तवनवासं च त्वया सार्धमरिंदमम्}
{अभिषिक्तमयोध्यायां क्षिप्रं द्रक्ष्यसि राघवम्} %||5-66-28||

\fourlineindentedshloka
{ततो मया वाग्भिरदीनभाषिणी}
{ शिवाभिरिष्टाभिरभिप्रसादिता}
{जगाम शान्तिं मम मैथिलात्मजा}
{ तवापि शोकेन तथाभिपीडिता} %||5-67-29||


{॥इत्यार्षे श्रीमद्रामायणे वाल्मीकीये आदिकाव्ये सुन्दरकाण्डे षट्षष्टितमः सर्गः॥}

